\documentclass[fs9,seriesRus]{korfarm-base}
\usepackage{korfarm-russell}

% ============================================================
% passage-note.tex (v11)
% 지문 박스 + 첫 페이지 우측 메모란 (동적 높이)
% 정책: PAGE_BREAK_POLICY.md §3, §4 참조
%
% 변경 (v11):
%  · 지문 박스 폭 확대: enlarge right by=-3.7cm  (메모 3.4cm + gap 3mm)
%  · 메모란 동적 높이: frame.north east ~ frame.south east 사이 자동 채움
%  · 메모 줄선: clip 영역으로 박스 내부만 채움 (박스 높이 모르고도 작동)
%  · 문단 들여쓰기 활성화: tcolorbox 내부 \parindent=1em, \noindent 제거
%
% 사용:
%   \kfPassageNote{제목}{지문 본문}      → 메모란 포함
%   \kfPassageNoteNT{지문 본문}          → 메모란 포함, 제목 없음
%   \kfPassageFull{제목}{본문}           → 메모란 없음 (활동 내 보기)
% ============================================================

\makeatletter

% 메모란 규격 (cm 고정)
\newlength{\kf@memoW}
\setlength{\kf@memoW}{3.4cm}
\newlength{\kf@memoGap}
\setlength{\kf@memoGap}{3mm}

% --- 메인: 제목 있는 버전 ---
\newcommand{\kf@passagenote@with}[2]{%
  \par\ifkfPassageNewPage\ifdim\pagetotal>0.4\textheight\clearpage\fi\fi\smallskip\needspace{6\baselineskip}%
  \begin{tcolorbox}[
    enhanced, breakable,
    width=\dimexpr\linewidth-4cm\relax,
    colback=kfPassageBg, colframe=kfPrimary!70,
    boxrule=0.6pt, arc=3pt,
    left=\kfBoxPad, right=\kfBoxPad, top=\dimexpr\kfBoxPad+8pt\relax, bottom=\kfBoxPad,
    title={\small\bfseries\color{kfPrimary} #1},
    coltitle=kfPrimary,
    colbacktitle=kfPrimaryLight!60,
    attach boxed title to top left={yshift=-2.4mm, xshift=8pt},
    boxed title style={colframe=kfPrimary!50, boxrule=0.5pt, arc=2pt},
    parbox=false,
    overlay unbroken and first={%
      \begin{scope}
        \draw[kfRule, line width=0.4pt, fill=kfNoteBg, rounded corners=2pt]
          ([xshift=\kf@memoGap]frame.north east) rectangle
          ([xshift=\kf@memoGap+\kf@memoW]frame.south east);
        \node[anchor=north east, font=\footnotesize\bfseries\color{kfMute}, inner sep=3pt]
          at ([xshift=\kf@memoGap+\kf@memoW, yshift=-2pt]frame.north east) {메모};
        \node[anchor=north east, fill=kfPrimary!15, text=kfPrimary,
              font=\scriptsize\bfseries, rounded corners=2pt,
              inner xsep=4pt, inner ysep=2pt]
          at ([xshift=\kf@memoGap+\kf@memoW-3pt, yshift=-19pt]frame.north east) {\kf@memocat};
        \begin{scope}
          \path[clip]
            ([xshift=\kf@memoGap]frame.north east) rectangle
            ([xshift=\kf@memoGap+\kf@memoW]frame.south east);
          \foreach \i in {1,...,40}{%
            \draw[kfRule, line width=0.25pt]
              ([xshift=\kf@memoGap+2mm, yshift=-7mm - \i*5mm]frame.north east) --
              ([xshift=\kf@memoGap+\kf@memoW-2mm, yshift=-7mm - \i*5mm]frame.north east);%
          }
        \end{scope}
      \end{scope}
    }
  ]
    \setlength{\parindent}{1em}%
    \hspace*{1em}\ignorespaces#2%
  \end{tcolorbox}\par\smallskip
}

% --- 제목 없는 버전 ---
\newcommand{\kf@passagenote@notitle}[1]{%
  \par\ifkfPassageNewPage\ifdim\pagetotal>0.4\textheight\clearpage\fi\fi\smallskip\needspace{6\baselineskip}%
  \begin{tcolorbox}[
    enhanced, breakable,
    width=\dimexpr\linewidth-4cm\relax,
    colback=kfPassageBg, colframe=kfPrimary!70,
    boxrule=0.6pt, arc=3pt,
    left=\kfBoxPad, right=\kfBoxPad, top=\dimexpr\kfBoxPad+8pt\relax, bottom=\kfBoxPad,
    parbox=false,
    overlay unbroken and first={%
      \begin{scope}
        \draw[kfRule, line width=0.4pt, fill=kfNoteBg, rounded corners=2pt]
          ([xshift=\kf@memoGap]frame.north east) rectangle
          ([xshift=\kf@memoGap+\kf@memoW]frame.south east);
        \node[anchor=north east, font=\footnotesize\bfseries\color{kfMute}, inner sep=3pt]
          at ([xshift=\kf@memoGap+\kf@memoW, yshift=-2pt]frame.north east) {메모};
        \node[anchor=north east, fill=kfPrimary!15, text=kfPrimary,
              font=\scriptsize\bfseries, rounded corners=2pt,
              inner xsep=4pt, inner ysep=2pt]
          at ([xshift=\kf@memoGap+\kf@memoW-3pt, yshift=-19pt]frame.north east) {\kf@memocat};
        \begin{scope}
          \path[clip]
            ([xshift=\kf@memoGap]frame.north east) rectangle
            ([xshift=\kf@memoGap+\kf@memoW]frame.south east);
          \foreach \i in {1,...,40}{%
            \draw[kfRule, line width=0.25pt]
              ([xshift=\kf@memoGap+2mm, yshift=-7mm - \i*5mm]frame.north east) --
              ([xshift=\kf@memoGap+\kf@memoW-2mm, yshift=-7mm - \i*5mm]frame.north east);%
          }
        \end{scope}
      \end{scope}
    }
  ]
    \setlength{\parindent}{1em}%
    \hspace*{1em}\ignorespaces#1%
  \end{tcolorbox}\par\smallskip
}

\newcommand{\kfPassageNote}[2]{%
  \def\kf@tmptitle{#1}%
  \ifx\kf@tmptitle\@empty
    \kf@passagenote@notitle{#2}%
  \else
    \kf@passagenote@with{#1}{#2}%
  \fi
}
\newcommand{\kfPassageNoteNT}[1]{\kf@passagenote@notitle{#1}}
\makeatother

% ============================================================
% 풀폭 지문 박스 (메모란 없음) — 활동 내 보기 박스
% PAGE_BREAK_POLICY.md §3
% ============================================================
\makeatletter
\newcommand{\kf@passagefull@with}[2]{%
  \par\ifkfPassageNewPage\ifdim\pagetotal>0.4\textheight\clearpage\fi\fi\smallskip\needspace{6\baselineskip}%
  \begin{tcolorbox}[
    enhanced, breakable,
    colback=kfPassageBg, colframe=kfPrimary!70,
    boxrule=0.6pt, arc=3pt,
    left=\kfBoxPad, right=\kfBoxPad, top=\dimexpr\kfBoxPad+8pt\relax, bottom=\kfBoxPad,
    title={\small\bfseries\color{kfPrimary} #1},
    coltitle=kfPrimary,
    colbacktitle=kfPrimaryLight!60,
    attach boxed title to top left={yshift=-2.4mm, xshift=8pt},
    boxed title style={colframe=kfPrimary!50, boxrule=0.5pt, arc=2pt},
    parbox=false
  ]
    \setlength{\parindent}{1em}%
    \hspace*{1em}\ignorespaces#2%
  \end{tcolorbox}\par\smallskip
}
\newcommand{\kf@passagefull@notitle}[1]{%
  \par\ifkfPassageNewPage\ifdim\pagetotal>0.4\textheight\clearpage\fi\fi\smallskip\needspace{6\baselineskip}%
  \begin{tcolorbox}[
    enhanced, breakable,
    colback=kfPassageBg, colframe=kfPrimary!70,
    boxrule=0.6pt, arc=3pt,
    left=\kfBoxPad, right=\kfBoxPad, top=\dimexpr\kfBoxPad+8pt\relax, bottom=\kfBoxPad,
    parbox=false
  ]
    \setlength{\parindent}{1em}%
    \hspace*{1em}\ignorespaces#1%
  \end{tcolorbox}\par\smallskip
}
\newcommand{\kfPassageFull}[2]{%
  \def\kf@tmptitle{#1}%
  \ifx\kf@tmptitle\@empty
    \kf@passagefull@notitle{#2}%
  \else
    \kf@passagefull@with{#1}{#2}%
  \fi
}
\newcommand{\kfPassageFullNT}[1]{\kf@passagefull@notitle{#1}}
\makeatother

% ============================================================
% 작은 문장 박스 (문장 독해용) — PAGE_BREAK_POLICY.md §2.5
% ============================================================
\newcommand{\kfSentenceBox}[2]{%
  \par\smallskip\needspace{3\baselineskip}%
  \noindent\begin{tcolorbox}[
    enhanced, breakable=false,
    colback=kfPassageBg, colframe=kfMute,
    boxrule=0.4pt, arc=2pt,
    left=8pt, right=6pt, top=4pt, bottom=4pt,
    title={\footnotesize\bfseries\color{kfPrimary} #1},
    coltitle=kfPrimary, colbacktitle=kfPaper,
    attach boxed title to top left={yshift=-1.6mm, xshift=4pt},
    boxed title style={colframe=kfMute, boxrule=0.3pt, arc=1pt}
  ]%
    \setstretch{1.3}#2%
  \end{tcolorbox}\par\smallskip%
}

% ============================================================
% 지문 본문 아래 안내/정리 박스 (v16 신규)
% 사용: \kfPassageHint{한 줄 정리}{본문 안내 텍스트}
% ============================================================
\newcommand{\kfPassageHint}[2]{%
  \par\smallskip\needspace{4\baselineskip}%
  \begin{tcolorbox}[
    enhanced, breakable=false,
    colback=kfPrimary!5, colframe=kfPrimary!40,
    boxrule=0.4pt, arc=2pt,
    left=8pt, right=6pt, top=4pt, bottom=4pt,
    title={\footnotesize\bfseries\color{kfPrimary} #1},
    coltitle=kfPrimary, colbacktitle=kfPaper,
    attach boxed title to top left={yshift=-1.5mm, xshift=4pt},
    boxed title style={colframe=kfPrimary!40, boxrule=0.3pt, arc=1pt}
  ]%
    \small\setstretch{1.3}#2%
  \end{tcolorbox}\par\smallskip%
}

% ============================================================
% 환경 버전 (호환성) — 메모란 포함
% ============================================================
\makeatletter
\newenvironment{kfPassageNoteEnv}[1]%
  {\par\needspace{6\baselineskip}%
   \def\kf@psrc{#1}%
   \begin{tcolorbox}[
     enhanced, breakable,
     width=\dimexpr\linewidth-4cm\relax,
     colback=kfPassageBg, colframe=kfPrimary!70,
     boxrule=0.6pt, arc=3pt,
     left=\kfBoxPad, right=\kfBoxPad, top=\dimexpr\kfBoxPad+8pt\relax, bottom=\kfBoxPad,
     title={\small\bfseries\color{kfPrimary} \kf@psrc},
     coltitle=kfPrimary, colbacktitle=kfPrimaryLight!60,
     attach boxed title to top left={yshift=-2.4mm, xshift=8pt},
     boxed title style={colframe=kfPrimary!50, boxrule=0.5pt, arc=2pt},
     parbox=false,
     overlay unbroken and first={%
       \begin{scope}
         \draw[kfRule, line width=0.4pt, fill=kfNoteBg, rounded corners=2pt]
           ([xshift=\kf@memoGap]frame.north east) rectangle
           ([xshift=\kf@memoGap+\kf@memoW]frame.south east);
         \node[anchor=north east, font=\footnotesize\bfseries\color{kfMute}, inner sep=3pt]
           at ([xshift=\kf@memoGap+\kf@memoW, yshift=-2pt]frame.north east) {메모};
         \begin{scope}
           \path[clip]
             ([xshift=\kf@memoGap]frame.north east) rectangle
             ([xshift=\kf@memoGap+\kf@memoW]frame.south east);
           \foreach \i in {1,...,40}{%
             \draw[kfRule, line width=0.25pt]
               ([xshift=\kf@memoGap+2mm, yshift=-7mm - \i*5mm]frame.north east) --
               ([xshift=\kf@memoGap+\kf@memoW-2mm, yshift=-7mm - \i*5mm]frame.north east);%
           }
         \end{scope}
       \end{scope}
     }
   ]%
   \setlength{\parindent}{1em}%
  }%
  {\end{tcolorbox}\par\smallskip}
\makeatother

% ============================================================
% condition-box.tex
% <조건> 박스 컴포넌트 (tcolorbox)
% 사용:
%   \begin{kfCondition}
%     \item 첫째 조건
%     \item 둘째 조건
%   \end{kfCondition}
% ============================================================

\newenvironment{kfCondition}%
  {\par\smallskip\needspace{4\baselineskip}%
   \begin{tcolorbox}[
     enhanced, breakable=false,
     colback=kfPrimaryLight!40, colframe=kfPrimary,
     boxrule=0.5pt, arc=2pt,
     left=\kfBoxPad, right=\kfBoxPad, top=\kfBoxPad, bottom=\kfBoxPad,
     title={\small\bfseries\color{kfPrimary} \,조건\,},
     coltitle=kfPrimary, colbacktitle=kfPaper,
     attach boxed title to top left={yshift=-2mm, xshift=6pt},
     boxed title style={colframe=kfPrimary, boxrule=0.4pt, arc=1pt}
   ]%
   \begin{enumerate}[leftmargin=16pt, itemsep=1pt, topsep=1pt,
                    label=\textbf{\arabic*.}]}%
  {\end{enumerate}\end{tcolorbox}\par\smallskip}

% ============================================================
% evidence-box.tex
% <보기> 박스 컴포넌트
% 사용:
%   \begin{kfEvidence}
%     보기 본문 ...
%   \end{kfEvidence}
% ============================================================

\newenvironment{kfEvidence}%
  {\par\smallskip\needspace{4\baselineskip}%
   \begin{tcolorbox}[
     enhanced, breakable=false,
     colback=kfPaper, colframe=kfMute,
     boxrule=0.5pt, arc=2pt,
     left=\kfBoxPad, right=\kfBoxPad, top=\kfBoxPad, bottom=\kfBoxPad,
     title={\small\bfseries\color{kfInk} \,보기\,},
     coltitle=kfInk, colbacktitle=kfPrimaryLight!50,
     attach boxed title to top left={yshift=-2mm, xshift=6pt},
     boxed title style={colframe=kfMute, boxrule=0.4pt, arc=1pt}
   ]%
   \leavevmode\mbox{}\ignorespaces%
  }%
  {\end{tcolorbox}\par\smallskip}

% ============================================================
% choice-list.tex
% 5선지 (① ② ③ ④ ⑤) 자동 들여쓰기 컴포넌트
% 사용:
%   \begin{kfChoices}
%     \kfChoice 첫째 선택지
%     \kfChoice 둘째 선택지
%     \kfChoice 셋째 선택지
%     \kfChoice 넷째 선택지
%     \kfChoice 다섯째 선택지
%   \end{kfChoices}
% ============================================================

\newcounter{kfChoiceCnt}
\newcommand{\kfChoiceMark}[1]{%
  \ifcase#1\relax%
  \or ①\or ②\or ③\or ④\or ⑤\or ⑥\or ⑦\or ⑧\or ⑨%
  \fi
}

% 환경: 자동 번호
\newenvironment{kfChoices}%
  {\setcounter{kfChoiceCnt}{0}%
   \par\smallskip%
   \begin{list}{}{%
     \setlength{\leftmargin}{20pt}%
     \setlength{\itemindent}{0pt}%
     \setlength{\labelwidth}{18pt}%
     \setlength{\labelsep}{6pt}%
     \setlength{\itemsep}{2pt}%
     \setlength{\topsep}{1pt}%
     \setlength{\parsep}{0pt}%
     \renewcommand{\makelabel}[1]{##1\hss}%
   }}%
  {\end{list}\par\smallskip}

\newcommand{\kfChoice}{%
  \stepcounter{kfChoiceCnt}%
  \item[\textbf{\kfChoiceMark{\value{kfChoiceCnt}}}]\ignorespaces%
}

% --- 한 줄 5선지 (짧은 선택지용) ---
\newcommand{\kfChoicesInline}[5]{%
  \par\smallskip%
  \noindent\textbf{①}~#1\quad\textbf{②}~#2\quad\textbf{③}~#3\quad\textbf{④}~#4\quad\textbf{⑤}~#5%
  \par\smallskip%
}

% ============================================================
% answer-note.tex
% 답안 작성란 (노트 스타일, 줄친 영역)
% 사용:
%   \kfAnswerLines{5}        % 5줄 답안란
%   \kfAnswerNote{서술형 답안란}{8} % 라벨+8줄
% ============================================================

% 단일 줄 (필요 시 인라인)
\newcommand{\kfAnswerLine}{%
  \par\noindent\textcolor{kfRule}{\rule{\linewidth}{0.4pt}}\par
}

% N줄 답안란 (간단)
\newcommand{\kfAnswerLines}[1]{%
  \par\smallskip%
  \begin{minipage}{\linewidth}%
  \foreach \i in {1,...,#1}{%
    \textcolor{kfRule}{\rule{\linewidth}{0.4pt}}\\[6pt]%
  }%
  \end{minipage}%
  \par\smallskip%
}

% 라벨이 있는 노트형 답안란
\newcommand{\kfAnswerNote}[2]{%
  \par\smallskip\needspace{#2\baselineskip}%
  \begin{tcolorbox}[
    enhanced, breakable=false,
    colback=kfPaper, colframe=kfRule,
    boxrule=0.4pt, arc=2pt,
    left=\kfBoxPad, right=\kfBoxPad, top=\kfBoxPad, bottom=\kfBoxPad,
    title={\footnotesize\bfseries\color{kfMute} #1},
    coltitle=kfMute, colbacktitle=kfPaper,
    attach boxed title to top left={yshift=-1.5mm, xshift=6pt},
    boxed title style={colframe=kfRule, boxrule=0.3pt, arc=1pt}
  ]%
  \foreach \i in {1,...,#2}{%
    \textcolor{kfRule}{\rule{\linewidth}{0.3pt}}\\[6pt]%
  }%
  \end{tcolorbox}\par\smallskip%
}

% 짧은 빈칸 (단답형)
\newcommand{\kfBlank}[1][3cm]{\underline{\hspace{#1}}}
% 넓은 빈칸 (서술 한 줄)
\newcommand{\kfBlankWide}[1][6cm]{\underline{\hspace{#1}}}
% 괄호 빈칸
\newcommand{\kfBlankParen}{(\,\hspace{1cm}\,)}
% O/X 빈칸
\newcommand{\kfBlankOX}{(\,\hspace{1.5cm}\,)}

% ============================================================
% question-block.tex
% 문제 블록 (번호 배지 + 발문 + 보기/조건 + 선택지 + 답안란)
% 모든 구성을 하나의 minipage로 묶어 단/페이지 단절 금지
%
% 사용:
%   \begin{kfQBlock}{1}{객관식}
%     \kfQStem{다음 글에 대한 설명으로 가장 적절한 것은?}
%     \begin{kfEvidence} ... \end{kfEvidence}    % 선택
%     \begin{kfCondition} ... \end{kfCondition}  % 선택
%     \begin{kfChoices}
%       \kfChoice ...
%     \end{kfChoices}
%     \kfAnswerLines{2}                            % 선택
%   \end{kfQBlock}
% ============================================================

% 무단절 minipage로 묶고, 발문/보기/조건/선택지/답안란을 자유롭게 배치
% 사용자 피드백 #4: [객관식]/[서술형] 등 텍스트 라벨 제거 — 번호 배지만 표시
% #2(유형 라벨)는 호환성 인자로만 받고 출력하지 않음.
\newenvironment{kfQBlock}[2]%
  {\par\smallskip\needspace{6\baselineskip}%
   \noindent\begin{minipage}{\linewidth}%
   \samepage%
   \par\noindent
   \kfQNum{#1}\hspace{4pt}%
   \begin{minipage}[t]{\dimexpr\linewidth-30pt\relax}%
  }%
  {\end{minipage}\end{minipage}\par\smallskip}

% 문제 발문
\newcommand{\kfQStem}[1]{%
  \par\noindent#1\par\vspace{2pt}%
}

% 짧은 소문항 라벨 (1), (2), (3)
\newcounter{kfSubQ}
\newcommand{\kfSubQReset}{\setcounter{kfSubQ}{0}}
\newcommand{\kfSubQ}{%
  \stepcounter{kfSubQ}%
  \par\noindent\textbf{(\arabic{kfSubQ})}~\ignorespaces%
}

% ============================================================
% activity-table.tex
% 활동: 정리표 / 내용 확인표 (마크다운 표 → tabularx 변환)
% 사용:
%   % 일반 tabular (직접 컬럼 명시)
%   \begin{kfActTable}{|l|X|X|}
%     \kfTHead{항목} & \kfTHead{설명} & \kfTHead{학생 작성} \\\hline
%     ① & ... & \kfTBlank \\\hline
%   \end{kfActTable}
%
%   % adjustbox 자동 축소 (정해진 폭으로 강제)
%   \kfActTableWrap{
%     ... (\begin{tabular}...\end{tabular} 코드) ...
%   }
% ============================================================

% 사용 시 본문 마지막 행은 \\\hline 으로 끝나야 함.
% v17.1: 흰색 mask로 Canva 배경 본문 침범 차단
\newenvironment{kfActTable}[1]%
  {\par\smallskip\needspace{4\baselineskip}%
   \renewcommand{\arraystretch}{1.3}%
   \arrayrulecolor{kfRule}%
   \begin{tcolorbox}[blanker, colback=white, left=2pt, right=2pt, top=2pt, bottom=2pt]%
   \noindent\begin{adjustbox}{max width=\linewidth}%
   \begin{tabular}{#1}%
   \hline}%
  {\end{tabular}%
   \end{adjustbox}%
   \end{tcolorbox}%
   \arrayrulecolor{black}%
   \par\smallskip}

% 헤더 행 헬퍼
\newcommand{\kfTHead}[1]{\textbf{\color{kfPrimary} #1}}
\newcommand{\kfTHeadRow}[1]{\rowcolor{kfPrimaryLight!50}\textbf{\color{kfPrimary} #1}}

% 빈 셀 (학생 작성용)
\newcommand{\kfTBlank}{\rule[-1.6em]{0pt}{3.4em}}
\newcommand{\kfTBlankSm}{\rule[-1.0em]{0pt}{2.4em}}

% adjustbox 강제 축소 래퍼 — Python에서 변환된 raw tabular 블록을 감쌀 때 사용
\newcommand{\kfActTableWrap}[1]{%
  \par\smallskip\needspace{4\baselineskip}%
  \begin{tcolorbox}[blanker, colback=white, left=2pt, right=2pt, top=2pt, bottom=2pt]%
  \noindent\begin{adjustbox}{max width=\linewidth, center}%
  #1%
  \end{adjustbox}%
  \end{tcolorbox}\par\smallskip%
}

% ============================================================
% activity-vocab.tex
% 어휘 학습 표 (3열: 번호 - 뜻 - 빈칸)
% 사용:
%   \begin{kfVocab}
%     \kfVocabItem{1}{눈으로 보고 마음으로 깨달음}
%     \kfVocabItem{2}{사물의 이치를 따져 헤아림}
%   \end{kfVocab}
% ============================================================

% adjustbox로 폭 강제 + 일반 tabular + p{} 열 사용
\newenvironment{kfVocab}%
  {\par\smallskip\needspace{4\baselineskip}%
   \renewcommand{\arraystretch}{1.35}%
   \arrayrulecolor{kfRule}%
   \noindent\begin{adjustbox}{max width=\linewidth, center}%
   \begin{tabular}{|>{\centering\arraybackslash}m{1.0cm}|m{8.5cm}|>{\centering\arraybackslash}m{4.0cm}|}%
   \hline
   \rowcolor{kfPrimaryLight!50}%
   \multicolumn{1}{|c|}{\textbf{\color{kfPrimary}번호}} &
   \multicolumn{1}{c|}{\textbf{\color{kfPrimary}뜻풀이}} &
   \multicolumn{1}{c|}{\textbf{\color{kfPrimary}어휘 (정답 작성)}} \\\hline
   \ignorespaces}%
  {\end{tabular}%
   \end{adjustbox}%
   \arrayrulecolor{black}%
   \par\smallskip}

\newcommand{\kfVocabItem}[2]{%
  \textbf{#1} & #2 & \rule[-1.0em]{0pt}{2.4em}\\\hline%
}

% ============================================================
% activity-blank.tex
% 빈칸 채우기 활동
% 사용:
%   \begin{kfBlankList}
%     \kfBlankItem 첫 번째 문장에서 (\kfBlank)은(는) ...
%     \kfBlankItem 둘째 문장에서 ...
%   \end{kfBlankList}
% ============================================================

\newenvironment{kfBlankList}%
  {\par\smallskip%
   \begin{enumerate}[leftmargin=20pt, itemsep=4pt, topsep=2pt,
                    label=\textbf{\arabic*.}, labelsep=6pt]}%
  {\end{enumerate}\par\smallskip}

\newcommand{\kfBlankItem}{\item}

% 텍스트 안에 박스형 빈칸 (단어 1개 채우기)
\newcommand{\kfBlankBox}[1][2.4cm]{%
  \tikz[baseline=-0.5ex]{%
    \draw[kfRule, line width=0.4pt] (0,0) -- ++(#1,0);%
  }%
}

% ============================================================
% activity-ox.tex (v15)
% O/X 표기 활동 — 그래픽 디자인
% 사용:
%   \begin{kfOXList}
%     \kfOXItem 글쓴이는 자연을 예찬하고 있다.\kfOX
%     \kfOXItem 1연과 4연은 수미상관 구조이다.\kfOX
%   \end{kfOXList}
%
% \kfOX = 우측에 [O] [X] 두 박스 (학생이 하나 동그라미)
% ============================================================

\newenvironment{kfOXList}%
  {\par\smallskip%
   \begin{enumerate}[leftmargin=20pt, itemsep=5pt, topsep=2pt,
                    label=\textbf{\arabic*.}, labelsep=6pt]}%
  {\end{enumerate}\par\smallskip}

\newcommand{\kfOXItem}{\item}

% v16: O/X 그래픽 박스 키움 (18×14 → 24×20pt) + 영역색 강조
\newcommand{\kfOX}{%
  \hfill\nolinebreak
  \tikz[baseline=-0.9ex]{%
    \node[draw=kfPrimary, fill=kfPrimary!5, line width=0.6pt, rounded corners=3pt,
          minimum width=24pt, minimum height=20pt, inner sep=0pt,
          font=\normalsize\bfseries\color{kfPrimary}] (ob) {O};%
    \node[draw=kfMute, fill=kfMute!5, line width=0.6pt, rounded corners=3pt,
          minimum width=24pt, minimum height=20pt, inner sep=0pt,
          font=\normalsize\bfseries\color{kfMute},
          right=4pt of ob] (xb) {X};%
  }%
}

% ============================================================
% activity-connect.tex
% 좌우 선 긋기 활동 (2단 양식 + 중앙 여백)
% 사용:
%   \begin{kfConnect}
%     \kfPair{사르트르}{실존은 본질에 앞선다}
%     \kfPair{니체}{초인 사상}
%     \kfPair{하이데거}{현존재}
%   \end{kfConnect}
% ============================================================

\newenvironment{kfConnect}%
  {\par\smallskip%
   \renewcommand{\arraystretch}{1.6}%
   \begin{tabular}{@{}>{\raggedleft\arraybackslash}m{4.5cm}
                     @{\hspace{2.5cm}}
                     >{\raggedright\arraybackslash}m{6.5cm}@{}}}%
  {\end{tabular}\par\smallskip}

\newcommand{\kfPair}[2]{%
  \tikz[baseline]{\node[draw=kfRule, fill=kfPrimaryLight!30, rounded corners=2pt,
        inner sep=4pt, font=\small]{#1};} \tikz[baseline]{\node[circle, draw=kfMute, fill=kfPaper, inner sep=1.5pt]{};}
  &
  \tikz[baseline]{\node[circle, draw=kfMute, fill=kfPaper, inner sep=1.5pt]{};} \tikz[baseline]{\node[draw=kfRule, fill=kfNoteBg, rounded corners=2pt,
        inner sep=4pt, font=\small]{#2};} \\
}

% ============================================================
% activity-writing.tex
% 글쓰기 활동 (큰 노트 영역)
% 사용:
%   \kfWriting{독후감}{12}     % 라벨 + 12줄 큰 노트
%   \kfWritingPrompt{프롬프트 텍스트}
% ============================================================

\newcommand{\kfWritingPrompt}[1]{%
  \par\smallskip%
  \begin{tcolorbox}[
    enhanced, breakable=false,
    colback=kfPrimaryLight!30, colframe=kfPrimary,
    boxrule=0.4pt, arc=3pt,
    left=\kfBoxPad, right=\kfBoxPad, top=\kfBoxPad, bottom=\kfBoxPad,
  ]%
  \small\textbf{\color{kfPrimary}글쓰기 안내}\\[2pt]%
  #1
  \end{tcolorbox}\par\smallskip%
}

\newcommand{\kfWriting}[2]{%
  \par\smallskip\needspace{\numexpr#2+2\relax\baselineskip}%
  \begin{tcolorbox}[
    enhanced, breakable=false,
    colback=kfPaper, colframe=kfRule,
    boxrule=0.4pt, arc=2pt,
    left=\kfBoxPad, right=\kfBoxPad, top=\kfBoxPad, bottom=\kfBoxPad,
    title={\footnotesize\bfseries\color{kfMute} #1},
    coltitle=kfMute, colbacktitle=kfPaper,
    attach boxed title to top left={yshift=-1.5mm, xshift=6pt},
    boxed title style={colframe=kfRule, boxrule=0.3pt, arc=1pt}
  ]%
  \foreach \i in {1,...,#2}{%
    \textcolor{kfRule}{\rule{\linewidth}{0.3pt}}\\[7pt]%
  }%
  \end{tcolorbox}\par\smallskip%
}

% ============================================================
% activity-structure.tex
% 구조도 (트리 + 빈칸)
% 사용:
%   \begin{kfStructure}
%     \kfNode{0}{서론}{문제 제기}
%     \kfNode{1}{본론}{(\,\quad\,)}
%     \kfNode{1}{본론}{근거 2}
%     \kfNode{0}{결론}{주장 강화}
%   \end{kfStructure}
% ============================================================

\newenvironment{kfStructure}%
  {\par\smallskip\needspace{6\baselineskip}%
   \begin{tcolorbox}[
     enhanced, breakable=false,
     colback=kfPaper, colframe=kfRule,
     boxrule=0.4pt, arc=2pt,
     left=\kfBoxPad, right=\kfBoxPad, top=\kfBoxPad, bottom=\kfBoxPad,
     title={\footnotesize\bfseries\color{kfMute} 글의 구조},
     coltitle=kfMute, colbacktitle=kfPrimaryLight!40,
     attach boxed title to top left={yshift=-1.5mm, xshift=6pt},
     boxed title style={colframe=kfRule, boxrule=0.3pt, arc=1pt}
   ]%
   \begin{tabbing}
     \hspace{1.4cm}\=\hspace{1.4cm}\=\hspace{8cm}\kill}%
  {\end{tabbing}\end{tcolorbox}\par\smallskip}

% 노드: #1=들여쓰기 단계 (0~3), #2=라벨, #3=내용
\newcommand{\kfNode}[3]{%
  \ifcase#1\relax%
    \>\>\textbf{\color{kfPrimary} ▶ #2}\quad #3\\
  \or
    \>\>\quad\textbf{\color{kfMute} ◦ #2}\quad #3\\
  \or
    \>\>\quad\quad\textbf{\color{kfMute} - #2}\quad #3\\
  \or
    \>\>\quad\quad\quad\textbf{\color{kfMute} \textperiodcentered\ #2}\quad #3\\
  \fi
}

% 빈칸 노드
\newcommand{\kfNodeBlank}[2]{\kfNode{#1}{#2}{(\hspace{2.5cm})}}

% ============================================================
% activity-sentence-reading.tex
% 문장 독해 활동 — 문장 박스 1개 + 그에 묶인 문제 N개를 한 minipage로
% 사용:
%   \begin{kfSentenceUnit}{1}{"바람은 내 귀에 속삭이며..."}
%     \kfSentQ{(1)}{서술어 '속삭이며'의 주어는?}
%     \begin{kfChoices}
%       \kfChoice 바람은
%       ...
%     \end{kfChoices}
%   \end{kfSentenceUnit}
% ============================================================

% kfSentenceBox 매크로는 passage-note.tex에 정의됨

\newenvironment{kfSentenceUnit}[2]%
  {\par\smallskip\needspace{6\baselineskip}%
   \noindent\begin{minipage}{\linewidth}%
   \samepage%
   \kfSentenceBox{문장 #1}{#2}%
   \par\smallskip%
  }%
  {\end{minipage}\par\smallskip}

% 같은 문장에 묶인 소문항 (문장 안내 + 발문)
\newcommand{\kfSentQ}[2]{%
  \par\noindent\textbf{#1}~#2\par\smallskip%
}

% 헤더 없이 문장만 있는 묶음 (sentence_ref가 없는 일반 문항)
\newenvironment{kfSentencelessUnit}%
  {\par\smallskip\noindent\begin{minipage}{\linewidth}\samepage}%
  {\end{minipage}\par\smallskip}

% ============================================================
% chapter-cover.tex (v6)
% 챕터 진입 페이지 (표지) — 하단에 영역 목차 추가
% 사용:
%   \kfChapterCover{1}{근대성과 자아}{Chapter 1}{이번 챕터에서는 ...}
%   \begin{kfChapterAreaList}
%     \kfChapterAreaItem{어휘}{kfAreaVocab}{핵심 어휘 12개}
%     ...
%   \end{kfChapterAreaList}
%
%   영역 목차가 없을 때는 \kfChapterCoverEnd 호출
% ============================================================

\newcommand{\kfChapterCover}[4]{%
  \clearpage%
  \thispagestyle{kfBlank}%
  \kfMarkChapter{Chapter #1. #2}%
  \kfChapterCoverBg%
  % v18.5: 챕터 표지 우상단에 로고
  \begin{tikzpicture}[remember picture, overlay]
    \node[anchor=north east, inner sep=0pt] at ($(current page.north east)+(-18mm,-18mm)$) {\kfLogoMd};
  \end{tikzpicture}%
  \vspace*{20mm}%
  \noindent
  \begin{tikzpicture}[remember picture, overlay]
    % 좌측 액센트 바
    \fill[kfPrimary] (current page.west) ++(12mm,25mm) rectangle ++(5mm, -50mm);
    % 챕터 번호 배지 (이중 원, 약간 작게)
    \node[anchor=north west, inner sep=0pt] at ($(current page.north west)+(24mm,-45mm)$) {%
      \begin{tikzpicture}
        \draw[kfPrimary, line width=1pt] (0,0) circle (10mm);
        \draw[kfPrimary, line width=0.4pt] (0,0) circle (8mm);
        \node[font=\normalsize\bfseries\color{kfPrimary}] at (0,2mm) {CHAPTER};
        \node[font=\LARGE\bfseries\color{kfPrimary}] at (0,-3mm) {#1};
      \end{tikzpicture}%
    };
  \end{tikzpicture}%
  \vspace{42mm}%
  \par\noindent\hspace{32mm}%
  \begin{minipage}{0.65\linewidth}%
    {\footnotesize\color{kfMute} #3}\\[4pt]
    {\LARGE\bfseries\color{kfPrimary} #2}\\[8pt]
    {\color{kfRule}\rule{50mm}{0.5pt}}\\[6pt]
    {\small\color{kfInk} #4}
  \end{minipage}%
  \par\vspace{14mm}%
}

% v6 / v18.5 / v18.6: 챕터 표지의 영역 목차 — 흰색 반투명 박스
% v18.6: 배경 가로선/도형과 안 겹치도록 TikZ overlay로 우하단 절대 위치
\makeatletter
% 영역 항목들을 모아 둘 박스
\newsavebox{\kf@arealistbox}
\newenvironment{kfChapterAreaList}{%
  \begin{lrbox}{\kf@arealistbox}%
  \begin{minipage}{0.62\linewidth}%
    \begin{tcolorbox}[%
      enhanced, colback=white, colframe=white,
      opacityback=0.92, opacityframe=0,
      boxrule=0pt, arc=4pt,
      width=\linewidth,
      top=8pt, bottom=8pt, left=10pt, right=10pt,
    ]%
      {\footnotesize\bfseries\color{kfMute} 이 챕터의 영역}\par\vspace{4pt}%
      {\color{kfRule}\hrule height 0.3pt}\par\vspace{4pt}%
      \begin{itemize}[leftmargin=14pt, itemsep=2pt, topsep=0pt, label={\color{kfPrimary!70}\textbullet}]%
}{%
      \end{itemize}%
    \end{tcolorbox}%
  \end{minipage}%
  \end{lrbox}%
  % v18.7: 배경 상단 가로선 ~ 하단 가로선 사이의 중간 영역에 배치
  % 가로선이 내용을 가리지 않도록 박스 중심이 두 선 사이 중점에 위치
  \begin{tikzpicture}[remember picture, overlay]%
    \node[anchor=center, inner sep=0pt]
      at ($(current page.south)+(0mm,90mm)$)
      {\usebox{\kf@arealistbox}};%
  \end{tikzpicture}%
  \vfill%
  \noindent\hfill\kfSeriesBadge\hspace{12mm}%
  \clearpage%
}
\makeatother

% 영역 항목: \kfChapterAreaItem{영역명}{kfArea색}{부제}
\makeatletter
\newcommand{\kfChapterAreaItem}[3]{%
  \item {\small\bfseries\color{#2} #1}%
  \def\kf@cci@sub{#3}%
  \ifx\kf@cci@sub\@empty\else
    {\small\color{kfMute}\, \textemdash\, #3}%
  \fi
}
\makeatother

% 영역 목록 없이 cover만 (legacy)
\newcommand{\kfChapterCoverEnd}{%
  \vfill%
  \noindent\hfill\kfSeriesBadge\hspace{12mm}%
  \clearpage%
}

% ============================================================
% area-header.tex (v16)
% 영역 시작 표제 — 시리즈별 차별화 라벨 + 큰 영역명 + 부제
% PAGE_BREAK_POLICY.md §1.4
%
% v16 (디자인 고도화 Phase 1):
%   - 시리즈별 라벨 박스 추가:
%     소쉬르 = AREA START (영역색 채움 박스)
%     프레게 = SECTION OPEN (영역색 테두리)
%     러셀   = RUSSELL (영역색 라이트 박스)
%     비트   = WITTGENSTEIN (회색 라이트 박스)
%   - 라벨 위치: 영역명 위 좌측
% ============================================================

\makeatletter

% 시리즈 라벨 텍스트
\newcommand{\kf@serieslabel}{%
  \ifcase\kf@series\relax
    AREA START%
  \or
    AREA START%
  \or
    SECTION OPEN%
  \or
    RUSSELL%
  \or
    WITTGENSTEIN%
  \fi
}

% 시리즈 라벨 박스 스타일 (#1 = 영역색)
\newcommand{\kf@labelbox}[1]{%
  \ifcase\kf@series\relax
    % 0/1 = 소쉬르: 영역색 채움
    \tikz[baseline=-2pt]{\node[fill=#1, text=white,
      rounded corners=3pt, inner xsep=6pt, inner ysep=2pt,
      font=\scriptsize\bfseries]{\kf@serieslabel};}%
  \or
    % 1 = 소쉬르: 영역색 채움
    \tikz[baseline=-2pt]{\node[fill=#1, text=white,
      rounded corners=3pt, inner xsep=6pt, inner ysep=2pt,
      font=\scriptsize\bfseries]{\kf@serieslabel};}%
  \or
    % 2 = 프레게: 영역색 테두리
    \tikz[baseline=-2pt]{\node[draw=#1, line width=0.6pt, text=#1,
      rounded corners=3pt, inner xsep=6pt, inner ysep=2pt,
      font=\scriptsize\bfseries]{\kf@serieslabel};}%
  \or
    % 3 = 러셀: 영역색 라이트 박스
    \tikz[baseline=-2pt]{\node[fill=#1!12, text=#1,
      rounded corners=2pt, inner xsep=5pt, inner ysep=1.5pt,
      font=\scriptsize\bfseries]{\kf@serieslabel};}%
  \or
    % 4 = 비트: 회색 미니 박스
    \tikz[baseline=-2pt]{\node[fill=kfMute!10, text=kfMute,
      rounded corners=2pt, inner xsep=5pt, inner ysep=1.5pt,
      font=\scriptsize\bfseries]{\kf@serieslabel};}%
  \fi
}

% v17: 시리즈+영역별 Canva 배경 자동 매핑
% \kf@hasbg=1로 set되면 \kfAreaStart에서 라벨 박스 출력 생략
\def\kf@areaVocab{어휘}\def\kf@areaGrammar{문법}\def\kf@areaConcept{개념}
\def\kf@areaLit{문학}\def\kf@areaNonlit{비문학}
\newif\ifkf@bgon
% 파일 있으면 배경 적용 + boolean true. 없으면 nothing.
\newcommand{\kfBgSet}[1]{%
  \IfFileExists{#1.pdf}{\kfPageBg{#1}\kf@bgontrue}{}%
}
\newcommand{\kf@trybg}[1]{%
  % v18.4: 영역 배경 PDF 비활성화 — 큰 도형이 본문 영역을 침범해
  % "영역 헤더 후 빈 페이지" 문제를 일으킴. 라벨 박스로만 영역 표시.
  \kf@bgonfalse%
}

% Note: 러셀(rus_*)/비트(wit_*) 배경 PDF 미존재 시 \kfPageBg가 에러 → \IfFileExists로 가드
% (현재는 파일이 모두 존재한다고 가정. 부분 제공 시 cls 수정 필요)

\newcommand{\kfAreaStart}[3]{%
  \clearpage%
  \thispagestyle{fancy}%
  \kfMarkArea{#1}%
  \kf@trybg{#1}%
  \vspace*{1mm}%
  % 배경 없으면 라텍스 라벨 박스 출력 (배경 있으면 일러스트가 라벨 역할)
  \ifkf@bgon
    \vspace*{4mm}%
  \else
    \noindent\kf@labelbox{#2}\par\vspace{4pt}%
  \fi
  % 영역명 + 부제
  \noindent
  \begin{tikzpicture}
    \node[anchor=west, inner sep=0pt] (bar) at (0,0) {\color{#2}\rule{3pt}{22pt}};
    \node[anchor=base west, inner sep=0pt, xshift=8pt, yshift=2pt] (lbl) at (bar.east |- 0,0) {%
      {\LARGE\bfseries\color{#2} #1}%
    };
    \def\kf@areasub{#3}%
    \ifx\kf@areasub\@empty\else
      \node[anchor=base west, inner sep=0pt, xshift=8pt, font=\normalsize\color{kfMute}]
        at (lbl.base east) {\,\textperiodcentered\, #3};%
    \fi
  \end{tikzpicture}%
  \par\vspace{2pt}
  {\color{#2!70}\hrule height 1pt}%
  \par\vspace{1pt}%
  {\color{kfRule}\hrule height 0.3pt}%
  \par\vspace{4pt}%
  \needspace{6\baselineskip}\nopagebreak[4]%
  \global\kf@areaJustOpenedtrue% v18.5: 첫 컨텐츠 needspace 무시
}
\makeatother

% 시리즈/영역 단축 매크로
\newcommand{\kfStartVocab}[1]{\kfAreaStart{어휘}{kfAreaVocab}{#1}}
\newcommand{\kfStartGrammar}[1]{\kfAreaStart{문법}{kfAreaGrammar}{#1}}
\newcommand{\kfStartConcept}[1]{\kfAreaStart{개념}{kfAreaConcept}{#1}}
\newcommand{\kfStartLit}[1]{\kfAreaStart{문학}{kfAreaLiterature}{#1}}
\newcommand{\kfStartNonlit}[1]{\kfAreaStart{비문학}{kfAreaNonlit}{#1}}
\newcommand{\kfStartWeekly}[1]{\kfAreaStart{실력 확인}{kfAreaWeekly}{#1}}

% ============================================================
% section-header.tex
% 섹션 시작 라벨 헤더 — [지문] / [활동] / [문제] 라벨
% 좌측 액센트 바 + 라벨 + 부제 (영역 액센트 색)
%
% 사용:
%   \kfSecHeader{kfAreaLiterature}{지문}{빼앗긴 들에도 봄은 오는가 / 이상화}
%   \kfSecHeader{kfAreaConcept}{활동}{내용 확인}
%   \kfSecHeader{kfAreaNonlit}{문제}{}
%
% 헤더 직후 페이지 분할 금지: \needspace{4\baselineskip} + \nopagebreak
% ============================================================

\providecommand{\kfSecHeader}[3]{%
  \par\needspace{10\baselineskip}\filbreak%
  \vspace{3pt}%
  \noindent
  \begin{tikzpicture}
    \node[anchor=west, inner sep=0pt] (bar) at (0,0) {\color{#1}\rule{3pt}{14pt}};
    \node[anchor=west, inner sep=0pt, xshift=6pt, font=\normalsize\bfseries\color{#1}]
       (lbl) at (bar.east) {[\,#2\,]};
    \node[anchor=west, inner sep=0pt, xshift=4pt, font=\small\color{kfMute}]
       at (lbl.east) {#3};
  \end{tikzpicture}%
  \par\vspace{1pt}%
  {\color{#1!50}\hrule height 0.4pt}%
  \par\vspace{2pt}\nopagebreak%
}

% 영역별 단축 매크로
\newcommand{\kfSecPassageVocab}[1]{\kfSecHeader{kfAreaVocab}{지문}{#1}}
\newcommand{\kfSecPassageGrammar}[1]{\kfSecHeader{kfAreaGrammar}{지문}{#1}}
\newcommand{\kfSecPassageConcept}[1]{\kfSecHeader{kfAreaConcept}{지문}{#1}}
\newcommand{\kfSecPassageLit}[1]{\kfSecHeader{kfAreaLiterature}{지문}{#1}}
\newcommand{\kfSecPassageNonlit}[1]{\kfSecHeader{kfAreaNonlit}{지문}{#1}}
\newcommand{\kfSecPassageWeekly}[1]{\kfSecHeader{kfAreaWeekly}{지문}{#1}}

\newcommand{\kfSecActVocab}[1]{\kfSecHeader{kfAreaVocab}{활동}{#1}}
\newcommand{\kfSecActGrammar}[1]{\kfSecHeader{kfAreaGrammar}{활동}{#1}}
\newcommand{\kfSecActConcept}[1]{\kfSecHeader{kfAreaConcept}{활동}{#1}}
\newcommand{\kfSecActLit}[1]{\kfSecHeader{kfAreaLiterature}{활동}{#1}}
\newcommand{\kfSecActNonlit}[1]{\kfSecHeader{kfAreaNonlit}{활동}{#1}}
\newcommand{\kfSecActWeekly}[1]{\kfSecHeader{kfAreaWeekly}{활동}{#1}}

\newcommand{\kfSecQVocab}[1]{\kfSecHeader{kfAreaVocab}{문제}{#1}}
\newcommand{\kfSecQGrammar}[1]{\kfSecHeader{kfAreaGrammar}{문제}{#1}}
\newcommand{\kfSecQConcept}[1]{\kfSecHeader{kfAreaConcept}{문제}{#1}}
\newcommand{\kfSecQLit}[1]{\kfSecHeader{kfAreaLiterature}{문제}{#1}}
\newcommand{\kfSecQNonlit}[1]{\kfSecHeader{kfAreaNonlit}{문제}{#1}}
\newcommand{\kfSecQWeekly}[1]{\kfSecHeader{kfAreaWeekly}{문제}{#1}}


\begin{document}

\thispagestyle{kfBlank}
\kfBookCoverBg
\vspace*{\kfBookCoverVspace}
\begin{center}
{\kfLogoLg}\\[14pt]
{\fontsize{72}{84}\selectfont\bfseries\kfBookCoverColor 러셀}\\[18pt]
{\fontsize{28}{32}\selectfont\kfBookCoverSubColor \textsc{Level} \textbf{3}}\\[22pt]
{\large\kfBookCoverSubColor 제 1 권 \quad\textbar\quad 챕터 1\textendash{} 4}
\end{center}
\vfill
\hfill\kfSeriesBadge\hspace{15mm}
\clearpage
\kfChapterCover{1}{러셀3 (챕터1)}{Chapter 1}{이번 챕터의 학습 내용을 시작합니다.}
\begin{kfChapterAreaList}
\kfChapterAreaItem{문법}{kfAreaGrammar}{단어의 이름표, 품사 체계의 이해}
\kfChapterAreaItem{개념}{kfAreaConcept}{구수한 사투리의 맛, 방언의 효과}
\kfChapterAreaItem{문학}{kfAreaLiterature}{「금 따는 콩밭」 — 김유정}
\kfChapterAreaItem{비문학}{kfAreaNonlit}{[인문] 용서와 추서, 정약용의 실천 윤리}
\end{kfChapterAreaList}

\kfStartGrammar{단어의 이름표, 품사 체계의 이해}


\kfPassageNoteNT{%
　우리가 사용하는 수많은 단어는 그 성질에 따라 몇 개의 갈래로 나눌 수 있는데, 이렇게 단어의 성격에 따라 붙인 이름표를 품사(品詞)라고 한다. 품사는 단어를 분류하는 기준이며, 이 분류는 단어의 형태가 변하는지, 문장에서 어떤 기능을 하는지, 그리고 어떤 공통된 의미를 가지는지에 따라 체계적으로 이루어진다.

　가장 큰 분류 기준인 형태에 따라, 단어는 문장 속에서 모습이 변하는 가변어와 형태가 고정된 불변어로 나뉜다. '신다'가 '신고, 신으니'처럼 변하는 것은 가변어이고, '하늘'이나 '매우'처럼 항상 같은 모습인 것은 불변어이다.

　이 기준을 바탕으로 단어의 기능과 의미에 따라 국어의 품사를 아홉 가지로 나눌 수 있다.

　첫째, 문장에서 주체적인 역할을 하여 뼈대를 이루는 체언이 있다. 체언은 문장에서 주로 주어나 목적어로 쓰이며, 여기에는 '학교', '우정'처럼 대상의 이름을 나타내는 명사, '나', '우리', '이것', '저기'처럼 명사를 대신하는 대명사, 그리고 '하나', '첫째'처럼 수량이나 순서를 나타내는 수사가 속한다. 이들은 모두 형태가 변하지 않는 불변어이다.

　둘째, 주로 체언 뒤에 붙어 문법적 관계를 표시하거나 특별한 뜻을 더하는 관계언이 있다. 여기에는 조사가 유일하게 속한다. '수호가', '운동화를', '학교에서', '너만'의 '가, 를, 에서, 만' 등이 모두 조사이다. 조사는 불변어이지만, '이것은 행복이다'의 '이다' 하나만은 예외적으로 형태가 변하는 서술격 조사이다.

　셋째, 문장의 서술어 역할을 하여 주체의 움직임이나 상태를 풀이하는 용언이 있다. 용언은 '형태가 변하는 말'이라는 뜻처럼, 어간에 다양한 어미가 붙어 '신다'가 '신고, 신었다' 등으로 변하는 활용을 한다는 것이 가장 큰 특징이다. 용언에는 '가다', '살다'처럼 움직임을 나타내는 동사와, '멋지다', '부드럽다'처럼 성질이나 상태를 나타내는 형용사가 있다. 동사와 형용사는 현재형 '-(ㄴ)는다'나 명령형 '-아라', 청유형 '-자'와 결합할 수 있는지 여부로 쉽게 구분할 수 있다.

　넷째, 문장의 다른 말을 꾸며주어 의미를 풍부하게 하는 수식언이 있다. 수식언은 용언과 달리 활용하지 않는 불변어이다. 여기에는 '새' 학년, '모든' 친구처럼 오직 체언만을 꾸미는 관형사와, '방긋' 웃는, '참' 좋다, '벌써' 가니처럼 주로 용언을 꾸미지만 다른 부사나 문장 전체를 꾸미기도 하는 부사가 있다.

　마지막으로, 문장의 다른 성분과 직접적인 관계없이 독립적으로 쓰이는 독립언이 있다. 여기에는 '어머나', '네', '야'처럼 놀람, 대답, 부름을 나타내는 감탄사가 속한다. 이 9가지 품사 체계를 이해하는 것은 단어의 성격을 파악하고 문장을 정확하게 분석하는 첫걸음이다.
\par\medskip\begin{itemize}[leftmargin=18pt, itemsep=2pt, topsep=2pt, label=\textbullet]
\item 가변/불변 — 동·형·서술격 조사(이다)=가변 / 나머지=불변
\item 5대 기능 — 체언·용언·수식언·관계언·독립언
\item 서술격 조사 예외 — '이다'는 조사지만 활용(이고·이니·이다)
\end{itemize}
}

\kfActivityBg \par\kfMaybeNeedspace{32\baselineskip}지문의 핵심 내용을 떠올리며 아래 빈칸에 알맞은 말을 채워 보세요.

\#\#\# 1. 품사 분류 기준 

\par\medskip\needspace{4\baselineskip}
\noindent\begin{adjustbox}{max width=\linewidth, center}
\renewcommand{\arraystretch}{1.45}
\arrayrulecolor{kfRule}
\begin{tabular}{|>{\raggedright\arraybackslash}p{\dimexpr 0.1398\linewidth-12pt\relax}|>{\raggedright\arraybackslash}p{\dimexpr 0.8602\linewidth-12pt\relax}|}
\hline
\rowcolor{kfPrimaryLight!50}\textbf{\color{kfPrimary} 기준} & \textbf{\color{kfPrimary} 분류} \\\hline
1. 형태 (형태 변화 여부) & ∙ [ ① ]: 형태 고정 (예: 하늘, 매우)<br>∙ [ ② ]: 형태 변화 (예: 신다 → 신고, 신으니) \\\hline
2. 기능 (문장 속 역할) & ∙ [ ③ ]: 문장의 뼈대 (주어, 목적어 등)<br>∙ [ ④ ]: 문법적 관계 표시<br>∙ [ ⑤ ]: 문장의 서술어 기능 (활용 O)<br>∙ [ ⑥ ]: 다른 말을 꾸며 줌 (활용 X)<br>∙ [ ⑦ ]: 독립적으로 쓰임 \\\hline
3. 의미 (단어의 공통된 뜻) & ∙ 9품사 (명사, 대명사, 수사, 조사, 동사, 형용사, 관형사, 부사, 감탄사) \\\hline
\end{tabular}
\arrayrulecolor{black}
\end{adjustbox}\par\medskip



\#\#\# 2. 9품사 분류 (기능과 의미 중심) 

\par\medskip\needspace{4\baselineskip}
\noindent\begin{adjustbox}{max width=\linewidth, center}
\renewcommand{\arraystretch}{1.45}
\arrayrulecolor{kfRule}
\begin{tabular}{|>{\raggedright\arraybackslash}p{\dimexpr 0.0646\linewidth-12pt\relax}|>{\raggedright\arraybackslash}p{\dimexpr 0.0682\linewidth-12pt\relax}|>{\raggedright\arraybackslash}p{\dimexpr 0.8673\linewidth-12pt\relax}|}
\hline
\rowcolor{kfPrimaryLight!50}\textbf{\color{kfPrimary} 기능 (5)} & \textbf{\color{kfPrimary} 품사 (9)} & \textbf{\color{kfPrimary} 의미 및 특징} \\\hline
\multirow{3}{*}{[ ③ ]} & 명사 & 대상의 이름을 나타냄 (예: 학교, 우정) \\\cline{2-3}
 & [ ⑧ ] & 명사를 대신하여 가리킴 (예: 나, 우리, 이것) \\\cline{2-3}
 & [ ⑨ ] & 수량이나 순서를 나타냄 (예: 하나, 첫째) \\\cline{2-3}
[ ④ ] & [ ⑩ ] & 주로 체언 뒤에 붙어 관계를 표시함 (예: 가, 를, 에서)<br>\textit{(단, [ ⑪ ] '이다'는 유일하게 가변어에 속함)} \\\hline
\multirow{2}{*}{[ ⑤ ]} & [ ⑫ ] & 대상의 움직임을 나타냄 (예: 가다, 살다) \\\cline{2-3}
 & [ ⑬ ] & 대상의 성질/상태를 나타냄 (예: 멋지다, 부드럽다) \\\cline{2-3}
\multirow{2}{*}{[ ⑥ ]} & [ ⑭ ] & 오직 체언 (명사, 대명사, 수사)만 꾸밈 (예: 새, 모든) \\\cline{2-3}
 & [ ⑮ ] & 주로 용언 (동사, 형용사)을 꾸밈 (예: 방긋, 참) \\\cline{2-3}
[ ⑦ ] & [ ⑯ ] & 놀람, 대답, 부름을 나타냄 (예: 어머나, 네) \\\hline
\end{tabular}
\arrayrulecolor{black}
\end{adjustbox}\par\medskip



\#\#\# 3. 심화 확인 (형태 기준)

* 가변어에 속하는 3가지 품사: [ ⑰ ], [ ⑱ ], 그리고 [ ⑲ ] '이다'   * 동사와 형용사의 구분 방법: 현재형 어미 [ ⑳ ], 명령형 '-아라', 청유형 '-자'와 결합할 수 있으면 동사, 없으면 형용사이다.

---\nopagebreak[4]\par\nopagebreak[4]

\kfActivityBg \kfSecHeader{kfAreaGrammar}{활동}{분석 훈련}
\par\kfMaybeNeedspace{24\baselineskip}\noindent{\kfTipMark\,\,}{\small\color{kfInk} 다음 활용 사례를 분석하여 빈칸을 채워 보세요.}\par\nopagebreak[4]\smallskip\nopagebreak[4]
\begin{enumerate}[leftmargin=22pt, itemsep=6pt, topsep=4pt, label=\textbf{\arabic*.}]
\item 다음 문장에서 형태가 변하는 '가변어'를 모두 찾아 쓰시오.
\begin{kfEvidence}저 아이는 학생이고, 매우 착하다.\end{kfEvidence}
\item 다음 문장에서 형태가 변하지 않는 '불변어'를 모두 찾아 쓰시오.
\begin{kfEvidence}저 아이는 학생이고, 매우 착하다.\end{kfEvidence}
\item 다음 문장에서 문장의 뼈대가 되는 '체언'을 모두 찾아 쓰시오.
\begin{kfEvidence}우리 셋이 저것을 다 먹었다.\end{kfEvidence}
\item 다음 문장에서 문법적 관계를 나타내는 '관계언(조사)'을 모두 찾아 쓰시오.
\begin{kfEvidence}너만이 나에게 희망이다.\end{kfEvidence}
\item 다음 문장에서 주체의 움직임이나 상태를 풀이하는 '용언'을 찾아 기본형을 쓰고, 각각 '동사', '형용사'로 구분하시오.
\begin{kfEvidence}하늘이 푸르고, 새가 힘차게 난다.\end{kfEvidence}
\item 다음 문장에서 다른 말을 꾸며주는 '수식언'을 모두 찾아 '관형사', '부사'로 구분하시오.
\begin{kfEvidence}그 헌 책을 아주 빨리 읽었다.\end{kfEvidence}
\item 다음 문장에서 '-는-' 또는 '-ㄴ다'를 붙여 보고, 이를 근거로 밑줄 친 단어의 품사가 '동사'인지 '형용사'인지 밝히시오.
\begin{kfEvidence}길이 (가)\uline{밝다}. \\\relax
날이 (나)\uline{밝는다}.\end{kfEvidence}
\item 다음 문장의 밑줄 친 단어들의 품사를 각각 쓰시오.
\begin{kfEvidence}(가) \uline{저} 하늘을 보아라. \\\relax
(나) \uline{저기}에 우리 집이 있다.\end{kfEvidence}
\item 다음 문장의 밑줄 친 단어들의 품사를 각각 쓰시오.
\begin{kfEvidence}(가) 학생 \uline{다섯}이 교실에 있다 \\\relax
(나) \uline{다섯} 학생이 교실에 있다.\end{kfEvidence}
\item 다음 문장에 쓰인 모든 단어를 9품사로 분류하시오.
\begin{kfEvidence}와, 모든 학생이 정말 열심히 공부한다!\end{kfEvidence}
\end{enumerate}

\kfSecHeader{kfAreaGrammar}{문제}{}

\par\medskip
\kfBeginMC
\begin{kfQBlock}{01}{객관식}
\kfQStem{품사 분류 기준에 대한 설명으로 적절하지 \uline{\underline{않은}} 것은?}
\begin{kfChoices}
\kfChoice '형태'는 단어가 문장에서 변하는지 여부이다.
\kfChoice '기능'은 단어가 문장 속에서 하는 역할이다.
\kfChoice '의미'는 단어가 가진 공통된 뜻이다.
\kfChoice '문장 성분'은 품사 분류의 가장 중요한 기준이다.
\kfChoice 품사는 이 세 가지 기준을 종합적으로 고려하여 분류한다.
\end{kfChoices}
\end{kfQBlock}
\begin{kfQBlock}{02}{객관식}
\kfQStem{다음 중 형태에 따른 분류가 나머지와 \uline{다른} 하나는?}
\begin{kfChoices}
\kfChoice 하늘
\kfChoice 매우
\kfChoice 에게
\kfChoice 모든
\kfChoice 멋지다
\end{kfChoices}
\end{kfQBlock}
\begin{kfQBlock}{03}{객관식}
\kfQStem{다음 밑줄 친 단어 중 유일하게 가변어에 속하는 것은?}
\begin{kfChoices}
\kfChoice \uline{새} 옷을 입었다.
\kfChoice \uline{저} 사람은 학생이다.
\kfChoice \uline{어머나}, 눈이 온다.
\kfChoice \uline{다섯}이 모였다.
\kfChoice 저 사람은 학생\uline{이다}.
\end{kfChoices}
\end{kfQBlock}
\begin{kfQBlock}{04}{객관식}
\kfQStem{다음 단어들을 기능에 따라 묶을 때, 성격이 \uline{다른} 하나는?}
\begin{kfChoices}
\kfChoice 학교
\kfChoice 우리
\kfChoice 첫째
\kfChoice 매우
\kfChoice 사랑
\end{kfChoices}
\end{kfQBlock}
\begin{kfQBlock}{05}{객관식}
\kfQStem{‘용언’에 대한 설명으로 적절하지 \uline{\underline{않은}} 것은?}
\begin{kfChoices}
\kfChoice 동사와 형용사를 함께 묶어 부르는 품사 분류이다.
\kfChoice 문장에서 주로 서술어 자리에 쓰여 의미를 마무리한다.
\kfChoice 어간에 어미가 결합하여 다양한 형태로 활용된다.
\kfChoice 주어의 동작이나 상태를 풀이해 주는 기능을 한다.
\kfChoice 형태가 변하지 않고 체언 뒤에 붙어 문법 관계를 나타낸다.
\end{kfChoices}
\end{kfQBlock}
\begin{kfQBlock}{06}{객관식}
\kfQStem{동사와 형용사를 구분하는 방법에 대한 설명으로 적절하지 \uline{\underline{않은}} 것은?}
\begin{kfChoices}
\kfChoice 현재형 ‘-는-/-ㄴ다’ 어미와 결합이 가능한지 살펴본다.
\kfChoice 명령형·청유형 어미와 자연스럽게 결합 가능한지 살펴본다.
\kfChoice 동사는 ‘먹는다·간다’처럼 현재형 결합이 가능한 부류이다.
\kfChoice 형용사는 ‘예쁜다·착한다’처럼 결합이 자연스럽지 않다.
\kfChoice 단어에 받침이 있는지 없는지로 동사와 형용사를 구분한다.
\end{kfChoices}
\end{kfQBlock}
\begin{kfQBlock}{07}{객관식}
\kfQStem{다음 중 '-는-' 또는 '-ㄴ다'와 결합이 불가능하여 형용사로 분류되는 것은?}
\begin{kfChoices}
\kfChoice 읽다
\kfChoice 자다
\kfChoice 슬프다
\kfChoice 뛰다
\kfChoice 걷다
\end{kfChoices}
\end{kfQBlock}
\begin{kfQBlock}{08}{객관식}
\kfQStem{<보기>의 ㉠, ㉡에 들어갈 품사로 바르게 짝지어진 것은?}
\begin{kfEvidence}( \\\relax
㉠ )은/는 체언을 꾸며주며, ( \\\relax
㉡ )은/는 주로 용언을 꾸며준다.\end{kfEvidence}
\begin{kfChoices}
\kfChoice ㉠ 관형사 / ㉡ 부사
\kfChoice ㉠ 부사 / ㉡ 관형사
\kfChoice ㉠ 명사 / ㉡ 동사
\kfChoice ㉠ 동사 / ㉡ 형용사
\kfChoice ㉠ 조사 / ㉡ 감탄사
\end{kfChoices}
\end{kfQBlock}
\begin{kfQBlock}{09}{객관식}
\kfQStem{다음 밑줄 친 단어 중 품사가 관형사인 것은?}
\begin{kfChoices}
\kfChoice \uline{하늘}이 맑다.
\kfChoice 밥을 \uline{빨리} 먹어라.
\kfChoice \uline{저} 사람은 누구인가?
\kfChoice \uline{다섯}이 모였다.
\kfChoice \uline{아}, 깜짝이야!
\end{kfChoices}
\end{kfQBlock}
\begin{kfQBlock}{10}{객관식}
\kfQStem{다음 밑줄 친 단어 중 품사가 부사인 것은?}
\begin{kfChoices}
\kfChoice \uline{새} 신발을 샀다.
\kfChoice \uline{우리}는 친구다.
\kfChoice \uline{첫째}도 건강이다.
\kfChoice \uline{어머나}, 꽃이 피었네.
\kfChoice \uline{정말} 멋지다.
\end{kfChoices}
\end{kfQBlock}
\begin{kfQBlock}{11}{객관식}
\kfQStem{다음 문장에 사용되지 \uline{\underline{않은}} 품사는?}
\begin{kfEvidence}와, 저 새가 아주 높이 나는구나!\end{kfEvidence}
\begin{kfChoices}
\kfChoice 감탄사
\kfChoice 명사
\kfChoice 관형사
\kfChoice 부사
\kfChoice 형용사
\end{kfChoices}
\end{kfQBlock}
\begin{kfQBlock}{12}{객관식}
\kfQStem{<보기>의 단어들을 바르게 분류한 것은?}
\begin{kfEvidence}하늘, 나, 첫째, 가, 먹다, 예쁘다, 새, 매우, 아\end{kfEvidence}
\begin{kfChoices}
\kfChoice 가변어: 하늘, 나, 첫째, 가, 새, 매우, 아
\kfChoice 불변어: 먹다, 예쁘다
\kfChoice 체언: 하늘, 나, 가
\kfChoice 용언: 먹다, 예쁘다
\kfChoice 수식언: 첫째, 새, 매우
\end{kfChoices}
\end{kfQBlock}
\begin{kfQBlock}{13}{객관식}
\kfQStem{다음 밑줄 친 단어의 품사가 나머지와 \uline{다른} 하나는?}
\begin{kfChoices}
\kfChoice \uline{그} 책을 다오.
\kfChoice \uline{새} 옷을 샀다.
\kfChoice \uline{헌} 집을 고쳤다.
\kfChoice \uline{모든} 사람이 웃었다.
\kfChoice \uline{그}는 학생이다.
\end{kfChoices}
\end{kfQBlock}
\begin{kfQBlock}{14}{객관식}
\kfQStem{다음 밑줄 친 단어의 품사가 나머지와 \uline{다른} 하나는?}
\begin{kfChoices}
\kfChoice \uline{하나}만 먹어라.
\kfChoice \uline{첫째}는 건강이다.
\kfChoice \uline{다섯}이 모였다.
\kfChoice \uline{한} 사람이 왔다.
\kfChoice \uline{둘}은 되고 셋은 안 된다.
\end{kfChoices}
\end{kfQBlock}
\begin{kfQBlock}{15}{객관식}
\kfQStem{기능에 따른 품사 분류가 \uline{잘못} 짝지어진 것은?}
\begin{kfChoices}
\kfChoice 체언 - 명사, 대명사, 수사
\kfChoice 관계언 - 조사
\kfChoice 용언 - 동사, 형용사
\kfChoice 수식언 - 관형사, 부사
\kfChoice 독립언 - 부사
\end{kfChoices}
\end{kfQBlock}
\begin{kfQBlock}{16}{단답형}
\kfQStem{품사를 분류하는 세 가지 기준은 무엇인가?}
\kfAnswerLines{1}
\end{kfQBlock}
\begin{kfQBlock}{17}{단답형}
\kfQStem{문장 속에서 형태가 변하는 단어를 '가변어'라 하고, 형태가 고정된 단어를 '이것'이라 한다. '이것'은 무엇인가?}
\kfAnswerLines{1}
\end{kfQBlock}
\begin{kfQBlock}{18}{단답형}
\kfQStem{국어의 9품사 중 가변어에 속하는 품사 3가지를 모두 쓰시오.}
\kfAnswerLines{1}
\end{kfQBlock}
\begin{kfQBlock}{19}{단답형}
\kfQStem{문장에서 뼈대 역할을 하는 '체언'에 속하는 품사 3가지를 쓰시오.}
\kfAnswerLines{1}
\end{kfQBlock}
\begin{kfQBlock}{20}{단답형}
\kfQStem{주로 체언 뒤에 붙어 문법적 관계를 나타내는 '관계언'에 속하는 품사는 무엇인가?}
\kfAnswerLines{1}
\end{kfQBlock}
\begin{kfQBlock}{21}{단답형}
\kfQStem{문장에서 서술어 역할을 하며 '활용'을 하는 '용언'에 속하는 품사 2가지를 쓰시오.}
\kfAnswerLines{1}
\end{kfQBlock}
\begin{kfQBlock}{22}{단답형}
\kfQStem{다른 말을 꾸며주는 '수식언'에 속하는 품사 2가지를 쓰시오.}
\kfAnswerLines{1}
\end{kfQBlock}
\begin{kfQBlock}{23}{단답형}
\kfQStem{'새', '헌', '모든'처럼 오직 체언만을 꾸며주는 품사는 무엇인가?}
\kfAnswerLines{1}
\end{kfQBlock}
\begin{kfQBlock}{24}{단답형}
\kfQStem{'아주', '매우', '방긋'처럼 주로 용언을 꾸며주는 품사는 무엇인가?}
\kfAnswerLines{1}
\end{kfQBlock}
\begin{kfQBlock}{25}{단답형}
\kfQStem{놀람, 부름, 대답을 나타내며 독립적으로 쓰이는 '독립언'에 속하는 품사는 무엇인가?}
\kfAnswerLines{1}
\end{kfQBlock}
\begin{kfQBlock}{26}{서술형}
\kfQStem{'용언'이 다른 품사와 구별되는 가장 큰 특징 2가지를 서술하시오.}
\kfAnswerNote{답안}{4}
\end{kfQBlock}
\kfEndMC


\kfStartConcept{구수한 사투리의 맛, 방언의 효과}


\kfPassageNoteNT{%
　우리가 사용하는 언어는 지역에 따라 조금씩 다른 특징을 보이는데, 이처럼 특정 지역에서만 사용하는 고유한 말을 방언 또는 사투리라고 한다. 표준어는 공식적인 자리에서 의사소통의 효율성을 위해 사용하는 공통어이지만, 방언은 특정 지역 사람들의 삶과 정서, 문화가 고스란히 녹아 있는 살아있는 언어이다. 그렇다면 문학 작품에서는 방언을 왜 사용할까?

　문학 작품에서 방언이 사용될 때 나타나는 효과는 여러 가지가 있다. 첫째, 작품에 향토성을 더해준다. '향토성'이란 그 지역 고유의 정겨운 느낌, 즉 '고향의 냄새'를 뜻한다. 작품의 배경이 되는 지역의 방언을 사용하면, 독자는 낯선 지명 대신 구수한 사투리를 통해 그 장소를 더 구체적이고 현실적으로 느끼게 된다. 

　둘째, 현장감과 사실성을 높여준다. 인물들이 실제 그 지역 사람들처럼 방언을 사용하면, 독자는 그 이야기가 꾸며낸 것이 아니라 '정말로 일어난 일'처럼 생생하게 느끼게 된다. 

　셋째, 인물의 성격을 효과적으로 드러내고 친근감을 준다. 특정 지역의 방언은 그 지역 사람들의 분위기나 성격을 보여주는 경우가 많다. 예를 들어, 김유정의 「동백꽃」에서 강원도 방언을 쓰는 '점순이'의 말투는 순박하면서도 앙큼한 성격을 잘 보여준다. 마찬가지로 이효석의 「메밀꽃 필 무렵」에서는 강원도 산골 사람들의 정겹고 소박한 삶이 방언을 통해 잘 드러난다. 독자는 표준어를 쓰는 인물보다, 방언을 쓰는 인물에게서 더 큰 친근감을 느끼고 이야기에 깊이 빠져들게 된다.

　이처럼 방언은 단순히 지역의 말을 옮겨 적는 것을 넘어선다. 방언은 작품의 배경을 생생하게 만들고, 인물에게 실제 살아 숨 쉬는 듯한 생명력을 불어넣는다. 독자는 방언을 통해 작품을 더욱 현실적으로 받아들이며 특별한 문학적 재미를 느낄 수 있다.
\par\medskip\begin{itemize}[leftmargin=18pt, itemsep=2pt, topsep=2pt, label=\textbullet]
\item 향토성 — 김유정 「동백꽃」의 강원도 방언
\item 현장감 — 이효석 「메밀꽃 필 무렵」의 산골 말투
\item 성격 제시 — 점순이의 앙큼한 말투를 통한 인물 성격 형상화
\end{itemize}
}

\kfSecHeader{kfAreaConcept}{문제}{}

\par\medskip
\kfBeginMC
\begin{kfQBlock}{01}{객관식}
\kfQStem{'방언'에 대한 설명으로 적절하지 \uline{\underline{않은}} 것은?}
\begin{kfChoices}
\kfChoice 사투리라고도 부른다.
\kfChoice 특정 지역에서만 사용하는 고유한 말이다.
\kfChoice 공식적인 자리에서 의미를 전달하기 위해 사용한다.
\kfChoice 그 지역 사람들의 삶과 문화가 녹아 있다.
\kfChoice 문학 작품에서 특별한 효과를 위해 자주 사용된다.
\end{kfChoices}
\end{kfQBlock}
\begin{kfQBlock}{02}{객관식}
\kfQStem{‘향토성’에 대한 설명으로 적절하지 \uline{\underline{않은}} 것은?}
\begin{kfChoices}
\kfChoice 그 지역 고유의 정겨운 분위기와 느낌을 가리키는 개념이다.
\kfChoice 흔히 ‘고향의 냄새’와 같은 정서로 풀이되기도 하는 개념이다.
\kfChoice 방언·지명·풍속 등이 작품의 향토성을 살리는 데 크게 기여한다.
\kfChoice 작품의 배경에 지역적 색채를 짙게 불어넣어 주는 요소이다.
\kfChoice 독자에게 단순히 문학적 재미만 주는 표면적인 효과이다.
\end{kfChoices}
\end{kfQBlock}
\begin{kfQBlock}{03}{객관식}
\kfQStem{방언의 효과로 적절하지 \uline{\underline{않은}} 것은?}
\begin{kfChoices}
\kfChoice 향토성을 더해준다.
\kfChoice 현장감과 사실성을 높여준다.
\kfChoice 인물의 성격을 효과적으로 드러낸다.
\kfChoice 독자에게 친근감을 준다.
\kfChoice 작품의 배경을 추상적으로 만든다.
\end{kfChoices}
\end{kfQBlock}
\begin{kfQBlock}{04}{객관식}
\kfQStem{위 글을 바탕으로 할 때, 방언을 사용한 소설을 읽는 독자의 반응으로 적절하지 \uline{\underline{않은}} 것은?}
\begin{kfChoices}
\kfChoice "주인공의 사투리가 참 구수해서 정감이 가네."
\kfChoice "이야기가 꼭 실제로 있었던 일처럼 생생하게 느껴져."
\kfChoice "소설의 배경이 되는 시골 마을이 눈앞에 그려지는 것 같아."
\kfChoice "주인공의 말투를 들으니 어떤 성격인지 짐작이 가."
\kfChoice "모든 인물이 표준어만 쓰니까 작품의 품격이 높아지는군."
\end{kfChoices}
\end{kfQBlock}
\begin{kfQBlock}{05}{객관식}
\kfQStem{<보기>는 김유정의 「동백꽃」의 한 구절이다. 밑줄 친 방언의 효과로 가장 적절한 것은?}
\begin{kfEvidence}"느 집엔 이거 없지?"   \\\relax
   하고 섰다. 그 다음에는 “느 아버지가 고자라지?” 하며 놀렸다.\end{kfEvidence}
\begin{kfChoices}
\kfChoice 인물의 성격을 앙큼하게 드러낸다.
\kfChoice 작품의 배경을 도시로 느끼게 한다.
\kfChoice 인물의 지적인 면모를 강조한다.
\kfChoice 독자가 인물과 거리감을 느끼게 한다.
\kfChoice 사건의 비극성을 강조한다.
\end{kfChoices}
\end{kfQBlock}
\begin{kfQBlock}{06}{객관식}
\kfQStem{<보기>는 이효석의 「메밀꽃 필 무렵」의 한 구절이다. 이 구절에 나타난 방언의 효과로 가장 적절한 것은?}
\begin{kfEvidence}"허 생원은 걸핏하면 이 길을 기억하니께."   \\\relax
   "그렇지, 밤길을 걸으면서도 이 길을 걸은 지가 이십 년이나 됐네."\end{kfEvidence}
\begin{kfChoices}
\kfChoice 인물의 냉철한 성격을 드러낸다.
\kfChoice 정겹고 소박한 삶의 모습을 보여준다.
\kfChoice 긴박한 사건의 전개를 암시한다.
\kfChoice 인물의 부유한 생활을 짐작하게 한다.
\kfChoice 작품의 배경이 서울임을 암시한다.
\end{kfChoices}
\end{kfQBlock}
\begin{kfQBlock}{07}{객관식}
\kfQStem{소설에서 방언을 사용하여 얻는 효과 중 <보기>가 설명하는 것은 무엇인가?}
\begin{kfEvidence}소설을 읽으니 마치 내가 강원도 봉평 장터에 직접 가 있는 듯한 생생한 느낌이 들었다.\end{kfEvidence}
\begin{kfChoices}
\kfChoice 향토성
\kfChoice 현장감
\kfChoice 친근감
\kfChoice 해학성
\kfChoice 교훈성
\end{kfChoices}
\end{kfQBlock}
\begin{kfQBlock}{08}{객관식}
\kfQStem{소설에서 방언을 사용하여 얻는 효과 중 <보기>가 설명하는 것은 무엇인가?}
\begin{kfEvidence}주인공이 사용하는 구수한 사투리를 들으니 왠지 모르게 정겹고 그 마을 특유의 분위기가 느껴졌다.\end{kfEvidence}
\begin{kfChoices}
\kfChoice 향토성
\kfChoice 비판성
\kfChoice 주관성
\kfChoice 객관성
\kfChoice 긴장감
\end{kfChoices}
\end{kfQBlock}
\begin{kfQBlock}{09}{단답형}
\kfQStem{특정 지역에서만 사용하는 고유한 말을 무엇이라 하는가?}
\kfAnswerLines{1}
\end{kfQBlock}
\begin{kfQBlock}{10}{단답형}
\kfQStem{공식적인 자리에서 의사소통을 위해 사용하는 공통어는 무엇인가?}
\kfAnswerLines{1}
\end{kfQBlock}
\begin{kfQBlock}{11}{단답형}
\kfQStem{방언이 그 지역 고유의 정겨운 느낌, 즉 '고향의 냄새'를 더해주는 효과를 무엇이라 하는가?}
\kfAnswerLines{1}
\end{kfQBlock}
\begin{kfQBlock}{12}{단답형}
\kfQStem{방언을 사용해 이야기가 꾸며낸 것이 아니라 '정말로 일어난 일'처럼 생생하게 느끼게 하는 효과는 무엇인가?}
\kfAnswerLines{1}
\end{kfQBlock}
\begin{kfQBlock}{13}{단답형}
\kfQStem{독자가 방언을 쓰는 인물에게 정이 가고 이야기 속으로 몰입하게 되는 감정은 무엇인가?}
\kfAnswerLines{1}
\end{kfQBlock}
\begin{kfQBlock}{14}{단답형}
\kfQStem{「동백꽃」에서 강원도 방언을 쓰는 '점순이'의 말투가 보여주는 성격적 특징 두 가지는 무엇인가?}
\kfAnswerLines{1}
\end{kfQBlock}
\begin{kfQBlock}{15}{단답형}
\kfQStem{「메밀꽃 필 무렵」에서 방언을 통해 드러나는 강원도 산골 사람들의 삶의 모습은 어떠한가?}
\kfAnswerLines{1}
\end{kfQBlock}
\begin{kfQBlock}{16}{서술형}
\kfQStem{'방언'과 '표준어'의 개념을 지문에 근거하여 각각 서술하시오.}
\kfAnswerNote{답안}{4}
\end{kfQBlock}
\kfEndMC


\kfStartLit{「금 따는 콩밭」 — 김유정}


\kfPassageNote{}{%
[앞부분 줄거리] 성실한 농부였던 영식은 친구 수재로부터 자신의 콩밭 밑에 금맥이 있다는 말을 듣고 아내와 함께 일확천금의 꿈에 부푼다. 그는 애써 가꾼 콩밭을 파헤치기 시작하지만 금은 나오지 않고 농사마저 망쳐버린다. 결국 땅 주인인 마름에게 밭을 원래대로 돌려놓으라는 큰 꾸지람을 듣고 쫓겨날 위기에 처한다. 시간이 갈수록 초조해진 영식은 금이 나올 거라는 말만 반복하는 수재와 다투고 절망감에 빠진다. 고된 노동 끝에 아무런 수확도 없이 진흙투성이가 되어 집으로 돌아온다.

　밤은 캄캄하게 어두웠다. 어디선가 개 떼가 요란하게 짖어댄다.

　남편은 진흙투성이가 되어 돌아왔다. 기운이 빠져 몸도 제대로 가누지 못하고 아랫목에 축 늘어진다.

　이 꼴을 보니 아내는 다시 기운이 빠진다. 오늘도 또 실패했구나. 금이 나오면 집을 한 채 산다고 자랑하고 왔는데 결국 헛일이었다. 이제는 창피해서 얼굴을 들고 나갈 수도 없게 되었다.

　남편에게 저녁을 차려 주고 안타깝게 바라본다.

"이제 빌려 온 쌀도 다 떨어졌는데……."

"새벽에 산에 제사를 좀 지낼 텐데 한 번만 더 빌려 와."

　아내의 말에는 대답도 없고 느릿느릿 힘없는 소리뿐이다. 그리고 드러누운 채 눈을 지그시 감아 버린다.

"죽 끓일 쌀도 없는데 산에 지내는 제사는 무슨……."

"듣기 싫어, 재수 없는 소리 하지 마."

　이 고함에 아내는 그만 움찔했다. 요즘 들어 남편은 아무 이유 없이 화만 내서 딱했다. 미친 사람처럼 밤잠도 안 자고 소리만 빽빽 지르며 덤벼들려고 한다. 심지어 어린아이가 좀 울어도 내다 버리라고 난리를 피우는 것이다.

　저녁을 안 먹기에 그냥 치워 버렸다. 남편의 명령을 거스르기 어려워 양근댁에게 또다시 안 갈 수가 없다. 그동안 쌀은 계속 빌려다 먹고 갚지도 못했는데 또 무슨 얼굴로 입을 벌릴지 곤란했다.

　그는 생각 끝에 부끄러움을 무릅쓰고 다시 한번 찾아가는 것이지만, 막상 마주 보고 입을 열어,

"내일 산에 제사를 지낸다는데 쌀이 있어야지요."

　라고 말하려니 얼굴이 화끈거렸다.

　그러나 그들은 무척 착한 사람들이었다.

"암, 그렇지요. 산신령님이 화내면 아무것도 안 됩니다."

　하고 말을 받으며 그 집 남편은 빙그레 웃는다. 원래 금 캐는 일에 경험이 많은 사람이라 이런 일에는 이해심이 깊었다. 직접 쌀 닷 되를 퍼다 주며,

"산에 지내는 제사란 안 지내면 몰라도 이왕 지내려면 아주 정성껏 해야 합니다. 산신이 화를 잘 내니까요."

　하고 그 방법까지 알려준다.

　쌀을 받아 들고 나오며 영식이 아내는 고마움보다 먼저 미안함에 얼굴이 다시 빨개졌다. 그리고 그들 부부가 살아가는 모습이 정말 몹시 부러웠다. 양근댁 남편은 날마다 금 캐는 곳을 돌아다니며 버려진 돌멩이 더미를 뒤져 금이 섞인 돌을 주워 온다. 그걸 온종일 돌판에 갈면 운이 좋으면 이삼 원, 적어도 칠팔십 전은 매일 버는 것이었다. 그러면 쌀을 사고, 옷감을 사고, 떡을 하고, 돈놀이도 하고……. 그런데 우리는 왜 늘 이 모양인지 생각만 해도 가슴이 답답해서 한숨이 계속 나왔다.

　아내는 집에 돌아와 떡 만들 쌀을 물에 담갔다. 내일은 뭘로 죽을 쑤어 먹을지. 윗목에 웅크리고 앉아서 건너편에 누워 있는 남편을 곁눈으로 살짝 쳐다본다. 남들은 돌아다니며 잘도 금을 주워 오는데, 저 한심한 인간은 제 밭 하나를 다 망쳐도 금 한 톨 못 주워 오나. 에휴, 변변치도 못한 사내. 저도 모르게 얕은 한숨이 두 번 터져 나온다.

　밤이 깊어 그들 부부는 떡을 하러 나왔다. 남편은 절구에 쿵쿵 쌀을 빻았다. 그러나 체가 없다. 동네로 돌아다니며 빌려 오느라 아내는 다리가 퉁퉁 부었다.

"왜 그냥 앉아 있어요, 불 좀 때지 않고."

　떡을 찧다가 넋이 빠져 멍하니 앉아 있는 남편이 얄밉다. 나는 이렇게 애를 쓰는데 저 사람은 무슨 생각을 하고 저러고 있는지. 낫으로 마른 나뭇가지를 탁탁 잘라 던져 주며 아내는 은근히 남편을 다그쳤다.

　닭이 두 번째 울고 나서야 떡은 완성되었다.

　아내는 시루를 이고 남편은 겨드랑이에 돗자리를 꼈다. 그리고 캄캄한 산길을 올라간다.

　비탈길을 얼마 올라가서야 콩밭이 나타났다. 앞이 우뚝 솟은 검은 산에 둘러싸여 막힌 곳이었다. 가장자리로 느티나무와 대추나무들이 서 있었다.

　밭머리에 조금 못 미처 남편은 걸음을 멈추고 뒤따라오던 아내를 돌아본다.

"이리 내. 그리고 여기 가만히 서 있어."

　시루를 받아 한 팔로 껴안고 그는 혼자서 콩밭으로 올라섰다. 앞에 쌓인 것이 모두 흙더미, 그 흙더미를 막 돌아 서려 할 때 아마 돌을 찼나 보다. 몸이 쓰러지려고 비틀거리자 아내가 깜짝 놀라 뛰어오르며 그를 부축했다.

"부정 타게 왜 올라와, 재수 없는 여자 같으니."

　남편은 몸을 바로 잡자마자 소리를 빽 지르며 아내의 뺨을 때린다. 가뜩이나 힘들어 죽겠는데 불길하게도 여자가. 그는 못마땅한 듯 투덜거리며 밭으로 들어간다.

　밭 한가운데에 돗자리를 펴고 그 위에 시루를 놓았다. 그리고 시루 앞에 공손하고 정성스럽게 큰절을 한다.

"우리를 살려 주십시오. 산신께서 도와주지 않으시면 저희는 죽을 수밖에 없습니다."

　그는 손을 모으고 이렇게 빌었다.

　아내는 이 꼴을 바라보며 화가 치밀어 올랐다. 금을 캔다고 하더니 금 한 톨 못 캐고 버릇만 점점 나빠져 간다. 전에는 안 그랬는데 요즘은 걸핏하면 때리는 못된 버릇이 생긴 것이다. 금을 캐랬지, 뺨을 때리랬나. 제발 그놈의 금 좀 나오지 말았으면. 그는 뺨 맞은 것에 앙심을 품고 마음껏 반항적인 생각을 했다.

[뒷부분 줄거리]   산신에게 제사를 지낸 후에도 금이 나오지 않자 영식의 절망은 깊어만 간다. 아내가 빌려온 쌀 문제를 언급하며 다툼이 벌어지고, 분노가 폭발한 영식은 아내에게 심한 폭력을 행사한다. 갈등이 최고조에 이른 순간, 영식의 폭력이 자신에게 향할 것을 두려워한 수재는 위기를 모면하기 위해 거짓으로 금맥을 발견했다고 소리친다. 이에 영식과 아내는 모든 갈등을 잊고 기쁨에 휩싸이지만, 수재는 그날 밤 마을을 떠나기로 결심한다.
}


\kfPassageNote{문장 독해 — 발췌}{%
1. 남편은 진흙투성이가 되어 돌아왔다. 기운이 빠져 몸도 제대로 가누지 못하고 아랫목에 축 늘어진다. 2. 그는 생각 끝에 부끄러움을 무릅쓰고 다시 한번 찾아가는 것이지만... 3. 그러나 그들은 무척 착한 사람들이었다. 4. “산신이 화를 잘 내니까요.” 하고 그 방법까지 알려준다. 5. 그런데 우리는 왜 늘 이 모양인지 생각만 해도 가슴이 답답해서 한숨이 계속 나왔다. 6. 저 한심한 인간은 제 밭 하나를 다 망쳐도 금 한 톨 못 주워 오나. 7. 나는 이렇게 애를 쓰는데 저 사람은 무슨 생각을 하고 저러고 있는지. 8. 우리를 살려 주십시오. 산신께서 도와주지 않으시면 저희는 죽을 수밖에 없습니다. 9. 그는 손을 모으고 이렇게 빌었다. 10. 아내는 이 꼴을 바라보며 화가 치밀어 올랐다.
}


\kfPassageNote{해설 학습 1}{%
소설가 김유정은 우리나라 문학에서 웃음과 슬픔을 가장 잘 표현한 작가로 유명하다. 그의 작품에는 '해학'과 '비애'가 함께 담겨 있다. 해학이란 인물의 어리석은 행동이나 모습을 재미있게 표현하여 웃음을 불러일으키는 것을 말한다. 반면 비애는 마음속 깊은 곳에서 느껴지는 슬픔을 의미한다. 「금 따는 콩밭」은 이 웃음과 슬픔이 절묘하게 섞여 있는 작품이다. 금이 나온다는 말에 속아 멀쩡한 콩밭을 파헤치는 영식 부부의 모습은 너무나 어리석어서 웃음이 나온다. 특히 다 망해버린 콩밭 한가운데에 떡시루를 가져다 놓고 살려달라고 비는 영식의 모습은 너무나 진지해서 오히려 우스꽝스럽게 느껴진다. 이것이 바로 해학이다. 하지만 이 웃음 뒤에는 깊은 슬픔이 숨어 있다. 영식 부부가 이처럼 헛된 꿈에 매달린 이유는 바로 지독한 가난 때문이다. 자기 땅 한 평 없이 하루하루 힘들게 살아가던 1930년대 농민들에게 금은 유일한 희망처럼 보였다. 가난에서 벗어나고 싶은 절박한 마음이 그들을 비상식적인 행동으로 내몰았던 것이다. 결국 영식은 농사도 망치고 빚만 지게 되는데, 이는 당시 농촌의 가난한 현실을 보여주는 비애라 할 수 있다. 이처럼 독자는 인물들의 어리석은 모습에 웃으면서도, 동시에 그들이 처한 비참한 현실에 연민과 안타까움을 느끼게 된다. 이처럼 웃음과 슬픔을 동시에 느끼게 하는 것이 바로 김유정 소설의 뛰어난 점이다.
}


\kfPassageNote{해설 학습 2}{%
김유정의 「금 따는 콩밭」은 성실했던 한 농부가 ‘금’이라는 헛된 욕망에 빠져 점차 망가져가는 과정을 그리고 있다. 주인공 영식은 원래 자신의 땅을 소중히 여기며 땀 흘려 일하던 성실한 농부였다. 하지만 친구 수재의 “콩밭에 금이 묻혀 있다”는 말 한마디에 넘어가면서 그의 삶은 흔들리기 시작한다. 이 작품에서 '콩밭'과 '금'은 서로 반대되는 의미를 지닌다. 콩밭은 땀 흘려 일하는 정직한 노동을 상징한다. 반면 금은 노력 없이 한 번에 큰돈을 벌려는 헛된 꿈, 즉 '일확천금'을 상징한다. 가난에 지친 영식에게 땅속의 금은 모든 문제를 단번에 해결해 줄 것처럼 보였다. 그래서 그는 자신의 유일한 생계 수단인 콩밭을 망가뜨리면서까지 금 캐기에 매달렸다. 금이 나오지 않자 그는 점점 더 초조해지고, 미신에 의지하며 올바른 판단을 하지 못하게 된다. 결국 성실했던 농부는 아내를 폭행하고, 일도 가정도 모두 잃은 채 절망에 빠지는 비참한 인물로 변하고 만다. 이 작품은 일확천금이라는 헛된 욕망이 한 사람을 얼마나 완전히 망가뜨릴 수 있는지를 보여준다. 이를 통해 작가는 눈에 보이지 않는 헛된 꿈보다, 땀 흘려 일하는 성실한 삶의 가치가 훨씬 소중하다는 교훈을 전달한다.
}


\kfPassageNote{해설 학습 3}{%
이 소설은 헛된 꿈이 가족과 이웃 사이에 어떤 갈등을 일으키는지를 자세히 보여준다. 특히 금 캐기를 둘러싼 인물들의 심리 변화와 갈등 과정이 생생하게 드러난다. 처음에 아내는 남편과 함께 금에 대한 희망을 가졌다. 그러나 금이 나오지 않고 가난이 계속되자 점차 남편을 원망하게 된다. 아내는 빚을 내어 쌀을 빌려오는 현실적인 고통 속에서 남편의 무능함을 탓하며 '변변치도 못한 사내'라고 생각한다. 심지어 남편의 폭력에 시달리자 '제발 그놈의 금 좀 나오지 말았으면' 하고 저주를 퍼붓기도 한다. 이는 헛된 희망이 부부 사이의 믿음과 사랑마저 깨뜨린다는 것을 보여준다. 영식 역시 금이 나오지 않는 초조함과 아내의 잔소리 때문에 점점 더 폭력적으로 변해간다. 한편 이 모든 사건의 원인을 제공한 수재는 영식과 계속해서 갈등을 겪는다. 갈등이 가장 심해졌을 때, 수재는 위기를 벗어나기 위해 아무 돌멩이를 주워들고 금맥을 발견했다고 거짓말을 한다. 절망에 빠져 있던 영식과 아내는 수재의 거짓말 한마디에 모든 갈등을 잊고 다시 기쁨에 빠진다. 이처럼 인물들은 절망과 희망 사이를 오가며 어리석은 모습을 반복하는 것이다.
}


\kfSecHeader{kfAreaLiterature}{문제}{}

\par\medskip
\kfBeginMC
\begin{kfQBlock}{01}{객관식}
\kfQStem{다음 문장에서 서술어 ‘늘어진다’의 주어로 적절한 것은?}
\begin{kfChoices}
\kfChoice 남편
\kfChoice 수재
\kfChoice 아내
\kfChoice 개 떼
\end{kfChoices}
\end{kfQBlock}
\begin{kfQBlock}{02}{객관식}
\kfQStem{다음 문장에서 ‘그’가 가리키는 인물로 적절한 것은?}
\begin{kfChoices}
\kfChoice 남편
\kfChoice 수재
\kfChoice 영식
\kfChoice 아내
\end{kfChoices}
\end{kfQBlock}
\begin{kfQBlock}{03}{객관식}
\kfQStem{다음 문장에서 ‘그들’이 가리키는 인물들로 적절한 것은?}
\begin{kfChoices}
\kfChoice 영식 부부
\kfChoice 양근댁 부부
\kfChoice 마을 사람들
\kfChoice 수재와 친구들
\end{kfChoices}
\end{kfQBlock}
\begin{kfQBlock}{04}{객관식}
\kfQStem{다음 문장에서 서술어 ‘알려준다’의 주체로 적절한 것은?}
\begin{kfChoices}
\kfChoice 수재
\kfChoice 아내
\kfChoice 영식
\kfChoice 양근댁 남편
\end{kfChoices}
\end{kfQBlock}
\begin{kfQBlock}{05}{객관식}
\kfQStem{다음 문장에서 ‘우리’가 가리키는 인물들로 적절한 것은?}
\begin{kfChoices}
\kfChoice 양근댁 부부
\kfChoice 수재와 영식
\kfChoice 영식과 아내
\kfChoice 마을 사람들
\end{kfChoices}
\end{kfQBlock}
\begin{kfQBlock}{06}{객관식}
\kfQStem{다음 문장에서 ‘저 한심한 인간’이 가리키는 인물로 적절한 것은?}
\begin{kfChoices}
\kfChoice 영식
\kfChoice 수재
\kfChoice 마름
\kfChoice 양근댁
\end{kfChoices}
\end{kfQBlock}
\begin{kfQBlock}{07}{객관식}
\kfQStem{다음 문장에서 ‘저 사람’이 가리키는 인물로 적절한 것은?}
\begin{kfChoices}
\kfChoice 수재
\kfChoice 영식
\kfChoice 마름
\kfChoice 양근댁 남편
\end{kfChoices}
\end{kfQBlock}
\begin{kfQBlock}{08}{객관식}
\kfQStem{다음 문장에서 ‘우리’가 가리키는 인물들로 적절한 것은?}
\begin{kfChoices}
\kfChoice 양근댁 부부
\kfChoice 수재와 영식
\kfChoice 영식과 아내
\kfChoice 마을 사람들
\end{kfChoices}
\end{kfQBlock}
\begin{kfQBlock}{09}{객관식}
\kfQStem{다음 문장에서 ‘그’가 가리키는 인물로 적절한 것은?}
\begin{kfChoices}
\kfChoice 수재
\kfChoice 영식
\kfChoice 아내
\kfChoice 양근댁 남편
\end{kfChoices}
\end{kfQBlock}
\begin{kfQBlock}{10}{객관식}
\kfQStem{다음 문장에서 ‘이 꼴’이 가리키는 모습으로 가장 적절한 것은?}
\begin{kfChoices}
\kfChoice 남편이 밭에서 산신에게 절을 하는 모습
\kfChoice 남편이 금이 나오지 않아 아내에게 화를 내는 모습
\kfChoice 남편이 밭 한가운데서 멍하니 하늘을 쳐다보는 모습
\kfChoice 남편이 수재와 함께 금을 캐기 위해 땅을 파는 모습
\end{kfChoices}
\end{kfQBlock}
\begin{kfQBlock}{11}{객관식}
\kfQStem{「금 따는 콩밭」의 작가는 누구인가?}
\kfAnswerLines{2}
\end{kfQBlock}
\begin{kfQBlock}{12}{객관식}
\kfQStem{김유정 작가가 작품 속에서 가장 잘 표현한 두 가지 감정은 무엇인가?}
\kfAnswerLines{2}
\end{kfQBlock}
\begin{kfQBlock}{13}{객관식}
\kfQStem{「금 따는 콩밭」 은 몇 년대의 시대적 배경을 바탕으로 하는가?}
\kfAnswerLines{2}
\end{kfQBlock}
\begin{kfQBlock}{14}{객관식}
\kfQStem{인물의 어리석은 행동을 재미있게 표현하여 웃음을 만드는 것을 무엇이라고 하는가?}
\kfAnswerLines{2}
\end{kfQBlock}
\begin{kfQBlock}{15}{객관식}
\kfQStem{마음속 깊은 곳에서 느껴지는 슬픔을 의미하는 단어는 무엇인가?}
\kfAnswerLines{2}
\end{kfQBlock}
\begin{kfQBlock}{16}{객관식}
\kfQStem{영식 부부가 허황된 꿈에 매달리게 된 근본적인 이유는 무엇인가?}
\kfAnswerLines{2}
\end{kfQBlock}
\begin{kfQBlock}{17}{객관식}
\kfQStem{주인공이 결국 농사를 망치고 빚만 지게 되는 모습은 당시 무엇을 보여주는가?}
\kfAnswerLines{2}
\end{kfQBlock}
\begin{kfQBlock}{18}{객관식}
\kfQStem{이 글의 주인공 이름은 무엇인가?}
\kfAnswerLines{2}
\end{kfQBlock}
\begin{kfQBlock}{19}{객관식}
\kfQStem{영식의 원래 직업은 무엇이었는가?}
\kfAnswerLines{2}
\end{kfQBlock}
\begin{kfQBlock}{20}{객관식}
\kfQStem{영식에게 금맥이 있다고 말해 준 친구의 이름은 무엇인가?}
\kfAnswerLines{2}
\end{kfQBlock}
\begin{kfQBlock}{21}{객관식}
\kfQStem{이 글에서 ‘콩밭’이 상징하는 것은 무엇인가?}
\kfAnswerLines{2}
\end{kfQBlock}
\begin{kfQBlock}{22}{객관식}
\kfQStem{‘콩밭’과 반대되는 의미를 지닌 ‘금’이 상징하는 것은 무엇인가?}
\kfAnswerLines{2}
\end{kfQBlock}
\begin{kfQBlock}{23}{객관식}
\kfQStem{헛된 욕망에 빠진 영식은 결국 어떤 인물로 변했는가?}
\kfAnswerLines{2}
\end{kfQBlock}
\begin{kfQBlock}{24}{객관식}
\kfQStem{작가는 이 작품을 통해 무엇의 가치를 버리면 삶이 비참하게 무너질 수 있음을 보여주는가?}
\kfAnswerLines{2}
\end{kfQBlock}
\begin{kfQBlock}{25}{객관식}
\kfQStem{처음에 아내는 남편과 함께 금에 대해 어떤 마음을 가졌는가?}
\kfAnswerLines{2}
\end{kfQBlock}
\begin{kfQBlock}{26}{객관식}
\kfQStem{아내가 남편을 원망하게 된 이유는 무엇인가?}
\kfAnswerLines{2}
\end{kfQBlock}
\begin{kfQBlock}{27}{객관식}
\kfQStem{아내는 현실적인 고통 속에서 남편을 어떤 사내라고 생각했는가?}
\kfAnswerLines{2}
\end{kfQBlock}
\begin{kfQBlock}{28}{객관식}
\kfQStem{헛된 희망은 부부 사이의 무엇을 깨뜨렸는가?}
\kfAnswerLines{2}
\end{kfQBlock}
\begin{kfQBlock}{29}{객관식}
\kfQStem{영식의 성격이 폭력적으로 변해간 이유는 무엇인가?}
\kfAnswerLines{2}
\end{kfQBlock}
\begin{kfQBlock}{30}{객관식}
\kfQStem{이 모든 사건의 원인을 제공한 인물은 누구인가?}
\kfAnswerLines{2}
\end{kfQBlock}
\begin{kfQBlock}{31}{객관식}
\kfQStem{수재는 위기를 벗어나기 위해 어떤 거짓말을 했는가?}
\kfAnswerLines{2}
\end{kfQBlock}
\begin{kfQBlock}{32}{객관식}
\kfQStem{이 글에 대한 설명으로 적절하지 \uline{\underline{않은}} 것은?}
\begin{kfChoices}
\kfChoice 어리석은 인물의 행동을 통해 당대 현실을 신랄하게 비판한다.
\kfChoice 일확천금이라는 헛된 욕망에 빠져드는 인물을 비중 있게 그린다.
\kfChoice 농촌의 어려운 현실과 인간의 탐욕을 함께 깊이 있게 다룬다.
\kfChoice 주인공의 비합리적이고 무모한 행동을 통해 풍자가 이루어진다.
\kfChoice 모든 갈등이 행복한 결말로 해소되어 희망을 보여 주는 작품이다.
\end{kfChoices}
\end{kfQBlock}
\begin{kfQBlock}{33}{객관식}
\kfQStem{영식의 성격 변화로 적절하지 \uline{\underline{않은}} 것은 무엇인가?}
\begin{kfChoices}
\kfChoice 현실적이었지만 미신을 믿게 되었다.
\kfChoice 이기적인 모습에서 아내를 이해하게 되었다.
\kfChoice 순박했지만 신경질적이고 난폭하게 변했다.
\kfChoice 성실한 농부에서 요행을 바라는 사람이 되었다.
\kfChoice 아내를 아꼈지만 폭력을 사용하는 사람으로 변했다.
\end{kfChoices}
\end{kfQBlock}
\begin{kfQBlock}{34}{객관식}
\kfQStem{이 글에 등장하는 ‘수재’에 대한 설명으로 가장 알맞은 것은 무엇인가?}
\begin{kfChoices}
\kfChoice 금을 캐는 능력이 매우 뛰어난 분야의 전문가이다.
\kfChoice 영식을 진심으로 돕는 의리 있고 따뜻한 친구이다.
\kfChoice 영식의 불행을 진심으로 안타까워하는 착한 이웃이다.
\kfChoice 늘 성실하게 일해서 스스로의 힘으로 성공한 인물이다.
\kfChoice 자기 이익을 위해 친구를 속이는 교활한 인물이다.
\end{kfChoices}
\end{kfQBlock}
\begin{kfQBlock}{35}{객관식}
\kfQStem{남편이 산신에게 제사를 지내려는 태도에 대한 설명으로 적절하지 \uline{\underline{않은}} 것은?}
\begin{kfChoices}
\kfChoice 현실적인 노력을 포기하고 초자연적 존재에 의지하려 한다.
\kfChoice 헛된 믿음에 매달리는 절박한 모습을 보여 준다.
\kfChoice 합리적 판단보다 미신적 발상에 더 의존한다.
\kfChoice 일확천금에 대한 욕망이 인물의 행동을 부추긴다.
\kfChoice 이웃들의 도움에 보답하려는 따뜻한 마음을 표현하려 한다.
\end{kfChoices}
\end{kfQBlock}
\begin{kfQBlock}{36}{객관식}
\kfQStem{아내가 ‘양근댁 부부’를 부러워한 이유로 가장 적절한 것은 무엇인가?}
\begin{kfChoices}
\kfChoice 자식들이 부모님께 효도하며 걱정을 끼치지 않아서
\kfChoice 마을 사람들로부터 인품이 훌륭하다고 존경을 받아서
\kfChoice 부부 사이가 좋아서 한 번도 큰 소리로 싸우지 않아서
\kfChoice 금맥을 발견해 마을에서 제일가는 부자로 살고 있어서
\kfChoice 매일 꾸준히 소득을 올리며 안정적으로 살아가서
\end{kfChoices}
\end{kfQBlock}
\begin{kfQBlock}{37}{객관식}
\kfQStem{㉠\textasciitilde{}㉢에 대한 설명으로 적절하지 \uline{\underline{않은}} 것은 무엇인가?}
\begin{kfChoices}
\kfChoice ㉠은 인물들을 갈등하게 만드는 원인이다.
\kfChoice ㉢은 합리적 노동의 가치를 의미한다.
\kfChoice ㉡은 인물의 비합리적인 믿음을 보여준다.
\kfChoice ㉠은 헛된 꿈을, ㉢은 묵묵한 노동의 가치를 각각 상징적으로 보여 준다.
\kfChoice ㉠과 ㉡은 농사 대신 금을 캐려는 인물의 합리적 계획을 함께 드러낸다.
\end{kfChoices}
\end{kfQBlock}
\begin{kfQBlock}{38}{객관식}
\kfQStem{이 글의 결말에 대한 설명으로 적절하지 \uline{\underline{않은}} 것은?}
\begin{kfChoices}
\kfChoice 등장인물은 잠시 기뻐하지만 곧 더 큰 비극이 예고된다.
\kfChoice 수재의 거짓말 때문에 인물의 기쁨이 곧 절망으로 바뀔 수 있다.
\kfChoice 독자에게 안타까움을 자아내는 비극적 결말이다.
\kfChoice 풍자적 비판이 마지막까지 이어지는 결말이다.
\kfChoice 모든 갈등이 해소되고 인물이 행복을 얻는 권선징악적 결말이다.
\end{kfChoices}
\end{kfQBlock}
\begin{kfQBlock}{39}{객관식}
\kfQStem{<보기>의 관점에서 윗글을 감상한 내용으로 적절하지 \uline{\underline{않은}} 것은 무엇인가?}
\begin{kfEvidence}김유정의 소설에는 당대 농촌의 어려운 현실이 잘 드러나 있다. 작가는 가난 때문에 비극적인 상황에 놓인 인물들의 어리석은 행동을 웃음으로 그려내지만, 그 이면에는 현실에 대한 날카로운 비판과 인간에 대한 깊은 연민을 담고 있다.\end{kfEvidence}
\begin{kfChoices}
\kfChoice 수재의 거짓말은 가난한 시대 현실과도 관련이 있겠군.
\kfChoice 영식의 폭력적인 모습은 비판받아야 하지만 안타까운 마음이 들어.
\kfChoice 영식이 헛된 꿈에 빠진 것은 가난에서 벗어나고 싶었기 때문이겠군.
\kfChoice 망가진 밭에서 제사를 지내는 모습은 어리석지만 그만큼 절박했다는 뜻이겠지.
\kfChoice 작가는 금을 발견한 영식이를 통해 독자들에게 희망을 주려고 하는군.
\end{kfChoices}
\end{kfQBlock}
\begin{kfQBlock}{40}{서술형}
\kfQStem{영식이의 성격은 금을 캐기 전과 후 어떻게 달라졌는지 서술하시오.}
\kfAnswerNote{답안}{4}
\end{kfQBlock}
\begin{kfQBlock}{41}{서술형}
\kfQStem{콩밭을 파는 영식 부부의 모습에서 웃음이 나는 이유와 슬픔이 느껴지는 이유를 각각 서술하시오.}
\kfAnswerNote{답안}{4}
\end{kfQBlock}
\kfEndMC


\kfStartNonlit{[인문] 용서와 추서, 정약용의 실천 윤리}


\kfPassageNote{[인문] 용서와 추서, 정약용의 실천 윤리}{%
　정약용은 조선 후기의 실학자로, 인간의 본성에 대한 연구를 통해 인간의 선한 행위를 설명하려고 했다. 그는 이전까지 절대적 권위를 가지고 있던 주희의 주자학을 비판하며 인간의 본성에 대한 자신만의 이론을 만들었다. 주희는 인간의 본성을 '본연지성'과 '기질지성'으로 설명했다. '본연지성'은 인간이 하늘로부터 받은 순수하고 선한 본성이고, '기질지성'은 본연지성에 사람마다 다른 기질이 더해진 것이다. 그래서 주희는 인간의 기질이 맑으면 선한 행위를 하고 탁하면 악한 행위를 할 수 있다고 보았다. 하지만 정약용은 주희의 관점을 비판했다. 선한 행위와 악한 행위의 원인을 기질이라는 태어날 때부터 정해진 요인으로 본다면 행위에 인간의 의지가 개입되지 않는다. 그렇다면 악한 행위를 한 사람에게 도덕적 책임을 물을 수 없다는 것이 정약용의 생각이었다. 정약용은 인간의 본성을 '기호'라고 보았다. 기호란 즐기고 좋아한다는 뜻으로, 생명이 있는 모든 존재는 각각의 기호를 본성으로 갖는다고 보았다. 꿩은 산을 좋아하고 벼는 물을 좋아하는 것처럼, 인간도 어떤 성향을 갖는다는 것이다. 정약용은 인간에게 '감각적 욕구에서 나온 기호'와 '도덕적 욕구에서 나온 기호'가 있다고 보았다. 감각적 욕구에서 나온 기호는 맛있는 것을 좋아하는 것 같은 몸의 성향이고, 도덕적 욕구에서 나온 기호는 선을 좋아하거나 악을 싫어하는 것 같은 영혼의 성향이다. 정약용은 감각적 욕구에서 나온 기호를 제어하지 못할 경우 악한 행위가 나타날 수 있고, 도덕적 욕구에서 나온 기호를 따를 경우 선한 행위가 나타난다고 보았다. 선한 행위를 하거나 악한 행위를 하는 것이 온전히 인간의 자유 의지에 달려 있으므로, 악한 행위를 한 사람에게 도덕적 책임을 물을 수 있다고 보았다. 정약용은 자유 의지로 선한 행위를 선택하고 이를 실천하는 것이 중요하다고 보았는데, 구체적인 실천 방법으로 '서(恕)'를 강조했다. 그는 '서'를 용서(容恕)와 추서(推恕)로 구분하고, 추서를 특히 강조했다. 용서(容恕)는 '타인의 악을 너그럽게 보아줌'을 의미하고, 추서(推恕)는 '내가 대접받고 싶은 대로 타인을 대우함'을 의미한다. 추서는 자신의 마음을 다른 사람의 처지에 미루어 생각해보는 것으로, 내가 어떤 대우를 받고 싶은지 먼저 생각해보고 그대로 타인을 대접하는 것이다. 용서는 타인의 악한 행위를 그냥 넘어가는 문제가 생길 수 있지만, 추서는 자신의 마음을 바탕으로 타인의 마음을 이해할 수 있으므로, 정약용은 추서에 따라 선한 행위를 실천해야 한다고 보았다.
}


\kfPassageNote{문장 독해 — 본문}{%
1. 정약용은 조선 후기의 실학자로, 인간의 본성에 대한 연구를 통해 인간의 선한 행위를 설명하려고 했다. 2. 그는 이전까지 절대적 권위를 가지고 있던 주희의 주자학을 비판하며 인간의 본성에 대한 자신만의 이론을 만들었다. 3. 주희는 인간의 본성을 '본연지성'과 '기질지성'으로 설명했다. 4. '본연지성'은 인간이 하늘로부터 받은 순수하고 선한 본성이고, '기질지성'은 본연지성에 사람마다 다른 기질이 더해진 것이다. 5. 그래서 주희는 인간의 기질이 맑으면 선한 행위를 하고 탁하면 악한 행위를 할 수 있다고 보았다. 6. 하지만 정약용은 주희의 관점을 비판했다. 7. 선한 행위와 악한 행위의 원인을 기질이라는 태어날 때부터 정해진 요인으로 본다면 행위에 인간의 의지가 개입되지 않는다. 8. 그렇다면 악한 행위를 한 사람에게 도덕적 책임을 물을 수 없다는 것이 정약용의 생각이었다. 9. 정약용은 인간의 본성을 '기호'라고 보았다. 10. 기호란 즐기고 좋아한다는 뜻으로, 생명이 있는 모든 존재는 각각의 기호를 본성으로 갖는다고 보았다. 11. 꿩은 산을 좋아하고 벼는 물을 좋아하는 것처럼, 인간도 어떤 성향을 갖는다는 것이다. 12. 정약용은 인간에게 '감각적 욕구에서 나온 기호'와 '도덕적 욕구에서 나온 기호'가 있다고 보았다. 13. 감각적 욕구에서 나온 기호는 맛있는 것을 좋아하는 것 같은 몸의 성향이고, 도덕적 욕구에서 나온 기호는 선을 좋아하거나 악을 싫어하는 것 같은 영혼의 성향이다. 14. 정약용은 감각적 욕구에서 나온 기호를 제어하지 못할 경우 악한 행위가 나타날 수 있고, 도덕적 욕구에서 나온 기호를 따를 경우 선한 행위가 나타난다고 보았다. 15. 선한 행위를 하거나 악한 행위를 하는 것이 온전히 인간의 자유 의지에 달려 있으므로, 악한 행위를 한 사람에게 도덕적 책임을 물을 수 있다고 보았다. 16. 정약용은 자유 의지로 선한 행위를 선택하고 이를 실천하는 것이 중요하다고 보았는데, 구체적인 실천 방법으로 '서(恕)'를 강조했다. 17. 그는 '서'를 용서(容恕)와 추서(推恕)로 구분하고, 추서를 특히 강조했다. 18. 용서(容恕)는 '타인의 악을 너그럽게 보아줌'을 의미하고, 추서(推恕)는 '내가 대접받고 싶은 대로 타인을 대우함'을 의미한다. 19. 추서는 자신의 마음을 다른 사람의 처지에 미루어 생각해보는 것으로, 내가 어떤 대우를 받고 싶은지 먼저 생각해보고 그대로 타인을 대접하는 것이다. 20. 용서는 타인의 악한 행위를 그냥 넘어가는 문제가 생길 수 있지만, 추서는 자신의 마음을 바탕으로 타인의 마음을 이해할 수 있으므로, 정약용은 추서에 따라 선한 행위를 실천해야 한다고 보았다.
}


\kfActivityBg \par\kfMaybeNeedspace{32\baselineskip}\noindent{\kfTipMark\,\,}{\small\color{kfInk} 다음 뜻풀이에 해당하는 어휘를 지문에서 찾아 쓰시오.}\par\nopagebreak[4]\smallskip\nopagebreak[4]
\begin{kfVocab}
\kfVocabItem{1}{조선 후기에 등장한, 실생활에 도움이 되는 학문을 연구한 학자}
\kfVocabItem{2}{사람이 본래부터 가지고 있는 성질이나 품성}
\kfVocabItem{3}{중국 송나라의 학자 주희가 집대성한 유학 사상}
\kfVocabItem{4}{주희가 말한, 하늘로부터 부여받은 순수하고 선한 본성}
\kfVocabItem{5}{주희가 말한, 본연지성에 사람마다 다른 기질이 더해진 본성}
\kfVocabItem{6}{사람마다 타고나는 마음이나 성질}
\kfVocabItem{7}{어떤 행위에 대해 마땅히 져야 할 책임}
\kfVocabItem{8}{정약용이 인간의 본성이라고 본, 즐기고 좋아하는 성향}
\kfVocabItem{9}{맛있는 것을 좋아하거나 편안함을 추구하는 몸의 욕구}
\kfVocabItem{10}{선을 좋아하고 악을 싫어하는 영혼의 욕구}
\kfVocabItem{11}{외부의 제약 없이 스스로의 의지에 따라 선택하는 능력}
\kfVocabItem{12}{정약용이 선한 행위의 실천 방법으로 강조한 것}
\kfVocabItem{13}{타인의 악한 행위를 너그럽게 보아주는 것}
\kfVocabItem{14}{내가 대접받고 싶은 대로 타인을 대우하는 것}
\end{kfVocab}

\kfActivityBg \par\kfMaybeNeedspace{32\baselineskip}\noindent{\kfTipMark\,\,}{\small\color{kfInk} 다음 글의 구조를 정리한 표의 빈칸을 채워 보세요.}\par\nopagebreak[4]\smallskip\nopagebreak[4]
\par\medskip\needspace{8\baselineskip}
\begin{tcolorbox}[enhanced, breakable=true,
  colback=kfPaper, colframe=kfRule, boxrule=0.4pt, arc=2pt,
  left=8pt, right=8pt, top=6pt, bottom=6pt,
  title={\footnotesize\bfseries\color{kfMute} 글의 구조도},
  coltitle=kfMute, colbacktitle=kfPrimaryLight!40,
  attach boxed title to top left={yshift=-1.5mm, xshift=6pt},
  boxed title style={colframe=kfRule, boxrule=0.3pt, arc=1pt}]
\footnotesize\setstretch{1.3}
\renewcommand{\arraystretch}{1.35}
\arrayrulecolor{kfRule}
\noindent\begin{tabularx}{\linewidth}{|>{\centering\arraybackslash}p{1.6cm}|>{\raggedright\arraybackslash}p{3.4cm}|>{\raggedright\arraybackslash}X|}
\hline
\rowcolor{kfPrimaryLight!40}\textbf{\color{kfPrimary}단} & \textbf{\color{kfPrimary}항목} & \textbf{\color{kfPrimary}내용} \\\hline
연구 동기: 인간의 ( ① ) 행위를 설명하고자 함 & 연구 동기: 인간의 ( ① ) 행위를 설명하고자 함 &  \\\hline
연구 방법: ( ② )을 비판하며 자신만의 이론 정립 & 연구 방법: ( ② )을 비판하며 자신만의 이론 정립 &  \\\hline
본연지성: 하늘로부터 받은 순수하고 ( ③ ) 본성 & 본연지성: 하늘로부터 받은 순수하고 ( ③ ) 본성 &  \\\hline
기질지성: 본연지성에 각기 다른 ( ④ )이 더해진 것 & 기질지성: 본연지성에 각기 다른 ( ④ )이 더해진 것 &  \\\hline
비판: 기질에 따라 선악이 결정된다면, 인간의 ( ⑤ )가 개입되지 않음 & 비판: 기질에 따라 선악이 결정된다면, 인간의 ( ⑤ )가 개입되지 않음 &  \\\hline
결론: 악한 행위에 ( ⑥ )을 물을 수 없음 & 결론: 악한 행위에 ( ⑥ )을 물을 수 없음 &  \\\hline
본성의 정의: 즐기고 좋아하는 성향, 즉 ‘( ⑦ )’ & 본성의 정의: 즐기고 좋아하는 성향, 즉 ‘( ⑦ )’ &  \\\hline
( ⑧ )에서 나온 기호 (몸의 성향) & ( ⑧ )에서 나온 기호 (몸의 성향) &  \\\hline
( ⑨ )에서 나온 기호 (영혼의 성향) & ( ⑨ )에서 나온 기호 (영혼의 성향) &  \\\hline
과정: 감각적 욕구 제어 실패 시 악, ( ⑩ ) 욕구 선택 시 선 & 과정: 감각적 욕구 제어 실패 시 악, ( ⑩ ) 욕구 선택 시 선 &  \\\hline
핵심: 선악의 선택은 온전히 인간의 ( ⑪ )에 달려 있음 & 핵심: 선악의 선택은 온전히 인간의 ( ⑪ )에 달려 있음 &  \\\hline
방법 제시: ( ⑫ )를 강조 & 방법 제시: ( ⑫ )를 강조 &  \\\hline
구분: 용서(容恕)와 ( ⑬ )(推恕)로 구분 & 구분: 용서(容恕)와 ( ⑬ )(推恕)로 구분 &  \\\hline
( ⑭ ): 타인의 악을 너그럽게 보아줌 & ( ⑭ ): 타인의 악을 너그럽게 보아줌 &  \\\hline
( ⑮ ): 내가 대접받고 싶은 대로 타인을 대우함 & ( ⑮ ): 내가 대접받고 싶은 대로 타인을 대우함 &  \\\hline
추서의 실천: ( ⑯ )의 마음을 미루어 ( ⑰ )의 처지를 생각함 & 추서의 실천: ( ⑯ )의 마음을 미루어 ( ⑰ )의 처지를 생각함 &  \\\hline
( ⑱ )의 문제점: 타인의 악한 행위를 그냥 넘어갈 수 있음 & ( ⑱ )의 문제점: 타인의 악한 행위를 그냥 넘어갈 수 있음 &  \\\hline
( ⑲ )의 장점: 타인을 이해하고 선한 행위를 실천하게 함 & ( ⑲ )의 장점: 타인을 이해하고 선한 행위를 실천하게 함 &  \\\hline
\end{tabularx}
\arrayrulecolor{black}
\end{tcolorbox}\par\medskip
\kfSecHeader{kfAreaNonlit}{문제}{}

\par\medskip
\kfBeginMC
\begin{kfQBlock}{01}{객관식}
\kfQStem{정약용이 인간의 선한 행위를 설명하기 위해 수행한 연구로 적절한 것은?}
\begin{kfChoices}
\kfChoice 서양의 과학 기술에 대한 연구
\kfChoice 농업 생산성 향상에 대한 연구
\kfChoice 법률 제도의 개혁에 대한 연구
\kfChoice 인간의 본성에 대한 연구
\end{kfChoices}
\end{kfQBlock}
\begin{kfQBlock}{02}{객관식}
\kfQStem{정약용이 비판했던 학문으로 적절한 것은?}
\begin{kfChoices}
\kfChoice 주희의 주자학
\kfChoice 왕양명의 양명학
\kfChoice 이황의 성리학
\kfChoice 이이의 성리학
\end{kfChoices}
\end{kfQBlock}
\begin{kfQBlock}{03}{객관식}
\kfQStem{정약용이 자신만의 이론을 정립한 분야로 적절한 것은?}
\begin{kfChoices}
\kfChoice 농업 기술
\kfChoice 상업 활동
\kfChoice 국방 강화
\kfChoice 인간의 본성
\end{kfChoices}
\end{kfQBlock}
\begin{kfQBlock}{04}{객관식}
\kfQStem{주희가 인간의 본성을 설명하기 위해 제시한 두 가지 개념으로 적절한 것은?}
\begin{kfChoices}
\kfChoice '본연지성'과 '기질지성'
\kfChoice '인심'과 '도심'
\kfChoice '사단'과 '칠정'
\kfChoice '이'와 '기'
\end{kfChoices}
\end{kfQBlock}
\begin{kfQBlock}{05}{객관식}
\kfQStem{주희의 이론에서 ‘본연지성’에 대한 설명으로 가장 적절한 것은?}
\begin{kfChoices}
\kfChoice 후천적인 학습과 경험을 통해 얻어진 본성
\kfChoice 상황이나 감정의 변화에 따라 수시로 변하는 본성
\kfChoice 인간이 하늘로부터 받은 순수하고 선한 본성
\kfChoice 자신이 속한 사회적 환경에 의해 형성된 본성
\end{kfChoices}
\end{kfQBlock}
\begin{kfQBlock}{06}{객관식}
\kfQStem{주희의 이론에서 ‘기질지성’이 형성되는 과정으로 적절한 것은?}
\begin{kfChoices}
\kfChoice 본연지성에 사람마다 다른 기질이 더해져서
\kfChoice 성장 과정에서 본래 가지고 있던 본연지성을 잃어버려서
\kfChoice 올바른 가치관을 심어주기 어려운 나쁜 환경에서 자라서
\kfChoice 선악을 구별할 수 있는 교육을 제대로 받지 못해서
\end{kfChoices}
\end{kfQBlock}
\begin{kfQBlock}{07}{객관식}
\kfQStem{주희에 따르면, 인간의 선악 행위를 결정하는 요인으로 적절한 것은?}
\begin{kfChoices}
\kfChoice 부모의 경제력
\kfChoice 사회적 지위
\kfChoice 기질이 맑고 탁함
\kfChoice 학문의 깊이
\end{kfChoices}
\end{kfQBlock}
\begin{kfQBlock}{08}{객관식}
\kfQStem{정약용은 주희의 관점에 대해 어떤 입장을 취했는가?}
\begin{kfChoices}
\kfChoice 동조하는 입장
\kfChoice 무관심한 입장
\kfChoice 절충하는 입장
\kfChoice 비판하는 입장
\end{kfChoices}
\end{kfQBlock}
\begin{kfQBlock}{09}{객관식}
\kfQStem{정약용에 따르면, 선악의 원인을 ‘기질’로 볼 때 발생하는 문제점으로 적절한 것은?}
\begin{kfChoices}
\kfChoice 인간의 운명이 신의 뜻을 거스르게 된다는 치명적인 문제
\kfChoice 행위에 인간의 의지가 개입되지 않는다는 문제
\kfChoice 기존의 사회적 질서를 무너뜨리고 큰 혼란을 야기할 수 있다는 문제
\kfChoice 사회 전체의 생산성을 떨어뜨려 경제적 이익을 저해하는 문제
\end{kfChoices}
\end{kfQBlock}
\begin{kfQBlock}{10}{객관식}
\kfQStem{정약용이 주희의 이론을 비판한 가장 큰 이유로 적절한 것은?}
\begin{kfChoices}
\kfChoice 타고난 기질은 인위적으로 바꿀 수 없어 교육이 불가능하기 때문이다
\kfChoice 도덕적인 문제는 강제성을 가진 법으로 처벌할 수 없기 때문이다
\kfChoice 선한 행위를 하려는 의도나 동기가 순수하지 않고 불순하기 때문이다
\kfChoice 악한 행위를 한 사람에게 도덕적 책임을 물을 수 없게 되기 때문이다
\end{kfChoices}
\end{kfQBlock}
\begin{kfQBlock}{11}{객관식}
\kfQStem{정약용이 정의한 인간의 본성으로 적절한 것은?}
\begin{kfChoices}
\kfChoice 이성
\kfChoice 감정
\kfChoice 기호
\kfChoice 본능
\end{kfChoices}
\end{kfQBlock}
\begin{kfQBlock}{12}{객관식}
\kfQStem{‘기호’의 의미로 가장 적절한 것은?}
\begin{kfChoices}
\kfChoice 판단하고 결정한다는 뜻
\kfChoice 즐기고 좋아한다는 뜻
\kfChoice 싫어하고 피한다는 뜻
\kfChoice 배우고 익힌다는 뜻
\end{kfChoices}
\end{kfQBlock}
\begin{kfQBlock}{13}{객관식}
\kfQStem{이 문장이 앞 문장의 내용을 설명하는 방식으로 가장 적절한 것은?}
\begin{kfChoices}
\kfChoice 추상적인 개념으로 정의하고 있다
\kfChoice 반대 사례를 들어 비판하고 있다
\kfChoice 전문가의 견해를 인용하고 있다
\kfChoice 구체적인 예시를 들어 설명하고 있다
\end{kfChoices}
\end{kfQBlock}
\begin{kfQBlock}{14}{객관식}
\kfQStem{정약용이 구분한 인간의 두 가지 기호로 적절한 것은?}
\begin{kfChoices}
\kfChoice 사고와 판단의 이성적 욕구와 감정적 반응의 감성적 욕구
\kfChoice 개인의 이익을 위한 욕구와 사회 공동체를 위한 욕구
\kfChoice 감각적 욕구에서 나온 기호와 도덕적 욕구에서 나온 기호
\kfChoice 풍요를 위한 물질적 욕구와 내적 평화를 위한 정신적 욕구
\end{kfChoices}
\end{kfQBlock}
\begin{kfQBlock}{15}{객관식}
\kfQStem{‘감각적 욕구에서 나온 기호’에 해당하는 성향으로 적절한 것은?}
\begin{kfChoices}
\kfChoice 감각적 욕구에서 나온 기호
\kfChoice 도덕적 욕구에서 나온 기호
\kfChoice 이성적 판단에서 나온 기호
\kfChoice 사회적 관습에서 나온 기호
\end{kfChoices}
\end{kfQBlock}
\begin{kfQBlock}{16}{객관식}
\kfQStem{‘도덕적 욕구에서 나온 기호’에 해당하는 성향으로 적절한 것은?}
\begin{kfChoices}
\kfChoice 감각적 욕구에서 나온 기호
\kfChoice 본능적 욕구에서 나온 기호
\kfChoice 생존 욕구에서 나온 기호
\kfChoice 도덕적 욕구에서 나온 기호
\end{kfChoices}
\end{kfQBlock}
\begin{kfQBlock}{17}{객관식}
\kfQStem{정약용에 따르면 악한 행위가 나타나는 경우로 적절한 것은?}
\begin{kfChoices}
\kfChoice 감각적 욕구에서 나온 기호를 제어하지 못할 경우
\kfChoice 영혼의 성향인 도덕적 욕구에서 나온 기호를 따를 경우
\kfChoice 감정에 치우치지 않고 냉철하게 이성적인 판단을 할 경우
\kfChoice 자신의 고집을 버리고 타인의 진심 어린 충고를 들을 경우
\end{kfChoices}
\end{kfQBlock}
\begin{kfQBlock}{18}{객관식}
\kfQStem{정약용에 따르면 선한 행위가 나타나는 경우로 적절한 것은?}
\begin{kfChoices}
\kfChoice 육체의 편안함을 추구하는 감각적 욕구에서 나온 기호를 따를 경우
\kfChoice 선을 지향하는 도덕적 욕구에서 나온 기호를 따를 경우
\kfChoice 타인의 권리보다 자신의 이익을 최우선으로 고려할 경우
\kfChoice 자신의 소신보다는 타인의 시선을 지나치게 의식할 경우
\end{kfChoices}
\end{kfQBlock}
\begin{kfQBlock}{19}{객관식}
\kfQStem{정약용에 따르면 선악 행위의 선택은 궁극적으로 무엇에 달려 있는가?}
\begin{kfChoices}
\kfChoice 타고난 기질
\kfChoice 신의 섭리
\kfChoice 인간의 자유 의지
\kfChoice 사회적 환경
\end{kfChoices}
\end{kfQBlock}
\begin{kfQBlock}{20}{객관식}
\kfQStem{정약용이 악한 행위에 도덕적 책임을 물을 수 있다고 본 근거로 적절한 것은?}
\begin{kfChoices}
\kfChoice 선악의 선택이 온전히 인간의 자유 의지에 달려 있기 때문에
\kfChoice 타고난 기질은 후천적인 노력으로 얼마든지 변할 수 있기 때문에
\kfChoice 악한 행위에 대한 처벌 규정이 법률로 명확히 정해져 있기 때문에
\kfChoice 아무리 악한 사람이라도 교육을 통해 충분히 교화될 수 있기 때문에
\end{kfChoices}
\end{kfQBlock}
\begin{kfQBlock}{21}{객관식}
\kfQStem{정약용이 강조한 선한 행위의 구체적인 실천 방법으로 적절한 것은?}
\begin{kfChoices}
\kfChoice 효(孝)
\kfChoice 충(忠)
\kfChoice 예(禮)
\kfChoice 서(恕)
\end{kfChoices}
\end{kfQBlock}
\begin{kfQBlock}{22}{객관식}
\kfQStem{정약용이 구분한 ‘서(恕)’의 두 가지 종류로 적절한 것은?}
\begin{kfChoices}
\kfChoice 인과 의
\kfChoice 용서와 추서
\kfChoice 지와 행
\kfChoice 리와 기
\end{kfChoices}
\end{kfQBlock}
\begin{kfQBlock}{23}{객관식}
\kfQStem{정약용이 ‘용서’와 ‘추서’ 중 특히 강조한 것은?}
\begin{kfChoices}
\kfChoice 추서
\kfChoice 용서
\kfChoice 인
\kfChoice 의
\end{kfChoices}
\end{kfQBlock}
\begin{kfQBlock}{24}{객관식}
\kfQStem{‘용서(容恕)’의 의미로 가장 적절한 것은?}
\begin{kfChoices}
\kfChoice 타인의 악을 엄격하게 처벌함
\kfChoice 타인의 악을 무시함
\kfChoice 타인의 악을 널리 알림
\kfChoice 타인의 악을 너그럽게 보아줌
\end{kfChoices}
\end{kfQBlock}
\begin{kfQBlock}{25}{객관식}
\kfQStem{‘추서(推恕)’의 의미로 가장 적절한 것은?}
\begin{kfChoices}
\kfChoice 법에 따라 공정하게 타인을 대우함
\kfChoice 내가 대접받고 싶은 대로 타인을 대우함
\kfChoice 이익이 되는 방향으로 타인을 대우함
\kfChoice 감정에 따라 즉흥적으로 타인을 대우함
\end{kfChoices}
\end{kfQBlock}
\begin{kfQBlock}{26}{객관식}
\kfQStem{‘추서’는 누구의 마음을 미루어 생각하는 것인가?}
\begin{kfChoices}
\kfChoice 인간의 모든 길흉화복을 주관하는 신의 뜻을 따르는 것
\kfChoice 공동체의 질서 유지를 위해 사회적 규칙을 준수하는 것
\kfChoice 자신의 마음을 다른 사람의 처지에 미루어 생각하는 것
\kfChoice 자기 자신의 주관보다는 다수의 객관적인 의견을 따르는 것
\end{kfChoices}
\end{kfQBlock}
\begin{kfQBlock}{27}{객관식}
\kfQStem{‘추서’를 실천하기 위해 가장 먼저 해야 할 일로 적절한 것은?}
\begin{kfChoices}
\kfChoice 내가 어떤 대우를 받고 싶은지 먼저 생각해보는 것
\kfChoice 상대방이 무엇을 원하는지 물어보는 것
\kfChoice 이 행동이 나에게 이익이 되는지 계산하는 것
\kfChoice 법률적으로 문제가 없는지 확인하는 것
\end{kfChoices}
\end{kfQBlock}
\begin{kfQBlock}{28}{객관식}
\kfQStem{‘용서(容恕)’가 가질 수 있는 문제점으로 적절한 것은?}
\begin{kfChoices}
\kfChoice 악한 사람에게 이용당하여 자신이 큰 손해를 볼 수 있다
\kfChoice 타인에게 자신의 선의가 부당하게 이용당할 수 있다
\kfChoice 도덕적 권고일 뿐 법적 효력이 없어 강제할 수 없다
\kfChoice 타인의 악한 행위를 그냥 넘어가는 문제가 생길 수 있다
\end{kfChoices}
\end{kfQBlock}
\begin{kfQBlock}{29}{객관식}
\kfQStem{정약용이 ‘용서’보다 ‘추서’를 더 강조한 이유로 적절한 것은?}
\begin{kfChoices}
\kfChoice 추서가 용서보다 실천하기가 훨씬 더 쉽고 간편하기 때문이다
\kfChoice 추서는 자신의 마음을 바탕으로 타인의 마음을 이해할 수 있기 때문이다
\kfChoice 추서가 예로부터 전해 내려오는 전통적인 도덕 실천 방법이기 때문이다
\kfChoice 추서가 개인의 도덕성뿐만 아니라 법률적인 강제성이 있기 때문이다
\end{kfChoices}
\end{kfQBlock}
\begin{kfQBlock}{30}{객관식}
\kfQStem{‘추서’는 무엇을 바탕으로 타인의 마음을 이해하는 것인가?}
\begin{kfChoices}
\kfChoice 사회적 통념
\kfChoice 타인의 조언
\kfChoice 객관적 사실
\kfChoice 자신의 마음
\end{kfChoices}
\end{kfQBlock}
\begin{kfQBlock}{31}{객관식}
\kfQStem{정약용은 무엇에 따라 선한 행위를 실천해야 한다고 보았는가?}
\begin{kfChoices}
\kfChoice 용서
\kfChoice 추서
\kfChoice 기질
\kfChoice 본능
\end{kfChoices}
\end{kfQBlock}
\begin{kfQBlock}{32}{객관식}
\kfQStem{윗글의 내용과 일치하지 \uline{않는} 것은?}
\begin{kfChoices}
\kfChoice 주희는 인간의 본성을 두 가지로 나누어 설명했다.
\kfChoice 정약용은 인간의 행위가 자유 의지에 따른 선택의 결과라고 보았다.
\kfChoice 정약용은 감각적 욕구를 따르는 것을 항상 악한 행위라고 규정했다.
\kfChoice 정약용은 선한 행위를 실천하는 구체적인 방법으로 '추서'를 제시했다.
\kfChoice 정약용은 주희의 이론이 도덕적 책임의 문제를 해결할 수 없다고 보았다.
\end{kfChoices}
\end{kfQBlock}
\begin{kfQBlock}{33}{객관식}
\kfQStem{주희와 정약용의 견해를 비교한 것으로 적절하지 \uline{\underline{않은}} 것은?}
\begin{kfChoices}
\kfChoice 본성의 정의: 주희는 '성(性)', 정약용은 '기호(嗜好)'로 보았다.
\kfChoice 선악의 근원: 주희는 '기질'의 맑고 탁함, 정약용은 '자유 의지'의 선택으로 보았다.
\kfChoice 도덕적 책임: 주희는 묻기 어렵다고, 정약용은 물을 수 있다고 보았다.
\kfChoice 인간의 의지: 주희는 중시하지 않았고, 정약용은 핵심적인 역할로 보았다.
\kfChoice 선한 본성: 주희는 긍정했고, 정약용은 선한 본성의 존재 자체를 부정했다.
\end{kfChoices}
\end{kfQBlock}
\begin{kfQBlock}{34}{객관식}
\kfQStem{㉠추서(推恕)에 대한 설명으로 적절하지 \uline{\underline{않은}} 것은?}
\begin{kfChoices}
\kfChoice 내 마음을 기준으로 타인에게 어떻게 행동할지를 정하는 적극적 실천이다.
\kfChoice ‘내가 대접받고 싶은 대로 남을 대하라’는 적극적 도덕률을 담는다.
\kfChoice 정약용 사상에서 강조되는 도덕 실천 방법이다.
\kfChoice 타인의 입장을 생각해 행동하는 윤리적 태도이다.
\kfChoice 사회적 규칙이나 법률을 통해 강제적으로 실현되어야 하는 외부적 의무이다.
\end{kfChoices}
\end{kfQBlock}
\begin{kfQBlock}{35}{객관식}
\kfQStem{<보기>의 스승이 제자에게 강조할 정약용의 사상으로 가장 적절한 것은?}
\begin{kfEvidence}한 제자가 스승에게 물었다. "사람들은 왜 착하게 살지 못합니까?"   \\\relax
   스승이 답했다. "타고난 기질 때문이라고 핑계를 대서는 안 된다. 맛있는 음식을 먹고 싶은 마음과 어려운 사람을 돕고 싶은 마음 사이에서 무엇을 따를지 결정하는 것은 너 자신에게 달려있기 때문이다."\end{kfEvidence}
\begin{kfChoices}
\kfChoice 인간의 본성은 하늘이 부여한 '본연지성'이다.
\kfChoice 선한 행위와 악한 행위는 모두 '기호'에 따른 것이다.
\kfChoice 악한 행위에 대해서는 너그러운 '용서'의 자세가 필요하다.
\kfChoice 선악의 문제는 인간의 '자유 의지'에 따른 선택의 문제이다.
\kfChoice 사람마다 '기질'이 다르므로 도덕적 수준도 다를 수밖에 없다.
\end{kfChoices}
\end{kfQBlock}
\begin{kfQBlock}{36}{객관식}
\kfQStem{윗글의 중심 사상에 대한 설명으로 적절하지 \uline{\underline{않은}} 것은?}
\begin{kfChoices}
\kfChoice 인간에게는 본래 선을 좋아하고 악을 싫어하는 영혼의 ‘도덕적 기호’가 갖추어져 있다.
\kfChoice 그와 동시에 욕구가 통제되지 않으면 악으로 이어질 수 있는 ‘감각적 기호’도 함께 있다.
\kfChoice 두 기호 사이에서 어느 쪽 욕구를 따를지는 인간 스스로의 ‘자유 의지’에 달려 있다.
\kfChoice 자신이 한 그 선택에 대해 스스로 ‘도덕적 책임’을 져야 한다고 강조한다.
\kfChoice 인간에게는 선을 좋아하는 마음이 없으므로 외부의 가르침과 규율을 통해서만 도덕을 배울 수 있다.
\end{kfChoices}
\end{kfQBlock}
\begin{kfQBlock}{37}{단답형}
\kfQStem{주희가 설명한, 하늘로부터 받은 순수하고 선한 본성은 무엇인가?}
\kfAnswerLines{1}
\end{kfQBlock}
\begin{kfQBlock}{38}{단답형}
\kfQStem{정약용이 주희의 기질론을 비판한 이유는 악한 행위를 한 사람에게 무엇을 물을 수 없다고 보았기 때문인가?}
\kfAnswerLines{1}
\end{kfQBlock}
\begin{kfQBlock}{39}{단답형}
\kfQStem{정약용이 인간의 본성이라고 주장한, '즐기고 좋아하는 성향'을 뜻하는 개념은 무엇인가?}
\kfAnswerLines{1}
\end{kfQBlock}
\begin{kfQBlock}{40}{단답형}
\kfQStem{정약용이 제시한 두 가지 기호를 쓰시오.}
\kfAnswerLines{1}
\end{kfQBlock}
\begin{kfQBlock}{41}{단답형}
\kfQStem{선을 좋아하고 악을 싫어하는 것과 같은 영혼의 성향은 어떤 기호에 해당하는가?}
\kfAnswerLines{1}
\end{kfQBlock}
\begin{kfQBlock}{42}{서술형}
\kfQStem{정약용이 주희의 '기질지성' 이론을 비판한 핵심적인 이유를 <조건>에 맞게 서술하시오.}
\kfAnswerNote{답안}{4}
\end{kfQBlock}
\kfEndMC



\clearpage

\kfChapterCover{2}{러셀3 (챕터2)}{Chapter 2}{이번 챕터의 학습 내용을 시작합니다.}
\begin{kfChapterAreaList}
\kfChapterAreaItem{문법}{kfAreaGrammar}{단어의 구조를 분석하는 틀: 형태소, 어근, 접사}
\kfChapterAreaItem{개념}{kfAreaConcept}{판소리에서 태어난 소설}
\kfChapterAreaItem{문학}{kfAreaLiterature}{「심청전」}
\kfChapterAreaItem{비문학}{kfAreaNonlit}{[사회] 국가가 처벌하는 형사법 vs 개인이 다투는 민사법}
\end{kfChapterAreaList}

\kfStartGrammar{단어의 구조를 분석하는 틀: 형태소, 어근, 접사}


\kfPassageNoteNT{%
　우리가 사용하는 '책가방' 같은 단어는 '책'과 '가방'이라는 더 작은 단위로 쪼갤 수 있다. 하지만 '책'을 다시 'ㅊ, ㅐ, ㄱ'으로 쪼개면 '글을 읽는 물건'이라는 의미가 사라지고 만다. 이처럼 뜻을 가진 채로 더는 쪼갤 수 없는 가장 작은 말의 단위를 형태소(形態素)라고 한다. 형태소는 단어의 구조를 분석하는 기본 단위이다.

　형태소는 두 가지 기준을 동시에 적용하여 분류할 수 있다. 첫째, 문장에서 홀로 쓰일 수 있는지에 따라 자립 형태소와 의존 형태소로 나뉜다. '하늘, 책, 매우'처럼 혼자 쓸 수 있으면 자립 형태소이고, '가다'의 '가-'나 '-다', 조사 '-이/가'처럼 반드시 다른 말에 기대어야만 쓰이면 의존 형태소이다.

　둘째, 실질적인 의미를 가졌는지에 따라 실질 형태소와 형식 형태소로 나뉜다. '하늘, 책, 가-'처럼 구체적인 대상이나 동작, 상태를 나타내면 실질 형태소이고, '-다, -이/가'처럼 문법적인 관계를 표시하거나 보조적인 의미만 더하면 형식 형태소이다. 형식 형태소에는 조사, 어미, 접사가 속한다.

　이 두 기준을 적용할 때, 학생들이 가장 혼동하기 쉬운 부분이 바로 '가다'의 '가-'와 같은 용언의 어간이다. 용언의 어간은 '가다'라는 동작의 실질적 의미를 담고 있으므로 실질 형태소에 속하지만, 결코 혼자 쓰일 수 없고 반드시 '-다'와 같은 어미에 기대어야 하므로 의존 형태소에 속한다. 이처럼 용언의 어간은 국어에서 유일하게 '의존 형태소이면서 동시에 실질 형태소'라는 독특한 지위를 가진다.

　이러한 형태소 개념은 단어의 내부 구조를 분석하는 두 가지 관점으로 이어진다. 첫 번째는 단어의 '형성'을 분석하는 어근(語根)과 접사(接辭)의 개념이다. 어근은 단어에서 실질적인 의미를 나타내는 핵심 부분이다. '덮개'의 '덮-', '맨손'의 '손'이 모두 어근이며 실질 형태소이다. 접사는 이 어근에 붙어 뜻을 더하거나 제한하는 부분이다. '덮개'의 '-개'처럼 어근 뒤에 붙으면 접미사, '맨손'의 '맨-'처럼 어근 앞에 붙으면 접두사라고 한다. 접사는 모두 의존 형태소이면서 형식 형태소이다.

　두 번째는 용언의 '활용'을 분석하는 어간(語幹)과 어미(語尾)의 개념이다. 어간은 '먹다, 먹고, 먹으니'에서 변하지 않는 부분인 '먹-'을 말하고, 어미는 '-다, -고, -으니'처럼 뒤에 붙어 시제나 문장 연결 등 문법적 기능을 하며 바뀌는 부분을 말한다. 어간은 실질 형태소이고, 어미는 형식 형태소이다.

　그렇다면 '어근'과 '어간'은 어떻게 구별할까? '먹다'는 어근도 '먹-'이고 어간도 '먹-'이라서 같아 보이지만, '치솟다'라는 단어를 보면 차이가 명확해진다. '치솟다'의 실질적 의미 중심인 어근은 '솟-' 하나이다. 하지만 이 단어가 활용할 때는 '치솟고, 치솟으니'처럼 변하므로, 변하지 않는 부분인 어간은 '치솟-' 전체가 된다. 이처럼 어근과 접사는 새로운 단어를 만드는 '형성'의 재료이고, 어간과 어미는 완성된 용언이 문장에서 쓰이는 '활용'의 도구라고 이해해야 한다.
\par\medskip\begin{itemize}[leftmargin=18pt, itemsep=2pt, topsep=2pt, label=\textbullet]
\item 용언 어간의 특수성 — '가-'(실질+의존)
\item 어근·접사(형성) — '덮개=덮-(어근)+-개(접사)'
\item 어근 vs 어간 — '치솟다': 어근='솟-' / 어간='치솟-'
\end{itemize}
}

\kfActivityBg \par\kfMaybeNeedspace{32\baselineskip}지문의 핵심 내용을 떠올리며 아래 빈칸에 알맞은 말을 채워 보세요.

1. 형태소의 개념과 분류 

\par\medskip\needspace{4\baselineskip}
\noindent\begin{adjustbox}{max width=\linewidth, center}
\renewcommand{\arraystretch}{1.45}
\arrayrulecolor{kfRule}
\begin{tabular}{|>{\raggedright\arraybackslash}p{\dimexpr 0.1185\linewidth-12pt\relax}|>{\raggedright\arraybackslash}p{\dimexpr 0.8815\linewidth-12pt\relax}|}
\hline
\rowcolor{kfPrimaryLight!50}\textbf{\color{kfPrimary} 구분} & \textbf{\color{kfPrimary} 내용} \\\hline
형태소 & ∙ [ ① ]을/를 가진 채로 더는 쪼갤 수 없는 가장 작은 말의 단위. \\\hline
분류 기준 1: [ ② ] & ∙ 자립 형태소: 문장에서 [ ③ ] 쓸 수 있음. (예: 하늘, 책)<br>∙ 의존 형태소: 반드시 다른 말에 [ ④ ] 쓰임. (예: -다, -이/가, 가-) \\\hline
분류 기준 2: [ ⑤ ] & ∙ [ ⑥ ]: 구체적인 대상, 동작, 상태 등 실질적 의미를 가짐. (예: 하늘, 가-)<br>∙ [ ⑦ ]: 문법적인 관계를 표시하거나 보조적인 의미만 더함. \\\hline
[ ⑦ ]의 종류 & ∙ 문법적 관계 표시: [ ⑧ ] (예: -이/가)<br>∙ 활용 시 변하는 부분: [ ⑨ ] (예: -다, -고)<br>∙ 어근에 붙어 뜻을 더함: [ ⑩ ] (예: 맨-, -개) \\\hline
\end{tabular}
\arrayrulecolor{black}
\end{adjustbox}\par\medskip



2. 형태소 2x2 분류  

\par\medskip\needspace{4\baselineskip}
\noindent\begin{adjustbox}{max width=\linewidth, center}
\renewcommand{\arraystretch}{1.45}
\arrayrulecolor{kfRule}
\begin{tabular}{|>{\raggedright\arraybackslash}p{\dimexpr 0.1212\linewidth-12pt\relax}|>{\raggedright\arraybackslash}p{\dimexpr 0.3918\linewidth-12pt\relax}|>{\raggedright\arraybackslash}p{\dimexpr 0.4870\linewidth-12pt\relax}|}
\hline
\rowcolor{kfPrimaryLight!50}\textbf{\color{kfPrimary} (기준)} & \textbf{\color{kfPrimary} 실질 형태소 (뜻 O)} & \textbf{\color{kfPrimary} 형식 형태소 (뜻 X)} \\\hline
자립 형태소 (홀로 O) & ∙ 명사, 대명사, 수사, 관형사, 부사, 감탄사<br>∙ (예: 하늘, 책, 매우) & ∙ (없음) \\\hline
의존 형태소 (홀로 X) & ∙ [ ⑪ ]<br>∙ (예: '먹다'의 '먹-') & ∙ [ ⑫ ], [ ⑬ ], [ ⑭ ]<br>∙ (예: -이/가, -다, -개) \\\hline
\end{tabular}
\arrayrulecolor{black}
\end{adjustbox}\par\medskip



 핵심 정리: 국어에서 유일하게  '의존 형태소이면서 동시에 실질 형태소' 인 것은 바로  [ ⑮ ] 이다.

3. 단어의 구조 분석 

\par\medskip\needspace{4\baselineskip}
\noindent\begin{adjustbox}{max width=\linewidth, center}
\renewcommand{\arraystretch}{1.45}
\arrayrulecolor{kfRule}
\begin{tabular}{|>{\raggedright\arraybackslash}p{\dimexpr 0.0582\linewidth-12pt\relax}|>{\raggedright\arraybackslash}p{\dimexpr 0.0582\linewidth-12pt\relax}|>{\raggedright\arraybackslash}p{\dimexpr 0.4931\linewidth-12pt\relax}|>{\raggedright\arraybackslash}p{\dimexpr 0.3905\linewidth-12pt\relax}|}
\hline
\rowcolor{kfPrimaryLight!50}\textbf{\color{kfPrimary} 분석 관점} & \textbf{\color{kfPrimary} 핵심 개념} & \textbf{\color{kfPrimary} 설명} & \textbf{\color{kfPrimary} 예시} \\\hline
\multirow{2}{*}{단어 형성} & [ ⑯ ] & ∙ 단어에서 실질적인 의미를 나타내는 핵심 부분 & ∙ '덮개'의 '덮-'<br>∙ '맨손'의 '손' \\\cline{2-4}
 & [ ⑰ ] & ∙ 어근에 붙어 뜻을 더하거나 제한하는 부분<br>-[ ⑱ ] : 어근 앞에 붙음<br>-[ ⑲ ] : 어근 뒤에 붙음 & ∙ '덮개'의 '-개'<br>∙ '맨손'의 '맨-' \\\cline{2-4}
\multirow{2}{*}{용언 활용} & [ ⑳ ] & ∙ 용언이 활용할 때 변하지 않는 부분 & ∙ '먹다, 먹고'의 '먹-' \\\cline{2-4}
 & [ ㉑ ] & ∙ 용언이 활용할 때 변하는 부분 & ∙ '먹다, 먹고'의 '-다, -고' \\\cline{2-4}
\hline
\end{tabular}
\arrayrulecolor{black}
\end{adjustbox}\par\medskip



4. 어근 vs 어간 심화 비교 

\par\medskip\needspace{4\baselineskip}
\noindent\begin{adjustbox}{max width=\linewidth, center}
\renewcommand{\arraystretch}{1.45}
\arrayrulecolor{kfRule}
\begin{tabular}{|>{\raggedright\arraybackslash}p{\dimexpr 0.0636\linewidth-12pt\relax}|>{\raggedright\arraybackslash}p{\dimexpr 0.9364\linewidth-12pt\relax}|}
\hline
\rowcolor{kfPrimaryLight!50}\textbf{\color{kfPrimary} 단어} & \textbf{\color{kfPrimary} 분석} \\\hline
먹다 & ∙ 어근 (실질적 핵심): [ ㉒ ]<br>∙ 어간 (활용 시 불변): [ ㉓ ] \\\hline
치솟다 & ∙ 어근 (실질적 핵심): [ ㉔ ]<br>∙ 어간 (활용 시 불변: '치솟고, 치솟으니'): [ ㉕ ] \\\hline
\end{tabular}
\arrayrulecolor{black}
\end{adjustbox}\par\medskip



---\nopagebreak[4]\par\nopagebreak[4]

\kfActivityBg \kfSecHeader{kfAreaGrammar}{활동}{분석 훈련}
\par\kfMaybeNeedspace{24\baselineskip}\noindent{\kfTipMark\,\,}{\small\color{kfInk} 다음 활용 사례를 분석하여 빈칸을 채워 보세요.}\par\nopagebreak[4]\smallskip\nopagebreak[4]
\begin{enumerate}[leftmargin=22pt, itemsep=6pt, topsep=4pt, label=\textbf{\arabic*.}]
\item 다음 문장을 형태소 단위로 모두 쪼개어 보시오.
\begin{kfEvidence}우리는 맨손으로 잡았다.\end{kfEvidence}
\item 1번 문장의 형태소 중 '자립 형태소'를 모두 찾아 쓰시오.
\item 1번 문장의 형태소 중 '실질 형태소'를 모두 찾아 쓰시오.
\item 1번 문장의 형태소 중 '형식 형태소'를 모두 찾아 쓰시오.
\item '먹었다'의 어간 '먹-'은 (자립/의존) 형태소이면서 (실질/형식) 형태소입니다. \kfFillBlank[12mm]{} 안에 알맞은 말에 O 표시 하시오.
\item 다음 단어를 '어근'과 '접사'로 나누어 쓰시오.
\begin{kfEvidence}덮개\end{kfEvidence}
\item 다음 단어를 '어근'과 '접사'로 나누어 쓰시오.
\begin{kfEvidence}맨손\end{kfEvidence}
\item 다음 단어를 '어간'과 '어미'로 나누어 쓰시오.
\begin{kfEvidence}예쁘고\end{kfEvidence}
\item 다음 단어의 '어근'을 찾아 쓰시오.
\begin{kfEvidence}지우개\end{kfEvidence}
\item 다음 단어의 '어근'과 '어간'을 각각 찾아 쓰시오.
\begin{kfEvidence}치솟았다\end{kfEvidence}
\end{enumerate}

\kfSecHeader{kfAreaGrammar}{문제}{}

\par\medskip
\kfBeginMC
\begin{kfQBlock}{01}{객관식}
\kfQStem{형태소에 대한 설명으로 적절하지 \uline{\underline{않은}} 것은?}
\begin{kfChoices}
\kfChoice 뜻을 가진 가장 작은 말의 단위이다.
\kfChoice 자립성 여부에 따라 자립 형태소와 의존 형태소로 나뉜다.
\kfChoice 모든 실질 형태소는 자립 형태소에 속한다.
\kfChoice 실질적 의미 여부에 따라 실질 형태소와 형식 형태소로 나뉜다.
\kfChoice 형식 형태소에는 조사, 어미, 접사가 있다.
\end{kfChoices}
\end{kfQBlock}
\begin{kfQBlock}{02}{객관식}
\kfQStem{다음 문장을 형태소로 분석한 개수가 \uline{잘못된} 것은?}
\begin{kfChoices}
\kfChoice 하늘이 맑다. (하늘 / 이 / 맑- / -다) - 4개
\kfChoice 꽃이 피었다. (꽃 / 이 / 피- / -었- / -다) - 5개
\kfChoice 밥을 먹는다. (밥 / 을 / 먹- / -는- / 다) - 5개
\kfChoice 새 책을 샀다. (새 / 책 / 을 / 샀다) - 4개
\kfChoice 아주 높이 난다. (아주 / 높- / -이 / 날- / -ㄴ- / 다) - 6개
\end{kfChoices}
\end{kfQBlock}
\begin{kfQBlock}{03}{객관식}
\kfQStem{형태소의 분류가 바르게 짝지어진 것은?}
\begin{kfEvidence}하늘, -이, 맑-, -다\end{kfEvidence}
\begin{kfChoices}
\kfChoice 자립 형태소: 하늘, 맑-
\kfChoice 의존 형태소: -이, -다
\kfChoice 실질 형태소: 하늘, -이, 맑-
\kfChoice 형식 형태소: -다
\kfChoice 의존 형태소이면서 실질 형태소: -다
\end{kfChoices}
\end{kfQBlock}
\begin{kfQBlock}{04}{객관식}
\kfQStem{<보기>의 ㉠\textasciitilde{}㉣에 대한 설명으로 적절하지 \uline{\underline{않은}} 것은?}
\begin{kfEvidence}㉠ 하늘 \\\relax
㉡ 이 \\\relax
㉢ 푸르 \\\relax
㉣ 다\end{kfEvidence}
\begin{kfChoices}
\kfChoice ㉠: 자립 형태소이면서 실질 형태소이다.
\kfChoice ㉡: 의존 형태소이면서 형식 형태소이다.
\kfChoice ㉢: 의존 형태소이면서 실질 형태소이다.
\kfChoice ㉣: 자립 형태소이면서 형식 형태소이다.
\kfChoice ㉢: 용언의 어간이다.
\end{kfChoices}
\end{kfQBlock}
\begin{kfQBlock}{05}{객관식}
\kfQStem{<보기>의 두 기준을 모두 만족하는 형태소는?}
\begin{kfEvidence}∙ 홀로 쓰일 수 없는 의존 형태소이다.   \\\relax
   ∙ 실질적인 의미를 가진 실질 형태소이다.\end{kfEvidence}
\begin{kfChoices}
\kfChoice 높-
\kfChoice 하늘
\kfChoice -이
\kfChoice -다
\kfChoice 헌
\end{kfChoices}
\end{kfQBlock}
\begin{kfQBlock}{06}{객관식}
\kfQStem{다음 밑줄 친 부분 중 성격이 나머지와 \uline{다른} 하나는?}
\begin{kfChoices}
\kfChoice \uline{새} 신
\kfChoice \uline{맨}주먹
\kfChoice \uline{풋}사랑
\kfChoice \uline{군}소리
\kfChoice \uline{참}기름
\end{kfChoices}
\end{kfQBlock}
\begin{kfQBlock}{07}{객관식}
\kfQStem{‘어근’과 ‘접사’에 대한 설명으로 적절하지 \uline{\underline{않은}} 것은?}
\begin{kfChoices}
\kfChoice 어근은 단어의 실질적 의미를 가지는 핵심 부분이다.
\kfChoice 접사는 어근의 앞이나 뒤에 붙어 의미를 더하거나 품사를 바꾼다.
\kfChoice 접사는 형식 형태소, 어근은 실질 형태소에 해당한다.
\kfChoice 어근과 접사는 단어의 ‘형성’을 분석하는 개념이다.
\kfChoice 접사는 어근의 뒤에만 붙으며, 앞에 붙는 경우는 없다.
\end{kfChoices}
\end{kfQBlock}
\begin{kfQBlock}{08}{객관식}
\kfQStem{다음 중 '어근'이 \uline{\underline{아닌}} 것은?}
\begin{kfChoices}
\kfChoice '덮개'의 '덮-'
\kfChoice '맨손'의 '손'
\kfChoice '치솟다'의 '솟-'
\kfChoice '짓밟다'의 '밟-'
\kfChoice '신발'의 '신-'
\end{kfChoices}
\end{kfQBlock}
\begin{kfQBlock}{09}{객관식}
\kfQStem{다음 중 '어간'과 '어미'의 구분이 \uline{불가능한} 단어는?}
\begin{kfChoices}
\kfChoice 맑다
\kfChoice 높이
\kfChoice 덮다
\kfChoice 웃고
\kfChoice 잡히다
\end{kfChoices}
\end{kfQBlock}
\begin{kfQBlock}{10}{객관식}
\kfQStem{다음 단어의 '어근'과 '어간'을 바르게 짝지은 것은?}
\begin{kfEvidence}짓밟히다\end{kfEvidence}
\begin{kfChoices}
\kfChoice 어근: 짓밟- / 어간: 짓밟히-
\kfChoice 어근: 밟- / 어간: 짓밟-
\kfChoice 어근: 밟- / 어간: 짓밟히-
\kfChoice 어근: 짓밟히- / 어간: 짓밟히-
\kfChoice 어근: 짓- / 어간: 밟-
\end{kfChoices}
\end{kfQBlock}
\begin{kfQBlock}{11}{객관식}
\kfQStem{‘어근’과 ‘어간’에 대한 설명으로 적절하지 \uline{\underline{않은}} 것은?}
\begin{kfChoices}
\kfChoice 어근은 단어 의미의 핵심을 가지는 부분이다.
\kfChoice 어간은 용언이 활용할 때 변하지 않는 부분이다.
\kfChoice 어근은 단어 ‘형성’과, 어간은 용언 ‘활용’과 관련된 개념이다.
\kfChoice 단어에 따라 어근과 어간이 일치하기도 하고 다르기도 하다.
\kfChoice 어근은 단어의 ‘활용’과, 어간은 단어의 ‘형성’과 관련된 개념이다.
\end{kfChoices}
\end{kfQBlock}
\begin{kfQBlock}{12}{객관식}
\kfQStem{다음 중 밑줄 친 부분이 '접사'가 \uline{\underline{아닌}} 것은?}
\begin{kfChoices}
\kfChoice \uline{새}파랗다
\kfChoice 잠\uline{꾸러기}
\kfChoice \uline{맨}발
\kfChoice \uline{새} 신발
\kfChoice 덮\uline{개}
\end{kfChoices}
\end{kfQBlock}
\begin{kfQBlock}{13}{객관식}
\kfQStem{다음 단어 중 어근이 둘 이상 결합한 단어는?}
\begin{kfChoices}
\kfChoice 덮개
\kfChoice 맨손
\kfChoice 짓밟다
\kfChoice 책가방
\kfChoice 치솟다
\end{kfChoices}
\end{kfQBlock}
\begin{kfQBlock}{14}{객관식}
\kfQStem{<보기>의 ㉠, ㉡에 모두 해당되는 단어는?}
\begin{kfEvidence}㉠ 의존 형태소이다. \\\relax
㉡ 실질 형태소이다.\end{kfEvidence}
\begin{kfChoices}
\kfChoice 하늘
\kfChoice 를
\kfChoice -다
\kfChoice 덮-
\kfChoice 맨-
\end{kfChoices}
\end{kfQBlock}
\begin{kfQBlock}{15}{객관식}
\kfQStem{다음 문장의 형태소 분석으로 옳은 것은?}
\begin{kfEvidence}아이가 웃는다.\end{kfEvidence}
\begin{kfChoices}
\kfChoice 총 3개의 형태소로 이루어져 있다.
\kfChoice 자립 형태소는 '아이, 웃-' 2개이다.
\kfChoice 형식 형태소는 '가' 1개이다.
\kfChoice 실질 형태소는 '아이, 웃-' 2개이다.
\kfChoice '-는다'는 하나의 형태소이다.
\end{kfChoices}
\end{kfQBlock}
\begin{kfQBlock}{16}{단답형}
\kfQStem{뜻을 가진 채로 더는 쪼갤 수 \underline{없는} 가장 작은 말의 단위를 무엇이라고 하는가?}
\kfAnswerLines{1}
\end{kfQBlock}
\begin{kfQBlock}{17}{단답형}
\kfQStem{형태소 중, '하늘', '책'처럼 홀로 쓰일 수 있는 형태소는 무엇인가?}
\kfAnswerLines{1}
\end{kfQBlock}
\begin{kfQBlock}{18}{단답형}
\kfQStem{형태소 중, '먹-', '-다', '-이/가'처럼 반드시 다른 말에 기대어 쓰이는 형태소는 무엇인가?}
\kfAnswerLines{1}
\end{kfQBlock}
\begin{kfQBlock}{19}{단답형}
\kfQStem{'하늘', '먹-'처럼 실질적인 의미를 가진 형태소는 무엇인가?}
\kfAnswerLines{1}
\end{kfQBlock}
\begin{kfQBlock}{20}{단답형}
\kfQStem{조사, 어미, 접사처럼 문법적인 의미만 가진 형태소는 무엇인가?}
\kfAnswerLines{1}
\end{kfQBlock}
\begin{kfQBlock}{21}{단답형}
\kfQStem{국어에서 유일하게 '의존 형태소'이면서 동시에 '실질 형태소'인 것은 무엇인가?}
\kfAnswerLines{1}
\end{kfQBlock}
\begin{kfQBlock}{22}{단답형}
\kfQStem{단어에서 실질적인 의미를 나타내는 핵심 부분을 무엇이라고 하는가?}
\kfAnswerLines{1}
\end{kfQBlock}
\begin{kfQBlock}{23}{단답형}
\kfQStem{어근에 붙어 뜻을 더하거나 제한하는 부분을 무엇이라고 하는가?}
\kfAnswerLines{1}
\end{kfQBlock}
\begin{kfQBlock}{24}{단답형}
\kfQStem{'맨손'의 '맨-'처럼 어근 앞에 붙는 접사를 무엇이라고 하는가?}
\kfAnswerLines{1}
\end{kfQBlock}
\begin{kfQBlock}{25}{단답형}
\kfQStem{'덮개'의 '-개'처럼 어근 뒤에 붙는 접사를 무엇이라고 하는가?}
\kfAnswerLines{1}
\end{kfQBlock}
\begin{kfQBlock}{26}{서술형}
\kfQStem{<보기>의 문장을 형태소로 분석하고, 각 형태소를 네 가지 기준 [자립/의존], [실질/형식]에 따라 분류하시오.}
\begin{kfEvidence}그가 새 신을 신었다.\end{kfEvidence}
\kfAnswerNote{답안}{4}
\end{kfQBlock}
\kfEndMC


\kfStartConcept{판소리에서 태어난 소설}


\kfPassageNoteNT{%
　조선 후기, 많은 사람들에게 사랑받던 공연 예술 판소리는 그 인기에 힘입어 새로운 문학 양식으로 발전하였다. 바로 판소리계 소설이다. 판소리계 소설은 판소리로 불리던 작품이 소설로 자리 잡은 것을 말하며, 「춘향전」, 「심청전」, 「흥보전」, 「토끼전」, 「적벽가」 등이 대표적이다. 본래 입으로 전하던 판소리가 글로 쓰인 소설로 변화하면서, 판소리계 소설은 판소리의 특징을 이어받는 동시에 소설로서의 새로운 특징을 지니게 되었다.

　판소리계 소설의 가장 큰 형식적 특징은 판소리 특유의 문체가 남아 있다는 점이다. 판소리는 노래인 창 부분과 말인 아니리 부분이 어우러진 형식이다. 그래서 판소리계 소설에도 운율이 느껴지는 운문체와 설명 중심의 산문체가 함께 쓰인다. 또한, 소리꾼이 관객에게 직접 말을 거는 듯한 '-습니다', '-옵네다'와 같은 높임말 투나, "이때 춘향이가\textasciitilde{}"처럼 이야기꾼인 서술자가 끼어들어 상황을 설명하기도 한다. 특히, 비슷한 문장 구조를 반복하여 리듬감을 살리는 대구법, 여러 가지를 늘어놓는 열거법 등은 판소리의 음악적 특성을 그대로 이어받은 표현이다.

　내용적으로는 판소리와 마찬가지로 평민의 삶을 다루며 재미있고 익살스러운 성격이 강하게 나타난다. 「춘향전」의 방자나 「흥보전」의 흥보처럼 평민 주인공들은 솔직하고 익살맞은 모습을 보이며, 이들의 활약을 통해 평민 의식의 성장을 보여준다. 또한, 백성을 괴롭히는 나쁜 관리(변학도)나 욕심 많은 양반(놀보)을 비꼬며 조선 시대 사회의 잘못된 점을 비판하는 주제를 담고 있다.

　이처럼 판소리계 소설은 '입으로 전하던 이야기'가 '글로 쓰는 소설'로 옮겨가는 중간 단계의 모습을 잘 보여주는 중요한 문학이다. 양반 중심의 문학이 아닌 평민들이 즐기던 이야기가 소설로 자리 잡았다는 점에서, 조선 후기 평민 문화의 성장을 보여준다. 또한 「춘향전」과 「심청전」은 영화·드라마·연극 등으로 각색되어 오늘날까지 사랑받고 있다. 이런 점에서 판소리계 소설은 한국인의 정서가 담긴 소중한 문학 유산이다.
\par\medskip\begin{itemize}[leftmargin=18pt, itemsep=2pt, topsep=2pt, label=\textbullet]
\item 대표작 — 「춘향전」·「심청전」·「흥보전」·「토끼전」·「적벽가」
\item 문체 — 운문체·산문체 혼용 / 서술자 개입('이때 춘향이가\textasciitilde{}')
\item 주제 — 변학도·놀보 풍자 → 평민 의식 성장
\end{itemize}
}

\kfSecHeader{kfAreaConcept}{문제}{}

\par\medskip
\kfBeginMC
\begin{kfQBlock}{01}{객관식}
\kfQStem{판소리계 소설의 형성 배경에 대한 설명으로 적절하지 \uline{\underline{않은}} 것은?}
\begin{kfChoices}
\kfChoice 공연 예술인 ‘판소리’가 그 바탕이 되었다.
\kfChoice 인기를 얻은 판소리 사설이 글로 정착되며 소설로 자리 잡았다.
\kfChoice 「춘향전」·「흥부전」·「심청전」 등이 대표적인 작품이다.
\kfChoice 노래(창)와 말(아니리)이 어우러진 형식의 흔적이 남아 있다.
\kfChoice 가면을 쓰고 춤추는 ‘탈춤’이 그 바탕이 된 공연 예술이다.
\end{kfChoices}
\end{kfQBlock}
\begin{kfQBlock}{02}{객관식}
\kfQStem{판소리계 소설에 운문체와 산문체가 섞여 있는 이유에 대한 설명으로 적절하지 \uline{\underline{않은}} 것은?}
\begin{kfChoices}
\kfChoice 판소리의 ‘창’이 운문체 부분으로 남았기 때문이다.
\kfChoice 판소리의 ‘아니리’가 산문체 부분으로 남았기 때문이다.
\kfChoice 노래로 부른 부분과 말로 풀어 간 부분이 함께 정착되었기 때문이다.
\kfChoice 공연 예술 판소리의 형식적 특성을 그대로 이어받은 결과이다.
\kfChoice 독자의 흥미를 끌려고 일부러 글을 어렵게 만들기 위해서이다.
\end{kfChoices}
\end{kfQBlock}
\begin{kfQBlock}{03}{객관식}
\kfQStem{판소리계 소설의 내용적 특징으로 적절하지 \uline{\underline{않은}} 것은?}
\begin{kfChoices}
\kfChoice 평민들의 삶과 정서를 솔직하게 작품에 담아낸다.
\kfChoice 익살과 해학적 요소가 작품 곳곳에서 두드러진다.
\kfChoice 사회 부조리에 대한 풍자가 작품 안에 자주 나타난다.
\kfChoice 권선징악의 교훈성과 결합되어 마무리되는 경우가 많다.
\kfChoice 전쟁터 영웅의 활약과 도술 같은 비현실적 사건이 중심이다.
\end{kfChoices}
\end{kfQBlock}
\begin{kfQBlock}{04}{객관식}
\kfQStem{<보기>의 (가)와 (나)에 공통으로 해당하는 판소리계 소설의 특징은?}
\begin{kfEvidence}(가) 「춘향전」에서 백성을 괴롭히는 나쁜 관리 '변학도'를 비판함   \\\relax
   (나) 「흥보전」에서 욕심 많고 심술궂은 양반 '놀보'를 비판함\end{kfEvidence}
\begin{kfChoices}
\kfChoice 풍자
\kfChoice 해학
\kfChoice 우연성
\kfChoice 서정성
\kfChoice 영웅성
\end{kfChoices}
\end{kfQBlock}
\begin{kfQBlock}{05}{객관식}
\kfQStem{다음 중 판소리계 소설에 해당하지 \uline{않는} 것은?}
\begin{kfChoices}
\kfChoice 「춘향전」
\kfChoice 「심청전」
\kfChoice 「흥보전」
\kfChoice 「토끼전」
\kfChoice 「홍길동전」
\end{kfChoices}
\end{kfQBlock}
\begin{kfQBlock}{06}{객관식}
\kfQStem{<보기>는 「흥보전」의 한 대목이다. 알 수 있는 표현상의 특징으로 가장 적절한 것은?}
\begin{kfEvidence}"돈 나온다, 돈! 돈! 돈! 엽전 한 꿰미, 은돈 한 꿰미, 황금 한 꿰미... (중략) ... 쌀 나온다, 쌀! 쌀! 쌀! 흰 쌀, 검은 쌀, 찹쌀, 멥쌀..."\end{kfEvidence}
\begin{kfChoices}
\kfChoice 서술자의 개입
\kfChoice 열거법과 반복법
\kfChoice 우연적인 사건 전개
\kfChoice 배경의 구체적인 묘사
\kfChoice 운문체와 산문체의 혼합
\end{kfChoices}
\end{kfQBlock}
\begin{kfQBlock}{07}{객관식}
\kfQStem{<보기>는 「춘향전」의 '십장가' 대목이다. 이 부분에서 드러나는 판소리계 소설의 특징은?}
\begin{kfEvidence}"일(一)자로 나를 치니 일편단심 굳은 마음 변할쏘냐. 이(二)자로 나를 치니 이부불경 굳은 절개 변할쏘냐. 삼(三)자로 나를 치니 삼청동 우리 도련님..."\end{kfEvidence}
\begin{kfChoices}
\kfChoice 시대 변화에 따른 평민 의식의 성장
\kfChoice 무능한 양반 계층에 대한 신랄한 풍자
\kfChoice 비합리적인 사건의 우연적 해결
\kfChoice 이야기 속 서술자의 직접 개입
\kfChoice 운율이 느껴지는 운문체
\end{kfChoices}
\end{kfQBlock}
\begin{kfQBlock}{08}{객관식}
\kfQStem{<보기>는 「심청전」의 한 대목이다. 이 부분에서 드러나는 판소리계 소설의 특징으로 가장 적절한 것은?}
\begin{kfEvidence}"아이고, 이게 누군가! 눈먼 심 봉사 아니시오? 내가 봉사님 눈 뜨게 해 줄 약을 가져왔으니... (중략) 어이쿠, 내가 속였네! 돈이나 내놓으시오!"   \\\relax
   (이 부분은 '뺑덕어멈'이라는 인물이 심 봉사를 속이는 장면이다.)\end{kfEvidence}
\begin{kfChoices}
\kfChoice 임금과 나라에 대한 변치 않는 충성심을 강조한다.
\kfChoice 평민들(뺑덕어멈)의 익살스럽고 재미있는 면모를 보여준다.
\kfChoice 사계절 자연의 아름다운 풍경을 섬세하게 묘사한다.
\kfChoice 사건이 반드시 합리적인 원인과 결과에 따라 일어난다.
\kfChoice 영웅적이고 비범한 인물의 활약상을 본격적으로 보여준다.
\end{kfChoices}
\end{kfQBlock}
\begin{kfQBlock}{09}{단답형}
\kfQStem{조선 후기, 공연 예술 판소리가 인기를 얻어 글로 자리 잡게 된 문학 양식을 무엇이라 하는가?}
\kfAnswerLines{1}
\end{kfQBlock}
\begin{kfQBlock}{10}{단답형}
\kfQStem{판소리처럼 본래 입으로 전하던 문학을 무엇이라 하는가?}
\kfAnswerLines{1}
\end{kfQBlock}
\begin{kfQBlock}{11}{단답형}
\kfQStem{판소리계 소설 5가지 중 한 가지를 쓰시오.}
\kfAnswerLines{1}
\end{kfQBlock}
\begin{kfQBlock}{12}{단답형}
\kfQStem{판소리에서 '노래'에 해당하는 부분을 무엇이라 하는가?}
\kfAnswerLines{1}
\end{kfQBlock}
\begin{kfQBlock}{13}{단답형}
\kfQStem{판소리에서 '말'에 해당하는 부분을 무엇이라 하는가?}
\kfAnswerLines{1}
\end{kfQBlock}
\begin{kfQBlock}{14}{단답형}
\kfQStem{판소리계 소설에서 시처럼 운율이 느껴지는 부분인 운문체는 판소리의 무엇에서 유래했는가?}
\kfAnswerLines{1}
\end{kfQBlock}
\begin{kfQBlock}{15}{단답형}
\kfQStem{판소리계 소설에서 줄글로 설명하는 부분인 산문체는 판소리의 무엇에서 유래했는가?}
\kfAnswerLines{1}
\end{kfQBlock}
\begin{kfQBlock}{16}{서술형}
\kfQStem{‘판소리’와 ‘판소리계 소설’의 관계를 설명하고, 판소리계 소설이 지니는 특징 두 가지를 서술하시오.}
\kfAnswerNote{답안}{4}
\end{kfQBlock}
\kfEndMC


\kfStartLit{「심청전」}


\kfPassageNote{}{%
　송나라 황주 도화동에 사는 심학규는 앞을 못 보는 맹인이었다. 그의 아내 곽씨 부인은 마음씨가 곱고 슬기로웠으며, 온갖 힘든 일을 하며 남편을 정성껏 모셨다. 마흔이 넘도록 자식이 없던 부부는 유명한 절에 정성을 다해 빌었고, 마침내 예쁜 딸을 낳았다.

　하지만 기쁨도 잠시, 곽씨 부인은 아이를 낳은 뒤 병을 얻어 세상을 떠나고 만다. 곽씨 부인은 죽기 전 아이의 이름을 '심청'이라 지어달라는 마지막 말을 남겼다. 홀로 남은 심 봉사는 갓난 딸을 품에 안고 동네를 다니며 아낙들에게 젖을 얻어 겨우 길러냈다.

　세월이 흘러 심청은 마음씨 고운 처녀로 자라, 이제는 아버지를 위해 헌 바가지를 들고 동냥을 하여 아버지를 모셨다.

　그러던 어느 날 저녁, 심청을 마중 나가던 심 봉사가 그만 개천에 빠지는 사고를 당했다. 마침 그곳을 지나던 몽운사의 스님이 그를 구해주고는 "우리 절 부처님께 바칠 쌀 삼백 석을 바치고 정성껏 기도를 올리면 틀림없이 눈을 뜨게 될 것입니다"라고 말했다. 심 봉사는 '눈을 뜬다'는 말에 앞뒤 가릴 것 없이 덜컥 약속을 하고 말았다. 스님이 돌아가자 심 봉사는 정신이 번쩍 들었다. 당장 먹을 쌀 한 톨 없는 형편에 약속을 지킬 길이 없어 자신의 어리석음을 탓하며 한숨을 쉬었다.

　집에 돌아온 심청은 아버지의 이야기를 듣고 위로했다. "아버지, 너무 걱정 마세요. 아버지께서 눈만 뜨실 수 있다면 부처님께 바칠 쌀 삼백 석은 제가 어떻게든 마련하겠습니다."

　며칠 후, 남경 상인들이 인당수라는 거친 바다를 건너기 위해 희생 제물로 바칠 열다섯 살 처녀를 구한다는 소문이 들렸다. 심청은 뱃사람들을 찾아가 "제가 그 희생 제물이 되겠습니다. 제 몸값으로 쌀 삼백 석을 저희 아버지가 약속한 절에 바쳐주십시오"라고 말했다. 뱃사람들은 심청의 효심에 감동하며 그 부탁을 들어주었다.

　심청은 아버지에게는 장 승상 댁 양딸로 가게 되었다고 거짓말을 했다. 배가 떠나기로 한 날, 심청은 마지막으로 아버지의 아침상을 차려드렸다. 식사를 마친 후, 심청은 아버지 앞에 엎드려 울며 모든 사실을 털어놓았다. "아버지, 용서하세요. 저는 인당수 희생 제물로 팔려 오늘 떠납니다. 부디 눈을 꼭 뜨셔서 밝은 세상을 보십시오."

　심 봉사는 딸의 말을 듣고 하늘이 무너지는 듯했다. 

"못 간다, 절대 못 간다! 자식을 죽여서 눈을 뜬들 그게 무슨 소용이냐!" 

　심 봉사는 딸을 붙들고 슬피 울었지만, 이미 약속은 끝난 뒤였다. 심청은 울부짖는 아버지를 뒤로하고 눈물을 머금은 채 배에 올랐다.

　몇 달을 항해한 배는 마침내 인당수에 도착했다. 거친 파도가 몰아치자 뱃사람들은 제사를 지내고 심청에게 물에 뛰어들라고 재촉했다. 심청은 마지막으로 "하늘이시여, 제 목숨을 거두는 대신, 부디 제 아버지의 눈을 뜨게 해주십시오"라고 빌고는 거친 파도 속으로 몸을 던졌다.

　그런데 물에 빠져 죽은 줄만 알았던 심청은 옥황상제의 명으로 용궁으로 가게 되어 정성스러운 대접을 받았다. 3년 후, 심청은 커다란 연꽃에 담겨 다시 인간 세상으로 보내졌다. 마침 그곳을 지나던 뱃사람들이 연꽃을 발견해 황제에게 바쳤고, 황제는 연꽃 속에서 나온 심청의 아름다움에 반해 황후로 맞이했다.

　황후가 된 심청은 늘 아버지를 그리워하다 황제에게 자신의 사연을 털어놓고 아버지를 찾게 해달라고 간절히 부탁했다. 황제는 심청의 효심에 감동하여 "전국의 모든 맹인들을 수도로 불러 큰 잔치를 열어주시오"라는 황후의 청을 따랐다.

　잔치가 끝날 무렵, 황후는 명단에도 없는 허름한 옷차림의 맹인 하나를 발견하고 가까이 불렀다. 황후가 누구인지 묻자 맹인은 눈물을 흘리며 자신의 슬픈 이야기를 했다.

"저는 황주 도화동에 살던 심학규라는 맹인이온데, 하나뿐인 딸 청이가 아비 눈을 뜨게 하려고 쌀 삼백 석에 몸을 팔아 인당수 희생 제물로 빠져 죽었습니다."

　그가 바로 아버지임을 확신한 심청은 급하게 달려 내려가 아버지를 끌어안았다.

"아버지! 제가 바로 인당수에 빠졌던 딸, 청이입니다!"

"뭐라고? 네가 청이란 말이냐!"

　딸의 목소리를 듣는 순간, 놀라움과 기쁨에 심 봉사의 두 눈에서 무언가 떨어지는 소리가 나더니, 어두웠던 눈이 번쩍 뜨였다.

"아, 보인다! 세상이 보여! 내 딸 얼굴이 보이는구나!"

　아버지와 딸은 서로를 부둥켜안고 기쁨의 눈물을 흘렸다. 이 광경을 본 잔치의 모든 맹인들도 동시에 눈을 뜨는 기적이 일어났다.
}


\kfPassageNote{문장 독해 — 발췌}{%
1. 마침 그곳을 지나던 몽운사의 스님이 그를 구해주고는... 2. 당장 먹을 쌀 한 톨 없는 형편에 약속을 지킬 길이 없어 자신의 어리석음을 탓하며 한숨을 쉬었다. 3. 아버지께서 눈만 뜨실 수 있다면 부처님께 바칠 쌀 삼백 석은 제가 어떻게든 마련하겠습니다. 4. 뱃사람들은 심청의 효심에 감동하며 그 부탁을 들어주었다. 5. 황후가 누구인지 묻자 맹인은 눈물을 흘리며 자신의 슬픈 이야기를 했다. 6. 전국의 모든 맹인들을 수도로 불러 큰 잔치를 열어주시오. 7. ...하나뿐인 딸 청이가 아비 눈을 뜨게 하려고 쌀 삼백 석에 몸을 팔아... 8. 그가 바로 아버지임을 확신한 심청은 급하게 달려 내려가 아버지를 끌어안았다.
}


\kfPassageNote{해설 학습 1}{%
심청전은 춘향전, 흥부전과 함께 우리나라 사람들이 가장 사랑하는 고전 소설 중 하나이다. 이 소설은 작가가 누구인지 정확히 알려져 있지 않지만, 원래는 판소리로 불리다가 소설로 만들어진 '판소리계 소설'로 추정된다.

이 이야기의 가장 핵심적인 주제는 바로 '효' 사상이다. 효란, 부모님을 사랑하고 공경하며, 부모님을 위해 자신의 모든 것을 바치는 마음을 말한다. 효는 조선 시대를 지배했던 유교에서 가장 중요하게 생각했던 최고의 가치였다. 주인공 심청은 눈먼 아버지를 위해 자신의 목숨까지 기꺼이 바치는 모습을 통해, 당시 사람들이 생각했던 가장 이상적인 효녀의 모습을 보여준다.

또한 심청전은 '인과응보'와 '권선징악'이라는 동양의 전통적인 가치관을 매우 뚜렷하게 보여준다. 인과응보란, 원인에 따라 결과가 응답하여 돌아온다는 뜻으로, 착한 일을 하면 좋은 결과가, 나쁜 일을 하면 나쁜 결과가 따른다는 불교적인 생각이다. 심청이 아버지를 위해 자신을 희생한 착한 행동이 결국 그녀가 황후가 되고 아버지가 눈을 뜨게 되는 행복한 결과를 가져온 것이 바로 인과응보이다. 권선징악은 착한 일을 권하고 악한 일을 벌한다는 뜻이다. 심청전은 심청이 큰 복을 받는 모습을 통해 효라는 착한 행동을 적극적으로 권장하고, 뺑덕어멈 같은 악한 인물이 벌을 받는 모습을 통해 나쁜 행동을 경계하게 한다.
}


\kfPassageNote{해설 학습 2}{%
「심청전」은 현실에서 일어날 법한 이야기와 함께, 현실을 뛰어넘는 신비롭고 환상적인 요소들이 어우러져 있다. 이러한 비현실적인 설정은 고전 소설의 중요한 특징 중 하나로, 이야기의 주제를 더욱 극적으로 만들어주는 역할을 한다.

심청이 인당수에 몸을 던진 후의 이야기가 대표적이다. 현실적으로는 사람이 물에 빠져 죽는 것이 당연하지만, 심청은 옥황상제의 도움으로 죽지 않고 용궁으로 가게 된다. 용궁은 인간 세상 아래에 있는 상상 속의 세계로, 이곳에서 심청은 극진한 대접을 받으며 행복한 시간을 보낸다. 이는 심청의 지극한 효심에 대한 하늘의 보상을 의미한다.

또한, 심청이 커다란 연꽃에 담겨 다시 인간 세상으로 돌아오는 장면 역시 매우 환상적이다. 연꽃은 불교에서 깨끗함과 고귀함을 상징하는 꽃으로, 심청의 순결하고 착한 마음을 나타내는 장치이다. 이처럼 작가는 용궁이나 연꽃 같은 비현실적인 장치들을 통해, 심청의 효행이 인간 세상을 넘어 천상의 세계까지 감동시킨 위대한 행동이었음을 보여준다. 이러한 신비로운 힘의 개입은 ‘착한 사람은 반드시 복을 받는다’는 주제를 더욱 분명하게 만들어준다.
}


\kfPassageNote{해설 학습 3}{%
「심청전」이 오랜 시간 동안 많은 사람들에게 사랑받을 수 있었던 이유 중 하나는, 이 이야기 속에 우리나라 사람들이 중요하게 생각했던 여러 가지 사상이 자연스럽게 녹아 있기 때문이다.

가장 바탕이 되는 것은 부모에 대한 효도를 강조하는 유교 사상이다. 심청이 아버지를 위해 자신을 희생하는 전체적인 이야기는 유교의 ‘효’ 사상을 그대로 보여준다.

여기에 부처님께 공양미를 바치면 소원이 이루어진다고 믿는 불교 사상이 더해진다. 심 봉사가 쌀 삼백 석을 바치기로 약속한 절이나, 심청을 상징하는 연꽃 등은 모두 불교와 깊은 관련이 있다.

마지막으로, 옥황상제나 용궁과 같은 신비로운 세계는 하늘의 뜻과 운명을 중시하는 도교 사상의 영향을 받은 것이다. 이처럼 「심청전」은 유교의 ‘효’를 중심으로 불교의 ‘인과응보’와 도교의 ‘신선 사상’ 등이 조화롭게 어우러져 있다. 이러한 다양한 사상의 결합은 이야기를 더욱 풍성하게 만들고, 당시 사람들의 세계관과 가치관을 잘 보여주는 역할을 한다.
}


\kfSecHeader{kfAreaLiterature}{문제}{}

\par\medskip
\kfBeginMC
\begin{kfQBlock}{01}{객관식}
\kfQStem{다음 문장에서 ‘그’가 가리키는 대상은 누구인가?}
\begin{kfChoices}
\kfChoice 심청
\kfChoice 심 봉사
\kfChoice 남경 상인들
\kfChoice 몽운사의 스님
\end{kfChoices}
\end{kfQBlock}
\begin{kfQBlock}{02}{객관식}
\kfQStem{다음 문장에서 서술어 ‘쉬었다’의 생략된 주어는 누구인가?}
\begin{kfChoices}
\kfChoice 심청
\kfChoice 곽씨 부인
\kfChoice 심 봉사
\kfChoice 몽운사의 스님
\end{kfChoices}
\end{kfQBlock}
\begin{kfQBlock}{03}{객관식}
\kfQStem{다음 대화에서 ‘저’가 가리키는 대상은 누구인가?}
\begin{kfChoices}
\kfChoice 심청
\kfChoice 심 봉사
\kfChoice 몽운사의 스님
\kfChoice 장 승상 댁 부인
\end{kfChoices}
\end{kfQBlock}
\begin{kfQBlock}{04}{객관식}
\kfQStem{다음 문장에서 ‘그 부탁’이 가리키는 내용으로 적절한 것은?}
\begin{kfChoices}
\kfChoice 쌀 삼백 석을 빌려 달라는 것
\kfChoice 아버지를 잔치에 초대해 달라는 것
\kfChoice 남경까지 안전하게 태워 달라는 것
\kfChoice 제물이 될 테니 쌀을 절에 바쳐 달라는 것
\end{kfChoices}
\end{kfQBlock}
\begin{kfQBlock}{05}{객관식}
\kfQStem{다음 문장에서 ‘황후’가 가리키는 대상은 누구인가?}
\begin{kfChoices}
\kfChoice 심청
\kfChoice 곽씨 부인
\kfChoice 뺑덕어멈
\kfChoice 용궁의 시녀
\end{kfChoices}
\end{kfQBlock}
\begin{kfQBlock}{06}{객관식}
\kfQStem{다음 대화에서 서술어 ‘열어주시오’의 생략된 주체는 누구인가?}
\begin{kfChoices}
\kfChoice 심청
\kfChoice 황제
\kfChoice 맹인들
\kfChoice 심 봉사
\end{kfChoices}
\end{kfQBlock}
\begin{kfQBlock}{07}{객관식}
\kfQStem{다음 대화에서 ‘아비’가 가리키는 대상은 누구인가?}
\begin{kfChoices}
\kfChoice 심 봉사
\kfChoice 남경 상인
\kfChoice 몽운사의 스님
\kfChoice 황주 도화동의 이웃
\end{kfChoices}
\end{kfQBlock}
\begin{kfQBlock}{08}{객관식}
\kfQStem{다음 문장에서 ‘그’가 가리키는 대상은 누구인가?}
\begin{kfChoices}
\kfChoice 이웃
\kfChoice 심청
\kfChoice 맹인
\kfChoice 황제
\end{kfChoices}
\end{kfQBlock}
\begin{kfQBlock}{09}{객관식}
\kfQStem{「심청전」처럼 판소리로 불리다가 소설로 만들어진 작품을 무엇이라고 하는가?}
\kfAnswerLines{2}
\end{kfQBlock}
\begin{kfQBlock}{10}{객관식}
\kfQStem{부모님을 사랑하고 공경하며 모든 것을 바치는 마음을 무엇이라고 하는가?}
\kfAnswerLines{2}
\end{kfQBlock}
\begin{kfQBlock}{11}{객관식}
\kfQStem{착한 일을 하면 좋은 결과가, 나쁜 일을 하면 나쁜 결과가 따른다는 불교적인 사상을 무엇이라고 하는가?}
\kfAnswerLines{2}
\end{kfQBlock}
\begin{kfQBlock}{12}{객관식}
\kfQStem{착한 일을 권하고 악한 일을 벌한다는 뜻의 사자성어는 무엇인가?}
\kfAnswerLines{2}
\end{kfQBlock}
\begin{kfQBlock}{13}{객관식}
\kfQStem{심청전에서 악한 인물은 누구인가?}
\kfAnswerLines{2}
\end{kfQBlock}
\begin{kfQBlock}{14}{객관식}
\kfQStem{고전 소설에서 이야기의 주제를 극적으로 만드는 비현실적인 설정을 무엇이라고 하는가?}
\kfAnswerLines{2}
\end{kfQBlock}
\begin{kfQBlock}{15}{객관식}
\kfQStem{물에 빠진 심청을 구해준 존재는 누구인가?}
\kfAnswerLines{2}
\end{kfQBlock}
\begin{kfQBlock}{16}{객관식}
\kfQStem{심청이 인간 세상으로 돌아오기 전 머물렀던 상상 속 공간은 어디인가?}
\kfAnswerLines{2}
\end{kfQBlock}
\begin{kfQBlock}{17}{객관식}
\kfQStem{심청이 용궁에서 대접을 받은 것은 무엇에 대한 보상을 의미하는가?}
\kfAnswerLines{2}
\end{kfQBlock}
\begin{kfQBlock}{18}{객관식}
\kfQStem{심청이 인간 세상으로 돌아올 때 담겨 있었던 것은 무엇인가?}
\kfAnswerLines{2}
\end{kfQBlock}
\begin{kfQBlock}{19}{객관식}
\kfQStem{이러한 비현실적인 요소들은 작품의 어떤 주제를 강조하는가?}
\kfAnswerLines{2}
\end{kfQBlock}
\begin{kfQBlock}{20}{객관식}
\kfQStem{부모에 대한 효도를 강조하는 것은 어떤 사상과 관련이 있는가?}
\kfAnswerLines{2}
\end{kfQBlock}
\begin{kfQBlock}{21}{객관식}
\kfQStem{심 봉사가 쌀 삼백 석을 바치려고 한 곳은 어디인가?}
\kfAnswerLines{2}
\end{kfQBlock}
\begin{kfQBlock}{22}{객관식}
\kfQStem{절에 공양미를 바치면 소원이 이루어진다는 믿음은 어떤 사상과 관련이 있는가?}
\kfAnswerLines{2}
\end{kfQBlock}
\begin{kfQBlock}{23}{객관식}
\kfQStem{옥황상제, 용궁과 같은 신비로운 세계는 어떤 사상의 영향을 받은 것인가?}
\kfAnswerLines{2}
\end{kfQBlock}
\begin{kfQBlock}{24}{객관식}
\kfQStem{「심청전」은 어떤 사상들이 조화롭게 어우러져 있는가?}
\kfAnswerLines{2}
\end{kfQBlock}
\begin{kfQBlock}{25}{객관식}
\kfQStem{이 글의 내용과 일치하지 \uline{않는} 것은?}
\begin{kfChoices}
\kfChoice 심청은 연꽃을 타고 세상으로 다시 돌아왔다.
\kfChoice 심청의 어머니는 심청을 낳은 직후에 세상을 떠났다.
\kfChoice 심청은 제물이 되는 대가로 쌀 삼백 석을 받기로 했다.
\kfChoice 심 봉사는 스님이 제안한 시주 액수를 흥정하여 줄인 뒤에 약속하였다.
\kfChoice 맹인 잔치에 모인 모든 맹인이 심 봉사와 더불어 그 자리에서 함께 눈을 떴다.
\end{kfChoices}
\end{kfQBlock}
\begin{kfQBlock}{26}{객관식}
\kfQStem{이 글의 중심 소재인 ‘쌀 삼백 석’의 기능으로 적절하지 \uline{\underline{않은}} 것은?}
\begin{kfChoices}
\kfChoice 심청의 효심을 시험하는 시련이 된다.
\kfChoice 심청이 인당수에 몸을 던지는 계기가 된다.
\kfChoice 심청 가족에게 부유하고 행복한 삶을 준다.
\kfChoice 이야기 속 새로운 사건을 만드는 갈등의 원인이 된다.
\kfChoice 심 봉사에게 희망을 주면서 동시에 절망의 원인이 된다.
\end{kfChoices}
\end{kfQBlock}
\begin{kfQBlock}{27}{객관식}
\kfQStem{이 글에 나타난 ‘심청’의 성격으로 가장 거리가 먼 것은?}
\begin{kfChoices}
\kfChoice 효심이 매우 지극하다.
\kfChoice 희생정신이 무척 강하다.
\kfChoice 아버지를 위하는 마음이 깊다.
\kfChoice 어려운 상황에도 슬기롭고 침착하다.
\kfChoice 자신의 운명을 스스로 개척하려 한다.
\end{kfChoices}
\end{kfQBlock}
\begin{kfQBlock}{28}{객관식}
\kfQStem{㉠\textasciitilde{}㉤에 대한 설명으로 적절하지 \uline{\underline{않은}} 것은?}
\begin{kfChoices}
\kfChoice ㉠: 딸을 향한 심 봉사의 지극한 사랑과 힘겨운 삶을 엿볼 수 있다.
\kfChoice ㉡: 눈을 뜨고 싶은 간절한 마음에 신중하게 생각하고 행동하고 있다.
\kfChoice ㉢: 아버지가 미리 슬퍼하거나 걱정하지 않도록 배려하는 마음이 드러난다.
\kfChoice ㉣: 딸을 잃는 슬픔이 눈을 뜨고 싶은 소망보다 크다는 것을 깨닫고 있다.
\kfChoice ㉤: 심청의 효심에 대한 보상이 개인을 넘어 사회 전체로 확대되고 있다.
\end{kfChoices}
\end{kfQBlock}
\begin{kfQBlock}{29}{객관식}
\kfQStem{이 글의 결말이 주는 교훈으로 적절하지 \uline{\underline{않은}} 것은?}
\begin{kfChoices}
\kfChoice 부모를 향한 지극한 효행은 결국 큰 복으로 돌아온다.
\kfChoice 권선징악의 가치가 마지막 결말에서 분명하게 드러난다.
\kfChoice 자기희생적 선행이 결국 보답으로 돌아온다는 메시지를 준다.
\kfChoice 인간의 도덕적 행동이 인생의 큰 흐름을 바꿀 수 있음을 강조한다.
\kfChoice 신중하지 못한 약속은 큰 불행만을 부른다는 점을 교훈으로 전한다.
\end{kfChoices}
\end{kfQBlock}
\begin{kfQBlock}{30}{객관식}
\kfQStem{이 소설의 핵심 주제에 대한 설명으로 적절하지 \uline{\underline{않은}} 것은?}
\begin{kfChoices}
\kfChoice 부모에 대한 자식의 지극한 효성을 핵심 가치로 삼는다.
\kfChoice 효의 실천이 큰 복으로 돌아온다는 권선징악을 보여 준다.
\kfChoice 자기희생적 사랑이 주제 형성에 중심이 된다.
\kfChoice 가족애와 효의 가치를 동시에 담는다.
\kfChoice 신분에 따른 사회적 차별과 계급 갈등이 핵심 주제이다.
\end{kfChoices}
\end{kfQBlock}
\begin{kfQBlock}{31}{객관식}
\kfQStem{이 소설이 오늘날까지도 많은 사람들에게 사랑받는 이유로 가장 알맞은 것은?}
\begin{kfChoices}
\kfChoice 바다 속 용궁의 모습을 신비롭게 묘사해서
\kfChoice 심 봉사의 말실수가 큰 웃음을 주기 때문에
\kfChoice 이 작품을 통해 조선 시대 평민들의 일상 풍습과 신분 제도, 효 사상 같은 역사적 사실을 사실 그대로 정확하고 자세하게 알 수 있어서
\kfChoice 당시 평민들이 겪었던 힘겨운 삶과 가난, 가족을 위해 자신을 기꺼이 희생해야 했던 비극적 처지를 매우 사실적으로 그려서 보여주기 때문에
\kfChoice 인당수에 몸을 던지는 자식의 자기희생, 죽음을 넘어 부녀가 다시 만나는 환상적 결말 등 시대를 초월한 인간 보편의 정서를 다루고 있어서
\end{kfChoices}
\end{kfQBlock}
\begin{kfQBlock}{32}{객관식}
\kfQStem{<보기>의 관점에서 이 글을 이해한 내용으로 적절하지 \uline{\underline{않은}} 것은?}
\begin{kfEvidence}고전 소설에는 현실에서 일어나기 어려운 비현실적이고 신비로운 사건들이 자주 등장한다. 이러한 장치들은 주인공이 겪는 시련의 고통을 극복하게 하고, ‘선한 인물은 복을 받는다’는 주제 의식을 효과적으로 드러내는 역할을 한다.\end{kfEvidence}
\begin{kfChoices}
\kfChoice 심청이 용궁에 간 것은 선행에 대한 보상을 항상 보여주는 장치이군.
\kfChoice 옥황상제의 등장은 하늘이 선인을 돕는다는 믿음을 보여주는군.
\kfChoice 심청이 연꽃을 타고 온 것은 고귀함을 상징하는 신비로운 장치이군.
\kfChoice 맹인들이 모두 눈을 뜬 것은 착한 효심이 만든 기적이라고 할 수 있군.
\kfChoice 남경 상인들이 심청을 제물로 바치는 것은 비현실적인 사건에 해당하군.
\end{kfChoices}
\end{kfQBlock}
\begin{kfQBlock}{33}{서술형}
\kfQStem{심청은 왜 인당수에 몸을 던지기로 결심했는지 서술하시오.}
\kfAnswerNote{답안}{4}
\end{kfQBlock}
\begin{kfQBlock}{34}{서술형}
\kfQStem{결말에서 심 봉사 외에 다른 맹인들도 모두 눈을 뜬 장면이 어떤 의미를 가지는지 서술하시오.}
\kfAnswerNote{답안}{4}
\end{kfQBlock}
\kfEndMC


\kfStartNonlit{[사회] 국가가 처벌하는 형사법 vs 개인이 다투는 민사법}


\kfPassageNote{[사회] 국가가 처벌하는 형사법 vs 개인이 다투는 민사법}{%
　우리 사회에서는 관계의 종류에 따라 적용되는 법이 달라진다. 법의 대표적인 두 분야는 형사법과 민사법이다. 형사법은 국가와 범죄자 사이의 관계를 다루고, 민사법은 개인과 개인 또는 회사 같은 법인과의 관계에 적용된다.

　형사법은 사회 질서를 유지하고 범죄를 처벌하기 위해 공익을 대표하는 국가가 범죄자에게 형벌을 가한다. 형벌이란 생명, 자유, 명예, 재산 등의 기본권을 빼앗는 것이다. 반면 민사법은 개인 간 분쟁을 해결하고 개인의 권리를 보호하는 것이 목적이다. 사건 당사자들이 평등한 관계라는 전제 하에 손해와 이익을 조정한다.

　형사 소송에서는 검사와 피고인이 당사자가 된다. 공익을 대표하는 검사가 범죄 혐의가 있는 사람을 피고인으로 기소하면서 소송이 시작된다. 검사가 피고인의 유죄를 입증해야 하고, 피고인은 변호인을 통해 반박할 수 있다. 법원은 이를 바탕으로 범죄 성립 여부와 잘못의 정도를 따져서 적당한 벌을 내린다. 같은 종류의 범죄라도 동기와 상황, 피해자와의 합의 여부 등을 고려해서 형량이 달라질 수 있다.

　민사 소송에서는 원고와 피고가 당사자가 된다. 피해를 당했다고 주장하며 소송을 제기한 사람이 원고가 되고, 가해자로 지목된 상대방이 피고가 된다. 민사 소송에서는 양쪽 모두 자기에게 유리한 사실을 입증해야 한다. 입증하지 못하면 법원은 해당 당사자에게 불리하게 판단한다. 형사 소송과 달리 두 당사자가 손해와 이익을 적절히 타협하면 바로 소송이 끝난다.

　하나의 사건도 두 분야 모두에서 문제가 될 수 있다. 예를 들어 갑이 을에게 맞아서 다쳤다면, 검사가 을을 상해죄로 처벌해 달라는 형사 소송을 제기할 수도 있고, 갑이 을에게 치료비와 위자료를 요구하는 민사 소송을 제기할 수도 있다. 하지만 같은 사건이라도 결론이 다를 수 있는데, 이는 입증 정도가 다르기 때문이다.

　입증의 정도에서 가장 큰 차이가 난다. 형사 소송은 '법관이 합리적인 의심을 할 여지가 없을 정도'의 강한 입증을 요구한다. 무죄추정의 원칙에 따라 "피고인은 유죄 판결이 확정될 때까지 무죄로 추정된다"고 보기 때문이다. 반면 민사 소송은 '일반인이라면 의심하지 않을 정도'의 입증만 요구한다. 원고와 피고의 입증 정도가 51 대 49라면 원고가 이기게 된다. 이처럼 형사법과 민사법은 목적과 절차, 입증 정도가 모두 다른 별개의 법 분야로 우리 사회의 질서를 각각 다른 방식으로 지키고 있다.
}


\kfPassageNote{문장 독해 — 본문}{%
1. 우리 사회에서는 관계의 종류에 따라 적용되는 법이 달라진다. 2. 법의 대표적인 두 분야는 형사법과 민사법이다. 3. 형사법은 국가와 범죄자 사이의 관계를 다루고, 민사법은 개인과 개인 또는 회사 같은 법인과의 관계에 적용된다. 4. 형사법은 사회 질서를 유지하고 범죄를 처벌하기 위해 공익을 대표하는 국가가 범죄자에게 형벌을 가한다. 5. 형벌이란 생명, 자유, 명예, 재산 등의 기본권을 빼앗는 것이다. 6. 반면 민사법은 개인 간 분쟁을 해결하고 개인의 권리를 보호하는 것이 목적이다. 7. 사건 당사자들이 평등한 관계라는 전제 하에 손해와 이익을 조정한다. 8. 형사 소송에서는 검사와 피고인이 당사자가 된다. 9. 공익을 대표하는 검사가 범죄 혐의가 있는 사람을 피고인으로 기소하면서 소송이 시작된다. 10. 검사가 피고인의 유죄를 입증해야 하고, 피고인은 변호인을 통해 반박할 수 있다. 11. 법원은 이를 바탕으로 범죄 성립 여부와 잘못의 정도를 따져서 적당한 벌을 내린다. 12. 같은 종류의 범죄라도 동기와 상황, 피해자와의 합의 여부 등을 고려해서 형량이 달라질 수 있다. 13. 민사 소송에서는 원고와 피고가 당사자가 된다. 14. 피해를 당했다고 주장하며 소송을 제기한 사람이 원고가 되고, 가해자로 지목된 상대방이 피고가 된다. 15. 민사 소송에서는 양쪽 모두 자기에게 유리한 사실을 입증해야 한다. 16. 입증하지 못하면 법원은 해당 당사자에게 불리하게 판단한다. 17. 형사 소송과 달리 두 당사자가 손해와 이익을 적절히 타협하면 바로 소송이 끝난다. 18. 하나의 사건도 두 분야 모두에서 문제가 될 수 있다. 19. 예를 들어 갑이 을에게 맞아서 다쳤다면, 검사가 을을 상해죄로 처벌해 달라는 형사 소송을 제기할 수도 있고, 갑이 을에게 치료비와 위자료를 요구하는 민사 소송을 제기할 수도 있다. 20. 하지만 같은 사건이라도 결론이 다를 수 있는데, 이는 입증 정도가 다르기 때문이다. 21. 입증의 정도에서 가장 큰 차이가 난다. 22. 형사 소송은 '법관이 합리적인 의심을 할 여지가 없을 정도'의 강한 입증을 요구한다. 23. 무죄추정의 원칙에 따라 "피고인은 유죄 판결이 확정될 때까지 무죄로 추정된다"고 보기 때문이다. 24. 반면 민사 소송은 '일반인이라면 의심하지 않을 정도'의 입증만 요구한다. 25. 원고와 피고의 입증 정도가 51 대 49라면 원고가 이기게 된다. 26. 이처럼 형사법과 민사법은 목적과 절차, 입증 정도가 모두 다른 별개의 법 분야로 우리 사회의 질서를 각각 다른 방식으로 지키고 있다.
}


\kfActivityBg \par\kfMaybeNeedspace{32\baselineskip}\noindent{\kfTipMark\,\,}{\small\color{kfInk} 다음 뜻풀이에 해당하는 어휘를 지문에서 찾아 쓰시오.}\par\nopagebreak[4]\smallskip\nopagebreak[4]
\begin{kfVocab}
\kfVocabItem{1}{범죄와 형벌에 대해 정해 놓은 법}
\kfVocabItem{2}{개인 간의 재산이나 신분 관계를 다루는 법}
\kfVocabItem{3}{법률에 따라 사람처럼 권리와 의무를 가지는 단체(회사 등)}
\kfVocabItem{4}{사회 전체의 이익}
\kfVocabItem{5}{범죄에 대한 대가로 국가가 가하는 제재}
\kfVocabItem{6}{사람들 사이에 일어나는 다툼}
\kfVocabItem{7}{법원에 재판을 신청하여 분쟁 해결을 요구하는 절차}
\kfVocabItem{8}{범죄 사건을 수사하고 법원에 재판을 청구하는 국가 공무원}
\kfVocabItem{9}{형사 재판에서 범죄 혐의를 받아 재판을 받는 사람}
\kfVocabItem{10}{검사가 법원에 특정 사건의 재판을 요구하는 일}
\kfVocabItem{11}{죄가 있다고 인정되는 것}
\kfVocabItem{12}{민사 재판에서 소송을 제기한 사람}
\kfVocabItem{13}{민사 재판에서 소송을 당한 상대방}
\kfVocabItem{14}{어떤 사실을 증거를 들어서 증명하는 것}
\kfVocabItem{15}{정신적 고통에 대한 손해 배상금}
\kfVocabItem{16}{유죄 판결이 확정되기 전까지는 무죄로 본다는 원칙}
\end{kfVocab}

\kfActivityBg \par\kfMaybeNeedspace{32\baselineskip}\noindent{\kfTipMark\,\,}{\small\color{kfInk} 다음 글의 구조를 정리한 표의 빈칸을 채워 보세요.}\par\nopagebreak[4]\smallskip\nopagebreak[4]
\par\medskip\needspace{8\baselineskip}
\begin{tcolorbox}[enhanced, breakable=true,
  colback=kfPaper, colframe=kfRule, boxrule=0.4pt, arc=2pt,
  left=8pt, right=8pt, top=6pt, bottom=6pt,
  title={\footnotesize\bfseries\color{kfMute} 글의 구조도},
  coltitle=kfMute, colbacktitle=kfPrimaryLight!40,
  attach boxed title to top left={yshift=-1.5mm, xshift=6pt},
  boxed title style={colframe=kfRule, boxrule=0.3pt, arc=1pt}]
\footnotesize\setstretch{1.3}
\renewcommand{\arraystretch}{1.35}
\arrayrulecolor{kfRule}
\noindent\begin{tabularx}{\linewidth}{|>{\centering\arraybackslash}p{1.6cm}|>{\raggedright\arraybackslash}p{3.4cm}|>{\raggedright\arraybackslash}X|}
\hline
\rowcolor{kfPrimaryLight!40}\textbf{\color{kfPrimary}단} & \textbf{\color{kfPrimary}항목} & \textbf{\color{kfPrimary}내용} \\\hline
\multirow{2}{*}{형사법} & 관계 & ( ① )와 범죄자 사이 \\\cline{2-3}
 & 목적 & 사회 질서 유지, 범죄 ( ② ) \\\cline{2-3}
\hline
\multirow{2}{*}{민사법} & 관계 & ( ③ ) 간의 관계 \\\cline{2-3}
 & 목적 & 개인 간 분쟁 해결, 개인의 ( ④ ) 보호 \\\cline{2-3}
\hline
당사자: ( ⑤ )(공익 대표) vs 피고인(범죄 혐의자) & 당사자: ( ⑤ )(공익 대표) vs 피고인(범죄 혐의자) &  \\\hline
입증 책임: 검사가 피고인의 ( ⑥ )를 입증해야 함 & 입증 책임: 검사가 피고인의 ( ⑥ )를 입증해야 함 &  \\\hline
판결: 법원이 범죄 성립 여부와 잘못의 정도를 따져 형벌 결정 & 판결: 법원이 범죄 성립 여부와 잘못의 정도를 따져 형벌 결정 &  \\\hline
당사자: ( ⑦ )(소송 제기) vs 피고(가해자로 지목) & 당사자: ( ⑦ )(소송 제기) vs 피고(가해자로 지목) &  \\\hline
입증 책임: 양쪽 모두 자신에게 ( ⑧ )이 없음을 입증해야 함 & 입증 책임: 양쪽 모두 자신에게 ( ⑧ )이 없음을 입증해야 함 &  \\\hline
특징: 당사자 간 ( ⑨ )이 이루어지면 소송 종료 가능 & 특징: 당사자 간 ( ⑨ )이 이루어지면 소송 종료 가능 &  \\\hline
관계: 하나의 사건에 두 소송 모두 가능하지만, 결론은 다를 수 있음 & 관계: 하나의 사건에 두 소송 모두 가능하지만, 결론은 다를 수 있음 &  \\\hline
\multirow{2}{*}{가장 큰 차이점: ( ⑩ )의 정도} & 형사 소송 & ‘합리적인 의심을 할 여지가 없을 정도’의 강한 입증 요구 ( ⑪ ) 원칙 때문 \\\cline{2-3}
 & 민사 소송 & ‘일반인이라면 의심하지 않을 정도’의 입증 요구 \\\cline{2-3}
\hline
형사법과 민사법은 목적, 절차, 입증 정도가 다른 ( ⑫ )의 법 분야임 & 형사법과 민사법은 목적, 절차, 입증 정도가 다른 ( ⑫ )의 법 분야임 &  \\\hline
\end{tabularx}
\arrayrulecolor{black}
\end{tcolorbox}\par\medskip
\kfSecHeader{kfAreaNonlit}{문제}{}

\par\medskip
\kfBeginMC
\begin{kfQBlock}{01}{객관식}
\kfQStem{우리 사회에서 적용되는 법은 무엇에 따라 달라지는가?}
\begin{kfChoices}
\kfChoice 개인의 의사
\kfChoice 국가의 명령
\kfChoice 사회의 규모
\kfChoice 관계의 종류
\end{kfChoices}
\end{kfQBlock}
\begin{kfQBlock}{02}{객관식}
\kfQStem{이 글에서 다루는 법의 대표적인 두 분야는 무엇인가?}
\begin{kfChoices}
\kfChoice 공법과 사법
\kfChoice 헌법과 형법
\kfChoice 상법과 민법
\kfChoice 형사법과 민사법
\end{kfChoices}
\end{kfQBlock}
\begin{kfQBlock}{03}{객관식}
\kfQStem{형사법은 누구와 누구 사이의 관계를 다루는가?}
\begin{kfChoices}
\kfChoice 개인과 개인
\kfChoice 국가와 국가
\kfChoice 국가와 범죄자
\kfChoice 기업과 소비자
\end{kfChoices}
\end{kfQBlock}
\begin{kfQBlock}{04}{객관식}
\kfQStem{민사법은 누구 사이의 관계에 적용되는가?}
\begin{kfChoices}
\kfChoice 국가와 국가
\kfChoice 국가와 범죄자
\kfChoice 정부와 공무원
\kfChoice 개인과 개인 또는 법인
\end{kfChoices}
\end{kfQBlock}
\begin{kfQBlock}{05}{객관식}
\kfQStem{형사법의 두 가지 목적은 무엇인가?}
\begin{kfChoices}
\kfChoice 국가 배상과 개인 구제
\kfChoice 개인 간의 화해와 조정
\kfChoice 사회 질서 유지와 범죄 처벌
\kfChoice 기업의 이익 추구와 보호
\end{kfChoices}
\end{kfQBlock}
\begin{kfQBlock}{06}{객관식}
\kfQStem{형사법에서 형벌을 가하는 주체는 누구인가?}
\begin{kfChoices}
\kfChoice 법원
\kfChoice 피해자 가족
\kfChoice 범죄 피해자
\kfChoice 공익을 대표하는 국가
\end{kfChoices}
\end{kfQBlock}
\begin{kfQBlock}{07}{객관식}
\kfQStem{'형벌'이란 무엇을 빼앗는 것인가?}
\begin{kfChoices}
\kfChoice 기본권
\kfChoice 투표권
\kfChoice 교육권
\kfChoice 노동 3권
\end{kfChoices}
\end{kfQBlock}
\begin{kfQBlock}{08}{객관식}
\kfQStem{다음 중 이 문장에서 예시로 든 기본권에 해당하는 것은?}
\begin{kfChoices}
\kfChoice 여행
\kfChoice 투표
\kfChoice 교육
\kfChoice 생명
\end{kfChoices}
\end{kfQBlock}
\begin{kfQBlock}{09}{객관식}
\kfQStem{'반면'이라는 표현을 볼 때, 이 문장은 앞 내용과 어떤 관계인가?}
\begin{kfChoices}
\kfChoice 인과 관계
\kfChoice 대조 관계
\kfChoice 병렬 관계
\kfChoice 예시 관계
\end{kfChoices}
\end{kfQBlock}
\begin{kfQBlock}{10}{객관식}
\kfQStem{민사법의 두 가지 목적은 무엇인가?}
\begin{kfChoices}
\kfChoice 범죄 예방과 교화
\kfChoice 사회 정의와 질서 유지
\kfChoice 국가의 권위와 위신 제고
\kfChoice 개인 간 분쟁 해결과 권리 보호
\end{kfChoices}
\end{kfQBlock}
\begin{kfQBlock}{11}{객관식}
\kfQStem{민사법이 손해와 이익을 조정할 때 전제하는 것은 무엇인가?}
\begin{kfChoices}
\kfChoice 가해자가 경제적 약자라는 것
\kfChoice 피해자가 도덕적 우위라는 것
\kfChoice 국가가 갈등을 중재한다는 것
\kfChoice 사건 당사자들이 평등한 관계라는 것
\end{kfChoices}
\end{kfQBlock}
\begin{kfQBlock}{12}{객관식}
\kfQStem{형사 소송의 당사자는 누구와 누구인가?}
\begin{kfChoices}
\kfChoice 원고와 피고
\kfChoice 판사와 변호사
\kfChoice 검사와 피고인
\kfChoice 경찰관과 목격자
\end{kfChoices}
\end{kfQBlock}
\begin{kfQBlock}{13}{객관식}
\kfQStem{형사 소송은 어떻게 시작되는가?}
\begin{kfChoices}
\kfChoice 피해자가 고소장을 제출하면서
\kfChoice 경찰이 범인을 체포하면서
\kfChoice 판사가 재판 개시를 선언하면서
\kfChoice 검사가 피고인을 기소하면서
\end{kfChoices}
\end{kfQBlock}
\begin{kfQBlock}{14}{객관식}
\kfQStem{형사 소송에서 '피고인'은 어떤 사람인가?}
\begin{kfChoices}
\kfChoice 범죄 피해를 입은 사람
\kfChoice 범죄 혐의가 있는 사람
\kfChoice 범죄를 목격한 사람
\kfChoice 재판을 진행하는 판사
\end{kfChoices}
\end{kfQBlock}
\begin{kfQBlock}{15}{객관식}
\kfQStem{형사 소송에서 피고인의 유죄를 입증해야 할 책임은 누구에게 있는가?}
\begin{kfChoices}
\kfChoice 검사
\kfChoice 판사
\kfChoice 변호인
\kfChoice 피고인
\end{kfChoices}
\end{kfQBlock}
\begin{kfQBlock}{16}{객관식}
\kfQStem{피고인은 누구를 통해 자신의 주장을 반박할 수 있는가?}
\begin{kfChoices}
\kfChoice 검사
\kfChoice 판사
\kfChoice 증인
\kfChoice 변호인
\end{kfChoices}
\end{kfQBlock}
\begin{kfQBlock}{17}{객관식}
\kfQStem{법원이 벌을 내리기 위해 따지는 두 가지는 무엇인가?}
\begin{kfChoices}
\kfChoice 피해 금액과 합의 여부
\kfChoice 범죄 성립 여부와 잘못의 정도
\kfChoice 피고인의 재산과 사회적 지위
\kfChoice 변호인의 능력과 검사의 태도
\end{kfChoices}
\end{kfQBlock}
\begin{kfQBlock}{18}{객관식}
\kfQStem{형량이 달라질 수 있는 요인으로 언급되지 \uline{\underline{않은}} 것은?}
\begin{kfChoices}
\kfChoice 동기
\kfChoice 상황
\kfChoice 날씨
\kfChoice 피해자와의 합의 여부
\end{kfChoices}
\end{kfQBlock}
\begin{kfQBlock}{19}{객관식}
\kfQStem{동기, 상황, 피해자와의 합의 여부 등은 무엇을 결정할 때 고려되는 요소인가?}
\begin{kfChoices}
\kfChoice 형량
\kfChoice 무죄
\kfChoice 벌금
\kfChoice 구속
\end{kfChoices}
\end{kfQBlock}
\begin{kfQBlock}{20}{객관식}
\kfQStem{민사 소송의 당사자는 누구와 누구인가?}
\begin{kfChoices}
\kfChoice 원고와 피고
\kfChoice 검사와 피고인
\kfChoice 판사와 배심원
\kfChoice 경찰과 용의자
\end{kfChoices}
\end{kfQBlock}
\begin{kfQBlock}{21}{객관식}
\kfQStem{민사 소송에서 '원고'는 어떤 사람인가?}
\begin{kfChoices}
\kfChoice 범죄 혐의를 받고 있는 사람
\kfChoice 가해자로 지목된 상대방
\kfChoice 재판을 진행하고 판결하는 사람
\kfChoice 피해를 당했다고 주장하며 소송을 제기한 사람
\end{kfChoices}
\end{kfQBlock}
\begin{kfQBlock}{22}{객관식}
\kfQStem{민사 소송에서 '피고'는 어떤 사람인가?}
\begin{kfChoices}
\kfChoice 소송을 제기한 사람
\kfChoice 가해자로 지목된 상대방
\kfChoice 재판을 참관하는 방청객
\kfChoice 법률 자문을 해주는 변호사
\end{kfChoices}
\end{kfQBlock}
\begin{kfQBlock}{23}{객관식}
\kfQStem{민사 소송에서 입증 책임은 누구에게 있는가?}
\begin{kfChoices}
\kfChoice 검사
\kfChoice 원고만
\kfChoice 피고만
\kfChoice 원고와 피고 모두
\end{kfChoices}
\end{kfQBlock}
\begin{kfQBlock}{24}{객관식}
\kfQStem{민사 소송에서 자신에게 책임이 없음을 입증하지 못하면 어떤 결과가 나타나는가?}
\begin{kfChoices}
\kfChoice 재판이 무효가 된다
\kfChoice 형량이 반으로 줄어든다
\kfChoice 법원이 해당 당사자에게 불리하게 판단한다
\kfChoice 국가가 대신 입증해 준다
\end{kfChoices}
\end{kfQBlock}
\begin{kfQBlock}{25}{객관식}
\kfQStem{민사 소송이 형사 소송과 달리 쉽게 끝날 수 있는 경우는 언제인가?}
\begin{kfChoices}
\kfChoice 범인이 자백했을 때
\kfChoice 증거가 너무 명확할 때
\kfChoice 두 당사자가 타협했을 때
\kfChoice 변호인이 사임했을 때
\end{kfChoices}
\end{kfQBlock}
\begin{kfQBlock}{26}{객관식}
\kfQStem{이 문장에서 '두 분야'가 가리키는 것은 무엇인가?}
\begin{kfChoices}
\kfChoice 헌법과 민법
\kfChoice 상법과 형법
\kfChoice 공법과 사법
\kfChoice 형사법과 민사법
\end{kfChoices}
\end{kfQBlock}
\begin{kfQBlock}{27}{객관식}
\kfQStem{폭행 사건의 경우, 가해자를 처벌하기 위해 제기하는 소송은 무엇인가?}
\begin{kfChoices}
\kfChoice 헌법 소원
\kfChoice 행정 소송
\kfChoice 형사 소송
\kfChoice 민사 소송
\end{kfChoices}
\end{kfQBlock}
\begin{kfQBlock}{28}{객관식}
\kfQStem{폭행 사건의 경우, 피해자가 치료비를 받기 위해 제기하는 소송은 무엇인가?}
\begin{kfChoices}
\kfChoice 형사 소송
\kfChoice 민사 소송
\kfChoice 행정 소송
\kfChoice 가사 소송
\end{kfChoices}
\end{kfQBlock}
\begin{kfQBlock}{29}{객관식}
\kfQStem{같은 사건에 대한 형사 소송과 민사 소송의 결론이 다를 수 있는 이유는 무엇인가?}
\begin{kfChoices}
\kfChoice 변호사의 실력 차이 때문
\kfChoice 입증 정도가 다르기 때문
\kfChoice 재판 기간이 다르기 때문
\kfChoice 법원의 위치가 다르기 때문
\end{kfChoices}
\end{kfQBlock}
\begin{kfQBlock}{30}{객관식}
\kfQStem{형사 소송과 민사 소송의 가장 큰 차이점은 무엇인가?}
\begin{kfChoices}
\kfChoice 판사의 수
\kfChoice 법정의 크기
\kfChoice 입증의 정도
\kfChoice 소송 비용
\end{kfChoices}
\end{kfQBlock}
\begin{kfQBlock}{31}{객관식}
\kfQStem{형사 소송에서 요구하는 입증의 정도는 어느 수준인가?}
\begin{kfChoices}
\kfChoice 약간의 의심은 허용되는 정도
\kfChoice 반반의 확률로 믿을 수 있는 정도
\kfChoice 일반인이 의심하지 않을 정도
\kfChoice 법관이 합리적인 의심을 할 여지가 없을 정도
\end{kfChoices}
\end{kfQBlock}
\begin{kfQBlock}{32}{객관식}
\kfQStem{형사 소송에서 강한 입증을 요구하는 이유는 어떤 원칙 때문인가?}
\begin{kfChoices}
\kfChoice 죄형법정주의
\kfChoice 일사부재리의 원칙
\kfChoice 무죄추정의 원칙
\kfChoice 과실 책임의 원칙
\end{kfChoices}
\end{kfQBlock}
\begin{kfQBlock}{33}{객관식}
\kfQStem{민사 소송에서 요구하는 입증의 정도는 어느 수준인가?}
\begin{kfChoices}
\kfChoice 전문가가 확신할 수 있는 정도
\kfChoice 일반인이라면 의심하지 않을 정도
\kfChoice 과학적으로 완벽하게 증명된 정도
\kfChoice 법관이 100\% 확신할 수 있는 정도
\end{kfChoices}
\end{kfQBlock}
\begin{kfQBlock}{34}{객관식}
\kfQStem{이 문장은 민사 소송의 입증 정도가 형사 소송보다 어떻다는 것을 보여주는가?}
\begin{kfChoices}
\kfChoice 훨씬 덜 엄격하다는 것
\kfChoice 훨씬 더 까다롭다는 것
\kfChoice 완전히 동일하다는 것
\kfChoice 입증이 불필요하다는 것
\end{kfChoices}
\end{kfQBlock}
\begin{kfQBlock}{35}{객관식}
\kfQStem{다음 중 형사법과 민사법의 차이점으로 언급되지 \uline{\underline{않은}} 것은?}
\begin{kfChoices}
\kfChoice 목적
\kfChoice 절차
\kfChoice 입증 정도
\kfChoice 법관의 성별
\end{kfChoices}
\end{kfQBlock}
\begin{kfQBlock}{36}{객관식}
\kfQStem{형사법과 민사법의 공통된 역할은 무엇인가?}
\begin{kfChoices}
\kfChoice 국가의 수입을 늘리는 것
\kfChoice 개인의 복수를 돕는 것
\kfChoice 우리 사회의 질서를 지키는 것
\kfChoice 기업의 성장을 촉진하는 것
\end{kfChoices}
\end{kfQBlock}
\begin{kfQBlock}{37}{객관식}
\kfQStem{윗글을 통해 알 수 있는 내용으로 적절하지 \uline{\underline{않은}} 것은?}
\begin{kfChoices}
\kfChoice 형사법은 공익을, 민사법은 개인의 이익을 보호하는 데 중점을 둔다.
\kfChoice 민사 소송에서는 당사자 간의 합의를 통해 소송을 끝낼 수 있다.
\kfChoice 형사 소송에서는 피고인의 유죄를 검사가 입증해야 한다.
\kfChoice 하나의 사건에 대해 형사 소송과 민사 소송이 동시에 진행될 수는 없다.
\kfChoice 형사 소송은 민사 소송보다 더 높은 수준의 입증을 요구한다.
\end{kfChoices}
\end{kfQBlock}
\begin{kfQBlock}{38}{객관식}
\kfQStem{형사 소송과 민사 소송을 비교한 설명으로 적절하지 \uline{\underline{않은}} 것은?}
\begin{kfChoices}
\kfChoice 형사 소송은 국가(검사) 대 개인(피고인)이 다투는 구도로 진행된다.
\kfChoice 민사 소송은 개인(원고) 대 개인(피고)이 서로 다투는 구도이다.
\kfChoice 형사는 범죄자 처벌을, 민사는 권리 관계 해결을 주된 목적으로 한다.
\kfChoice 형사 소송에서는 검사가 입증의 책임을 지는 비중이 매우 크다.
\kfChoice 형사는 51\% 입증으로 승소하고, 민사는 100\% 입증되어야 승소한다.
\end{kfChoices}
\end{kfQBlock}
\begin{kfQBlock}{39}{객관식}
\kfQStem{같은 사건에 대해 형사 소송과 민사 소송의 결론이 다를 수 있는 가장 큰 이유는?}
\begin{kfChoices}
\kfChoice 재판을 담당하는 법관이 다르기 때문이다.
\kfChoice 적용되는 법 조항이 완전히 다르기 때문이다.
\kfChoice 소송에 필요한 비용의 차이가 크기 때문이다.
\kfChoice 각 소송에서 요구하는 증거의 종류가 다르기 때문이다.
\kfChoice 유무죄나 책임을 판단하는 입증의 기준이 다르기 때문이다.
\end{kfChoices}
\end{kfQBlock}
\begin{kfQBlock}{40}{객관식}
\kfQStem{<보기>의 ㉠, ㉡에 들어갈 당사자로 올바르게 짝지어진 것은?}
\begin{kfEvidence}친구에게 빌려준 돈을 받지 못한 A는 법원에 소송을 제기했다. 이 소송에서 A는 ( \\\relax
㉠ )이/가 되고, 친구는 ( \\\relax
㉡ )이/가 된다.\end{kfEvidence}
\begin{kfChoices}
\kfChoice 원고 - 피고
\kfChoice 피고 - 원고
\kfChoice 검사 - 피고인
\kfChoice 피고인 - 검사
\kfChoice 원고 - 피고인
\end{kfChoices}
\end{kfQBlock}
\begin{kfQBlock}{41}{객관식}
\kfQStem{윗글을 바탕으로 <보기>의 상황을 이해한 것으로 가장 적절한 것은?}
\begin{kfEvidence}폭행 혐의로 기소된 연예인 B는 형사 재판에서 증거 불충분으로 무죄 판결을 받았다. 하지만 피해자가 제기한 민사 소송에서는 패소하여 치료비와 위자료를 지급하라는 판결을 받았다.\end{kfEvidence}
\begin{kfChoices}
\kfChoice 형사 재판의 판결이 잘못되었으므로 다시 재판해야 한다.
\kfChoice 민사 재판부는 형사 재판의 결과를 반드시 따라야 하므로 판결이 잘못되었다.
\kfChoice 형사상 유죄를 입증할 만큼 증거가 충분하지는 않았지만, 민사상 책임은 인정된 것이다.
\kfChoice 민사 소송에서 이겼으므로 B는 형사 재판 결과와 상관없이 전과 기록이 남게 된다.
\kfChoice B가 민사 소송에서 먼저 합의했다면 형사 재판은 열리지 않았을 것이다.
\end{kfChoices}
\end{kfQBlock}
\begin{kfQBlock}{42}{단답형}
\kfQStem{국가와 범죄자 사이의 관계를 다루며 사회 질서 유지를 목적으로 하는 법은 무엇인가?}
\kfAnswerLines{1}
\end{kfQBlock}
\begin{kfQBlock}{43}{단답형}
\kfQStem{개인 간의 분쟁 해결 및 권리 보호를 목적으로 하는 법은 무엇인가?}
\kfAnswerLines{1}
\end{kfQBlock}
\begin{kfQBlock}{44}{단답형}
\kfQStem{민사법이 개인 외에 적용되는 또 다른 대상은 무엇인가?}
\kfAnswerLines{1}
\end{kfQBlock}
\begin{kfQBlock}{45}{단답형}
\kfQStem{형벌을 통해 국가가 범죄자에게서 박탈하는 생명, 자유, 재산 등과 같은 권리를 통칭하여 무엇이라고 하는가?}
\kfAnswerLines{1}
\end{kfQBlock}
\begin{kfQBlock}{46}{단답형}
\kfQStem{형사 소송에서 범죄 혐의가 있는 사람을 기소하는 당사자는 누구인가?}
\kfAnswerLines{1}
\end{kfQBlock}
\begin{kfQBlock}{47}{서술형}
\kfQStem{형사 소송과 민사 소송에서 요구되는 '입증의 정도'가 어떻게 다른지 두 소송의 기준을 비교하여 서술하시오.}
\kfAnswerNote{답안}{4}
\end{kfQBlock}
\kfEndMC



\clearpage

\kfChapterCover{3}{러셀3 (챕터3)}{Chapter 3}{이번 챕터의 학습 내용을 시작합니다.}
\begin{kfChapterAreaList}
\kfChapterAreaItem{문법}{kfAreaGrammar}{새로운 단어의 탄생: 파생어와 합성어}
\kfChapterAreaItem{개념}{kfAreaConcept}{시간의 흐름에 담긴 의미}
\kfChapterAreaItem{문학}{kfAreaLiterature}{「비」 — 정지용}
\kfChapterAreaItem{비문학}{kfAreaNonlit}{[과학] 방사성 동위원소를 이용한 암석 연대 측정법}
\end{kfChapterAreaList}

\kfStartGrammar{새로운 단어의 탄생: 파생어와 합성어}


\kfPassageNoteNT{%
　우리가 사용하는 단어는 그 짜임에 따라 하나의 어근으로만 이루어진 단일어(單一語)와 둘 이상의 형태소가 결합한 복합어(複合語)로 나뉜다. '하늘', '바다', '먹다'처럼 더 이상 쪼갤 수 없는 단어가 단일어라면, '맨손'이나 '돌다리'처럼 둘 이상의 형태소가 결합한 단어는 복합어이다. 이 복합어는 새로운 단어를 만드는 방식, 즉 단어 형성법에 따라 파생어와 합성어로 다시 나뉜다.

　파생어(派生語)는 단어의 실질적 의미를 나타내는 어근(語根)에, 뜻을 더하거나 제한하는 접사(接辭)가 결합하여 만들어진 단어다. 접사는 붙는 위치에 따라 종류가 나뉜다. '맨손'의 '맨-', '헛기침'의 '헛-', '풋사과'의 '풋-', '참기름'의 '참-', '군소리'의 '군-'처럼 어근 앞에 붙는 접사를 접두사라고 한다. 이들은 각각 어근에 '아무것도 없는', '이유 없는', '덜 익은', '진짜의', '쓸데없는'과 같은 뜻을 더해준다.

　반면, 어근 뒤에 붙는 접미사는 어근에 '그런 사람', '도구', '행위'와 같은 구체적인 뜻을 더해준다. '장난꾸러기'의 '-꾸러기', '사냥꾼'의 '-꾼', '욕심쟁이'의 '-쟁이', '울보'의 '-보'는 모두 '그런 특성을 가진 사람'이라는 뜻을 더한다. '덮개'나 '지우개'의 '-개'는 '도구'의 뜻을, '낚시질'이나 '가위질'의 '-질'은 '행위'의 뜻을 더한다. 또한 '정답다'의 '-답-'이나 '평화롭다'의 '-롭-'처럼 어근에 붙어 '...같은' 또는 '...스러운 상태'라는 뜻을 더해주는 접미사도 있다.

　합성어(合成語)는 접사가 아닌, 실질적 의미를 지닌 어근과 어근이 결합하여 만들어진 단어다. 합성어는 결합하는 어근의 품사에 따라 다양하게 만들어진다. '돌다리'(돌+다리), '책가방'(책+가방), '집안'(집+안), '손발'(손+발), '눈사람'(눈+사람)처럼 명사와 명사가 결합하는 경우가 가장 많다. 또한 '오가다'(오다+가다), '뛰어가다'(뛰다+가다)처럼 동사 어근이 결합하거나, '검붉다'(검다+붉다), '높푸르다'(높다+푸르다)처럼 형용사 어근이 결합하기도 한다. '새해'(새+해), '큰집'(크다+집)처럼 관형사나 용언의 관형사형이 명사와 결합하여 만들어지는 합성어도 있다.

　이처럼 우리말은 어근에 접사를 붙이는 방식(파생)이나, 어근과 어근을 결합하는 방식(합성)을 통해 새로운 단어를 풍부하게 만들어 낸다.
\par\medskip\begin{itemize}[leftmargin=18pt, itemsep=2pt, topsep=2pt, label=\textbullet]
\item 접두 파생 — '맨손·헛기침·풋사과·참기름·군소리'
\item 접미 파생 — '장난꾸러기·욕심쟁이·지우개·낚시질'
\item 합성 — '돌다리(명+명)·오가다(동+동)·검붉다(형+형)·새해(관+명)'
\end{itemize}
}

\kfActivityBg \par\kfMaybeNeedspace{32\baselineskip}지문의 핵심 내용을 떠올리며 아래 빈칸에 알맞은 말을 채워 보세요.

1. 단어의 분류 

\par\medskip\needspace{4\baselineskip}
\noindent\begin{adjustbox}{max width=\linewidth, center}
\renewcommand{\arraystretch}{1.45}
\arrayrulecolor{kfRule}
\begin{tabular}{|>{\raggedright\arraybackslash}p{\dimexpr 0.1370\linewidth-12pt\relax}|>{\raggedright\arraybackslash}p{\dimexpr 0.4274\linewidth-12pt\relax}|>{\raggedright\arraybackslash}p{\dimexpr 0.4356\linewidth-12pt\relax}|}
\hline
\rowcolor{kfPrimaryLight!50}\textbf{\color{kfPrimary} 분류} & \textbf{\color{kfPrimary} 개념} & \textbf{\color{kfPrimary} 예시} \\\hline
[ ① ] & ∙ 하나의 어근으로만 이루어진 단어 & ∙ 하늘, 바다, 먹다 \\\hline
[ ② ] & ∙ 둘 이상의 형태소가 결합한 단어 & ∙ 맨손, 돌다리, 덮개 \\\hline
\end{tabular}
\arrayrulecolor{black}
\end{adjustbox}\par\medskip



2. 복합어의 종류 (단어 형성법) 

\par\medskip\needspace{4\baselineskip}
\noindent\begin{adjustbox}{max width=\linewidth, center}
\renewcommand{\arraystretch}{1.45}
\arrayrulecolor{kfRule}
\begin{tabular}{|>{\raggedright\arraybackslash}p{\dimexpr 0.0462\linewidth-12pt\relax}|>{\raggedright\arraybackslash}p{\dimexpr 0.0730\linewidth-12pt\relax}|>{\raggedright\arraybackslash}p{\dimexpr 0.1710\linewidth-12pt\relax}|>{\raggedright\arraybackslash}p{\dimexpr 0.7098\linewidth-12pt\relax}|}
\hline
\rowcolor{kfPrimaryLight!50}\textbf{\color{kfPrimary} 형성법} & \textbf{\color{kfPrimary} 개념 (짜임)} & \textbf{\color{kfPrimary} 세부 분류} & \textbf{\color{kfPrimary} 예시} \\\hline
\multirow{2}{*}{[ ③ ]} & ∙ 어근 + [ ④ ] & 접두 파생어<br>([ ⑤ ]가 어근 앞에 붙음) & ∙ [ ⑥ ] + 손 → 맨손<br>∙ [ ⑦ ] + 사과 → 풋사과<br>∙ 군- + 소리 → 군소리 \\\cline{2-4}
 & \rule[-1.0em]{0pt}{2.4em} & 접미 파생어<br>([ ⑧ ]가 어근 뒤에 붙음) & ∙ 장난 + [ ⑨ ] → 장난꾸러기<br>∙ 덮- + [ ⑩ ] → 덮개 (도구)<br>∙ 낚시 + -질 → 낚시질 (행위) \\\cline{2-4}
\multirow{3}{*}{[ ⑪ ]} & ∙ 어근 + [ ⑫ ] & 명사+명사 & ∙ 돌 + 다리 → [ ⑬ ]<br>∙ 책 + 가방 → 책가방 \\\cline{2-4}
 & \rule[-1.0em]{0pt}{2.4em} & 용언 어근 결합 & ∙ 오- + 가- → [ ⑭ ]<br>∙ 검- + 붉- → [ ⑮ ] \\\cline{2-4}
 & \rule[-1.0em]{0pt}{2.4em} & 기타 결합 & ∙ 새 + 해 → [ ⑯ ]<br>∙ 큰 + 집 → 큰집 \\\cline{2-4}
\hline
\end{tabular}
\arrayrulecolor{black}
\end{adjustbox}\par\medskip\nopagebreak[4]\par\nopagebreak[4]

\kfActivityBg \kfSecHeader{kfAreaGrammar}{활동}{분석 훈련}
\par\kfMaybeNeedspace{24\baselineskip}\noindent{\kfTipMark\,\,}{\small\color{kfInk} 다음 활용 사례를 분석하여 빈칸을 채워 보세요.}\par\nopagebreak[4]\smallskip\nopagebreak[4]
\kfBeginMC
\begin{enumerate}[leftmargin=22pt, itemsep=6pt, topsep=4pt, label=\textbf{\arabic*.}]
\item 하늘
\item 헛기침
\item 돌다리
\item 장난꾸러기
\item 뛰어가다
\item 지우개
\item 풋사과
\item 검붉다
\item 새해
\item 정답다
\end{enumerate}
\kfEndMC

\kfSecHeader{kfAreaGrammar}{문제}{}

\par\medskip
\kfBeginMC
\begin{kfQBlock}{01}{객관식}
\kfQStem{단어의 짜임에 대한 설명으로 적절하지 \uline{\underline{않은}} 것은?}
\begin{kfChoices}
\kfChoice 하나의 어근으로 이루어진 단어를 '단일어'라고 한다.
\kfChoice '복합어'는 둘 이상의 형태소가 결합한 단어이다.
\kfChoice '복합어'는 파생어와 합성어로 나뉜다.
\kfChoice '어근'은 단어의 실질적 의미를 나타내는 핵심 부분이다.
\kfChoice '접사'는 실질적 의미를 가지며 홀로 쓰일 수 있다.
\end{kfChoices}
\end{kfQBlock}
\begin{kfQBlock}{02}{객관식}
\kfQStem{<보기>의 단어들을 단어의 짜임에 따라 바르게 분류한 것은?}
\begin{kfEvidence}하늘, 손발, 덧신, 구름\end{kfEvidence}
\begin{kfChoices}
\kfChoice 단일어: 하늘, 손발 / 복합어: 덧신, 구름
\kfChoice 단일어: 하늘, 구름 / 복합어: 손발, 덧신
\kfChoice 단일어: 하늘, 덧신 / 복합어: 손발, 구름
\kfChoice 단일어: 덧신, 구름 / 복합어: 하늘, 손발
\kfChoice 단일어: 손발, 덧신 / 복합어: 하늘, 구름
\end{kfChoices}
\end{kfQBlock}
\begin{kfQBlock}{03}{객관식}
\kfQStem{다음 중 단어의 형성 방법이 나머지와 \uline{다른} 하나는?}
\begin{kfChoices}
\kfChoice 맨손
\kfChoice 돌다리
\kfChoice 풋사과
\kfChoice 장난꾸러기
\kfChoice 덮개
\end{kfChoices}
\end{kfQBlock}
\begin{kfQBlock}{04}{객관식}
\kfQStem{<보기>의 ㉠\textasciitilde{}㉢에 해당하는 예시가 바르게 짝지어진 것은?}
\begin{kfEvidence}㉠ 단일어 \\\relax
㉡ 파생어 \\\relax
㉢ 합성어\end{kfEvidence}
\begin{kfChoices}
\kfChoice ㉠ 바다 / ㉡ 덮개 / ㉢ 돌다리
\kfChoice ㉠ 덮개 / ㉡ 바다 / ㉢ 돌다리
\kfChoice ㉠ 바다 / ㉡ 돌다리 / ㉢ 덮개
\kfChoice ㉠ 돌다리 / ㉡ 덮개 / ㉢ 바다
\kfChoice ㉠ 덮개 / ㉡ 돌다리 / ㉢ 바다
\end{kfChoices}
\end{kfQBlock}
\begin{kfQBlock}{05}{객관식}
\kfQStem{밑줄 친 부분이 '접두사'가 \uline{\underline{아닌}} 것은?}
\begin{kfChoices}
\kfChoice \uline{헛}기침을 하다.
\kfChoice \uline{군}소리를 듣다.
\kfChoice \uline{새}파란 하늘이 보인다.
\kfChoice \uline{새} 신을 신었다.
\kfChoice \uline{참}기름을 샀다.
\end{kfChoices}
\end{kfQBlock}
\begin{kfQBlock}{06}{객관식}
\kfQStem{다음 밑줄 친 부분이 '접미사'인 것은?}
\begin{kfChoices}
\kfChoice 하늘\uline{이}  맑다.
\kfChoice 밥을 먹\uline{고} 잔다.
\kfChoice \uline{맨}손으로 잡다.
\kfChoice \uline{돌}다리가 무너졌다.
\kfChoice 장난\uline{꾸러기}가 웃는다.
\end{kfChoices}
\end{kfQBlock}
\begin{kfQBlock}{07}{객관식}
\kfQStem{단어의 짜임이 <보기>와 같은 것은?}
\begin{kfEvidence}어근 + 어근\end{kfEvidence}
\begin{kfChoices}
\kfChoice 헛수고
\kfChoice 잠꾸러기
\kfChoice 손발
\kfChoice 맑다
\kfChoice 덮개
\end{kfChoices}
\end{kfQBlock}
\begin{kfQBlock}{08}{객관식}
\kfQStem{<보기>의 ㉠, ㉡에 들어갈 말로 바르게 짝지어진 것은?}
\begin{kfEvidence}∙ '돌다리', '책가방', '오가다', '검붉다'는 ( \\\relax
㉠ )이다.   \\\relax
   ∙ '맨손', '풋사과', '지우개', '욕심쟁이'는 ( \\\relax
㉡ )이다.\end{kfEvidence}
\begin{kfChoices}
\kfChoice ㉠ 단일어 / ㉡ 파생어
\kfChoice ㉠ 파생어 / ㉡ 합성어
\kfChoice ㉠ 합성어 / ㉡ 파생어
\kfChoice ㉠ 단일어 / ㉡ 합성어
\kfChoice ㉠합성어 / ㉡ 단일어
\end{kfChoices}
\end{kfQBlock}
\begin{kfQBlock}{09}{객관식}
\kfQStem{밑줄 친 접사의 의미가 바르게 연결되지 \uline{\underline{않은}} 것은?}
\begin{kfChoices}
\kfChoice \uline{맨}손 - '아무것도 없는'
\kfChoice \uline{헛}기침 - '이유 없는'
\kfChoice \uline{풋}사과 - '덜 익은'
\kfChoice \uline{참}기름 - '질이 떨어지는'
\kfChoice \uline{군}소리 - '쓸데없는'
\end{kfChoices}
\end{kfQBlock}
\begin{kfQBlock}{10}{객관식}
\kfQStem{다음 중 파생어의 짜임이 나머지와 \uline{다른} 하나는?}
\begin{kfChoices}
\kfChoice 풋사과
\kfChoice 덮개
\kfChoice 욕심쟁이
\kfChoice 사냥꾼
\kfChoice 울보
\end{kfChoices}
\end{kfQBlock}
\begin{kfQBlock}{11}{객관식}
\kfQStem{다음 중 '파생어'가 \uline{\underline{아닌}} 것은?}
\begin{kfChoices}
\kfChoice 덮개
\kfChoice 맨손
\kfChoice 지우개
\kfChoice 군소리
\kfChoice 눈사람
\end{kfChoices}
\end{kfQBlock}
\begin{kfQBlock}{12}{객관식}
\kfQStem{다음 중 '합성어'가 \uline{\underline{아닌}} 것은?}
\begin{kfChoices}
\kfChoice 돌다리
\kfChoice 오가다
\kfChoice 새해
\kfChoice 높푸르다
\kfChoice 사냥꾼
\end{kfChoices}
\end{kfQBlock}
\begin{kfQBlock}{13}{객관식}
\kfQStem{다음 중 단어의 짜임이 '접두사 + 어근'인 것은?}
\begin{kfChoices}
\kfChoice 읽히다
\kfChoice 집안
\kfChoice 굳이
\kfChoice 군밤
\kfChoice 잠자다
\end{kfChoices}
\end{kfQBlock}
\begin{kfQBlock}{14}{객관식}
\kfQStem{‘새해’와 ‘큰집’이 합성어인 이유에 대한 설명으로 적절하지 \uline{\underline{않은}} 것은?}
\begin{kfChoices}
\kfChoice ‘새’와 ‘해’가 모두 실질적 의미를 가진 어근이다.
\kfChoice ‘크-’와 ‘집’이 모두 실질적 의미를 가진 어근이다.
\kfChoice 두 단어 모두 어근과 어근이 결합한 짜임이다.
\kfChoice 접사가 결합한 파생어와 구별되는 단어 형성 방식이다.
\kfChoice ‘새-’와 ‘크-’가 접사이기 때문에 합성어로 분류된다.
\end{kfChoices}
\end{kfQBlock}
\begin{kfQBlock}{15}{객관식}
\kfQStem{<보기>의 단어 형성법에 대한 설명으로 옳은 것은?}
\begin{kfEvidence}(가) 덮개, 지우개, 장난꾸러기    \\\relax
    (나) 돌다리, 책가방, 오가다\end{kfEvidence}
\begin{kfChoices}
\kfChoice (가)는 어근과 어근이 결합한 합성어이다.
\kfChoice (가)는 모두 어근 앞에 접두사가 붙은 파생어이다.
\kfChoice (나)는 어근에 접사가 결합한 파생어이다.
\kfChoice (나)는 모두 어근 뒤에 접미사가 붙은 파생어이다.
\kfChoice (가)는 파생어, (나)는 합성어이다.
\end{kfChoices}
\end{kfQBlock}
\begin{kfQBlock}{16}{단답형}
\kfQStem{하나의 어근으로만 이루어진 단어를 무엇이라고 하는가?}
\kfAnswerLines{1}
\end{kfQBlock}
\begin{kfQBlock}{17}{단답형}
\kfQStem{둘 이상의 형태소가 결합한 단어를 통틀어 무엇이라고 하는가?}
\kfAnswerLines{1}
\end{kfQBlock}
\begin{kfQBlock}{18}{단답형}
\kfQStem{복합어 중, 어근에 접사가 결합하여 만들어진 단어를 무엇이라고 하는가?}
\kfAnswerLines{1}
\end{kfQBlock}
\begin{kfQBlock}{19}{단답형}
\kfQStem{복합어 중, 어근과 어근이 결합하여 만들어진 단어를 무엇이라고 하는가?}
\kfAnswerLines{1}
\end{kfQBlock}
\begin{kfQBlock}{20}{단답형}
\kfQStem{'군소리', '참기름'의 '군-', '참-'처럼 어근 앞에 붙는 접사를 무엇이라고 하는가?}
\kfAnswerLines{1}
\end{kfQBlock}
\begin{kfQBlock}{21}{단답형}
\kfQStem{'사냥꾼', '덮개'의 '-꾼', '-개'처럼 어근 뒤에 붙는 접사를 무엇이라고 하는가?}
\kfAnswerLines{1}
\end{kfQBlock}
\begin{kfQBlock}{22}{단답형}
\kfQStem{'돌다리'를 직접 구성 성분으로 나누면 [　　　　　] + [　　　　　]이다.}
\kfAnswerLines{1}
\end{kfQBlock}
\begin{kfQBlock}{23}{단답형}
\kfQStem{'맨손'을 직접 구성 성분으로 나누면 [　　　　　] + [　　　　　]이다.}
\kfAnswerLines{1}
\end{kfQBlock}
\begin{kfQBlock}{24}{단답형}
\kfQStem{'지우개'의 실질적 의미를 나타내는 핵심 부분(어근)은 무엇인가?}
\kfAnswerLines{1}
\end{kfQBlock}
\begin{kfQBlock}{25}{단답형}
\kfQStem{'장난꾸러기'에서 '사람'의 뜻을 더해주는 접미사는 무엇인가?}
\kfAnswerLines{1}
\end{kfQBlock}
\begin{kfQBlock}{26}{서술형}
\kfQStem{'파생어'와 '합성어'는 모두 둘 이상의 형태소가 결합한 '복합어'입니다. 두 단어를 구분하는 기준을 '어근'과 '접사'라는 용어를 사용하여 서술하시오.}
\kfAnswerNote{답안}{4}
\end{kfQBlock}
\kfEndMC


\kfStartConcept{시간의 흐름에 담긴 의미}


\kfPassageNoteNT{%
　시를 읽다 보면, 시상이 펼쳐지는 순서가 작품마다 다르다는 것을 느낄 수 있다. 어떤 시는 처음부터 끝까지 시간 순서대로 흘러가고, 어떤 시는 현재 장면에서 시작하여 과거를 떠올리기도 한다. 이처럼 시인이 자신의 생각이나 감정을 전개하는 방식을 '시상 전개 방식'이라고 한다. 그중 대표적인 방식이 '시간의 흐름'에 따른 전개인데, 여기에는 크게 '순행적 구성'과 '역순행적 구성'이 있다.

　'순행적 구성'은 시간의 흐름에 따라 순서대로 시상이 진행되는 방식이다. 과거에서 현재로, 아침에서 저녁으로, 혹은 봄에서 겨울로 이어지는 계절의 변화처럼 시간의 순서를 따르는 것이 대표적이다. 이육사 「광야」가 좋은 예이다. 이 시는 "까마득한 날"이라는 과거의 시간에 하늘이 처음 열리고, "지금"이라는 현재의 상황에서 눈이 내리는 것을 거쳐, "다시 천고의 뒤"라는 미래의 시간에 초인이 올 것이라는 희망으로 이어진다. 이처럼 과거-현재-미래의 순서로 시상이 전개되므로 순행적 구성에 해당한다. 순행적 구성은 자연스러운 시간의 흐름을 따르기 때문에 독자에게 안정감을 준다.

　'순행적 구성'과 반대로 '역순행적 구성'은 현재의 상황에서 과거를 떠올리거나 회상하는 방식으로 시상을 전개하는 것이다. 기형도 「엄마 걱정」이 대표적인 예이다. 이 시는 어른이 된 "지금" 현재의 화자가 "내 유년의 윗목"이었던 "아주 먼 옛날"의 과거를 떠올리는 구조로 되어 있다. "열무 삼십 단을 이고 / 시장에 간" 엄마를 기다리며 "빈방에 혼자 엎드려 훌쩍거리던" 어린 시절의 모습을 과거 회상으로 구체적으로 보여 준다. 이처럼 현재에서 과거를 회상하는 방식은 현재까지 이어지는 화자의 정서를 독자에게 강렬하게 전달한다.

　시상 전개 방식은 단순히 생각을 늘어놓는 순서가 아니라, 시인이 전달하고자 하는 주제와 정서를 효과적으로 드러내기 위한 표현 기법이다. 같은 내용이라도 어떤 순서로 배치하느냐에 따라 독자가 받는 느낌은 완전히 달라진다.
\par\medskip\begin{itemize}[leftmargin=18pt, itemsep=2pt, topsep=2pt, label=\textbullet]
\item 순행 — 이육사 「광야」(과거-현재-미래 순)
\item 역순행 — 기형도 「엄마 걱정」(현재에서 유년 회상)
\item 효과 차이 — 순행=안정감 / 역순행=현재 정서 강조
\end{itemize}
}

\kfSecHeader{kfAreaConcept}{문제}{}

\par\medskip
\kfBeginMC
\begin{kfQBlock}{01}{객관식}
\kfQStem{윗글의 내용과 일치하지 \uline{않는} 것은?}
\begin{kfChoices}
\kfChoice 시상 전개 방식은 표현 기법의 하나이다.
\kfChoice 순행적 구성은 자연스러운 시간의 흐름을 따른다.
\kfChoice 역순행적 구성은 현재에서 과거를 회상하는 방식이다.
\kfChoice 〈광야〉는 과거-현재-미래의 시간 순서를 보여준다.
\kfChoice 〈엄마 걱정〉은 미래에서 현재를 바라보는 구성이다.
\end{kfChoices}
\end{kfQBlock}
\begin{kfQBlock}{02}{객관식}
\kfQStem{다음 중, ‘역순행적 구성’에 따른 시상 전개가 가장 두드러지게 나타나는 작품은?}
\begin{kfChoices}
\kfChoice 산봉우리에서 강물로 시선을 옮기며 풍경을 묘사한 시
\kfChoice 봄, 여름, 가을, 겨울의 순서로 계절의 변화를 노래한 시
\kfChoice 고향역에 도착한 화자가 정들었던 집을 떠나던 날을 회상하는 시
\kfChoice 사랑의 기쁨이 점차 이별의 슬픔으로 심화되는 과정을 그린 시
\kfChoice 화자가 거울을 보며 자신의 내면을 성찰하는 내용을 담은 시
\end{kfChoices}
\end{kfQBlock}
\begin{kfQBlock}{03}{객관식}
\kfQStem{이 시의 시상 전개 방식에 대한 설명으로 적절하지 \uline{\underline{않은}} 것은?}
\begin{kfEvidence}흔들리는 나뭇가지에 꽃 한번 피우려고

눈은 얼마나 많은 도전을 멈추지 않았을까

사각사각 두드려 보았겠지

펄펄 흩날리며 춤추었겠지

미끄러지고 미끄러지길 수백 번,

바람 한 번 불면 휙 날아갈 사랑을 위하여

새하얀 솜 같은 마음을 다 퍼부어 준 다음에야

마침내 피워 낸 저 눈부신 모습 보아라

봄이면 가지는 그 한번 시렸던 자리에

세상에서 가장 아름다운 상처를 피워낸다

- 고재종 「첫사랑」\end{kfEvidence}
\begin{kfChoices}
\kfChoice ‘겨울’에서 ‘봄’으로 자연스레 이어지는 계절의 흐름을 따라 시상이 전개된다.
\kfChoice 눈이 꽃을 피우려는 과정과 그 결과를 차례로 보여 준다.
\kfChoice 자연 현상을 시간의 흐름 속에서 의미 있게 재구성한다.
\kfChoice 시간이 흐르면서 의미가 점점 깊어지는 구조를 가진다.
\kfChoice 화자의 시선이 하늘에서 나뭇가지로 점차 이동하며 공간을 따라 시상이 전개된다.
\end{kfChoices}
\end{kfQBlock}
\begin{kfQBlock}{04}{객관식}
\kfQStem{이 시의 시상 전개 방식에 대한 설명으로 가장 적절한 것은?}
\begin{kfEvidence}아픈 몸 일으켜 혼자 찬밥을 먹는다

찬밥 속에 날카로운 차가움이 목을 찌른다

부엌에는 각종 전기 제품이 있어

일 분만 단추를 눌러도 따끈한 밥이 되는 세상

찬밥을 먹기도 쉽지 않지만

오늘 혼자 찬밥을 먹는다

가족에겐 따스한 밥 지어 먹이고

찬밥을 먹던 사람

이가 깨진 그릇에 찬밥 긁어모아

누가 남긴 무 조각에 생선 가시를 핥고

몸에서는 제일 따스한 사랑을 뿜던 그녀

깊은 밤에도

혼자 달그락거리던 그 손이 그리워

나 오늘 아픈 몸 일으켜 찬밥을 먹는다

집집마다 신을 보낼 수 없어

신 대신 보냈다는 설도 있지만

홀로 먹는 찬밥 속에서 그녀를 만난다

나 오늘

세상의 찬밥이 되어

-문정희「찬밥」\end{kfEvidence}
\begin{kfChoices}
\kfChoice 시간의 순서에 따라 과거에서 현재로 감정이 고조된다.
\kfChoice 화자가 방에서 부엌으로 공간을 이동하며 시상이 전개된다.
\kfChoice 현재의 상황에서 미래의 희망을 떠올리며 의지를 다진다.
\kfChoice 현재의 행동을 계기로 과거의 기억을 회상한 후 현재로 돌아온다.
\kfChoice 과거의 즐거웠던 기억과 현재의 외로운 처지를 대비하고 있다.
\end{kfChoices}
\end{kfQBlock}
\begin{kfQBlock}{05}{객관식}
\kfQStem{<보기>의 시상 전개 방식에 대한 설명으로 가장 적절한 것은?}
\begin{kfEvidence}죽는 날까지 하늘을 우러러   \\\relax
   한 점 부끄럼이 없기를,   \\\relax
   잎새에 이는 바람에도   \\\relax
   나는 괴로워했다.   \\\relax
   별을 노래하는 마음으로   \\\relax
   모든 죽어 가는 것을 사랑해야지   \\\relax
   그리고 오늘 밤에도 별이 바람에 스치운다.   \\\relax
   -윤동주 「서시」\end{kfEvidence}
\begin{kfChoices}
\kfChoice 현재에서 과거를 회상하고 있다.
\kfChoice 과거에서 현재로 시간이 흐르고 있다.
\kfChoice 공간의 이동에 따라 시상이 전개되고 있다.
\kfChoice 미래의 상황을 가정하며 시상이 전개되고 있다.
\kfChoice 시간의 순서와 관계없이 시상이 진행되고 있다.
\end{kfChoices}
\end{kfQBlock}
\begin{kfQBlock}{06}{객관식}
\kfQStem{<보기>의 시상 전개 방식의 효과에 대한 설명으로 적절하지 \uline{\underline{않은}} 것은?}
\begin{kfEvidence}어제는 아버지를 잃고, 오늘은 아들을 잃었다. 내일은 또 어떤 슬픔이 나를 기다리고 있을까.\end{kfEvidence}
\begin{kfChoices}
\kfChoice ‘어제 → 오늘 → 내일’의 시간 흐름을 따라 시상이 전개된다.
\kfChoice 화자의 슬픔이 시간의 흐름과 함께 점점 커진다.
\kfChoice 시간 흐름 속에서 정서가 누적되어 깊어진다.
\kfChoice 같은 정서가 반복되며 슬픔의 무게를 강조한다.
\kfChoice 장소를 옮겨 다니며 화자가 걸어온 길을 보여 준다.
\end{kfChoices}
\end{kfQBlock}
\begin{kfQBlock}{07}{객관식}
\kfQStem{다음 중, 시상 전개 방식이 나머지 넷과 가장 \uline{다른} 하나는?}
\begin{kfChoices}
\kfChoice 1월, 2월, 3월의 풍경을 차례로 묘사한다.
\kfChoice 아침 안개, 한낮 햇살, 저녁노을 순서로 노래한다.
\kfChoice 어린 시절, 현재의 모습, 미래의 다짐을 이어서 그린다.
\kfChoice 봄의 파종, 여름의 성장, 가을의 수확을 순서대로 보여준다.
\kfChoice 산 정상에서 마을로, 다시 강으로 시선을 옮기며 풍경을 그린다.
\end{kfChoices}
\end{kfQBlock}
\begin{kfQBlock}{08}{단답형}
\kfQStem{시인이 자신의 생각이나 감정을 전개하는 방식을 무엇이라고 하는가?}
\kfAnswerLines{1}
\end{kfQBlock}
\begin{kfQBlock}{09}{단답형}
\kfQStem{지문에서 다룬 시상 전개 방식의 대표적인 분류 기준은 무엇인가?}
\kfAnswerLines{1}
\end{kfQBlock}
\begin{kfQBlock}{10}{단답형}
\kfQStem{시간의 흐름에 따라 순서대로 시상이 진행되는 방식을 무엇이라고 하는가?}
\kfAnswerLines{1}
\end{kfQBlock}
\begin{kfQBlock}{11}{단답형}
\kfQStem{현재의 상황에서 과거를 떠올리거나 회상하는 방식으로 시상이 전개되는 방식을 무엇이라고 하는가?}
\kfAnswerLines{1}
\end{kfQBlock}
\begin{kfQBlock}{12}{단답형}
\kfQStem{「광야」에서 현재를 나타내는 시어 한 단어를 쓰시오.}
\kfAnswerLines{1}
\end{kfQBlock}
\begin{kfQBlock}{13}{단답형}
\kfQStem{지문에 따르면, 순행적 구성이 독자에게 주는 느낌은 무엇인가?}
\kfAnswerLines{1}
\end{kfQBlock}
\begin{kfQBlock}{14}{단답형}
\kfQStem{지문에서 역순행적 구성의 대표적인 예시로 제시한 작품명을 쓰시오.}
\kfAnswerLines{1}
\end{kfQBlock}
\begin{kfQBlock}{15}{서술형}
\kfQStem{‘시간의 흐름’에 따른 시상 전개 방식의 유형 두 가지를 들고, 각각의 개념을 서술하시오.}
\kfAnswerNote{답안}{4}
\end{kfQBlock}
\kfEndMC


\kfStartLit{「비」 — 정지용}


\kfPassageNote{}{%
돌에 \\[4pt]
그늘이 차고, \\[14pt]
한쪽으로 몰려드는 \\[4pt]
서늘한 바람. \\[14pt]
앞서거니 뒤서거니 \\[4pt]
꼬리 휙 치켜들고, \\[14pt]
종종 다리 까칠한 \\[4pt]
산새 걸음걸이. \\[14pt]
여울지어 흘러가는 \\[4pt]
수척한 흰 물살, \\[14pt]
갈래갈래 \\[4pt]
손가락을 펴는 듯. \\[14pt]
멎은 듯하더니 \\[4pt]
다시 돋는 빗방울, \\[14pt]
붉은 잎 잎 \\[4pt]
소란히 밟고 간다.
}


\kfPassageNote{문장 독해 — 발췌}{%
1. 꼬리 휙 치켜들고, \\[4pt]
2. 종종 다리 까칠한 산새 걸음걸이. \\[4pt]
3. 산새 걸음걸이. \\[4pt]
4. 갈래갈래 손가락을 펴는 듯. \\[4pt]
5. 멎은 듯하더니 다시 돋는 빗방울 \\[4pt]
6. 붉은 잎 잎 소란히 밟고 간다. \\[4pt]
7. 갈래갈래 손가락을 펴는 듯.
}


\kfPassageNote{해설 학습 1}{%
정지용의 시 「비」는 비 오는 날의 풍경을 마치 한 폭의 그림처럼 생생하게 그려낸다. 시인은 ‘비가 온다’고 직접 말하는 대신, 비가 내릴 때 일어나는 여러 현상을 생생한 감각으로 보여준다.

첫 부분에서는 비가 막 시작되려는 순간의 서늘한 분위기를 시각과 촉각으로 전달한다. ‘돌에 그늘이 차고’라는 표현은 비구름 때문에 돌의 그림자가 짙어지는 모습을 보여주면서 동시에 서늘한 기운이 느껴지게 한다. 이어서 부는 ‘서늘한 바람’은 으스스하고 차가운 느낌을 더하며 비가 곧 쏟아질 것임을 암시한다.

본격적으로 비가 내리는 모습은 시각적 이미지를 통해 역동적으로 표현된다. 빗줄기는 ‘꼬리’를 세우고 땅에 떨어지며, ‘산새 걸음걸이’처럼 총총거리며 움직인다. 시냇물은 ‘수척한 흰 물살’이 되어 흐르고, 빗방울은 떨어지며 ‘손가락’을 펴는 듯한 모양으로 퍼져나간다.

마지막 부분에서는 빗줄기가 잠시 멎는가 싶더니 다시 세차게 쏟아지는 모습을 그린다. ‘멎은 듯하더니 다시 돋는 빗방울’은 빗방울이 다시 선명해지는 것을 보여주고, ‘붉은 잎’ 위로 떨어지는 빗소리는 ‘소란히’라는 청각적 표현을 통해 굵어진 빗줄기를 짐작하게 한다. 이처럼 시인은 다양한 감각을 이용하여 비 오는 풍경을 선명하게 그려낸다.
}


\kfPassageNote{해설 학습 2}{%
이 시의 가장 큰 특징은 비유와 의인화를 통해 비가 내리는 풍경을 살아 움직이는 생명체처럼 표현했다는 점이다. 시인은 눈에 보이지 않는 바람이나 형태가 없는 빗줄기에 구체적인 동물의 모습을 빗대어 생동감을 불어넣는다.

먼저 빗줄기가 힘차게 뻗어 나가는 모습을 마치 동물이 ‘꼬리’를 바짝 세운 것에 비유했다. 이는 가늘고 날카로운 빗줄기의 이미지를 효과적으로 보여준다. 또한 땅에 통통 튀며 떨어지는 빗방울의 모습은 ‘종종 다리 까칠한 산새 걸음걸이’에 빗대었다. 이를 통해 독자는 마치 작고 빠른 산새가 걸어가는 듯한 빗방울의 움직임을 머릿속에 그릴 수 있다.

물살의 모습 역시 비유를 통해 자세하게 표현된다. 빗물이 모여 이룬 물살이 바닥에 부딪혀 여러 갈래로 퍼져나가는 모습을 ‘갈래갈래 손가락을 펴는 듯’이라고 표현했다. 마치 물살이 살아있는 생명체처럼 손가락을 펼치는 듯한 인상을 준다.

마지막 연에서 비가 ‘붉은 잎 잎 소란히 밟고 간다’고 표현한 것은 이 시의 의인법이 가장 잘 드러나는 부분이다. 비를 사람인 것처럼 여겨, 마치 사람이 발걸음을 옮기듯 나뭇잎을 밟고 지나가는 존재로 묘사했다. 이러한 표현들은 비가 내리는 자연 현상을 하나의 생명체가 움직이는 역동적인 사건으로 바꾸어 놓으며 시 전체에 활기를 불어넣는다.
}


\kfPassageNote{해설 학습 3}{%
「비」는 매우 짧고 간결한 시이지만, 독자에게 풍부한 상상력을 불러일으킨다. 그 이유는 시인이 단어를 아끼고 문장 성분을 과감하게 생략하는 ‘절제’의 방식을 사용했기 때문이다. 시인은 주어나 목적어, 그리고 문장을 이어주는 접속 부사 등을 거의 사용하지 않는다. \\[14pt]
시는 ‘돌에 그늘이 차고,’처럼 명사나 동사 중심으로 짧게 끊어서 제시된다. ‘비가 오기 때문에’ 돌에 그늘이 차가워지고 ‘차가운’ 바람이 분다는 식의 친절한 설명을 모두 생략한다. 대신 ‘그늘이 차고,’ ‘서늘한 바람’과 같이 핵심적인 이미지만을 나열하여 독자가 그 사이의 빈 공간을 스스로 채우도록 유도한다. \\[14pt]
이러한 생략과 절제는 시에 속도감과 긴장감을 부여한다. 예를 들어 ‘앞서거니 뒤서거니 \\[4pt]
꼬리 휙 치켜들고’와 ‘종종 다리 까칠한 \\[4pt]
산새 걸음걸이’는 마치 영상의 한 장면 한 장면이 빠르게 지나가는 듯한 느낌을 준다. 독자는 이 짧은 구절들을 연결하며 비가 내리는 전체 과정을 머릿속으로 재구성하게 된다. \\[14pt]
마지막 연의 ‘붉은 잎 잎 \\[4pt]
소란히 밟고 간다.’ 역시 주어인 ‘비가’가 생략되어 있다. 하지만 독자는 시 전체의 흐름을 통해 누가 붉은 잎을 밟고 가는지를 자연스럽게 알 수 있다. 이처럼 정지용 시인은 말을 아낌으로써 오히려 더 많은 것을 보여주는 언어의 마술을 보여준다. 짧은 형식 속에 깊은 여운을 남기는 것이 이 시의 진정한 매력이다.
}


\kfSecHeader{kfAreaLiterature}{문제}{}

\par\medskip
\kfBeginMC
\begin{kfQBlock}{01}{객관식}
\kfQStem{다음 시구는 비가 내리는 모습을 비유한 표현이다. 여기서 원관념은 무엇인가?}
\begin{kfChoices}
\kfChoice 꼬리
\kfChoice 산새
\kfChoice 바람
\kfChoice 빗줄기
\end{kfChoices}
\end{kfQBlock}
\begin{kfQBlock}{02}{객관식}
\kfQStem{다음 시구는 무엇이 종종걸음을 걷는 모습을 묘사한 것인가?}
\begin{kfChoices}
\kfChoice 산새가 걷는 모습
\kfChoice 바람이 부는 모습
\kfChoice 빗줄기가 떨어지는 모습
\kfChoice 시냇물이 흘러가는 모습
\end{kfChoices}
\end{kfQBlock}
\begin{kfQBlock}{03}{객관식}
\kfQStem{다음 시구에 사용된 비유 표현에서 ‘원관념’과 ‘보조 관념’을 올바르게 짝지은 것은?}
\begin{kfChoices}
\kfChoice 원관념: 산새 / 보조 관념: 빗방울
\kfChoice 원관념: 비가 내리는 모습 / 보조 관념: 산새 걸음걸이
\kfChoice 원관념: 바람 / 보조 관념: 산새 걸음걸이
\kfChoice 원관념: 빗소리 / 보조 관념: 산새 소리
\end{kfChoices}
\end{kfQBlock}
\begin{kfQBlock}{04}{객관식}
\kfQStem{다음 시구에 사용된 비유 표현에서 ‘원관념’과 ‘보조 관념’을 올바르게 짝지은 것은?}
\begin{kfChoices}
\kfChoice 원관념: 물살 / 보조 관념: 손가락
\kfChoice 원관념: 빗줄기 / 보조 관념: 손가락
\kfChoice 원관념: 손가락 / 보조 관념: 물살
\kfChoice 원관념: 나뭇가지 / 보조 관념: 손가락
\end{kfChoices}
\end{kfQBlock}
\begin{kfQBlock}{05}{객관식}
\kfQStem{다음 문장에서 서술어 ‘돋는’의 생략된 주어는 무엇인가?}
\begin{kfChoices}
\kfChoice 구름이
\kfChoice 햇살이
\kfChoice 바람이
\kfChoice 빗방울이
\end{kfChoices}
\end{kfQBlock}
\begin{kfQBlock}{06}{객관식}
\kfQStem{다음 문장에서 서술어 ‘밟고 간다’의 생략된 주어는 무엇인가?}
\begin{kfChoices}
\kfChoice 잎이
\kfChoice 비가
\kfChoice 발이
\kfChoice 님이
\end{kfChoices}
\end{kfQBlock}
\begin{kfQBlock}{07}{객관식}
\kfQStem{다음 문장에서 생략된 성분을 넣어 문장을 복원할 때 가장 적절한 것은?}
\begin{kfChoices}
\kfChoice 바람이 갈래갈래 손가락을 펴는 듯하다.
\kfChoice 산새가 갈래갈래 손가락을 펴는 듯하다.
\kfChoice 물살이 갈래갈래 손가락을 펴는 듯하다.
\kfChoice 햇살이 갈래갈래 손가락을 펴는 듯하다.
\end{kfChoices}
\end{kfQBlock}
\begin{kfQBlock}{08}{객관식}
\kfQStem{이 시는 비 오는 날의 풍경을 무엇처럼 생생하게 그려내는가?}
\kfAnswerLines{2}
\end{kfQBlock}
\begin{kfQBlock}{09}{객관식}
\kfQStem{‘돌에 그늘이 차고’라는 표현은 어떤 감각을 통해 분위기를 전달하는가?}
\kfAnswerLines{2}
\end{kfQBlock}
\begin{kfQBlock}{10}{객관식}
\kfQStem{‘서늘한 바람’은 어떤 느낌을 주는 시어인가?}
\kfAnswerLines{2}
\end{kfQBlock}
\begin{kfQBlock}{11}{객관식}
\kfQStem{땅에 떨어지는 빗방울의 움직임을 무엇에 빗대어 표현했는가?}
\kfAnswerLines{2}
\end{kfQBlock}
\begin{kfQBlock}{12}{객관식}
\kfQStem{시냇물이 된 빗물은 어떤 모습으로 흐르는가?}
\kfAnswerLines{2}
\end{kfQBlock}
\begin{kfQBlock}{13}{객관식}
\kfQStem{빗방울이 물 위로 퍼지는 모습을 무엇이 펴지는 것에 비유했는가?}
\kfAnswerLines{2}
\end{kfQBlock}
\begin{kfQBlock}{14}{객관식}
\kfQStem{‘소란히’라는 표현은 어떤 감각적 이미지에 해당하는가?}
\kfAnswerLines{2}
\end{kfQBlock}
\begin{kfQBlock}{15}{객관식}
\kfQStem{이 시에서 ‘소란히’라는 표현을 통해 무엇을 짐작할 수 있는가?}
\kfAnswerLines{2}
\end{kfQBlock}
\begin{kfQBlock}{16}{객관식}
\kfQStem{이 시는 어떤 표현 방법을 통해 비를 살아 움직이는 생명체처럼 만드는가?}
\kfAnswerLines{2}
\end{kfQBlock}
\begin{kfQBlock}{17}{객관식}
\kfQStem{힘차게 뻗어 나가는 빗줄기를 동물의 무엇에 비유했는가?}
\kfAnswerLines{2}
\end{kfQBlock}
\begin{kfQBlock}{18}{객관식}
\kfQStem{땅에 떨어지는 빗방울의 모습을 어떤 동물의 걸음걸이에 비유했는가?}
\kfAnswerLines{2}
\end{kfQBlock}
\begin{kfQBlock}{19}{객관식}
\kfQStem{‘산새 걸음걸이’라는 비유는 빗방울의 어떤 움직임을 표현하는가?}
\kfAnswerLines{2}
\end{kfQBlock}
\begin{kfQBlock}{20}{객관식}
\kfQStem{물살이 여러 갈래로 퍼져나가는 모습을 신체의 어느 부분이 펴지는 것에 비유했는가?}
\kfAnswerLines{2}
\end{kfQBlock}
\begin{kfQBlock}{21}{객관식}
\kfQStem{이 시에서 인격이 부여된 자연물은 무엇인가?}
\kfAnswerLines{2}
\end{kfQBlock}
\begin{kfQBlock}{22}{객관식}
\kfQStem{비가 나뭇잎을 때리는 모습을 사람이 무엇을 하는 행동에 비유했는가?}
\kfAnswerLines{2}
\end{kfQBlock}
\begin{kfQBlock}{23}{객관식}
\kfQStem{이러한 비유와 의인화는 시 전체에 어떤 효과를 주는가?}
\kfAnswerLines{2}
\end{kfQBlock}
\begin{kfQBlock}{24}{객관식}
\kfQStem{비유란 표현하려는 대상을 다른 대상에 빗대어 표현하는 방법이다. 그렇다면 의인화란 무엇인가?}
\kfAnswerLines{2}
\end{kfQBlock}
\begin{kfQBlock}{25}{객관식}
\kfQStem{이 시에서 동물의 이미지를 활용하여 표현한 대상은 무엇인가?}
\kfAnswerLines{2}
\end{kfQBlock}
\begin{kfQBlock}{26}{객관식}
\kfQStem{이 시가 독자에게 풍부한 상상력을 불러일으키는 이유는 어떤 표현 방식 때문인가?}
\kfAnswerLines{2}
\end{kfQBlock}
\begin{kfQBlock}{27}{객관식}
\kfQStem{이 시에서 주로 생략된 문장 성분은 무엇인가?}
\kfAnswerLines{2}
\end{kfQBlock}
\begin{kfQBlock}{28}{객관식}
\kfQStem{시인이 친절한 설명을 생략하고 핵심적인 이미지만을 나열하는 것은 독자의 무엇을 유도하기 위함인가?}
\kfAnswerLines{2}
\end{kfQBlock}
\begin{kfQBlock}{29}{객관식}
\kfQStem{이 시의 생략과 절제는 어떤 효과를 만들어 내는가?}
\kfAnswerLines{2}
\end{kfQBlock}
\begin{kfQBlock}{30}{객관식}
\kfQStem{시는 마치 무엇의 한 장면들이 빠르게 지나가는 듯한 느낌을 주는가?}
\kfAnswerLines{2}
\end{kfQBlock}
\begin{kfQBlock}{31}{객관식}
\kfQStem{독자는 짧게 끊어진 구절들을 연결하며 무엇을 하게 되는가?}
\kfAnswerLines{2}
\end{kfQBlock}
\begin{kfQBlock}{32}{객관식}
\kfQStem{마지막 연에서 생략된 주어는 무엇인가?}
\kfAnswerLines{2}
\end{kfQBlock}
\begin{kfQBlock}{33}{객관식}
\kfQStem{시인이 단어를 아껴 쓰는 것을 무엇이라고 표현하는가?}
\kfAnswerLines{2}
\end{kfQBlock}
\begin{kfQBlock}{34}{객관식}
\kfQStem{이 시는 짧은 형식 속에 깊은 무엇을 남기는가?}
\kfAnswerLines{2}
\end{kfQBlock}
\begin{kfQBlock}{35}{객관식}
\kfQStem{정지용 시인은 말을 아낌으로써 오히려 더 많은 것을 보여주는 무엇을 보여주는가?}
\kfAnswerLines{2}
\end{kfQBlock}
\begin{kfQBlock}{36}{객관식}
\kfQStem{이 시에 대한 설명으로 적절하지 \uline{\underline{않은}} 것은?}
\begin{kfChoices}
\kfChoice 비가 내리는 풍경을 감각적으로 묘사한다.
\kfChoice 비를 동물의 움직임에 빗대어 표현한다.
\kfChoice 시각·청각·촉각 등 여러 감각을 함께 활용한다.
\kfChoice 비가 시작해서 굵어지는 과정을 시간 흐름에 따라 보여 준다.
\kfChoice 화자의 깊은 슬픔과 절망적 정서를 직접적으로 토로한다.
\end{kfChoices}
\end{kfQBlock}
\begin{kfQBlock}{37}{객관식}
\kfQStem{이 시의 분위기에 대한 설명으로 적절하지 \uline{\underline{않은}} 것은?}
\begin{kfChoices}
\kfChoice 비를 살아 움직이는 듯한 이미지로 그려 생동감을 만든다.
\kfChoice 동적인 비유로 활기찬 분위기를 형성한다.
\kfChoice 변화하는 자연 풍경을 역동적으로 표현한다.
\kfChoice 산뜻하고 가벼운 어조가 시 전체에 깔린다.
\kfChoice 깊은 쓸쓸함·외로움이 시 전체를 무겁게 누른다.
\end{kfChoices}
\end{kfQBlock}
\begin{kfQBlock}{38}{객관식}
\kfQStem{㉠과 같이 대상을 다른 사물에 빗대어 표현하는 방법의 효과로 적절하지 \uline{\underline{않은}} 것은?}
\begin{kfChoices}
\kfChoice 대상을 구체적이고 생생하게 느끼게 한다.
\kfChoice 추상적 현상을 친근한 이미지로 떠올리게 한다.
\kfChoice 독자의 상상력을 자극해 시상을 입체적으로 만든다.
\kfChoice 시 전체의 분위기를 풍부하게 만들어 준다.
\kfChoice 말하고자 하는 내용을 숨겨 긴장감을 유발한다.
\end{kfChoices}
\end{kfQBlock}
\begin{kfQBlock}{39}{객관식}
\kfQStem{이 시의 표현상 특징으로 적절하지 \uline{\underline{않은}} 것은?}
\begin{kfChoices}
\kfChoice 간결한 시어를 사용한다.
\kfChoice 다양한 비유적 표현이 나타난다.
\kfChoice 시간의 흐름에 따라 시상을 전개한다.
\kfChoice 문장 성분을 생략하여 여운을 남긴다.
\kfChoice 영탄적인 어조로 화자의 감정을 드러낸다.
\end{kfChoices}
\end{kfQBlock}
\begin{kfQBlock}{40}{객관식}
\kfQStem{시의 흐름으로 보아, 시간 순서대로 가장 마지막에 해당하는 연은?}
\begin{kfChoices}
\kfChoice 1연
\kfChoice 2연
\kfChoice 3연
\kfChoice 4연
\kfChoice 8연
\end{kfChoices}
\end{kfQBlock}
\begin{kfQBlock}{41}{객관식}
\kfQStem{<보기>를 참고하여 이 시를 감상한 내용으로 적절하지 \uline{\underline{않은}} 것은?}
\begin{kfEvidence}정지용 시인은 언어를 섬세하게 다듬어 회화적인 이미지를 만들어 내는 데 뛰어난 시인이다. 그는 단어의 소리나 색채, 질감 등을 활용하여 독자가 시를 읽으며 마치 한 폭의 그림을 감상하는 듯한 느낌을 받게 한다.\end{kfEvidence}
\begin{kfChoices}
\kfChoice ‘그늘이 차고’에서 비 오기 전의 어두운 색채와 차가운 느낌이 동시에 느껴져.
\kfChoice ‘서늘한 바람’이라는 표현에서 비가 오기 직전의 서늘한 촉감이 느껴지는 것 같아.
\kfChoice ‘수척한 흰 물살’에서 가늘고 하얀 물줄기가 흘러가는 모습이 눈앞에 그려지는군.
\kfChoice ‘붉은 잎’과 ‘흰 물살’의 색채 대비를 통해 비 내리는 풍경을 더욱 선명하게 만들었어.
\kfChoice ‘밟고 간다’는 표현에서 인생의 시련과 고난을 꿋꿋이 이겨내려는 의지가 느껴져.
\end{kfChoices}
\end{kfQBlock}
\begin{kfQBlock}{42}{서술형}
\kfQStem{이 시에서 비가 내리는 모습을 비유적으로 표현한 구절을 두 가지 찾아 쓰시오.}
\kfAnswerNote{답안}{4}
\end{kfQBlock}
\begin{kfQBlock}{43}{서술형}
\kfQStem{이 시에 나타난 감각적 이미지를 두 가지 이상 찾아 쓰고, 각각 어떤 감각에 해당하는지 서술하시오.}
\kfAnswerNote{답안}{4}
\end{kfQBlock}
\kfEndMC


\kfStartNonlit{[과학] 방사성 동위원소를 이용한 암석 연대 측정법}


\kfPassageNote{[과학] 방사성 동위원소를 이용한 암석 연대 측정법}{%
　19세기 초 지질학자들은 전 세계의 지질학적 연구 결과를 종합하여 '중생대 쥐라기'와 같은 지질학적 시간 척도를 만들었다. 그러나 이런 척도는 한 지층이 다른 지층보다 오래되었는지는 알 수 있어도 정확히 얼마나 오래되었는지는 알 수 없었다.

　1905년 러더포드가 방사성 동위원소를 이용하여 지층 연대 측정에 성공했다. 그는 암석 속 우라늄의 양을 측정하여 암석의 연대를 계산해냈다. 이것이 동위원소 연대측정법의 시작이었다.

　지질학자들은 방사성 동위원소의 어떤 특성을 활용했을까? 원자 중심에는 양성자와 중성자로 이루어진 원자핵이 있다. 원자핵의 양성자 수에 따라 원소의 종류가 결정된다. 같은 원소라도 중성자 수가 다른 것들이 있는데 이를 동위원소라 한다. 탄소-12는 양성자 6개와 중성자 6개가 있고, 탄소-14는 양성자 6개와 중성자 8개가 있다.

　㉠동위원소 중에는 원자핵이 불안정한 것들이 있다. 이런 불안정한 원자핵은 스스로 방사선을 방출하고 안정된 상태로 변하는데, 이를 방사성 붕괴라 한다. 방사성 붕괴를 일으키는 동위원소를 방사성 동위원소라 한다. 탄소-14는 방사성 붕괴로 질소-14가 된다. 붕괴 전을 모원소, 붕괴 후를 자원소라고 한다.

　방사성 동위원소는 일정한 시간이 지나면 모원소의 개수가 절반으로 줄어드는 특성이 있다. 이 시간을 반감기라 한다. 첫 반감기 때 모원소는 절반으로 줄고, 두 번째 반감기에는 1/4로, 세 번째 반감기에는 1/8로 줄어든다. 모원소와 자원소의 비율은 첫 반감기에 1:1, 두 번째 반감기에 1:3, 세 번째 반감기에 1:7이 된다. 탄소-14는 5730년, 포타슘-40은 13억 년, 우라늄-238은 44억 년의 반감기를 갖는다.

　방사성 동위원소의 반감기는 온도나 압력에 영향을 받지 않는다. 따라서 암석에 포함된 모원소와 자원소의 비율을 알고 반감기를 이용하면 암석이 만들어진 연대를 추정할 수 있다. 예를 들어 모원소와 자원소의 비율이 1:3이라면 반감기를 두 번 거쳤으므로, 포타슘-40의 경우 26억 년 전에 생성되었다고 볼 수 있다.
}


\kfPassageNote{문장 독해 — 본문}{%
1. 19세기 초 지질학자들은 전 세계의 지질학적 연구 결과를 종합하여 '중생대 쥐라기'와 같은 지질학적 시간 척도를 만들었다. 2. 그러나 이런 척도는 한 지층이 다른 지층보다 오래되었는지는 알 수 있어도 정확히 얼마나 오래되었는지는 알 수 없었다. 3. 1905년 러더포드가 방사성 동위원소를 이용하여 지층 연대 측정에 성공했다. 4. 그는 암석 속 우라늄의 양을 측정하여 암석의 연대를 계산해냈다. 5. 이것이 동위원소 연대측정법의 시작이었다. 6. 지질학자들은 방사성 동위원소의 어떤 특성을 활용했을까? 7. 원자 중심에는 양성자와 중성자로 이루어진 원자핵이 있다. 8. 원자핵의 양성자 수에 따라 원소의 종류가 결정된다. 9. 같은 원소라도 중성자 수가 다른 것들이 있는데 이를 동위원소라 한다. 10. 탄소-12는 양성자 6개와 중성자 6개가 있고, 탄소-14는 양성자 6개와 중성자 8개가 있다. 11. 동위원소 중에는 원자핵이 불안정한 것들이 있다. 12. 이런 불안정한 원자핵은 스스로 방사선을 방출하고 안정된 상태로 변하는데, 이를 방사성 붕괴라 한다. 13. 방사성 붕괴를 일으키는 동위원소를 방사성 동위원소라 한다. 14. 탄소-14는 방사성 붕괴로 질소-14가 된다. 15. 붕괴 전을 모원소, 붕괴 후를 자원소라고 한다. 16. 방사성 동위원소는 일정한 시간이 지나면 모원소의 개수가 절반으로 줄어드는 특성이 있다. 17. 이 시간을 반감기라 한다. 18. 첫 반감기 때 모원소는 절반으로 줄고, 두 번째 반감기에는 1/4로, 세 번째 반감기에는 1/8로 줄어든다. 19. 모원소와 자원소의 비율은 첫 반감기에 1:1, 두 번째 반감기에 1:3, 세 번째 반감기에 1:7이 된다. 20. 탄소-14는 5730년, 포타슘-40은 13억 년, 우라늄-238은 44억 년의 반감기를 갖는다. 21. 방사성 동위원소의 반감기는 온도나 압력에 영향을 받지 않는다. 22. 따라서 암석에 포함된 모원소와 자원소의 비율을 알고 반감기를 이용하면 암석이 만들어진 연대를 추정할 수 있다. 23. 예를 들어 모원소와 자원소의 비율이 1:3이라면 반감기를 두 번 거쳤으므로, 포타슘-40의 경우 26억 년 전에 생성되었다고 볼 수 있다.
}


\kfActivityBg \par\kfMaybeNeedspace{32\baselineskip}\noindent{\kfTipMark\,\,}{\small\color{kfInk} 다음 뜻풀이에 해당하는 어휘를 지문에서 찾아 쓰시오.}\par\nopagebreak[4]\smallskip\nopagebreak[4]
\begin{kfVocab}
\kfVocabItem{1}{지질학적 사건들을 시간 순서대로 정리한 표}
\kfVocabItem{2}{원자핵이 불안정하여 방사선을 내면서 붕괴하는 동위원소}
\kfVocabItem{3}{방사성 동위원소를 이용해 암석 등의 절대 연대를 측정하는 방법}
\kfVocabItem{4}{원자핵을 구성하는 양(+)전하를 띤 입자}
\kfVocabItem{5}{원자핵을 구성하는 전하가 없는 입자}
\kfVocabItem{6}{양성자 수는 같지만 중성자 수가 다른 원소}
\kfVocabItem{7}{불안정한 원자핵이 방사선을 방출하며 안정한 원자핵으로 변하는 현상}
\kfVocabItem{8}{방사성 붕괴 전의 처음 원소}
\kfVocabItem{9}{방사성 붕괴 후에 새로 생성된 안정된 원소}
\kfVocabItem{10}{방사성 동위원소의 원자 수가 원래의 절반으로 줄어드는 데 걸리는 시간}
\end{kfVocab}

\kfActivityBg \par\kfMaybeNeedspace{32\baselineskip}\noindent{\kfTipMark\,\,}{\small\color{kfInk} 다음 글의 구조를 정리한 표의 빈칸을 채워 보세요.}\par\nopagebreak[4]\smallskip\nopagebreak[4]
\par\medskip\needspace{8\baselineskip}
\begin{tcolorbox}[enhanced, breakable=true,
  colback=kfPaper, colframe=kfRule, boxrule=0.4pt, arc=2pt,
  left=8pt, right=8pt, top=6pt, bottom=6pt,
  title={\footnotesize\bfseries\color{kfMute} 글의 구조도},
  coltitle=kfMute, colbacktitle=kfPrimaryLight!40,
  attach boxed title to top left={yshift=-1.5mm, xshift=6pt},
  boxed title style={colframe=kfRule, boxrule=0.3pt, arc=1pt}]
\footnotesize\setstretch{1.3}
\renewcommand{\arraystretch}{1.35}
\arrayrulecolor{kfRule}
\noindent\begin{tabularx}{\linewidth}{|>{\centering\arraybackslash}p{1.6cm}|>{\raggedright\arraybackslash}p{3.4cm}|>{\raggedright\arraybackslash}X|}
\hline
\rowcolor{kfPrimaryLight!40}\textbf{\color{kfPrimary}단} & \textbf{\color{kfPrimary}항목} & \textbf{\color{kfPrimary}내용} \\\hline
한계: 초기의 ( ① )은(는) 정확한 연대를 알 수 없는 상대적 척도였음 & 한계: 초기의 ( ① )은(는) 정확한 연대를 알 수 없는 상대적 척도였음 &  \\\hline
발전: 러더포드가 ( ② )을(를) 이용하여 암석의 절대 연대를 계산하는 ( ③ )을(를) 시작함 & 발전: 러더포드가 ( ② )을(를) 이용하여 암석의 절대 연대를 계산하는 ( ③ )을(를) 시작함 &  \\\hline
원소 결정: 원자핵 속 ( ④ ) 수에 따라 결정됨 & 원소 결정: 원자핵 속 ( ④ ) 수에 따라 결정됨 &  \\\hline
( ⑤ ): 양성자 수는 같지만 ( ⑥ ) 수가 다른 원소 & ( ⑤ ): 양성자 수는 같지만 ( ⑥ ) 수가 다른 원소 &  \\\hline
정의: 불안정한 원자핵이 방사선을 방출하며 ( ⑦ ) 상태로 변하는 과정 & 정의: 불안정한 원자핵이 방사선을 방출하며 ( ⑦ ) 상태로 변하는 과정 &  \\\hline
용어: 붕괴 전 원소는 ‘( ⑧ )’, 붕괴 후 원소는 ‘( ⑨ )’ & 용어: 붕괴 전 원소는 ‘( ⑧ )’, 붕괴 후 원소는 ‘( ⑨ )’ &  \\\hline
정의: ( ⑩ )의 개수가 원래의 절반으로 줄어드는 데 걸리는 시간 & 정의: ( ⑩ )의 개수가 원래의 절반으로 줄어드는 데 걸리는 시간 &  \\\hline
\multirow{3}{*}{비율 변화 (모원소:자원소):} & 1회 후 & ( ⑪ ) \\\cline{2-3}
 & 2회 후 & ( ⑫ ) \\\cline{2-3}
 & 3회 후 & ( ⑬ ) \\\cline{2-3}
\hline
방법: 암석 속 모원소와 자원소의 ( ⑭ )과(와) 원소의 ( ⑮ )을(를) 이용해 연대 추정 & 방법: 암석 속 모원소와 자원소의 ( ⑭ )과(와) 원소의 ( ⑮ )을(를) 이용해 연대 추정 &  \\\hline
특징: 반감기는 ( ⑯ )이나 ( ⑰ )에 영향을 받지 않음 & 특징: 반감기는 ( ⑯ )이나 ( ⑰ )에 영향을 받지 않음 &  \\\hline
\end{tabularx}
\arrayrulecolor{black}
\end{tcolorbox}\par\medskip
\kfSecHeader{kfAreaNonlit}{문제}{}

\par\medskip
\kfBeginMC
\begin{kfQBlock}{01}{객관식}
\kfQStem{'지질학적 시간 척도'는 언제 만들어졌는가?}
\begin{kfChoices}
\kfChoice 19세기 초
\kfChoice 18세기 말
\kfChoice 20세기 중반
\kfChoice 21세기 초
\end{kfChoices}
\end{kfQBlock}
\begin{kfQBlock}{02}{객관식}
\kfQStem{지질학자들은 무엇을 종합하여 지질학적 시간 척도를 만들었는가?}
\begin{kfChoices}
\kfChoice 암석의 종류
\kfChoice 화석 발견 기록
\kfChoice 기후 변화 기록
\kfChoice 전 세계의 지질학적 연구 결과
\end{kfChoices}
\end{kfQBlock}
\begin{kfQBlock}{03}{객관식}
\kfQStem{초기의 지질학적 시간 척도가 알려줄 수 있었던 것은 무엇인가?}
\begin{kfChoices}
\kfChoice 지층의 성분
\kfChoice 화석의 종류
\kfChoice 지층의 절대 연령
\kfChoice 한 지층이 다른 지층보다 오래되었는지 여부
\end{kfChoices}
\end{kfQBlock}
\begin{kfQBlock}{04}{객관식}
\kfQStem{이 척도가 가진 한계점은 무엇인가?}
\begin{kfChoices}
\kfChoice 지층의 순서를 알 수 없었다는 점
\kfChoice 암석의 종류를 알 수 없었다는 점
\kfChoice 화석의 연대를 알 수 없었다는 점
\kfChoice 정확히 얼마나 오래되었는지는 알 수 없었다는 점
\end{kfChoices}
\end{kfQBlock}
\begin{kfQBlock}{05}{객관식}
\kfQStem{1905년에 러더포드가 성공한 것은 무엇인가?}
\begin{kfChoices}
\kfChoice 지층 연대 측정
\kfChoice 원소 주기율표 완성
\kfChoice 방사능의 발견
\kfChoice 새로운 지질학적 척도 생성
\end{kfChoices}
\end{kfQBlock}
\begin{kfQBlock}{06}{객관식}
\kfQStem{그는 무엇을 이용하여 연대 측정에 성공했는가?}
\begin{kfChoices}
\kfChoice 지진파
\kfChoice 화석 분석
\kfChoice 방사성 동위원소
\kfChoice 지자기 역전 현상
\end{kfChoices}
\end{kfQBlock}
\begin{kfQBlock}{07}{객관식}
\kfQStem{러더포드는 어떤 원소의 양을 측정하여 암석의 연대를 계산했는가?}
\begin{kfChoices}
\kfChoice 탄소
\kfChoice 질소
\kfChoice 우라늄
\kfChoice 포타슘
\end{kfChoices}
\end{kfQBlock}
\begin{kfQBlock}{08}{객관식}
\kfQStem{'이것'이 가리키는 것은 무엇인가?}
\begin{kfChoices}
\kfChoice 우라늄의 양을 측정하는 것
\kfChoice 지질학적 척도를 수정하는 것
\kfChoice 새로운 동위원소를 발견하는 것
\kfChoice 방사성 동위원소를 이용하여 지층 연대를 측정하는 것
\end{kfChoices}
\end{kfQBlock}
\begin{kfQBlock}{09}{객관식}
\kfQStem{이 문장은 앞으로 어떤 내용이 이어질 것을 알려주는가?}
\begin{kfChoices}
\kfChoice 지층 연대 측정의 역사
\kfChoice 방사성 동위원소의 발견 과정
\kfChoice 지질학자들이 발견한 새로운 화석
\kfChoice 지질학자들이 활용한 방사성 동위원소의 특성
\end{kfChoices}
\end{kfQBlock}
\begin{kfQBlock}{10}{객관식}
\kfQStem{원자핵은 어디에 위치하는가?}
\begin{kfChoices}
\kfChoice 원자 껍질
\kfChoice 전자 궤도
\kfChoice 원자 중심
\kfChoice 원자 외부
\end{kfChoices}
\end{kfQBlock}
\begin{kfQBlock}{11}{객관식}
\kfQStem{원자 중심에 있는 '원자핵'은 무엇과 무엇으로 이루어져 있는가?}
\begin{kfChoices}
\kfChoice 전자와 원자핵
\kfChoice 양성자와 전자
\kfChoice 중성자와 전자
\kfChoice 양성자와 중성자
\end{kfChoices}
\end{kfQBlock}
\begin{kfQBlock}{12}{객관식}
\kfQStem{원소의 종류를 결정하는 것은 무엇의 수인가?}
\begin{kfChoices}
\kfChoice 전자 수
\kfChoice 중성자 수
\kfChoice 양성자 수
\kfChoice 원자량
\end{kfChoices}
\end{kfQBlock}
\begin{kfQBlock}{13}{객관식}
\kfQStem{'동위원소'란 무엇인가?}
\begin{kfChoices}
\kfChoice 중성자 수가 다른 같은 원소
\kfChoice 양성자 수가 다른 같은 원소
\kfChoice 전자 수가 다른 같은 원소
\kfChoice 질량수가 같은 다른 원소
\end{kfChoices}
\end{kfQBlock}
\begin{kfQBlock}{14}{객관식}
\kfQStem{동위원소들은 양성자 수와 중성자 수 중 무엇의 수가 \uline{다른}가?}
\begin{kfChoices}
\kfChoice 전자 수
\kfChoice 중성자 수
\kfChoice 양성자 수
\kfChoice 원자가전자 수
\end{kfChoices}
\end{kfQBlock}
\begin{kfQBlock}{15}{객관식}
\kfQStem{'탄소-12'는 양성자와 중성자를 각각 몇 개 가지고 있는가?}
\begin{kfChoices}
\kfChoice 양성자 6개, 중성자 6개
\kfChoice 양성자 6개, 중성자 7개
\kfChoice 양성자 6개, 중성자 8개
\kfChoice 양성자 8개, 중성자 6개
\end{kfChoices}
\end{kfQBlock}
\begin{kfQBlock}{16}{객관식}
\kfQStem{'탄소-14'는 양성자와 중성자를 각각 몇 개 가지고 있는가?}
\begin{kfChoices}
\kfChoice 양성자 6개, 중성자 6개
\kfChoice 양성자 6개, 중성자 8개
\kfChoice 양성자 8개, 중성자 6개
\kfChoice 양성자 8개, 중성자 8개
\end{kfChoices}
\end{kfQBlock}
\begin{kfQBlock}{17}{객관식}
\kfQStem{'탄소-12'와 '탄소-14'의 공통점은 무엇인가?}
\begin{kfChoices}
\kfChoice 전자 수가 같다.
\kfChoice 중성자 수가 같다.
\kfChoice 원자량이 같다.
\kfChoice 양성자 수가 6개로 같다.
\end{kfChoices}
\end{kfQBlock}
\begin{kfQBlock}{18}{객관식}
\kfQStem{'탄소-12'와 '탄소-14'의 차이점은 무엇인가?}
\begin{kfChoices}
\kfChoice 전자 수가 다르다.
\kfChoice 중성자 수가 다르다.
\kfChoice 양성자 수가 다르다.
\kfChoice 화학적 성질이 다르다.
\end{kfChoices}
\end{kfQBlock}
\begin{kfQBlock}{19}{객관식}
\kfQStem{동위원소 중에는 어떤 상태의 원자핵을 가진 것들이 있는가?}
\begin{kfChoices}
\kfChoice 안정한 상태
\kfChoice 불안정한 상태
\kfChoice 변화하지 않는 상태
\kfChoice 에너지가 없는 상태
\end{kfChoices}
\end{kfQBlock}
\begin{kfQBlock}{20}{객관식}
\kfQStem{'방사성 붕괴'란 무엇인가?}
\begin{kfChoices}
\kfChoice 원자가 결합하여 분자가 되는 것
\kfChoice 원자핵이 융합하여 더 큰 핵이 되는 것
\kfChoice 전자를 잃고 양이온이 되는 것
\kfChoice 불안정한 원자핵이 방사선을 방출하며 안정된 상태로 변하는 것
\end{kfChoices}
\end{kfQBlock}
\begin{kfQBlock}{21}{객관식}
\kfQStem{방사성 붕괴 과정에서 원자핵은 무엇을 방출하는가?}
\begin{kfChoices}
\kfChoice 전파
\kfChoice 초음파
\kfChoice 방사선
\kfChoice 가시광선
\end{kfChoices}
\end{kfQBlock}
\begin{kfQBlock}{22}{객관식}
\kfQStem{'방사성 동위원소'란 무엇인가?}
\begin{kfChoices}
\kfChoice 전기를 띠는 동위원소
\kfChoice 빛을 내는 동위원소
\kfChoice 화학 반응을 잘하는 동위원소
\kfChoice 방사성 붕괴를 일으키는 동위원소
\end{kfChoices}
\end{kfQBlock}
\begin{kfQBlock}{23}{객관식}
\kfQStem{'탄소-14'는 방사성 붕괴를 통해 어떤 원소로 변하는가?}
\begin{kfChoices}
\kfChoice 산소-14
\kfChoice 질소-14
\kfChoice 탄소-12
\kfChoice 탄소-13
\end{kfChoices}
\end{kfQBlock}
\begin{kfQBlock}{24}{객관식}
\kfQStem{'모원소'란 무엇인가?}
\begin{kfChoices}
\kfChoice 붕괴 중인 원소
\kfChoice 붕괴 후의 안정된 원소
\kfChoice 붕괴 전의 방사성 동위원소
\kfChoice 붕괴와 상관없는 원소
\end{kfChoices}
\end{kfQBlock}
\begin{kfQBlock}{25}{객관식}
\kfQStem{'자원소'란 무엇인가?}
\begin{kfChoices}
\kfChoice 붕괴 전의 원소
\kfChoice 붕괴 중인 원소
\kfChoice 붕괴 후의 안정된 원소
\kfChoice 새로 발견된 원소
\end{kfChoices}
\end{kfQBlock}
\begin{kfQBlock}{26}{객관식}
\kfQStem{방사성 동위원소가 가진 중요한 특성은 무엇인가?}
\begin{kfChoices}
\kfChoice 일정한 시간이 지나면 크기가 커지는 특성
\kfChoice 일정한 시간이 지나면 무게가 늘어나는 특성
\kfChoice 일정한 시간이 지나면 개수가 절반으로 줄어드는 특성
\kfChoice 일정한 시간이 지나면 색깔이 변하는 특성
\end{kfChoices}
\end{kfQBlock}
\begin{kfQBlock}{27}{객관식}
\kfQStem{이 문장에서 '이 시간'이 가리키는 것은 무엇인가?}
\begin{kfChoices}
\kfChoice 모원소가 모두 사라지는 시간
\kfChoice 자원소가 모원소로 변하는 시간
\kfChoice 암석이 생성되는 데 걸리는 시간
\kfChoice 모원소의 개수가 원래 개수의 절반으로 줄어드는 시간
\end{kfChoices}
\end{kfQBlock}
\begin{kfQBlock}{28}{객관식}
\kfQStem{모원소의 개수가 절반으로 줄어드는 데 걸리는 시간을 무엇이라고 하는가?}
\begin{kfChoices}
\kfChoice 반감기
\kfChoice 붕괴기
\kfChoice 생성기
\kfChoice 안정기
\end{kfChoices}
\end{kfQBlock}
\begin{kfQBlock}{29}{객관식}
\kfQStem{반감기를 두 번 거치면 모원소의 양은 처음의 몇 분의 몇으로 줄어드는가?}
\begin{kfChoices}
\kfChoice 1/2
\kfChoice 1/3
\kfChoice 1/4
\kfChoice 1/8
\end{kfChoices}
\end{kfQBlock}
\begin{kfQBlock}{30}{객관식}
\kfQStem{반감기를 세 번 거치면 모원소의 양은 처음의 몇 분의 몇으로 줄어드는가?}
\begin{kfChoices}
\kfChoice 1/4
\kfChoice 1/6
\kfChoice 1/8
\kfChoice 1/16
\end{kfChoices}
\end{kfQBlock}
\begin{kfQBlock}{31}{객관식}
\kfQStem{첫 번째 반감기가 지났을 때 모원소와 자원소의 비율은 어떻게 되는가?}
\begin{kfChoices}
\kfChoice 1 : 1
\kfChoice 1 : 2
\kfChoice 1 : 3
\kfChoice 1 : 4
\end{kfChoices}
\end{kfQBlock}
\begin{kfQBlock}{32}{객관식}
\kfQStem{세 번째 반감기가 지났을 때 모원소와 자원소의 비율은 어떻게 되는가?}
\begin{kfChoices}
\kfChoice 1 : 1
\kfChoice 1 : 3
\kfChoice 1 : 7
\kfChoice 1 : 15
\end{kfChoices}
\end{kfQBlock}
\begin{kfQBlock}{33}{객관식}
\kfQStem{이 문장에서 소개된 원소 중 반감기가 가장 긴 것은 무엇인가?}
\begin{kfChoices}
\kfChoice 탄소-14
\kfChoice 질소-14
\kfChoice 포타슘-40
\kfChoice 우라늄-238
\end{kfChoices}
\end{kfQBlock}
\begin{kfQBlock}{34}{객관식}
\kfQStem{이 문장에서 소개된 원소 중 반감기가 가장 짧은 것은 무엇인가?}
\begin{kfChoices}
\kfChoice 탄소-14
\kfChoice 질소-14
\kfChoice 포타슘-40
\kfChoice 우라늄-238
\end{kfChoices}
\end{kfQBlock}
\begin{kfQBlock}{35}{객관식}
\kfQStem{이 문장을 통해 알 수 있는, '반감기'의 중요한 특징은 무엇인가?}
\begin{kfChoices}
\kfChoice 모든 원소의 반감기는 같다.
\kfChoice 원소에 따라 반감기가 다르다.
\kfChoice 반감기는 시간이 지날수록 길어진다.
\kfChoice 반감기는 환경에 따라 변한다.
\end{kfChoices}
\end{kfQBlock}
\begin{kfQBlock}{36}{객관식}
\kfQStem{반감기가 연대 측정에 유용한 이유는 무엇인가?}
\begin{kfChoices}
\kfChoice 암석의 종류에 따라 달라지기 때문이다.
\kfChoice 측정이 간편하고 비용이 적게 들기 때문이다.
\kfChoice 온도나 압력 같은 외부 요인에 영향을 받지 않기 때문이다.
\kfChoice 모든 암석에 방사성 동위원소가 들어있기 때문이다.
\end{kfChoices}
\end{kfQBlock}
\begin{kfQBlock}{37}{객관식}
\kfQStem{암석의 연대를 추정하기 위해 알아야 할 정보 두 가지는 무엇인가?}
\begin{kfChoices}
\kfChoice 암석의 무게와 부피
\kfChoice 모원소와 자원소의 비율, 반감기
\kfChoice 암석 발견 장소와 깊이
\kfChoice 지층의 두께와 색깔
\end{kfChoices}
\end{kfQBlock}
\begin{kfQBlock}{38}{객관식}
\kfQStem{이 두 가지 정보를 이용해 무엇을 추정할 수 있는가?}
\begin{kfChoices}
\kfChoice 암석의 성분
\kfChoice 암석의 용도
\kfChoice 암석의 가격
\kfChoice 암석이 만들어진 연대
\end{kfChoices}
\end{kfQBlock}
\begin{kfQBlock}{39}{객관식}
\kfQStem{모원소와 자원소의 비율이 1:3일 때, 이는 반감기를 몇 번 거쳤다는 의미인가?}
\begin{kfChoices}
\kfChoice 한 번
\kfChoice 두 번
\kfChoice 세 번
\kfChoice 네 번
\end{kfChoices}
\end{kfQBlock}
\begin{kfQBlock}{40}{객관식}
\kfQStem{예시에서 암석의 나이를 26억 년으로 추정한 근거는 무엇인가?}
\begin{kfChoices}
\kfChoice 우라늄-238의 반감기를 이용했기 때문이다.
\kfChoice 탄소-14의 반감기를 이용했기 때문이다.
\kfChoice 포타슘-40의 반감기(13억 년)를 두 번 거쳤기 때문이다.
\kfChoice 모원소와 자원소의 비율이 1:1이었기 때문이다.
\end{kfChoices}
\end{kfQBlock}
\begin{kfQBlock}{41}{객관식}
\kfQStem{윗글의 내용과 일치하는 것은?}
\begin{kfChoices}
\kfChoice 반감기를 두 번 거친 암석의 모원소와 자원소 비율은 1:4이다.
\kfChoice 방사성 붕괴 후 생성된 안정된 원소를 모원소라고 한다.
\kfChoice 동위원소는 양성자 수는 다르지만 중성자 수는 같은 원소를 말한다.
\kfChoice 방사성 붕괴는 불안정한 원자핵이 안정한 상태로 변해가는 과정이다.
\kfChoice 러더포드는 탄소-14를 이용하여 암석의 연대를 최초로 계산했다.
\end{kfChoices}
\end{kfQBlock}
\begin{kfQBlock}{42}{객관식}
\kfQStem{㉠동위원소에 대한 설명으로 적절하지 \uline{\underline{않은}} 것은?}
\begin{kfChoices}
\kfChoice 같은 원소이지만 원자핵의 중성자 수가 서로 다르다.
\kfChoice 양성자 수는 같으므로 같은 원소로 분류된다.
\kfChoice 탄소-12와 탄소-14는 양성자 수가 같은 동위원소이다.
\kfChoice 일부 동위원소는 불안정해 방사성 붕괴를 일으킨다.
\kfChoice 양성자 수가 다르고 중성자 수는 같은 원소들의 모임이다.
\end{kfChoices}
\end{kfQBlock}
\begin{kfQBlock}{43}{객관식}
\kfQStem{윗글을 바탕으로 <보기>의 암석 나이를 계산한 것으로 옳은 것은?}
\begin{kfEvidence}어떤 화성암 속에 방사성 동위원소 X가 포함되어 있었다. X의 반감기는 4억 년이며, 현재 이 암석 속 X와 이것이 붕괴하여 생성된 Y의 비율은 1 : 7이다.\end{kfEvidence}
\begin{kfChoices}
\kfChoice 4억 년
\kfChoice 8억 년
\kfChoice 12억 년
\kfChoice 16억 년
\kfChoice 28억 년
\end{kfChoices}
\end{kfQBlock}
\begin{kfQBlock}{44}{객관식}
\kfQStem{연대 측정 방법에 대한 설명으로 적절하지 \uline{\underline{않은}} 것은?}
\begin{kfChoices}
\kfChoice 층서 원리는 지층의 생성 순서를 알려주는 상대적 척도이다.
\kfChoice 러더포드는 우라늄의 방사성 붕괴를 이용하여 암석의 나이를 계산했다.
\kfChoice 동위원소 연대측정법은 암석 생성 이후 외부로부터 원소 유입이 없었다고 가정한다.
\kfChoice 반감기가 긴 '우라늄-238'은 수천 년 전 유물의 연대를 측정하는 데 매우 유용하다.
\kfChoice 암석에 포함된 모원소와 자원소의 양을 정확히 측정하는 것이 중요하다.
\end{kfChoices}
\end{kfQBlock}
\begin{kfQBlock}{45}{객관식}
\kfQStem{윗글의 서술 방식에 대한 설명으로 적절하지 \uline{\underline{않은}} 것은?}
\begin{kfChoices}
\kfChoice 핵심이 되는 개념을 먼저 정의한 뒤에 그 원리를 자세히 설명한다.
\kfChoice 탄소-14, 포타슘-40 같은 구체적인 사례를 들어 독자의 이해를 돕는다.
\kfChoice 동위원소·방사성 붕괴·반감기 같은 개념을 단계별로 차근차근 설명한다.
\kfChoice 개념을 먼저 제시하고 원리, 사례의 순서로 전체 글을 전개한다.
\kfChoice 가설을 세운 뒤 직접 실험을 통해 그 타당성을 증명하는 형식이다.
\end{kfChoices}
\end{kfQBlock}
\begin{kfQBlock}{46}{단답형}
\kfQStem{19세기 지질학적 시간 척도의 한계는 무엇이었는가?}
\kfAnswerLines{1}
\end{kfQBlock}
\begin{kfQBlock}{47}{단답형}
\kfQStem{1905년 방사성 동위원소를 이용해 암석의 절대 연령 측정에 성공한 과학자는 누구인가?}
\kfAnswerLines{1}
\end{kfQBlock}
\begin{kfQBlock}{48}{단답형}
\kfQStem{양성자 수는 같지만 중성자 수가 다른 같은 원소의 원자들을 무엇이라고 하는가?}
\kfAnswerLines{1}
\end{kfQBlock}
\begin{kfQBlock}{49}{단답형}
\kfQStem{불안정한 원자핵이 방사선을 방출하며 안정된 상태로 변하는 과정은 무엇인가?}
\kfAnswerLines{1}
\end{kfQBlock}
\begin{kfQBlock}{50}{단답형}
\kfQStem{방사성 붕괴 전의 불안정한 동위원소를 무엇이라고 부르는가?}
\kfAnswerLines{1}
\end{kfQBlock}
\begin{kfQBlock}{51}{서술형}
\kfQStem{동위원소 연대측정법을 이용해 암석의 절대 연령을 알 수 있는 원리를 <조건>에 맞게 서술하시오.}
\kfAnswerNote{답안}{4}
\end{kfQBlock}
\kfEndMC



\clearpage

\kfChapterCover{4}{러셀3 (챕터4)}{Chapter 4}{이번 챕터의 학습 내용을 시작합니다.}
\begin{kfChapterAreaList}
\kfChapterAreaItem{문법}{kfAreaGrammar}{단어의 변신, 품사 통용의 이해}
\kfChapterAreaItem{개념}{kfAreaConcept}{여성 화자를 통한 정서 표현}
\kfChapterAreaItem{문학}{kfAreaLiterature}{「사미인곡」}
\kfChapterAreaItem{비문학}{kfAreaNonlit}{[기술] 항원-항체 반응을 이용한 검사용 키트의 원리}
\end{kfChapterAreaList}

\kfStartGrammar{단어의 변신, 품사 통용의 이해}


\kfPassageNoteNT{%
　우리는 9품사를 배우면서 단어들을 성질에 따라 분류했다. 하지만 어떤 단어들은 하나의 품사로만 고정되지 않고, 문장 속에서 맡은 역할에 따라 여러 품사의 성격을 동시에 지니기도 한다. 예를 들어 '오늘'이라는 단어가 어떤 문장에서는 명사로, 다른 문장에서는 부사로 쓰이는 것이 그 예이다. 이처럼 하나의 단어가 둘 이상의 품사로 쓰이는 현상을 ‘품사 통용(通用)’ 이라고 한다. 품사 통용 단어의 품사를 구별하는 것은 문장 구조를 정확히 분석하는 데 매우 중요하다.

　품사를 구별하는 가장 중요한 기준은 그 단어가 문장 속에서 어떤 기능을 하는지, 특히 조사와 결합할 수 있는지 또는 다른 말을 꾸며 주는지를 확인하는 것이다.

　가장 대표적인 예는 수(數)를 나타내는 말이다. "사과 다섯이 남았다."라는 문장에서 '다섯'은 뒤에 주격 조사 '이'가 붙었다. 조사와 결합할 수 있는 말은 체언이므로, 이때 '다섯'은 수사이다. 하지만 "사과 다섯 개가 남았다."에서는 '다섯'이 뒤에 오는 단위 의존 명사 '개'를 꾸며주고 있다. 체언을 꾸미는 것은 관형사의 역할이므로, 이때 '다섯'은 관형사이다.

　지시를 나타내는 '이, 그, 저'도 마찬가지이다. "이는 내가 가장 아끼는 물건이다."에서 '이'는 보조사 '는'과 결합했으므로 대명사이다. 반면 "이 물건은 내가 가장 아낀다."에서 '이'는 뒤따르는 명사 '물건'을 꾸며주는 관형사로 쓰였다.

　시간을 나타내는 일부 명사들도 품사 통용을 보인다. "오늘은 정말 잊을 수 없는 날이다."에서 '오늘'은 보조사 '은'과 결합하였으므로 명사이다. 하지만 "나는 오늘 그를 만났다."에서 '오늘'은 뒤따르는 서술어 '만났다'를 꾸며주는 부사로 기능한다.

　형태는 완전히 같지만 문맥에 따라 동사와 형용사로 나뉘는 경우도 있다. "길이 참 밝다."의 '밝다'는 '밝는다'로 쓸 수 없고 상태를 나타내므로 형용사이다. 하지만 "어느덧 날이 밝다."의 '밝다'는 '밝는다'처럼 현재형 어미와 결합할 수 있고 동작의 의미를 가지므로 동사이다. 마찬가지로 "그는 키가 크다."는 형용사이지만, "아이가 쑥쑥 큰다."는 동사이다.

　띄어쓰기와도 밀접하게 관련된 품사 통용이 있다. '만큼', '뿐', '대로'와 같은 말들은 앞에 체언이 오는지 용언이 오는지에 따라 품사가 달라진다. "너만큼 잘할 수 있다."에서 '만큼'은 대명사 '너' 뒤에 바로 붙었으므로 조사이다. 하지만 "아는 만큼 보인다."에서 '만큼'은 용언의 관형사형 '아는' 뒤에 쓰였으므로, 띄어 써야 하는 의존 명사이다.

　이처럼 한 단어의 품사는 그 형태만으로 결정되지 않는다. 반드시 문장 안에서 조사와 결합했는지, 어떤 말을 꾸며주는지, 서술어로 쓰일 때 '-(ㄴ)는다'와 결합할 수 있는지 등 그 문법적 기능을 꼼꼼히 따져 보아야 한다.
\par\medskip\begin{itemize}[leftmargin=18pt, itemsep=2pt, topsep=2pt, label=\textbullet]
\item 수(수사·관형사) — '다섯이 남았다(수사)' vs '다섯 개가(관형사)'
\item 밝다(형·동) — '길이 밝다(형)' vs '날이 밝는다(동)'
\item 만큼(조·의존명사) — '너만큼(조사, 붙임)' vs '아는 만큼(의존 명사, 띄움)'
\end{itemize}
}

\kfActivityBg \par\kfMaybeNeedspace{32\baselineskip}지문의 핵심 내용을 떠올리며 아래 빈칸에 알맞은 말을 채워 보세요.

1. 품사 통용의 개념 

\par\medskip\needspace{4\baselineskip}
\noindent\begin{adjustbox}{max width=\linewidth, center}
\renewcommand{\arraystretch}{1.45}
\arrayrulecolor{kfRule}
\begin{tabular}{|>{\raggedright\arraybackslash}p{\dimexpr 0.1586\linewidth-12pt\relax}|>{\raggedright\arraybackslash}p{\dimexpr 0.8414\linewidth-12pt\relax}|}
\hline
\rowcolor{kfPrimaryLight!50}\textbf{\color{kfPrimary} 구분} & \textbf{\color{kfPrimary} 내용} \\\hline
품사 [ ① ] (通用) & ∙ 하나의 단어가 [ ② ] 이상의 품사로 쓰이는 현상. <br>(예: '오늘'이 명사로도, 부사로도 쓰임) \\\hline
중요 구별 기준 & ∙ 문장 안에서 어떤 [ ③ ]을/를 하는지, 특히 [ ④ ]와/과 결합하는지, 다른 말을 꾸미는지 등을 확인해야 함. \\\hline
\end{tabular}
\arrayrulecolor{black}
\end{adjustbox}\par\medskip



2. 품사 통용의 주요 사례 및 구별 기준 

\par\medskip\needspace{4\baselineskip}
\noindent\begin{adjustbox}{max width=\linewidth, center}
\renewcommand{\arraystretch}{1.45}
\arrayrulecolor{kfRule}
\begin{tabular}{|>{\raggedright\arraybackslash}p{\dimexpr 0.0891\linewidth-12pt\relax}|>{\raggedright\arraybackslash}p{\dimexpr 0.0365\linewidth-12pt\relax}|>{\raggedright\arraybackslash}p{\dimexpr 0.3343\linewidth-12pt\relax}|>{\raggedright\arraybackslash}p{\dimexpr 0.0365\linewidth-12pt\relax}|>{\raggedright\arraybackslash}p{\dimexpr 0.5036\linewidth-12pt\relax}|}
\hline
\rowcolor{kfPrimaryLight!50}\textbf{\color{kfPrimary} 사례} & \textbf{\color{kfPrimary} 품사 1} & \textbf{\color{kfPrimary} 기준 1 (예시)} & \textbf{\color{kfPrimary} 품사 2} & \textbf{\color{kfPrimary} 기준 2 (예시)} \\\hline
수(數) 표현 & [ ⑤ ] & ∙  [ ⑥ ] 와/과 결합함.<br>∙ (예: "사과 다섯이 남았다.") & [ ⑦ ] & ∙ 뒤의  체언(명사) 을 꾸밈.<br>∙ (예: "다섯 개가 남았다.") \\\hline
지시 표현 & [ ⑧ ] & ∙ [ ⑨ ]와/과 결합함.<br>∙ (예: "이는 내가 아낀다.") & [ ⑩ ] & ∙ 뒤의  체언(명사)을 꾸밈.<br>∙ (예: "이 물건은 아낀다.") \\\hline
시간 표현<br> & [ ⑪ ] & ∙ [ ⑫ ]와/과 결합함.<br>∙ (예: "오늘은 잊을 수 없다.") & [ ⑬ ] & ∙ 뒤의  서술어(용언)를 꾸밈.<br>∙ (예: "오늘 만났다.") \\\hline
형태가 같은 용언 & [ ⑭ ] & ∙ 상태를 나타냄.<br>∙ '-(ㄴ)는다' 결합 불가능.<br>∙ (예: "길이 밝다." (X 밝는다)) & [ ⑮ ] & ∙ 동작/변화를 나타냄.<br>∙ '-(ㄴ)는다' 결합 가능.<br>∙ (예: "날이 밝는다." (O)) \\\hline
'-만큼', '-뿐', '-대로' & [ ⑯ ] & ∙ 체언 뒤에 붙여 씀.<br>∙ (예: "너만큼  잘한다.") & [ ⑰ ] & ∙ 용언의 관형사형 뒤에 띄어 씀.<br>∙ (예: "아는 만큼 보인다.") \\\hline
\end{tabular}
\arrayrulecolor{black}
\end{adjustbox}\par\medskip\nopagebreak[4]\par\nopagebreak[4]

\kfActivityBg \kfSecHeader{kfAreaGrammar}{활동}{분석 훈련}
\par\kfMaybeNeedspace{24\baselineskip}\noindent{\kfTipMark\,\,}{\small\color{kfInk} 다음 활용 사례를 분석하여 빈칸을 채워 보세요.}\par\nopagebreak[4]\smallskip\nopagebreak[4]
\begin{enumerate}[leftmargin=22pt, itemsep=6pt, topsep=4pt, label=\textbf{\arabic*.}]
\item 밑줄 친 단어의 품사를 쓰고, 그렇게 판단한 이유를 간략히 서술하시오.
\begin{kfEvidence}(가) \uline{셋}이 모였다. / (나) \uline{세} 사람이 모였다.\end{kfEvidence}
\item 밑줄 친 단어의 품사를 쓰고, 그렇게 판단한 이유를 간략히 서술하시오.
\begin{kfEvidence}(가) \uline{그}는 학생이다. / (나) \uline{그} 학생은 착하다.\end{kfEvidence}
\item 밑줄 친 단어의 품사를 쓰고, 그렇게 판단한 이유를 간략히 서술하시오.
\begin{kfEvidence}(가) \uline{내일}은 축제일이다. / (나) 우리는 \uline{내일} 만난다.\end{kfEvidence}
\item 밑줄 친 단어의 품사를 쓰고, 그렇게 판단한 이유를 간략히 서술하시오.
\begin{kfEvidence}(가) 아이가 쑥쑥 \uline{큰다}. / (나) 아이의 키가 \uline{크다}.\end{kfEvidence}
\item 밑줄 친 단어의 품사를 쓰고, 그렇게 판단한 이유를 간략히 서술하시오.
\begin{kfEvidence}(가) 법\uline{대로} 해라. / (나) 네가 말한 \uline{대로} 해라.\end{kfEvidence}
\item 밑줄 친 단어의 품사를 쓰고, 그렇게 판단한 이유를 간략히 서술하시오.
\begin{kfEvidence}(가) 그림이 \uline{밝다}. / (나) 날이 \uline{밝는다}.\end{kfEvidence}
\item 밑줄 친 단어의 품사를 쓰고, 그렇게 판단한 이유를 간략히 서술하시오.
\begin{kfEvidence}(가) \uline{이}것이 정답이다. / (나) \uline{이} 정답은 틀렸다.\end{kfEvidence}
\item 밑줄 친 단어의 품사를 쓰고, 그렇게 판단한 이유를 간략히 서술하시오.
\begin{kfEvidence}(가) 노력할 \uline{뿐}이다. / (나) 가진 것은 너\uline{뿐}이다.\end{kfEvidence}
\end{enumerate}

\kfSecHeader{kfAreaGrammar}{문제}{}

\par\medskip
\kfBeginMC
\begin{kfQBlock}{01}{객관식}
\kfQStem{품사 통용에 대한 설명으로 적절하지 \uline{\underline{않은}} 것은?}
\begin{kfChoices}
\kfChoice 한 단어가 문장에서의 기능에 따라 다른 품사로 쓰일 수 있다.
\kfChoice ‘오늘’은 명사·부사 두 품사로 쓰일 수 있는 대표적 예이다.
\kfChoice 품사 통용을 판단하려면 문장에서의 쓰임을 함께 살펴야 한다.
\kfChoice 단어 형태만 보고 품사를 단정해서는 안 된다.
\kfChoice 한 단어는 사전에 등재된 하나의 품사로만 쓰일 수 있다.
\end{kfChoices}
\end{kfQBlock}
\begin{kfQBlock}{02}{객관식}
\kfQStem{<보기>의 밑줄 친 단어의 품사를 구별하는 기준으로 가장 적절한 것은?}
\begin{kfEvidence}∙  \uline{다섯}이 왔다.   \\\relax
   ∙  \uline{다섯} 사람이 왔다.\end{kfEvidence}
\begin{kfChoices}
\kfChoice 서술어와 결합하는 방식을 비교해 본다.
\kfChoice 조사와의 결합 가능 여부와 수식 대상을 본다.
\kfChoice 현재형 어미를 사용할 수 있는지 따져 본다.
\kfChoice 띄어쓰기를 어떻게 했는지를 기준으로 삼는다.
\kfChoice 단어가 지닌 고유한 의미만을 비교해 본다.
\end{kfChoices}
\end{kfQBlock}
\begin{kfQBlock}{03}{객관식}
\kfQStem{<보기>의 ㉠, ㉡에 해당하는 품사를 바르게 짝지은 것은?}
\begin{kfEvidence}㉠ \uline{이것}은 내 책이다. \\\relax
㉡ \uline{이} 책은 내 것이다.\end{kfEvidence}
\begin{kfChoices}
\kfChoice ㉠ 대명사 / ㉡ 관형사
\kfChoice ㉠ 관형사 / ㉡ 대명사
\kfChoice ㉠ 수사 / ㉡ 관형사
\kfChoice ㉠ 관형사 / ㉡ 수사
\kfChoice ㉠ 대명사 / ㉡ 수사
\end{kfChoices}
\end{kfQBlock}
\begin{kfQBlock}{04}{객관식}
\kfQStem{다음 밑줄 친 '오늘'의 품사가 나머지와 \uline{다른} 하나는?}
\begin{kfChoices}
\kfChoice \uline{오늘}의 날씨는 맑다.
\kfChoice \uline{오늘}도 여전히 춥다.
\kfChoice \uline{오늘} 그가 떠났다.
\kfChoice \uline{오늘}은 기분이 좋다.
\kfChoice \uline{오늘}부터 시작이다.
\end{kfChoices}
\end{kfQBlock}
\begin{kfQBlock}{05}{객관식}
\kfQStem{다음 밑줄 친 단어 중 품사가 관형사인 것은?}
\begin{kfChoices}
\kfChoice \uline{하나}만 주세요.
\kfChoice \uline{첫째}로 도착했다.
\kfChoice \uline{저기}는 공사 중이다.
\kfChoice \uline{저} 건물이 우리 학교다.
\kfChoice \uline{모두}가 잠든 깊은 밤.
\end{kfChoices}
\end{kfQBlock}
\begin{kfQBlock}{06}{객관식}
\kfQStem{<보기>의 ㉠과 ㉡의 품사를 바르게 구별한 것은?}
\begin{kfEvidence}㉠ 아이가 쑥쑥 \uline{큰다}. \\\relax
㉡ 내 동생은 키가 \uline{크다}.\end{kfEvidence}
\begin{kfChoices}
\kfChoice ㉠: 동사, ㉡: 동사
\kfChoice ㉠: 형용사, ㉡: 형용사
\kfChoice ㉠: 동사, ㉡: 형용사
\kfChoice ㉠: 형용사, ㉡: 동사
\kfChoice ㉠: 부사, ㉡: 관형사
\end{kfChoices}
\end{kfQBlock}
\begin{kfQBlock}{07}{객관식}
\kfQStem{<보기>의 단어 중, 품사 통용에 해당하지 \uline{않는} 것끼리 묶인 것은?}
\begin{kfEvidence}㉠ 하늘 / \\\relax
㉡ 셋 / \\\relax
㉢ 이 / \\\relax
㉣ 밝다 / \\\relax
㉤ 바다\end{kfEvidence}
\begin{kfChoices}
\kfChoice ㉠, ㉡
\kfChoice ㉠, ㉤
\kfChoice ㉡, ㉢
\kfChoice ㉢, ㉣
\kfChoice ㉡, ㉢, ㉣
\end{kfChoices}
\end{kfQBlock}
\begin{kfQBlock}{08}{객관식}
\kfQStem{다음 밑줄 친 단어의 품사가 나머지와 \uline{다른} 하나는?}
\begin{kfChoices}
\kfChoice 아는 \uline{만큼} 보인다.
\kfChoice 먹을 \uline{뿐}이다.
\kfChoice 약속한 \uline{대로} 되었다.
\kfChoice 하나\uline{뿐}인 내 편.
\kfChoice 볼 \uline{사람}이 없다.
\end{kfChoices}
\end{kfQBlock}
\begin{kfQBlock}{09}{객관식}
\kfQStem{다음 문장의 띄어쓰기가 \uline{잘못된} 것은?}
\begin{kfChoices}
\kfChoice 너 만큼 나도 할 수 있다.
\kfChoice 법대로 처리하자.
\kfChoice 아는 만큼 말해라.
\kfChoice 너뿐만 아니라 나도 간다.
\kfChoice 노력할 뿐이다.
\end{kfChoices}
\end{kfQBlock}
\begin{kfQBlock}{10}{객관식}
\kfQStem{다음 밑줄 친 단어의 품사가 부사인 것은?}
\begin{kfChoices}
\kfChoice \uline{내일}은 맑겠다.
\kfChoice \uline{내일} 보자.
\kfChoice \uline{저리} 가거라.
\kfChoice \uline{셋}이 왔다.
\kfChoice \uline{셋}째 줄에 앉아라.
\end{kfChoices}
\end{kfQBlock}
\begin{kfQBlock}{11}{객관식}
\kfQStem{밑줄 친 단어의 품사가 수사인 것은?}
\begin{kfChoices}
\kfChoice \uline{두} 마리
\kfChoice \uline{세} 사람
\kfChoice \uline{네} 명
\kfChoice 큰 \uline{다섯} 자루
\kfChoice \uline{하나}를 안다.
\end{kfChoices}
\end{kfQBlock}
\begin{kfQBlock}{12}{객관식}
\kfQStem{<보기>의 ㉠, ㉡에 들어갈 품사를 바르게 짝지은 것은?}
\begin{kfEvidence}"\uline{이} 사람은 \uline{셋째} 아들이다." \\\relax
㉠ '이'의 품사: \kfFillBlank[12mm]{} \\\relax
㉡ '셋째'의 품사: \kfFillBlank[12mm]{}\end{kfEvidence}
\begin{kfChoices}
\kfChoice ㉠ 대명사 / ㉡ 수사
\kfChoice ㉠ 대명사 / ㉡ 관형사
\kfChoice ㉠ 관형사 / ㉡ 수사
\kfChoice ㉠ 관형사 / ㉡ 관형사
\kfChoice ㉠ 수사 / ㉡ 수사
\end{kfChoices}
\end{kfQBlock}
\begin{kfQBlock}{13}{객관식}
\kfQStem{<보기>의 ㉠, ㉡에 대한 설명으로 옳은 것은?}
\begin{kfEvidence}㉠ 길이 \uline{밝다}. \\\relax
㉡ 날이 \uline{밝다}.\end{kfEvidence}
\begin{kfChoices}
\kfChoice ㉠, ㉡ 모두 형용사이다.
\kfChoice ㉠, ㉡ 모두 동사이다.
\kfChoice ㉠은 동사, ㉡은 형용사이다.
\kfChoice ㉠은 형용사, ㉡은 동사이다.
\kfChoice ㉠, ㉡은 모두 품사 통용의 예가 아니다.
\end{kfChoices}
\end{kfQBlock}
\begin{kfQBlock}{14}{객관식}
\kfQStem{밑줄 친 단어의 품사 구별이 \uline{잘못된} 것은?}
\begin{kfChoices}
\kfChoice \uline{그}가 왔다.
\kfChoice \uline{그} 집은 크다.
\kfChoice \uline{하나}만 안다.
\kfChoice \uline{한} 사람이 왔다.
\kfChoice \uline{오늘} 떠난다.
\end{kfChoices}
\end{kfQBlock}
\begin{kfQBlock}{15}{객관식}
\kfQStem{<보기>의 ㉠, ㉡에 들어갈 단어로 가장 적절한 것은?}
\begin{kfEvidence}∙ ㉠: 조사와 결합하여 '체언'으로 쓰임.   \\\relax
    ∙ ㉡: 뒤의 체언을 꾸며 '관형사'로 쓰임.\end{kfEvidence}
\begin{kfChoices}
\kfChoice ㉠:  이는 / ㉡:  이 사람
\kfChoice ㉠:  다섯 사람 / ㉡:  다섯이
\kfChoice ㉠:  오늘 만났다 / ㉡:  오늘의
\kfChoice ㉠:  크다 / ㉡:  큰 사람
\kfChoice ㉠:  너만큼 / ㉡:  아는 만큼
\end{kfChoices}
\end{kfQBlock}
\begin{kfQBlock}{16}{단답형}
\kfQStem{하나의 단어가 둘 이상의 품사로 쓰이는 현상을 무엇이라고 하는가?}
\kfAnswerLines{1}
\end{kfQBlock}
\begin{kfQBlock}{17}{단답형}
\kfQStem{품사 통용을 구별할 때, 단어 뒤에 붙을 수 있으면 '체언'으로 보는 문법 요소를 무엇이라고 하는가?}
\kfAnswerLines{1}
\end{kfQBlock}
\begin{kfQBlock}{18}{단답형}
\kfQStem{"사과  넷을 샀다."에서 '넷'의 품사는 무엇인가?}
\kfAnswerLines{1}
\end{kfQBlock}
\begin{kfQBlock}{19}{단답형}
\kfQStem{" 네 사람이 걸어간다."에서 '네'의 품사는 무엇인가?}
\kfAnswerLines{1}
\end{kfQBlock}
\begin{kfQBlock}{20}{단답형}
\kfQStem{" 오늘 우리는 학교에 간다."에서 '오늘'의 품사는 무엇인가?}
\kfAnswerLines{1}
\end{kfQBlock}
\begin{kfQBlock}{21}{단답형}
\kfQStem{" 오늘의 날씨는 맑다."에서 '오늘'의 품사는 무엇인가?}
\kfAnswerLines{1}
\end{kfQBlock}
\begin{kfQBlock}{22}{단답형}
\kfQStem{" 저는 학생입니다."에서 '저'의 품사는 무엇인가?}
\kfAnswerLines{1}
\end{kfQBlock}
\begin{kfQBlock}{23}{단답형}
\kfQStem{" 저 학생은 모범생이다."에서 '저'의 품사는 무엇인가?}
\kfAnswerLines{1}
\end{kfQBlock}
\begin{kfQBlock}{24}{단답형}
\kfQStem{"아이가 큰다."의 '크다'는 (동사/형용사)이다.}
\kfAnswerLines{1}
\end{kfQBlock}
\begin{kfQBlock}{25}{단답형}
\kfQStem{"키가 크다 ."의 '크다'는 (동사/형용사)이다.}
\kfAnswerLines{1}
\end{kfQBlock}
\begin{kfQBlock}{26}{서술형}
\kfQStem{<보기>의 밑줄 친 '다섯'의 품사가 각각 무엇인지 쓰고, 왜 그렇게 판단했는지 '조사'와 '수식'이라는 용어를 사용하여 서술하시오.}
\begin{kfEvidence}(가) 아이 \uline{다섯}이 놀고 있다.   \\\relax
(나) \uline{다섯} 아이가 놀고 있다.\end{kfEvidence}
\kfAnswerNote{답안}{4}
\end{kfQBlock}
\kfEndMC


\kfStartConcept{여성 화자를 통한 정서 표현}


\kfPassageNoteNT{%
　시를 읽을 때 우리는 작품 속에서 말하는 사람인 '화자'를 만나게 된다. 화자는 시인이 자신의 생각과 감정을 효과적으로 전달하기 위해 내세우는 목소리의 주인이다. 시인은 때로 어린아이, 노인, 남성, 여성 등 다양한 인물을 화자로 설정하여 시를 쓴다. 그중에서도 많은 시인이 자신의 성별과 관계없이 '여성 화자'를 등장시키는 경우가 있다. 이는 여성의 목소리를 통해 특정한 정서를 더욱 섬세하고 절실하게 표현할 수 있기 때문이다.

　이러한 여성 화자 설정은 고전 시가에서도 깊은 전통을 가지고 있다. 조선 시대의 남성 신하였던 정철은 「사미인곡」과 「속미인곡」에서 자신을 임과 헤어진 여성에 비유했다. 이 작품들의 여성 화자는 임, 곧 임금을 향한 변함없는 그리움과 충성심을 애절하게 노래한다. 이처럼 남성 사대부가 여성의 목소리를 빌려 임금에 대한 충성을 표현한 사례는, 이별의 정서를 다루는 여성 화자의 전통이 얼마나 효과적이었는지를 잘 보여준다.

　여성 화자는 전통적으로 '이별의 슬픔'이나 '기다림'의 정서를 표현하는 데 자주 등장한다. 대표적인 예로 김소월의 「진달래꽃」이 있다. 이 시의 화자는 "나 보기가 역겨워 가실 때에는 말없이 고이 보내 드리우리다"라고 말한다. 자신을 떠나는 임을 원망하는 대신, 묵묵히 이별을 받아들이고 꽃을 뿌려 축복하겠다는 이 목소리는 슬픔을 참고 견디는 모습을 보여준다. 이러한 설정을 통해 '이별의 정한'이라는 슬픔의 감정이 더욱 애절하게 전달된다.

　이처럼 시에서 여성 화자를 설정하는 것은 시의 주제를 형상화하는 중요한 방법이다. 남성 시인이라도 여성 화자의 목소리를 빌리면, 그리움이나 기다림, 슬픔과 같은 감정을 더욱 부드럽고 섬세하게 표현하여 독자의 공감을 얻기 쉽다. 따라서 독자는 시 속의 화자가 누구인지, 어떤 목소리를 내고 있는지 파악함으로써 작품의 정서와 주제를 더 깊이 있게 이해할 수 있다.
\par\medskip\begin{itemize}[leftmargin=18pt, itemsep=2pt, topsep=2pt, label=\textbullet]
\item 남성 시인의 여성 화자 — 정철 「사미인곡」·「속미인곡」(임=임금에 대한 충성)
\item 이별의 정한 — 김소월 「진달래꽃」의 여성 화자
\item 효과 — 그리움·슬픔·기다림을 섬세하고 애절하게 전달
\end{itemize}
}

\kfSecHeader{kfAreaConcept}{문제}{}

\par\medskip
\kfBeginMC
\begin{kfQBlock}{01}{객관식}
\kfQStem{윗글의 내용과 일치하지 \uline{않는} 것은?}
\begin{kfChoices}
\kfChoice 화자는 시 속에서 항상 말하는 사람이다.
\kfChoice 정철은 남성임에도 여성 화자를 설정했다.
\kfChoice 여성 화자는 그리움의 정서를 표현하는 데 쓰인다.
\kfChoice 시인은 반드시 자신의 성별과 같은 화자만 설정해야 한다.
\kfChoice 「진달래꽃」의 화자는 임과의 이별을 묵묵히 받아들이고 있다.
\end{kfChoices}
\end{kfQBlock}
\begin{kfQBlock}{02}{객관식}
\kfQStem{시인이 여성 화자를 설정하는 이유로 적절하지 \uline{\underline{않은}} 것은?}
\begin{kfChoices}
\kfChoice 정서를 부드럽고 섬세하게 표현하기 위해서이다.
\kfChoice 이별·기다림 같은 정서를 절실하게 형상화하기 위해서이다.
\kfChoice 한국 시가 전통의 ‘한’의 정서와 자연스럽게 어울리기 때문이다.
\kfChoice 독자의 공감을 끌어내기에 효과적이기 때문이다.
\kfChoice 독자에게 강한 교훈과 비판을 직접 전달하기 위해서이다.
\end{kfChoices}
\end{kfQBlock}
\begin{kfQBlock}{03}{객관식}
\kfQStem{정철이 「사미인곡」·「속미인곡」에서 여성 화자를 설정한 목적에 대한 설명으로 적절하지 \uline{\underline{않은}} 것은?}
\begin{kfChoices}
\kfChoice 자신을 ‘임과 헤어진 여성’에 비유한 표현 전략이다.
\kfChoice 임금을 향한 신하의 충성심을 절실하게 드러내기 위해서이다.
\kfChoice 정치적 좌절과 그리움을 부드러운 어조로 담아내기 위해서이다.
\kfChoice 충신연주지사 전통과 결합해 효과를 얻기 위해서이다.
\kfChoice 여성의 권리 신장을 직접적으로 주장하기 위해서이다.
\end{kfChoices}
\end{kfQBlock}
\begin{kfQBlock}{04}{객관식}
\kfQStem{<보기>의 화자가 느끼는 주된 정서로 알맞은 것은?}
\begin{kfEvidence}동짓달 기나긴 밤을 한가운데 잘라내어   \\\relax
   따뜻한 이불 밑에다 차곡차곡 넣어 뒀다   \\\relax
   사랑하는 님 오신 날 밤에 길게길게 펴리라.   \\\relax
   * 황진이 「동짓달 기나긴 밤」\end{kfEvidence}
\begin{kfChoices}
\kfChoice 분노
\kfChoice 후회
\kfChoice 절망
\kfChoice 체념
\kfChoice 그리움
\end{kfChoices}
\end{kfQBlock}
\begin{kfQBlock}{05}{객관식}
\kfQStem{<보기>의 화자에 대한 설명으로 적절하지 \uline{\underline{않은}} 것은?}
\begin{kfEvidence}임이여, 물을 건너지 마오.   \\\relax
   임은 결국 물을 건너시네.   \\\relax
   물에 빠져 죽었으니,   \\\relax
   장차 임을 어이할꼬.   \\\relax
   * 백수광부의 아내 「공무도하가」\end{kfEvidence}
\begin{kfChoices}
\kfChoice 임의 죽음을 슬퍼하고 있다.
\kfChoice 임과의 이별을 안타까워하고 있다.
\kfChoice 임을 말렸지만 끝내 막지 못했다.
\kfChoice 임에 대한 깊은 원망과 체념이 드러난다.
\kfChoice 임을 잊고 새로운 만남을 다짐하고 있다.
\end{kfChoices}
\end{kfQBlock}
\begin{kfQBlock}{06}{객관식}
\kfQStem{<보기>에 대한 설명으로 가장 적절한 것은?}
\begin{kfEvidence}아아, 님은 갔지마는 나는 님을 보내지   \\\relax
   아니하였습니다.   \\\relax
   제 곡조를 못 이기는 사랑의 노래는   \\\relax
   님의 침묵을 휩싸고 돕니다.   \\\relax
   * 한용운 「님의 침묵」\end{kfEvidence}
\begin{kfChoices}
\kfChoice 떠나간 임을 매우 깊이 원망하며 슬퍼하고 있다.
\kfChoice 임과의 재회를 사실상 포기하고 체념했다.
\kfChoice 임과의 이별을 운명으로 완전히 받아들였다.
\kfChoice 떠난 임에 대한 사랑이 식었음을 고백한다.
\kfChoice 임은 떠났지만 이별을 거부하며 재회를 다짐한다.
\end{kfChoices}
\end{kfQBlock}
\begin{kfQBlock}{07}{단답형}
\kfQStem{시인이 자신의 생각이나 감정을 효과적으로 전달하기 위해 내세우는 목소리의 주인은 무엇인가?}
\kfAnswerLines{1}
\end{kfQBlock}
\begin{kfQBlock}{08}{단답형}
\kfQStem{시인이 자신의 성별과 관계없이 특정한 정서를 섬세하고 절실하게 표현하기 위해 자주 설정하는 화자의 유형은 무엇인가?}
\kfAnswerLines{1}
\end{kfQBlock}
\begin{kfQBlock}{09}{단답형}
\kfQStem{조선 시대 남성 신하인 정철이 임금을 향한 그리움과 충성심을 여성 화자의 목소리로 표현한 대표적인 작품 두 가지를 쓰시오.}
\kfAnswerLines{1}
\end{kfQBlock}
\begin{kfQBlock}{10}{단답형}
\kfQStem{정철의 「사미인곡」과 「속미인곡」에서 여성 화자가 그리워하는 '임'은 실제 누구를 의미하는가?}
\kfAnswerLines{1}
\end{kfQBlock}
\begin{kfQBlock}{11}{단답형}
\kfQStem{여성 화자가 전통적으로 '이별의 슬픔'과 함께 자주 표현하는 정서는 무엇인가?}
\kfAnswerLines{1}
\end{kfQBlock}
\begin{kfQBlock}{12}{단답형}
\kfQStem{「진달래꽃」의 화자 설정을 통해 더욱 애절하게 전달되는 감정은 무엇인가?}
\kfAnswerLines{1}
\end{kfQBlock}
\begin{kfQBlock}{13}{단답형}
\kfQStem{시인이 자신의 성별과 달리 여성 화자를 설정함으로써 독자의 무엇을 얻기 쉬운가?}
\kfAnswerLines{1}
\end{kfQBlock}
\begin{kfQBlock}{14}{서술형}
\kfQStem{시인이 자신의 성별과 관계없이 '여성 화자'를 설정하는 이유를 서술하시오.}
\kfAnswerNote{답안}{4}
\end{kfQBlock}
\kfEndMC


\kfStartLit{「사미인곡」}


\kfPassageNote{}{%
이 몸이 태어날 때 임을 따라 태어났으니 \\[4pt]
이것은 평생 인연이며 하늘이 모를 리 없는 일. \\[4pt]
나는 오직 젊었고 님은 오직 나를 사랑하시니 \\[4pt]
이 마음 이 사랑을 비교할 곳이 전혀 없어라. \\[4pt]
평생의 소원은 님과 함께 살자 하는 것이었는데 \\[4pt]
늙어서 무슨 일로 홀로 두고 그리워하는가. \\[4pt]
엊그제는 님을 모셔 달나라 궁궐에 올랐더니 \\[4pt]
그동안 어찌하여 세상에 내려왔는가. \\[4pt]
님과 헤어질 때 빗은 머리 헝클어진 지 3년일세. \\[4pt]
연지 화장품이 있지만 누구를 위해 곱게 단장할까. \\[4pt]
마음에 맺힌 시름 겹겹이 쌓여 있어 \\[4pt]
나오는 것은 한숨이요 흐르는 것은 눈물이라. \\[4pt]
인생은 끝이 있는데 나의 시름은 끝이 없구나. \\[4pt]
무심한 세월은 물 흐르듯 하는구나. \\[4pt]
더위와 추위가 때를 알아 갔다가 다시 오니 \\[4pt]
듣고 보는 것마다 느낄 일이 많기도 많구나. \\[14pt]
봄바람이 문득 불어 쌓인 눈을 헤쳐 내니 \\[4pt]
창밖에 심어 둔 매화 두세 가지 피었구나. \\[4pt]
가뜩이나 날씨가 차가운데 은은한 향기는 무슨 일일까. \\[4pt]
황혼에 달이 따라와 베개맡에 비치니 \\[4pt]
흐느끼는 듯 반가워하는 듯, 혹시 님이신가 아니신가. \\[4pt]
저 매화 꺾어 내어 님 계신 곳에 보내고 싶어라. \\[4pt]
님이 너를 보고 어떻게 생각하실까. \\[14pt]
꽃 지고 새 잎 나니 푸른 나무 그늘이 깔렸는데 \\[4pt]
비단 휘장은 외롭고 수놓은 장막은 텅 비어 있네. \\[4pt]
연꽃 장막 걷어 놓고 공작 병풍 둘러두니 \\[4pt]
가뜩이나 시름도 많은데 날은 어찌 이리도 길까. \\[4pt]
원앙새 수놓은 비단을 잘라 놓고 오색실을 풀어 내어 \\[4pt]
금자로 재단해서 님의 옷을 지어 내니 \\[4pt]
솜씨는 물론이고 격식도 잘 갖추었구나. \\[4pt]
산호수 지게 위에 백옥 상자에 담아 두고 \\[4pt]
님에게 보내려고 님 계신 곳을 바라보니 \\[4pt]
산인지 구름인지, 멀고도 험하구나. \\[4pt]
천 리 만 리 머나먼 길을 누가 찾아갈 수 있을까. \\[4pt]
가거든 열어 보고 나인가 반가워하실까. \\[14pt]
하룻밤 서리가 내릴 때 기러기가 울며 날아갈 제 \\[4pt]
높은 누각에 혼자 올라 수정 발을 걷어 올리니 \\[4pt]
동산에 달이 뜨고 북극성이 보이니 \\[4pt]
님이신가 반가워하니 눈물이 저절로 난다. \\[4pt]
저 맑은 달빛을 끌어내어 님이 계신 궁궐에 부치고 싶어라. \\[4pt]
누각 위에 걸어 두고 온 세상에 다 비추어 \\[4pt]
깊은 산골짜기까지 대낮처럼 환하게 비춰 주소서. \\[14pt]
온 세상이 얼어붙어 흰 눈으로 한 빛깔일 때 \\[4pt]
사람은 물론이고 날아다니는 새도 끊겨 있구나. \\[4pt]
내가 있는 이곳 남쪽도 추위가 이 정도인데 \\[4pt]
님이 계신 높은 궁궐이야 말해 무엇하겠는가. \\[4pt]
따뜻한 봄 햇살을 보내 님 계신 곳에 쬐게 해 드리고 싶어라. \\[4pt]
내 초가집 처마에 비친 해를 님 계신 궁궐에 올리고 싶어라. \\[4pt]
붉은 치마를 여며 입고 푸른 소매를 반쯤 걷어 \\[4pt]
해 질 녘 대나무에 기대니 생각과 걱정이 많기도 많구나. \\[4pt]
짧은 해가 금방 지고 긴 밤을 꼿꼿이 앉아 \\[4pt]
푸른 등불 걸어 둔 곁에 공후를 놓아두고 \\[4pt]
꿈에서라도 님을 보려 턱을 괴고 기대어 있으니 \\[4pt]
원앙 이불도 차디차구나, 이 밤은 언제나 샐까. \\[14pt]
하루도 열두 때, 한 달도 서른 날 \\[4pt]
잠시라도 님 생각 말고 이 시름을 잊으려 하니 \\[4pt]
어느새 마음속에 맺혀서 뼛속까지 사무치니 \\[4pt]
전설의 명의가 열 명 온다 해도 이 병을 어찌할까. \\[4pt]
아아, 내 병은 바로 님 때문에 생긴 것이로다. \\[4pt]
차라리 죽어서 호랑나비가 되리라. \\[4pt]
꽃나무 가지마다 가는 곳마다 앉아 다니다가 \\[4pt]
향기 묻은 날개로 님의 옷에 내려앉으리라. \\[4pt]
님께서야 나인 줄 모르셔도 나는 그저 님을 따르려 하노라.
}


\kfPassageNote{문장 독해 — 발췌}{%
1. 나는 오직 젊었고 님은 오직 나를 사랑하시니 2. 엊그제는 님을 모셔 달나라 궁궐에 올랐더니 그동안 어찌하여 세상에 내려왔는가. 3. 님이 너를 보고 어떻게 생각하실까. 4. 금자로 재단해서 님의 옷을 지어 내니 솜씨는 물론이고 격식도 잘 갖추었구나. 5. 가거든 열어 보고 나인가 반가워하실까. 6. 깊은 산골짜기까지 대낮처럼 환하게 비춰 주소서. 7. 내가 있는 이곳 남쪽도 추위가 이 정도인데 님이 계신 높은 궁궐이야 말해 무엇하겠는가. 8. 아아, 내 병은 바로 님 때문에 생긴 것이로다.
}


\kfPassageNote{해설 학습 1}{%
「사미인곡」은 제목 그대로 ‘임을 그리워하는 노래’라는 뜻으로, 표면적으로는 사랑하는 임과 헤어진 한 여인의 슬픈 사랑 노래처럼 보이지만, 그 속에는 신하가 임금을 그리워하는 절절한 마음이 담겨 있다. 이렇게 남성인 작가가 일부러 여성 화자를 내세우는 방법을 사용하면, 임을 향한 사랑과 그리움이라는 감정을 훨씬 더 애절하고 절실하게 전달하는 효과가 있다.

이 작품에서 슬픔에 잠긴 ‘여성 화자’는 바로 작가인 정철 자신을 의미한다. 그리고 화자가 하늘이 정해준 인연이라 믿으며 애타게 그리워하는 ‘임’은 당시의 왕이었던 선조 임금을 상징한다. 시의 공간적 배경 또한 작가의 실제 상황에 빗대어 표현되었는데, 화자가 임과 헤어져 내려온 ‘세상’은 작가가 관직에서 물러나 머물던 전남 창평을, 임이 계신 ‘달나라 궁궐’은 임금이 계신 한양의 궁궐을 뜻한다.

정철은 당쟁으로 인해 관직에서 물러나 고향에 머무르게 되자, 임금을 향한 자신의 변함없는 충성심과 그리움을 전달하기 위해 이 작품을 지었다. 사랑하는 사람을 그리워하는 마음과 신하가 임금을 그리워하는 마음은 그 절실함이 비슷하기에, 여성의 목소리를 빌려 표현했을 때 독자들에게 더 큰 공감과 감동을 줄 수 있었다. 이처럼 임금을 그리워하는 신하의 마음을 ‘연군지정’이라고 하며, 이는 조선 시대 사대부 문학의 중요한 주제 중 하나였다.
}


\kfPassageNote{해설 학습 2}{%
「사미인곡」의 가장 큰 특징 중 하나는 계절의 변화(봄-여름-가을-겨울)에 따라 시상을 전개한다는 점이다. 화자는 계절이 바뀔 때마다 그 계절에 맞는 자연물이나 사물을 통해 임에 대한 자신의 사랑과 정성을 표현한다. 하지만 그 마음을 임에게 전달할 수 없는 안타까운 상황이 반복되면서 그리움의 정서는 점점 더 깊어진다.

봄에는 추위를 뚫고 피어나는 ‘매화’를 보며 임에 대한 변함없는 절개를 다짐하고, 그 향기를 임에게 보내고 싶어 한다. 여름에는 임을 위해 정성껏 ‘옷’을 짓지만, 험한 산과 구름에 가로막혀 보내지 못하고 안타까워한다. 가을 밤에는 밝고 맑은 ‘달빛’을 보며, 임금의 은혜가 온 세상을 비추기를 소망하는 마음을 드러낸다. 겨울에는 혹독한 추위 속에서 임이 춥지는 않을까 걱정하며 따뜻한 ‘봄 햇살’을 보내드리고 싶어 한다.

이처럼 화자는 계절마다 임을 생각하며 자신의 사랑과 충성심을 각기 다른 방식으로 표현한다. 시간이 흘러도 변하지 않는, 오히려 깊어지기만 하는 그리움의 감정이 계절의 순환 구조를 통해 효과적으로 드러난다.
}


\kfPassageNote{해설 학습 3}{%
「사미인곡」의 마지막 부분은 임을 향한 그리움이 가장 커지는 부분으로, 화자의 간절한 소망이 인상 깊게 나타난다. 화자는 임을 너무 그리워한 나머지 깊은 병에 걸렸다고 말한다. 이 병은 어떤 훌륭한 의사도 고칠 수 없고, 오직 ‘임’만이 낫게 할 수 있는 ‘상사병’이다.

화자는 이 깊은 그리움의 병을 이기지 못해 차라리 죽겠다고 다짐한다. 하지만 그냥 죽는 것이 아니라, 죽어서 ‘호랑나비’가 되겠다고 말한다. 여기서 호랑나비는 곧 화자 자신을 의미하며, 살아있을 때의 한계를 뛰어넘어 임에게 다가가려는 소망이 담긴 존재이다. 즉, 살아있는 몸으로는 임에게 갈 수 없으니, 죽어서 자유롭게 날아다니는 나비가 되어 임의 곁에 머물고 싶다는 뜻이다. 죽어서 '호랑나비'가 되겠다는 말에는 죽음의 한계마저 극복하여 임을 따르려는 적극적이고 의지적인 태도가 나타난다. 이는 임을 향한 화자의 사랑이 영원하고 절대적임을 보여준다.

호랑나비가 된 화자는 꽃의 향기를 날개에 묻혀 임의 옷에 옮겨 놓겠다고 말한다. 이는 죽어서도 임을 향한 변치 않는 사랑과 충성심을 전하려는 굳은 다짐을 보여준다. 설령 임이 나비가 자신인 줄 알아보지 못해도, 그저 곁에서 임을 따르는 것만으로도 좋다는 순수한 마음도 드러낸다.

이처럼 「사미인곡」은 죽음도 막을 수 없는 화자의 영원한 사랑을 노래하며 마무리된다. 이러한 모습은 임을 향한 화자의 마음이 얼마나 깊고 간절한지를 보여주며 독자에게 깊은 감동을 준다.
}


\kfSecHeader{kfAreaLiterature}{문제}{}

\par\medskip
\kfBeginMC
\begin{kfQBlock}{01}{객관식}
\kfQStem{다음 문장에서 '나'와 '님'이 가리키는 실제 대상으로 가장 적절한 것은?}
\begin{kfChoices}
\kfChoice 나: 여성, 님: 남편
\kfChoice 나: 작가 자신, 님: 임금
\kfChoice 나: 선녀, 님: 옥황상제
\kfChoice 나: 지상의 선녀, 님: 임금
\end{kfChoices}
\end{kfQBlock}
\begin{kfQBlock}{02}{객관식}
\kfQStem{다음 문장에서 서술어 '올랐더니'와 '내려왔는가'의 생략된 주어로 가장 적절한 것은?}
\begin{kfChoices}
\kfChoice 올랐더니: 님, 내려왔는가: 님
\kfChoice 올랐더니: 님, 내려왔는가: 나
\kfChoice 올랐더니: 나, 내려왔는가: 님
\kfChoice 올랐더니: 나, 내려왔는가: 나
\end{kfChoices}
\end{kfQBlock}
\begin{kfQBlock}{03}{객관식}
\kfQStem{다음 시구에서 '너'가 가리키는 시어로 가장 적절한 것은?}
\begin{kfChoices}
\kfChoice 매화
\kfChoice 달빛
\kfChoice 기러기
\kfChoice 북극성
\end{kfChoices}
\end{kfQBlock}
\begin{kfQBlock}{04}{객관식}
\kfQStem{다음 문장에서 서술어 '지어 내니'와 '갖추었구나'의 주어로 가장 적절한 것은?}
\begin{kfChoices}
\kfChoice 지어 내니: 님, 갖추었구나: 옷
\kfChoice 지어 내니: 님, 갖추었구나: 나
\kfChoice 지어 내니: 나, 갖추었구나: 옷
\kfChoice 지어 내니: 나, 갖추었구나: 님
\end{kfChoices}
\end{kfQBlock}
\begin{kfQBlock}{05}{객관식}
\kfQStem{다음 문장에서 '반가워하실까'의 주체로 가장 적절한 것은?}
\begin{kfChoices}
\kfChoice 나
\kfChoice 임
\kfChoice 매화
\kfChoice 기러기
\end{kfChoices}
\end{kfQBlock}
\begin{kfQBlock}{06}{객관식}
\kfQStem{다음 문장에서 '깊은 산골짜기'가 상징하는 대상으로 가장 적절한 것은?}
\begin{kfChoices}
\kfChoice 속세
\kfChoice 유배지
\kfChoice 한양 궁궐
\kfChoice 소외된 백성
\end{kfChoices}
\end{kfQBlock}
\begin{kfQBlock}{07}{객관식}
\kfQStem{다음 문장에서 '이곳'과 '높은 궁궐'이 가리키는 실제 장소로 가장 적절한 것은?}
\begin{kfChoices}
\kfChoice 이곳: 창평, 높은 궁궐: 한양
\kfChoice 이곳: 한양, 높은 궁궐: 창평
\kfChoice 이곳: 속세, 높은 궁궐: 천상계
\kfChoice 이곳: 유배지, 높은 궁궐: 달나라
\end{kfChoices}
\end{kfQBlock}
\begin{kfQBlock}{08}{객관식}
\kfQStem{다음 문장에서 '이 병'의 원인으로 가장 적절한 것은?}
\begin{kfChoices}
\kfChoice 노환
\kfChoice 가난
\kfChoice 임에 대한 그리움
\kfChoice 유배 생활의 고단함
\end{kfChoices}
\end{kfQBlock}
\begin{kfQBlock}{09}{객관식}
\kfQStem{}
\kfAnswerLines{2}
\end{kfQBlock}
\begin{kfQBlock}{10}{객관식}
\kfQStem{}
\kfAnswerLines{2}
\end{kfQBlock}
\begin{kfQBlock}{11}{객관식}
\kfQStem{}
\kfAnswerLines{2}
\end{kfQBlock}
\begin{kfQBlock}{12}{객관식}
\kfQStem{이 작품에 대한 설명으로 적절하지 \uline{\underline{않은}} 것은?}
\begin{kfChoices}
\kfChoice 임금에 대한 신하의 변함없는 충성심을 노래한다.
\kfChoice 작가는 정철로, 유배지에서 멀리 임금을 그리워하는 정서를 담는다.
\kfChoice 임과 이별한 여성 화자의 목소리에 충성심을 담아 표현한다.
\kfChoice 충신연주지사의 전통을 잘 보여 주는 작품이다.
\kfChoice 자연에 묻혀 사는 사대부의 한가로운 즐거움만을 노래한 강호 가사이다.
\end{kfChoices}
\end{kfQBlock}
\begin{kfQBlock}{13}{객관식}
\kfQStem{이 시의 화자가 궁극적으로 바라는 것에 대한 설명으로 적절하지 \uline{\underline{않은}} 것은?}
\begin{kfChoices}
\kfChoice 임과 다시 만나기를 간절히 소망한다.
\kfChoice 임과 멀리 떨어진 슬픔을 견디며 그리움을 노래한다.
\kfChoice 죽어서라도 임의 곁에 머물고 싶어 한다.
\kfChoice 임에 대한 변함없는 사랑과 충성을 다짐한다.
\kfChoice 임을 잊고 새로운 인연을 만나 행복을 얻기를 바란다.
\end{kfChoices}
\end{kfQBlock}
\begin{kfQBlock}{14}{객관식}
\kfQStem{이 작품의 화자에 대한 설명으로 적절하지 \uline{\underline{않은}} 것은?}
\begin{kfChoices}
\kfChoice 임과 이별한 상황에 처해 있다.
\kfChoice 죽어서라도 임의 곁에 있고 싶어 한다.
\kfChoice 여성의 목소리로 자신의 심정을 드러내고 있다.
\kfChoice 시간이 흘러도 임에 대한 마음이 변하지 않는다.
\kfChoice 자신의 처지를 운명으로 받아들이고 임을 단념한다.
\end{kfChoices}
\end{kfQBlock}
\begin{kfQBlock}{15}{객관식}
\kfQStem{이 시의 표현상 특징으로 적절하지 \uline{\underline{않은}} 것은?}
\begin{kfChoices}
\kfChoice 아름다운 우리말을 효과적으로 사용하고 있다.
\kfChoice 밝고 희망적인 분위기의 시어를 사용하고 있다.
\kfChoice 계절의 변화에 따라 시상을 전개하고 있다.
\kfChoice 여성 화자를 설정하여 애절함을 더하고 있다.
\kfChoice 자연물에 감정을 이입하여 마음을 표현하고 있다.
\end{kfChoices}
\end{kfQBlock}
\begin{kfQBlock}{16}{객관식}
\kfQStem{다음 ㉠\textasciitilde{}㉤ 중, 화자와 임 사이를 가로막는 장애물을 상징하는 것은?}
\begin{kfChoices}
\kfChoice ㉠구름
\kfChoice ㉡봄바람
\kfChoice ㉢기러기
\kfChoice ㉣북극성
\kfChoice ㉤대나무
\end{kfChoices}
\end{kfQBlock}
\begin{kfQBlock}{17}{객관식}
\kfQStem{화자가 임에게 보내고 싶어 한 소재와 그에 담긴 정서가 바르게 연결되지 \uline{\underline{않은}} 것은?}
\begin{kfChoices}
\kfChoice 옷 - 사랑과 정성
\kfChoice 매화 - 변치 않는 절개와 충심
\kfChoice 봄 햇살 - 임에 대한 걱정과 사랑
\kfChoice 달빛 - 임금의 은혜가 퍼지길 기원함
\kfChoice 호랑나비 - 이별의 슬픔과 깊은 한
\end{kfChoices}
\end{kfQBlock}
\begin{kfQBlock}{18}{객관식}
\kfQStem{화자가 ‘호랑나비’가 되려는 이유에 대한 설명으로 적절하지 \uline{\underline{않은}} 것은?}
\begin{kfChoices}
\kfChoice 살아생전에는 임에게 갈 수 없다는 절망감 때문이다.
\kfChoice 죽어서라도 임의 곁에 머물고 싶은 간절함 때문이다.
\kfChoice 자기 사랑이 사후에도 이어지길 바라는 마음 때문이다.
\kfChoice 임의 옷에 내려앉아 자신의 사랑을 전하고 싶기 때문이다.
\kfChoice 향기로운 꿀을 모아 임에게 선물로 바치기 위해서이다.
\end{kfChoices}
\end{kfQBlock}
\begin{kfQBlock}{19}{객관식}
\kfQStem{<보기>를 참고하여 이 시를 감상한 내용으로 적절하지 \uline{\underline{않은}} 것은?}
\begin{kfEvidence}'충신연주지사'는 신하가 임금을 그리워하는 마음을 남녀 간의 사랑에 빗대어 표현한 노래이다. 임금에 대한 충성심을 효과적으로 드러내기 위해, 작가들은 주로 이별한 여성을 화자로 설정하여 임을 향한 변치 않는 사랑을 노래했다.\end{kfEvidence}
\begin{kfChoices}
\kfChoice 작품 속 ‘임’은 임금을, ‘나’는 신하의 처지를 상징하는 인물이겠군.
\kfChoice 임을 향한 화자의 그리움은 곧 임금에 대한 충성심을 의미한다고 볼 수 있겠군.
\kfChoice 여성을 화자로 설정한 것은 사랑을 더 절실하게 표현하려는 의도이겠군.
\kfChoice 결국 이 작품은 임금을 향한 충성심을 전하려는 의도를 담은 노래이겠군.
\kfChoice ‘달나라 궁궐’은 신하가 사는 곳, ‘세상’은 임금의 궁궐을 비유한 것이겠군.
\end{kfChoices}
\end{kfQBlock}
\begin{kfQBlock}{20}{객관식}
\kfQStem{<보기>의 관점에서 이 작품을 이해한 것으로 가장 적절한 것은?}
\begin{kfEvidence}조선 시대 사대부들은 유교적 이념에 따라 임금에 대한 충성을 매우 중요한 가치로 여겼다. 그들은 임금을 향한 마음을 남녀 간의 사랑에 빗대어 표현하는 경우가 많았는데, 이를 통해 변치 않는 충성심을 효과적으로 드러내고자 했다.\end{kfEvidence}
\begin{kfChoices}
\kfChoice 화자가 겪는 그리움은 임금을 향한 충성심의 표현이라고 할 수 있군.
\kfChoice 화자가 임과 이별한 것은 실제 남녀 간의 사랑이 깨졌기 때문이겠군.
\kfChoice 화자가 여성으로 설정된 것은 당시 여성의 슬픈 삶을 보여주기 위함이군.
\kfChoice 화자가 호랑나비가 되겠다는 것은 비현실적인 사랑을 꿈꾸는 것이겠군.
\kfChoice 이 작품은 유교적 충성심보다 남녀 간의 사랑을 더 중요하게 여기고 있군.
\end{kfChoices}
\end{kfQBlock}
\begin{kfQBlock}{21}{서술형}
\kfQStem{이 작품의 화자를 여성으로 설정하여 얻을 수 있는 효과를 서술하시오.}
\kfAnswerNote{답안}{4}
\end{kfQBlock}
\begin{kfQBlock}{22}{서술형}
\kfQStem{이 시의 마지막에 화자가 '호랑나비'가 되겠다고 말한 것에 담긴 화자의 태도를 서술하시오.}
\kfAnswerNote{답안}{4}
\end{kfQBlock}
\kfEndMC


\kfStartNonlit{[기술] 항원-항체 반응을 이용한 검사용 키트의 원리}


\kfPassageNote{[기술] 항원-항체 반응을 이용한 검사용 키트의 원리}{%
　건강 상태를 진단하거나 범죄 현장에서 혈흔을 조사하기 위해 검사용 키트가 널리 이용된다. 항원-항체 반응을 응용하여 시료에 존재하는 성분을 분석하는 다양한 형태의 키트가 개발되고 있다. 항원-항체 반응은 항원과 그 항원에만 특이적으로 반응하는 항체가 결합하는 면역 반응을 말한다. 이는 마치 자물쇠와 열쇠처럼 딱 맞는 것끼리만 결합하는 원리이다.

측면유동면역분석법(LFIA) 키트를 이용하면 키트에 나타나는 선을 통해 목표 성분의 유무를 간편하게 확인할 수 있다. LFIA 키트는 시료 패드, 결합 패드, 반응막, 흡수 패드가 순서대로 배열된 구조로 되어 있다. 시료 패드로 흡수된 시료는 마치 물이 종이 위를 흐르듯이 결합 패드를 거쳐 반응막으로 이동한다. 결합 패드에 있는 복합체는 표지 물질에 특정 물질이 붙어 이루어진다. 표지 물질은 발색 반응에 의해 색깔을 낸다. 이는 마치 형광펜이 종이에 표시를 남기는 것과 비슷하다.

반응막에는 항체들이 띠 모양으로 두 가닥 고정되어 있는데, 시료 패드와 가까운 쪽에 있는 가닥이 검사선이고 다른 가닥은 표준선이다. 표지 물질이 검사선이나 표준선에 놓이면 발색 반응에 의해 반응선이 나타난다. 검사선이 발색되어 나타나는 반응선을 통해 목표 성분의 유무를 판정할 수 있다. 표준선은 검사가 제대로 진행되었는지 확인하는 역할을 한다.

LFIA 키트는 주로 직접 방식 또는 경쟁 방식으로 제작된다. 직접 방식은 찾고자 하는 물질이 있으면 선이 나타나는 방식이다. 시료에 목표 성분이 포함되어 있다면 검사선이 발색된다. 반면 경쟁 방식은 찾고자 하는 물질이 많으면 오히려 선이 나타나지 않는 방식이다. 시료에 목표 성분이 충분히 많다면 검사선이 발색되지 않는다.

검사용 키트의 정확성을 측정하기 위해 여러 번의 검사를 실시하고 그 결과를 분석한다. 키트가 시료에 목표 성분이 들어 있다고 판정하면 이를 양성이라고 한다. 이때 실제로 존재하면 진양성, 없다면 위양성이라고 한다. 위양성은 거짓 경보와 같다. 반대로 키트가 목표 성분이 없다고 판정했는데 실제로는 있는 경우를 위음성이라고 하며, 이는 놓친 경보와 같다.

정확도는 민감도와 특이도로 나뉜다. 민감도는 시료에 목표 성분이 존재하는 경우에 대해 키트가 이를 양성으로 판정한 비율이다. 민감도가 높다는 것은 찾아야 할 것을 잘 찾아낸다는 의미이다. 특이도는 시료에 목표 성분이 없는 경우에 대해 키트가 이를 음성으로 판정한 비율이다. 특이도가 높다는 것은 없는 것을 없다고 정확히 판정한다는 의미이다. 민감도와 특이도가 모두 높은 키트가 가장 이상적이지만 현실에서는 그렇지 않은 경우가 많다. 따라서 상황에 따라 민감도나 특이도를 고려하여 키트를 선택해야 한다.
}


\kfPassageNote{문장 독해 — 본문}{%
1. 건강 상태를 진단하거나 범죄 현장에서 혈흔을 조사하기 위해 검사용 키트가 널리 이용된다. 2. 항원-항체 반응을 응용하여 시료에 존재하는 성분을 분석하는 다양한 형태의 키트가 개발되고 있다. 3. 항원-항체 반응은 항원과 그 항원에만 특이적으로 반응하는 항체가 결합하는 면역 반응을 말한다. 4. 이는 마치 자물쇠와 열쇠처럼 딱 맞는 것끼리만 결합하는 원리이다. 5. 측면유동면역분석법(LFIA) 키트를 이용하면 키트에 나타나는 선을 통해 목표 성분의 유무를 간편하게 확인할 수 있다. 6. LFIA 키트는 시료 패드, 결합 패드, 반응막, 흡수 패드가 순서대로 배열된 구조로 되어 있다. 7. 시료 패드로 흡수된 시료는 마치 물이 종이 위를 흐르듯이 결합 패드를 거쳐 반응막으로 이동한다. 8. 결합 패드에 있는 복합체는 표지 물질에 특정 물질이 붙어 이루어진다. 9. 표지 물질은 발색 반응에 의해 색깔을 낸다. 10. 이는 마치 형광펜이 종이에 표시를 남기는 것과 비슷하다. 11. 반응막에는 항체들이 띠 모양으로 두 가닥 고정되어 있는데, 시료 패드와 가까운 쪽에 있는 가닥이 검사선이고 다른 가닥은 표준선이다. 12. 표지 물질이 검사선이나 표준선에 놓이면 발색 반응에 의해 반응선이 나타난다. 13. 검사선이 발색되어 나타나는 반응선을 통해 목표 성분의 유무를 판정할 수 있다. 14. 표준선은 검사가 제대로 진행되었는지 확인하는 역할을 한다. 15. LFIA 키트는 주로 직접 방식 또는 경쟁 방식으로 제작된다. 16. 직접 방식은 찾고자 하는 물질이 있으면 선이 나타나는 방식이다. 17. 시료에 목표 성분이 포함되어 있다면 검사선이 발색된다. 18. 반면 경쟁 방식은 찾고자 하는 물질이 많으면 오히려 선이 나타나지 않는 방식이다. 19. 시료에 목표 성분이 충분히 많다면 검사선이 발색되지 않는다. 20. 검사용 키트의 정확성을 측정하기 위해 여러 번의 검사를 실시하고 그 결과를 분석한다. 21. 키트가 시료에 목표 성분이 들어 있다고 판정하면 이를 양성이라고 한다. 22. 이때 실제로 존재하면 진양성, 없다면 위양성이라고 한다. 23. 위양성은 거짓 경보와 같다. 24. 반대로 키트가 목표 성분이 없다고 판정했는데 실제로는 있는 경우를 위음성이라고 하며, 이는 놓친 경보와 같다. 25. 정확도는 민감도와 특이도로 나뉜다. 26. 민감도는 시료에 목표 성분이 존재하는 경우에 대해 키트가 이를 양성으로 판정한 비율이다. 27. 민감도가 높다는 것은 찾아야 할 것을 잘 찾아낸다는 의미이다. 28. 특이도는 시료에 목표 성분이 없는 경우에 대해 키트가 이를 음성으로 판정한 비율이다. 29. 특이도가 높다는 것은 없는 것을 없다고 정확히 판정한다는 의미이다. 30. 민감도와 특이도가 모두 높은 키트가 가장 이상적이지만 현실에서는 그렇지 않은 경우가 많다. 31. 따라서 상황에 따라 민감도나 특이도를 고려하여 키트를 선택해야 한다.
}


\kfActivityBg \par\kfMaybeNeedspace{32\baselineskip}\noindent{\kfTipMark\,\,}{\small\color{kfInk} 다음 뜻풀이에 해당하는 어휘를 지문에서 찾아 쓰시오.}\par\nopagebreak[4]\smallskip\nopagebreak[4]
\begin{kfVocab}
\kfVocabItem{1}{항체 생성을 유도하는 물질}
\kfVocabItem{2}{특정 항원에만 결합하여 면역 작용을 하는 단백질}
\kfVocabItem{3}{액체 시료가 모세관 현상에 의해 이동하며 성분을 분석하는 방법}
\kfVocabItem{4}{검사 결과가 색으로 나타나도록 표시하는 물질}
\kfVocabItem{5}{화학 반응을 통해 색깔이 나타나는 현상}
\kfVocabItem{6}{키트에서 목표 성분의 유무에 따라 색이 나타나는 선}
\kfVocabItem{7}{키트가 정상적으로 작동했는지 보여주는 선}
\kfVocabItem{8}{검사 결과 목표 성분이 있다고 나온 판정}
\kfVocabItem{9}{실제 양성인 것을 양성으로 올바르게 판정한 결과}
\kfVocabItem{10}{실제 음성인 것을 양성으로 잘못 판정한 결과 (거짓 경보)}
\kfVocabItem{11}{실제 양성인 것을 음성으로 잘못 판정한 결과 (놓친 경보)}
\kfVocabItem{12}{목표 성분이 있는 경우, 이를 양성으로 판정하는 비율 (찾아야 할 것을 잘 찾아내는 정도)}
\kfVocabItem{13}{목표 성분이 없는 경우, 이를 음성으로 판정하는 비율 (없는 것을 없다고 잘 판정하는 정도)}
\end{kfVocab}

\kfActivityBg \par\kfMaybeNeedspace{32\baselineskip}\noindent{\kfTipMark\,\,}{\small\color{kfInk} 다음 글의 구조를 정리한 표의 빈칸을 채워 보세요.}\par\nopagebreak[4]\smallskip\nopagebreak[4]
\par\medskip\needspace{8\baselineskip}
\begin{tcolorbox}[enhanced, breakable=true,
  colback=kfPaper, colframe=kfRule, boxrule=0.4pt, arc=2pt,
  left=8pt, right=8pt, top=6pt, bottom=6pt,
  title={\footnotesize\bfseries\color{kfMute} 글의 구조도},
  coltitle=kfMute, colbacktitle=kfPrimaryLight!40,
  attach boxed title to top left={yshift=-1.5mm, xshift=6pt},
  boxed title style={colframe=kfRule, boxrule=0.3pt, arc=1pt}]
\footnotesize\setstretch{1.3}
\renewcommand{\arraystretch}{1.35}
\arrayrulecolor{kfRule}
\noindent\begin{tabularx}{\linewidth}{|>{\centering\arraybackslash}p{1.6cm}|>{\raggedright\arraybackslash}p{3.4cm}|>{\raggedright\arraybackslash}X|}
\hline
\rowcolor{kfPrimaryLight!40}\textbf{\color{kfPrimary}단} & \textbf{\color{kfPrimary}항목} & \textbf{\color{kfPrimary}내용} \\\hline
활용: 건강 진단, 범죄 현장 혈흔 조사 등 & 활용: 건강 진단, 범죄 현장 혈흔 조사 등 &  \\\hline
핵심 원리: ( ① )과 그에 특이적으로 반응하는 ( ② )가 결합하는 면역 반응 (자물쇠와 열쇠 비유) & 핵심 원리: ( ① )과 그에 특이적으로 반응하는 ( ② )가 결합하는 면역 반응 (자물쇠와 열쇠 비유) &  \\\hline
구조: 시료 패드 → 결합 패드 → 반응막 → 흡수 패드 & 구조: 시료 패드 → 결합 패드 → 반응막 → 흡수 패드 &  \\\hline
결합 패드: ( ③ )을 일으키는 표지 물질이 포함된 복합체가 있음 & 결합 패드: ( ③ )을 일으키는 표지 물질이 포함된 복합체가 있음 &  \\\hline
( ④ ): 목표 성분의 유무를 판정하는 선 & ( ④ ): 목표 성분의 유무를 판정하는 선 &  \\\hline
( ⑤ ): 검사가 제대로 진행되었는지 확인하는 선 & ( ⑤ ): 검사가 제대로 진행되었는지 확인하는 선 &  \\\hline
직접 방식: 목표 성분이 있으면 검사선이 ( ⑥ )됨 & 직접 방식: 목표 성분이 있으면 검사선이 ( ⑥ )됨 &  \\\hline
경쟁 방식: 목표 성분이 충분히 많으면 검사선이 ( ⑦ )되지 않음 & 경쟁 방식: 목표 성분이 충분히 많으면 검사선이 ( ⑦ )되지 않음 &  \\\hline
\multirow{2}{*}{( ⑧ ): 키트가 목표 성분이 있다고 판정} & 진양성 & 실제로 목표 성분이 있는 경우 (올바른 경보) \\\cline{2-3}
 & ( ⑨ ) & 실제로는 목표 성분이 없는 경우 (거짓 경보) \\\cline{2-3}
\hline
( ⑩ ): 키트가 목표 성분이 없다고 판정했는데 실제로는 있는 경우 (놓친 경보) & ( ⑩ ): 키트가 목표 성분이 없다고 판정했는데 실제로는 있는 경우 (놓친 경보) &  \\\hline
( ⑪ ): 목표 성분이 있는 경우, 양성으로 판정한 비율 & ( ⑪ ): 목표 성분이 있는 경우, 양성으로 판정한 비율 &  \\\hline
( ⑫ ): 목표 성분이 없는 경우, 음성으로 판정한 비율 & ( ⑫ ): 목표 성분이 없는 경우, 음성으로 판정한 비율 &  \\\hline
\end{tabularx}
\arrayrulecolor{black}
\end{tcolorbox}\par\medskip
\kfSecHeader{kfAreaNonlit}{문제}{}

\par\medskip
\kfBeginMC
\begin{kfQBlock}{01}{객관식}
\kfQStem{검사용 키트의 이용 예시로 가장 적절한 것은?}
\begin{kfChoices}
\kfChoice 식품 영양 성분 분석, 토양 오염 측정
\kfChoice 대기질 오염도 측정, 수질 오염 측정
\kfChoice 건강 상태 진단, 범죄 현장 혈흔 조사
\kfChoice 지하수 수질 검사, 전염병 확산 예측
\end{kfChoices}
\end{kfQBlock}
\begin{kfQBlock}{02}{객관식}
\kfQStem{검사용 키트가 응용한 원리로 가장 적절한 것은?}
\begin{kfChoices}
\kfChoice 중화 반응
\kfChoice 연소 반응
\kfChoice 산화-환원 반응
\kfChoice 항원-항체 반응
\end{kfChoices}
\end{kfQBlock}
\begin{kfQBlock}{03}{객관식}
\kfQStem{검사용 키트가 분석하는 대상으로 가장 적절한 것은?}
\begin{kfChoices}
\kfChoice 시료의 색깔
\kfChoice 시료의 온도
\kfChoice 시료의 무게
\kfChoice 시료의 성분
\end{kfChoices}
\end{kfQBlock}
\begin{kfQBlock}{04}{객관식}
\kfQStem{'항원-항체 반응'에 대한 설명으로 가장 적절한 것은?}
\begin{kfChoices}
\kfChoice 항원과 항체가 서로 분리되는 반응
\kfChoice 항원이 항체를 공격하여 없애는 반응
\kfChoice 항체가 스스로 복제하여 늘어나는 반응
\kfChoice 항원과 항체가 특이적으로 결합하는 반응
\end{kfChoices}
\end{kfQBlock}
\begin{kfQBlock}{05}{객관식}
\kfQStem{항체가 결합하는 항원으로 가장 적절한 것은?}
\begin{kfChoices}
\kfChoice 크기가 자신보다 작은 항원
\kfChoice 체내에 들어온 모든 종류의 항원
\kfChoice 자신에게만 특이적으로 반응하는 항원
\kfChoice 모양이 둥근 형태를 가진 모든 항원
\end{kfChoices}
\end{kfQBlock}
\begin{kfQBlock}{06}{객관식}
\kfQStem{항원-항체 반응의 원리를 비유한 것으로 가장 적절한 것은?}
\begin{kfChoices}
\kfChoice 물과 기름
\kfChoice 펜과 종이
\kfChoice 자석과 철
\kfChoice 자물쇠와 열쇠
\end{kfChoices}
\end{kfQBlock}
\begin{kfQBlock}{07}{객관식}
\kfQStem{LFIA 키트에서 목표 성분의 유무를 확인하는 방법으로 가장 적절한 것은?}
\begin{kfChoices}
\kfChoice 키트의 무게 변화
\kfChoice 키트의 온도 변화
\kfChoice 키트에 나타나는 선
\kfChoice 키트의 진동 발생
\end{kfChoices}
\end{kfQBlock}
\begin{kfQBlock}{08}{객관식}
\kfQStem{LFIA 키트의 구성 요소를 순서대로 바르게 나열한 것은?}
\begin{kfChoices}
\kfChoice 결합 패드 → 시료 패드 → 반응막 → 흡수 패드
\kfChoice 시료 패드 → 결합 패드 → 반응막 → 흡수 패드
\kfChoice 반응막 → 결합 패드 → 시료 패드 → 흡수 패드
\kfChoice 흡수 패드 → 반응막 → 결합 패드 → 시료 패드
\end{kfChoices}
\end{kfQBlock}
\begin{kfQBlock}{09}{객관식}
\kfQStem{시료가 반응막까지 이동하는 경로로 가장 적절한 것은?}
\begin{kfChoices}
\kfChoice 시료 패드 → 반응막 → 결합 패드
\kfChoice 시료 패드 → 결합 패드 → 반응막
\kfChoice 반응막 → 결합 패드 → 시료 패드
\kfChoice 결합 패드 → 시료 패드 → 반응막
\end{kfChoices}
\end{kfQBlock}
\begin{kfQBlock}{10}{객관식}
\kfQStem{시료가 이동하는 모습을 비유한 것으로 가장 적절한 것은?}
\begin{kfChoices}
\kfChoice 새가 하늘을 날아가는 것
\kfChoice 물이 종이 위를 흐르는 것
\kfChoice 차가 도로 위를 달리는 것
\kfChoice 공이 비탈길을 굴러가는 것
\end{kfChoices}
\end{kfQBlock}
\begin{kfQBlock}{11}{객관식}
\kfQStem{결합 패드에 있는 '복합체'의 구성으로 가장 적절한 것은?}
\begin{kfChoices}
\kfChoice 표지 물질, 항원
\kfChoice 항체, 반응막
\kfChoice 표지 물질, 특정 물질
\kfChoice 시료 패드, 흡수 패드
\end{kfChoices}
\end{kfQBlock}
\begin{kfQBlock}{12}{객관식}
\kfQStem{'표지 물질'의 역할로 가장 적절한 것은?}
\begin{kfChoices}
\kfChoice 시료를 빠르게 이동시키는 역할
\kfChoice 발색 반응에 의해 색깔을 내는 역할
\kfChoice 항원과 항체를 분리하는 역할
\kfChoice 키트의 온도를 일정하게 유지하는 역할
\end{kfChoices}
\end{kfQBlock}
\begin{kfQBlock}{13}{객관식}
\kfQStem{표지 물질이 색깔을 내는 현상을 비유한 것으로 가장 적절한 것은?}
\begin{kfChoices}
\kfChoice 풀로 종이를 붙이는 것
\kfChoice 칼로 종이를 자르는 것
\kfChoice 지우개로 글씨를 지우는 것
\kfChoice 형광펜이 종이에 표시를 남기는 것
\end{kfChoices}
\end{kfQBlock}
\begin{kfQBlock}{14}{객관식}
\kfQStem{반응막에 고정된 두 가닥의 선으로 가장 적절한 것은?}
\begin{kfChoices}
\kfChoice 검사선, 표준선
\kfChoice 시료선, 결합선
\kfChoice 반응선, 흡수선
\kfChoice 양성선, 음성선
\end{kfChoices}
\end{kfQBlock}
\begin{kfQBlock}{15}{객관식}
\kfQStem{시료 패드와 비교했을 때 '검사선'의 위치로 가장 적절한 것은?}
\begin{kfChoices}
\kfChoice 먼 쪽
\kfChoice 가까운 쪽
\kfChoice 중간 지점
\kfChoice 가장자리
\end{kfChoices}
\end{kfQBlock}
\begin{kfQBlock}{16}{객관식}
\kfQStem{'반응선'이 나타나는 경우로 가장 적절한 것은?}
\begin{kfChoices}
\kfChoice 표지 물질이 물과 반응했을 때
\kfChoice 표지 물질이 공기 중에 노출되었을 때
\kfChoice 표지 물질이 시료 패드에 흡수되었을 때
\kfChoice 표지 물질이 검사선이나 표준선에 놓였을 때
\end{kfChoices}
\end{kfQBlock}
\begin{kfQBlock}{17}{객관식}
\kfQStem{목표 성분의 유무를 판정하는 기준으로 가장 적절한 것은?}
\begin{kfChoices}
\kfChoice 표준선의 색깔 변화
\kfChoice 시료 패드의 색깔 변화
\kfChoice 흡수 패드의 색깔 변화
\kfChoice 검사선이 발색되어 나타나는 반응선
\end{kfChoices}
\end{kfQBlock}
\begin{kfQBlock}{18}{객관식}
\kfQStem{'표준선'의 역할로 가장 적절한 것은?}
\begin{kfChoices}
\kfChoice 목표 성분의 종류를 구별하는 역할
\kfChoice 검사 시간을 단축시키는 역할
\kfChoice 시료의 오염 여부를 확인하는 역할
\kfChoice 검사가 제대로 진행되었는지 확인하는 역할
\end{kfChoices}
\end{kfQBlock}
\begin{kfQBlock}{19}{객관식}
\kfQStem{LFIA 키트의 두 가지 제작 방식으로 가장 적절한 것은?}
\begin{kfChoices}
\kfChoice 직접 방식, 경쟁 방식
\kfChoice 수동 방식, 자동 방식
\kfChoice 개별 방식, 통합 방식
\kfChoice 화학 방식, 물리 방식
\end{kfChoices}
\end{kfQBlock}
\begin{kfQBlock}{20}{객관식}
\kfQStem{'직접 방식'에서 선이 나타나는 경우로 가장 적절한 것은?}
\begin{kfChoices}
\kfChoice 시료의 양이 충분할 때
\kfChoice 찾고자 하는 물질이 있을 때
\kfChoice 찾고자 하는 물질이 없을 때
\kfChoice 찾고자 하는 물질이 적을 때
\end{kfChoices}
\end{kfQBlock}
\begin{kfQBlock}{21}{객관식}
\kfQStem{'직접 방식'에서 검사선이 발색되는 경우로 가장 적절한 것은?}
\begin{kfChoices}
\kfChoice 검사가 실패했을 때
\kfChoice 시료가 오염되었을 때
\kfChoice 시료에 목표 성분이 없을 때
\kfChoice 시료에 목표 성분이 포함되어 있을 때
\end{kfChoices}
\end{kfQBlock}
\begin{kfQBlock}{22}{객관식}
\kfQStem{'경쟁 방식'에서 선이 나타나지 않는 경우로 가장 적절한 것은?}
\begin{kfChoices}
\kfChoice 시료가 부족할 때
\kfChoice 찾고자 하는 물질이 많을 때
\kfChoice 찾고자 하는 물질이 없을 때
\kfChoice 찾고자 하는 물질이 적을 때
\end{kfChoices}
\end{kfQBlock}
\begin{kfQBlock}{23}{객관식}
\kfQStem{'경쟁 방식'에서 검사선이 발색되지 않는 경우로 가장 적절한 것은?}
\begin{kfChoices}
\kfChoice 키트가 불량일 때
\kfChoice 시료에 불순물이 섞여 있을 때
\kfChoice 시료에 목표 성분이 전혀 없을 때
\kfChoice 시료에 목표 성분이 충분히 많을 때
\end{kfChoices}
\end{kfQBlock}
\begin{kfQBlock}{24}{객관식}
\kfQStem{검사용 키트의 정확성을 측정하기 위한 방법으로 가장 적절한 것은?}
\begin{kfChoices}
\kfChoice 더 비싼 키트를 사용해야 한다.
\kfChoice 검사 비용을 줄이는 방법을 찾아야 한다.
\kfChoice 한 번의 검사로 신속하게 결론을 내려야 한다.
\kfChoice 여러 번의 검사를 실시하고 그 결과를 분석해야 한다.
\end{kfChoices}
\end{kfQBlock}
\begin{kfQBlock}{25}{객관식}
\kfQStem{키트가 '양성'으로 판정하는 경우로 가장 적절한 것은?}
\begin{kfChoices}
\kfChoice 검사 결과가 불확실한 경우
\kfChoice 재검사가 필요한 경우
\kfChoice 시료에 목표 성분이 없다고 판정하는 경우
\kfChoice 시료에 목표 성분이 들어 있다고 판정하는 경우
\end{kfChoices}
\end{kfQBlock}
\begin{kfQBlock}{26}{객관식}
\kfQStem{'진양성'에 대한 설명으로 가장 적절한 것은?}
\begin{kfChoices}
\kfChoice 음성으로 판정했는데 실제로 목표 성분이 없는 경우
\kfChoice 양성으로 판정했는데 실제로 목표 성분이 없는 경우
\kfChoice 음성으로 판정했는데 실제로 목표 성분이 존재하는 경우
\kfChoice 양성으로 판정했는데 실제로 목표 성분이 존재하는 경우
\end{kfChoices}
\end{kfQBlock}
\begin{kfQBlock}{27}{객관식}
\kfQStem{'위양성'에 대한 설명으로 가장 적절한 것은?}
\begin{kfChoices}
\kfChoice 양성으로 판정했는데 실제로 목표 성분이 없는 경우
\kfChoice 음성으로 판정했는데 실제로 목표 성분이 없는 경우
\kfChoice 양성으로 판정했는데 실제로 목표 성분이 존재하는 경우
\kfChoice 음성으로 판정했는데 실제로 목표 성분이 존재하는 경우
\end{kfChoices}
\end{kfQBlock}
\begin{kfQBlock}{28}{객관식}
\kfQStem{'위양성'을 비유한 표현으로 가장 적절한 것은?}
\begin{kfChoices}
\kfChoice 늦은 경보
\kfChoice 거짓 경보
\kfChoice 진짜 경보
\kfChoice 놓친 경보
\end{kfChoices}
\end{kfQBlock}
\begin{kfQBlock}{29}{객관식}
\kfQStem{'위음성'에 대한 설명으로 가장 적절한 것은?}
\begin{kfChoices}
\kfChoice 목표 성분이 있다고 판정했는데 실제로는 없는 경우
\kfChoice 목표 성분이 있다고 판정했는데 실제로도 있는 경우
\kfChoice 목표 성분이 없다고 판정했는데 실제로는 있는 경우
\kfChoice 목표 성분이 없다고 판정했는데 실제로도 없는 경우
\end{kfChoices}
\end{kfQBlock}
\begin{kfQBlock}{30}{객관식}
\kfQStem{'위음성'을 비유한 표현으로 가장 적절한 것은?}
\begin{kfChoices}
\kfChoice 놓친 경보
\kfChoice 거짓 경보
\kfChoice 빠른 경보
\kfChoice 안전 경보
\end{kfChoices}
\end{kfQBlock}
\begin{kfQBlock}{31}{객관식}
\kfQStem{키트의 정확도를 나누는 두 가지 기준으로 가장 적절한 것은?}
\begin{kfChoices}
\kfChoice 정확도, 신속도
\kfChoice 민감도, 특이도
\kfChoice 경제성, 효율성
\kfChoice 안정성, 내구성
\end{kfChoices}
\end{kfQBlock}
\begin{kfQBlock}{32}{객관식}
\kfQStem{'민감도'에 대한 설명으로 가장 적절한 것은?}
\begin{kfChoices}
\kfChoice 시료가 오염된 경우
\kfChoice 검사가 실패한 경우
\kfChoice 시료에 목표 성분이 없는 경우
\kfChoice 시료에 목표 성분이 존재하는 경우
\end{kfChoices}
\end{kfQBlock}
\begin{kfQBlock}{33}{객관식}
\kfQStem{'민감도가 높다'는 것의 의미로 가장 적절한 것은?}
\begin{kfChoices}
\kfChoice 검사 속도가 빠르다는 의미
\kfChoice 검사 비용이 저렴하다는 의미
\kfChoice 찾아야 할 것을 잘 찾아낸다는 의미
\kfChoice 없는 것을 없다고 판정한다는 의미
\end{kfChoices}
\end{kfQBlock}
\begin{kfQBlock}{34}{객관식}
\kfQStem{'특이도'에 대한 설명으로 가장 적절한 것은?}
\begin{kfChoices}
\kfChoice 시료의 양이 많은 경우
\kfChoice 시료에 목표 성분이 없는 경우
\kfChoice 시료에 목표 성분이 있는 경우
\kfChoice 반응이 빠르게 일어나는 경우
\end{kfChoices}
\end{kfQBlock}
\begin{kfQBlock}{35}{객관식}
\kfQStem{'특이도가 높다'는 것의 의미로 가장 적절한 것은?}
\begin{kfChoices}
\kfChoice 있는 것을 있다고 판정한다는 의미
\kfChoice 미세한 성분까지 찾아낸다는 의미
\kfChoice 다양한 성분을 동시에 분석한다는 의미
\kfChoice 없는 것을 없다고 정확히 판정한다는 의미
\end{kfChoices}
\end{kfQBlock}
\begin{kfQBlock}{36}{객관식}
\kfQStem{가장 이상적인 검사용 키트의 조건으로 가장 적절한 것은?}
\begin{kfChoices}
\kfChoice 가격이 저렴한 키트
\kfChoice 민감도와 특이도가 모두 높은 키트
\kfChoice 민감도는 낮고 특이도는 높은 키트
\kfChoice 민감도는 높고 특이도는 낮은 키트
\end{kfChoices}
\end{kfQBlock}
\begin{kfQBlock}{37}{객관식}
\kfQStem{키트를 선택할 때 고려해야 할 점으로 가장 적절한 것은?}
\begin{kfChoices}
\kfChoice 제조사의 인지도
\kfChoice 키트의 디자인과 색상
\kfChoice 포장 재질의 내구성
\kfChoice 상황에 따른 민감도나 특이도
\end{kfChoices}
\end{kfQBlock}
\begin{kfQBlock}{38}{객관식}
\kfQStem{윗글의 내용과 일치하지 \uline{않는} 것은?}
\begin{kfChoices}
\kfChoice LFIA 키트는 항원-항체 반응을 응용한 것이다.
\kfChoice 직접 방식 키트는 목표 성분이 있으면 검사선에 색이 나타난다.
\kfChoice 경쟁 방식 키트는 목표 성분이 없으면 검사선에 색이 나타나지 않는다.
\kfChoice 민감도는 키트가 '병이 있는' 경우를 놓치지 않고 찾아내는 비율이다.
\kfChoice 특이도는 키트가 '병이 없는' 경우를 병이 없다고 정확히 판정하는 비율이다.
\end{kfChoices}
\end{kfQBlock}
\begin{kfQBlock}{39}{객관식}
\kfQStem{LFIA 키트의 작동 원리에 대한 설명으로 적절하지 \uline{\underline{않은}} 것은?}
\begin{kfChoices}
\kfChoice 결합 패드의 복합체가 시료액과 함께 반응막을 거쳐 흡수 패드로 이동한다.
\kfChoice 시료가 패드를 따라 흐르는 ‘측면 유동’ 원리를 이용한다.
\kfChoice 반응막의 검사선·표준선은 색 변화로 결과를 보여 준다.
\kfChoice 표지 물질은 결합 패드에서 목표 성분과 반응해 색을 만들 수 있다.
\kfChoice 흡수 패드는 시료 속 불순물을 골라내어 키트 밖으로 내보내는 정수 장치이다.
\end{kfChoices}
\end{kfQBlock}
\begin{kfQBlock}{40}{객관식}
\kfQStem{‘민감도’가 특히 중요하게 요구되는 상황에 대한 설명으로 적절하지 \uline{\underline{않은}} 것은?}
\begin{kfChoices}
\kfChoice 전염성이 매우 강한 질병을 조기에 선별해 내야 하는 검사 상황
\kfChoice 양성 환자를 한 명도 놓치지 않는 것이 중요한 검사 상황
\kfChoice 위음성을 줄이는 것이 사회적으로 큰 가치를 갖는 검사 상황
\kfChoice 질병 의심자를 폭넓게 걸러 내야 하는 1차 선별 검사 상황
\kfChoice 부작용이 큰 약을 처방하기 전 ‘병이 없음’을 확인하는 정밀 검사
\end{kfChoices}
\end{kfQBlock}
\begin{kfQBlock}{41}{객관식}
\kfQStem{윗글을 바탕으로 <보기>의 코로나19 자가 검사 키트를 이해한 것으로 적절하지 \uline{\underline{않은}} 것은?}
\begin{kfEvidence}코로나19 자가 검사 키트는 검사 결과 창에 C(표준선)와 T(검사선)가 있다. 바이러스가 시료에 포함된 경우, C와 T에 모두 선이 나타나 '양성'으로 판정한다. 바이러스가 없는 경우 C에만 선이 나타난다.\end{kfEvidence}
\begin{kfChoices}
\kfChoice 항원-항체 반응을 이용한 LFIA 키트의 일종이다.
\kfChoice 바이러스가 있으면 검사선이 나타나므로 '직접 방식'에 해당한다.
\kfChoice C선은 키트가 정상적으로 작동했는지를 보여주는 역할을 한다.
\kfChoice 실제 감염자를 검사했는데 T선이 나타나지 않았다면, 이 키트는 민감도가 100\%는 아닌 것이다.
\kfChoice T선이 나타났지만 실제로는 감염되지 않았다면, 이를 '위음성'이라고 한다.
\end{kfChoices}
\end{kfQBlock}
\begin{kfQBlock}{42}{객관식}
\kfQStem{윗글의 서술 방식에 대한 설명으로 적절하지 \uline{\underline{않은}} 것은?}
\begin{kfChoices}
\kfChoice ‘LFIA 키트’라는 개념을 정의한 뒤 그 구조와 작동을 설명한다.
\kfChoice 직접·경쟁 방식을 비교하며 작동 원리를 차례로 다룬다.
\kfChoice 정확도를 판단하는 ‘민감도·특이도’ 기준을 함께 소개한다.
\kfChoice 개념·구조·평가 기준을 체계적으로 연결하여 설명한다.
\kfChoice 시간 흐름에 따라 LFIA 기술이 발전해 온 역사를 추적하는 형식이다.
\end{kfChoices}
\end{kfQBlock}
\begin{kfQBlock}{43}{단답형}
\kfQStem{항원과 항체가 자물쇠와 열쇠처럼 결합하는 면역 반응은 무엇인가?}
\kfAnswerLines{1}
\end{kfQBlock}
\begin{kfQBlock}{44}{단답형}
\kfQStem{LFIA 키트를 구성하는 네 가지 요소를 순서대로 쓰시오.}
\kfAnswerLines{1}
\end{kfQBlock}
\begin{kfQBlock}{45}{단답형}
\kfQStem{결합 패드에 있으며, 형광펜처럼 발색 반응으로 색깔을 내는 물질은 무엇인가?}
\kfAnswerLines{1}
\end{kfQBlock}
\begin{kfQBlock}{46}{단답형}
\kfQStem{반응막에서 검사가 제대로 진행되었는지 확인하는 역할을 하는 선은 무엇인가?}
\kfAnswerLines{1}
\end{kfQBlock}
\begin{kfQBlock}{47}{단답형}
\kfQStem{목표 성분이 있을 때 검사선이 발색되는 방식은 무엇인가?}
\kfAnswerLines{1}
\end{kfQBlock}
\begin{kfQBlock}{48}{서술형}
\kfQStem{LFIA 키트의 '직접 방식'과 '경쟁 방식'은 목표 성분의 유무에 따라 검사선의 발색 결과가 어떻게 다른지 서술하시오.}
\kfAnswerNote{답안}{4}
\end{kfQBlock}
\kfEndMC



\end{document}
